% !TeX program = xelatex
\documentclass[11pt,a4paper]{ctexart}
\usepackage[margin=25mm,headheight=15pt]{geometry}
\usepackage{amsmath,amssymb,amsthm,mathtools,booktabs,longtable,array,fancyhdr,xcolor,pdflscape}
\usepackage[unicode,colorlinks=true,linkcolor=blue!45!black,citecolor=blue!45!black,urlcolor=blue!45!black]{hyperref}
\newtheoremstyle{cnplain}{7pt}{7pt}{\normalfont}{}{\bfseries}{。}{0.5em}{}
\theoremstyle{cnplain}\newtheorem{theorem}{定理}[section]
\newtheorem{proposition}[theorem]{命题}
\newtheorem{lemma}[theorem]{引理}
\newtheorem{corollary}[theorem]{推论}
\theoremstyle{definition}\newtheorem{definition}[theorem]{定义}
\theoremstyle{remark}\newtheorem{remark}[theorem]{注}
\renewcommand{\proofname}{证明}
\newcommand{\Q}{\mathbb Q}\newcommand{\cC}{\mathcal C}
\newcommand{\cS}{\mathcal S}\newcommand{\dR}{\mathrm{dR}}
\newcommand{\rk}{\operatorname{rk}}\newcommand{\im}{\operatorname{im}}
\newcommand{\cA}{\mathcal A}
\newcommand{\dd}{\mathrm d}
\setlength{\emergencystretch}{3em}\allowdisplaybreaks[2]
\pagestyle{fancy}\fancyhf{}
\fancyhead[L]{\small 费曼周期：统一主定理与结构预测}
\fancyhead[R]{\small 研究主线}\fancyfoot[C]{\thepage}
\title{费曼周期的结构预测与物理可见性\\[4pt]
\large 几何上下界、重整化的结构保持与主线闭合}
\author{费曼周期研究项目\quad 理论主线稿}
\date{2026年10月14日}
\begin{document}
\maketitle
\begin{abstract}
给定带质量、运动学、分子、积分链、方案和可观测量的费曼图，
本文研究在不先求完整积分时能够严格预测的周期结构。
二十三行可核查图族分类表进一步纳入原物理标架比较、严格误差状态及先行工作，把MZV肯定证据、非Tate结构、物理可见性、完整扩张和特殊点分别列出，并给每行的证明入口。
分类对象是实际周期的最小标架子商；纯几何、算术下降、混合扩张和物理可见性分别记录。
核心比较定理证明：若两个形式周期之差由给定Serre背景承载，
则其最小对象在相应Serre商中同构，背景外简单重数相同；
在socle分离条件下，完整混合对象还受到商意义的保护。
几何载体给上界，实际管变分和链的全局轨道给下界，
相对留数把非齐次源与固定环境分开。

对明确接受的覆盖、串联、拼接和局部插入图类，
本文构造实际下降与比较，而非只赋予几何标签。
阈值分离的两圈theta族由分子网格值决定椭圆支撑；
若保留$r$个核，其完整函数在$\mathbb C(s)$上的最小齐次阶数为$2r+1$，
微分Galois群为$\mathbb G_a^{2r}\rtimes(\mathrm{SL}_2)^r$。
固定圈数虽给有限裸骨架，同一骨架的可见秩仍可随质量列表增长。
二维四传播子banana的已构造秩八对象作为核心混合实例，
其固定纤维算术结论与函数单值结论保持区分。
实际四重商对应进一步将banana纯核与四粒子谱曲面识别。
在指定非等质量射线上，每个正偶数维的单位四粒子谱都恰保留一份六维K3因子：
升维改变代数微分，正常阈值的秩一变分排除物理读出将其消去。
四维三次减除的二点色散读出亦保留该因子，有限局部多项式不能移除它。
对只依赖两对不变质量、传播子极点位于相应阈值以下的完整实树振幅，
非零合并分子当且仅当正谱保留一份K3因子；
有限阈值喷射精确决定首项，新增极点的位置给椭圆边界候选。
接触项与交换图的首项相消只提高阈值幂次的例子，以及非零有符号权积分恒零的反例，
分别说明该规则的物理用途与必要范围。
进一步对一列稳定的简单交换极点，实际末端夹挤把物理变分识别为留数周期。
平方极部的首项不能被接触干涉抵消，故保留$r$个不同极点时，
非Artin--Tate半单部分恰为K3核加$r$个不同的椭圆因子；
完整函数的混合扩张与一般角依赖振幅仍不由这一判据决定。
进一步在稳定质量窗口内引入自由旁观场和局部源，
以固定一粒子矩阵元的在壳条件完成有限阶全部中间态求和。
所考察三圈源关联函数的完整非局部谱恰为上述四粒子谱，
因此全部可见核成为真实未割系数的下界。
有限子反项插回后可以产生椭圆sunset背景；
正常阈值的局部变分将该背景消去，从而保护割所生成的全部组成。

物理出口需要额外的结构规则。
有限森林差由完整系数张量与装饰商图承载；
非线性读出扩大背景，相空间积分又可能产生新的几何。
本文给有理树权的实际质量壳载体、二体有理角权的Kummer上界，
以及三体正链的椭圆可见性。
进一步对既有四维$\phi^4$两圈基准，
将阶$g^4$弹性总截面的交叉泡角积分构造成双Kummer系数的属一周期；
两次顺序续接给非零纯椭圆读出，排除其它项相消，
并证明固定极点质量的有限耦合方案变化不能消去该因子。
带系数角色过滤进一步给一般单变量对数投影的纯核上界，
并确定该弹性系数恰含一份椭圆简单因子；完整混合扩张另行保留。

RG出口进一步按完整森林方向、固定点漂移和临界谱分三步比较；
简单谱时有限迹检验给精确一阶不可见核，
共振及参数正则性则界定了“指数不变”与“仅换方案”的差别。
实际多标量乘积固定点又给全阶交叉尺度组合律；
观察图的奇圈条件精确判定能否恢复各块热指数，
已有六圈形式MZV方向在全部正整数分量的交叉指数中保留。
有质量欧氏BPHZ完整森林再给一般有限圈图的实际周期入口；
有限局部喷射构造相容方案比较，二次发散必须消去整个局部多项式空间后才获得低圈背景保护。
合并森林进一步给保持形式周期的有理关联空间、有限有理参数表达及光滑Tate纤维；
贪心扇区与动量同步缩放进一步将所有秩亏面化为正则边界，
在正质量欧氏绝对收敛接受域内，任意有限圈数的完整标架比较成立。
这些结果组成有实际模型、适用条件和输出强度的有限推导规则。
有质量接受域以外的森林、任意圈振幅的完整物理推前及全部固定点算术仍有明确缺口。
本文不把接受类的定理冒充任意图自动分类，
也不把动机障碍改写为未经证明的数值超越性。
进一步把完整减除积分放入固定三维代数空间，
实际留数正合列给全部运动学组成的上界，物理管链给匹配下界。
在所声明的源归一化下，完整三圈系数随运动学变化的无限单值简单因子
恰为一份K3及每个保留交换极点的一份椭圆核。
环境常数和混合扩张仍保留，几何局部系统分类不冒充固定纤维算术分类。

对一般标量积多项式分子，协变球面偶矩给精确有理角下降，
没有增加两对几何之外的根式或有限极点。
质量壳上的多项式非零性因而成为四维正谱保留一份K3的有限充要判据。
任意有限角导数四场源具有真实的完整三圈二点出口，
其物理零方向由有限多项式系数矩阵给出。
对于含角传播子的真实$\phi^4$四体产生，
共同方向上的正线性分母又给全部干涉的精确旋转积分核。
该核由一次对数和代数根式组成；合并极部给严格零核判据，
并证明四体产生在整个阈值以上具有正的首项系数。
剩余五变量形状推前仍有独立的几何分类缺口。
对于核心二维banana，原正参数与正谱色散载体进一步由实际有理商和有向链比较。
八张正分支给出形式周期关系$I^{\mathfrak m}=8J^{\mathfrak m}$，
从而指定固定纤维的完整秩八对象及coaction转移到色散读出。
实际能量留数固定欧氏动量积分的Lefschetz归一化；
任意有限边数的二叉分组又给不依赖分组树的形式周期比较规则。
这一接受类的载体构造不代替高维纯核的完整分类。
偶数维的能量MOM代表现经实际代数球面推前比较到同一谱双覆盖。
最少减除数为$(D/2-1)(n-1)$，角纤维的Lefschetz因子被显式消去；
四维四线比较系数恰为一。其自由复合源的物理MOM实现得到证明，
独立森林代表之间的全部形式周期比较仍单独保留。


原假设11.2的几何输入通过同伦纤维推出、共轭固定正链与形式周期态射关系，
连接到原积分的实际余作用；指定秩八最小性保留独立Picard与非分裂输入。
等质量四体20点的角截断可严格控制到$6\times10^{-17}$以内，
但剩余积分尚无窄区间，三圈前向全切割亦单独保留。
文献对照将已知结果、应用、差异候选及开放项分开，不以独立推导替代查新。

\end{abstract}
\tableofcontents

\section{总目标的数学对象}
\subsection{输入与分类输出}
输入为
\[
 \mathfrak P=(G,D,\nu,\mu,\kappa,N,\sigma,k,R,\mathcal O),
\]
分别包含图、维数、传播子幂次、质量、运动学、分子、链、基域、减除方案及可观测量。
输出首先是带标架的实际对象，不是仅由边表决定的“MPL/椭圆/K3”单选标签。

在已构造的有理Nori范畴或明确指定的实现范畴$\cC$中，采用同调方向：
$E$承载链$\sigma\in E_B$，形式为读出$\lambda\in E_{\dR}^{\vee}$。
de Rham系数可在数域$k$，实际子对象仍在原范畴中求，保留下降条件。
周期记作$p=[E,\sigma,\lambda]$，此记号省略了与上同调约定之间的对偶。定义
\[
 R_\sigma=\bigcap_{\sigma\in J_B}J,\quad
 I_\lambda=\sum_{\lambda|J_{\dR}=0}J,\quad
 M(p)=R_\sigma/(R_\sigma\cap I_\lambda).
\]
有限长度与实现正合性保证存在。最小矩阵系数对象的表示无关性是经典结果
\cite[\S2.4]{brown}，不是本项目新创的定义定理。

完整分类包含三组数据：简单组成、权重、Hodge数、极化与算术下降；
各层间的实际粘合与多层相容关系；标架及每个操作的源、靶、实际态射。
只列简单组成会漏掉Kummer和regulator；只列相邻$\operatorname{Ext}^1$
也未必确定多层对象。精确输出仍是带过滤和标架的实际对象或等价的已证明表示。

\subsection{结构表的分层语义}
\begin{center}\small
\begin{tabular}{p{24mm}p{47mm}p{72mm}}
\toprule 名称&合法结论的层次&不能由名称单独推出\\\midrule
MPL&指定字母、路径和端点的迭代积分表示&任意代数字母的值都是MZV\\
MZV&指定归一化的常数及其形式周期&全部mixed Tate周期属于该子类\\
mixed Tate&实际对象的组成均为相应Tate对象&扩张分裂或自动有指定MPL表达式\\
椭圆/高亏格&曲线$H^1$及可见因子&同亏格等于同对象或全部因子物理可见\\
K3&指定超越$H^2$及边界粘合&秩六方程必为K3或整个周期只有纯层\\
较高维CY&指定几何、次数与上同调证书&任意高阶微分方程都是CY\\
Kummer&单位及其实际扩张类&不需基域、分支和比较数据\\\bottomrule
\end{tabular}
\end{center}
例如质量单位$2$的Kummer扩张与特定MZV端点对象有不同的算术输入。
这一说明不需要断言$\log2$与某组普通数值之间的未知代数独立性。
本表规定输出语义；充分判据仍须由实际几何给出。

\section{主定理一：加上已知结构，不能消去什么}
\subsection{背景子范畴与正确的商}
设$\cS\subset\cC$为Serre子范畴，即对子对象、商与扩张封闭。
记精确商函子$\pi:\cC\to\cC/\cS$，采用标准存在性定理
\cite[Tag 02MS]{stacks}。该商只比较哪些组成超出背景结构，
一般不是带纤维函子的动机范畴。

若背景含单位对象，就不能同时要求它是对任意环境对象张量封闭的
“张量理想”：$1\otimes E=E$会使其等于全部$\cC$。
下文用给定背景张量子范畴的\emph{Serre闭包}，没有这一要求。
在有限长度范畴中，它恰由全部简单组成属于背景清单的对象构成。

\begin{theorem}[背景周期的加性稳定性]\label{thm:addition}
设$p$为实际形式周期，$\ell=[L,\tau,\eta]$，$L\in\cS$，
$q=p+\ell$。则
\begin{equation}\label{eq:quotient}
 \pi M(q)\simeq\pi M(p).
\end{equation}
对任何不属于$\cS$的简单对象$S$，有
\[
 [M(q):S]=[M(p):S].
\]
另有与给定扰动载体有关的秩界
\[
 |\rk M(q)-\rk M(p)|\leq\rk L .
\]
无需计算$p,\ell$的复数值，也不需要周期映射单射性。
\end{theorem}
\begin{proof}
取$P=M(p)$的最小标架$(\sigma,\lambda)$。
于是$\sigma$生成$P$，且$P$中没有被$\lambda$消去的非零子对象。
$q$在$P\oplus L$上由链$(\sigma,\tau)$和读出$(\lambda,\eta)$表示。
令$R$为该链生成的子对象，$I\subset R$为该读出消去的最大子对象，
则$M(q)=R/I$。

投影$R\to P$的像承载$\sigma$，故满射；其核$K$是$L$的子对象，
所以$K\in\cS$。另一方面，
\[
 I\cap(P\oplus0)=0,
\]
因为该交是$P$的实际子对象，在其上$\lambda$为零，违背最小性。
因此投影$I\to L$为单射，故$I\in\cS$。两条短正合列
\[
 0\to K\to R\to P\to0,\qquad
 0\to I\to R\to M(q)\to0
\]
在$\pi$下成为同构，证明\eqref{eq:quotient}。
组成重数的可加性给
$[M(q):S]-[P:S]=[K:S]-[I:S]$；
对背景外的$S$右侧为零。
秩差为$\rk K-\rk I$，二者都介于$0$和$\rk L$，得秩界。
这些均是实际子商，不是数值相加后猜测对象。
\end{proof}

该结论强于“加入背景后生成同一张量范畴”：
后者不决定新简单因子在对象中出现几次，本定理保留精确重数。
背景可代表低圈图、已知边界或其它允许结构，
但须先证明扰动确由该背景内的实际对象承载。

\subsection{为何仍必须单独研究混合扩张}
若$\cS$为所有组成均为Tate的对象，则非分裂Kummer对象和其分裂替代都映为零。
例如$p=\log^{\mathfrak m}2$、$\ell=-p$允许$q=0$：
它们在商中都是零，原最小对象不同。
所以该商不替代regulator、单位和完整coaction，
只精确回答背景之外的新结构能否被加性变化消掉。
完整分类要把这一稳定纯组成与实际粘合并列保留。

\section{主定理二：锚定重整化保留首次出现的几何}
\subsection{实际图类与低圈背景}
采用本项目已证明的正质量欧氏$D=4$、单位传播子、至多对数发散图类，
使用已构造的零动量或相容代数非零MOM参考方案。
令$a,b$为两方案，$h_G=h>0$，固定共同物理锚点$\kappa_*$：
\[
 p_a=\widehat\Phi_a^{\mathfrak m}(G;\kappa,\kappa_*)
 =\Phi_a^{\mathfrak m}(G;\kappa)-\Phi_a^{\mathfrak m}(G;\kappa_*).
\]
当前点、锚点及参考数据在已证正质量代数欧氏范围。
改变基域或代数规范时，所需算术对象须显式纳入背景。

前批实际有限方案定理给
\begin{equation}\label{eq:forest}
 p_b-p_a=\sum_{\varnothing\ne\gamma\subsetneq G}
 g_{ba}^{\mathfrak m}(\gamma)\,
 \widehat\Phi_a^{\mathfrak m}
 (G/\gamma;\pi_\gamma\kappa,\pi_\gamma\kappa_*).
\end{equation}
子图可为森林余积允许的不连通并，每个因子均来自严格低圈实际载体。
经典森林背景见\cite{BK}；实际周期提升与截断物理比较由前批证明承担。

令$\mathcal B_{<G}$为本次转换所需有限子图--收缩闭包、
参考点和当前/锚点评价的全部严格低圈载体生成的满张量子范畴，
包含所需Tate因子；$\cS_{<G}$为其在共同环境中的Serre闭包。
“新几何”总是相对于这个明确背景，不能只按圈数口头指定。

\begin{theorem}[新结构的方案不变性]\label{thm:ren}
在上述范围，
\[
 \pi_{<G}M(p_a)\simeq\pi_{<G}M(p_b).
\]
每个背景外简单对象$S$的重数相同。
若全部低圈背景mixed Tate，而某方案的锚定最小对象含非Tate简单组成，
另一方案仍含同一组成及同样重数。
\end{theorem}
\begin{proof}
式\eqref{eq:forest}每项乘积由实际载体的张量积承载，
有限和由这些张量积的直和$L$承载。
故$L\in\mathcal B_{<G}\subset\cS_{<G}$。
应用定理\ref{thm:addition}即可。
若低圈均mixed Tate，其张量、子商及扩张仍只有Tate组成，
因此非Tate简单因子在背景外。
\end{proof}

这直接接合结构表与物理操作：一旦几何证明某个新纯因子进入指定标架，
允许的锚定方案变化不能将其改成背景类型。
没有求完整积分，也没有把局部性误解为所有常数消失。

\subsection{锚点条件及物理求和的区别}
未锚定差含整图常数$g_{ba}(G)$，
可含同圈数参考周期，不能归到低圈。
本原图已经有
\[
 \Phi_b(G;\kappa)-\Phi_a(G;\kappa)
 =\Phi_0(G;\rho_a)-\Phi_0(G;\rho_b).
\]
bubble在自己的减除点为零，改参考后该点可严格非零，
所以孤立最小对象可由零变为非零。
物理分子或完整图和的改变也不是方案变化；
它们可能在同一新几何内部真实相消，本定理没有禁止这种现象。

\section{主定理三：几何上界与物理下界的夹逼}
\subsection{有限实际操作的上界}
\begin{proposition}[结构上界]\label{prop:upper}
给定有限组实际种子，允许有限直和、张量、对偶、Tate扭转、
实际态射的核像商及已给出的短正合列。
若其组成属于对子商、张量、对偶及扩张封闭的背景类，
则全部输出及物理最小对象的组成仍属于该类。
\end{proposition}
\begin{proof}
直和与正合列的重数可加，子商只删组成。
张量积对有限长度过滤正合，其相继商是组成的张量积；
再取组成，连同对偶和扭转即可归纳。
物理最小化最后也是子商。
\end{proof}
扩张仍可增加regulator与幺幂深度。
切割、收缩及Gauss--Manin导数也不都是固定对象的自同态：
新几何必须给其实际对象与比较，而非从形式相似推断。

\subsection{链相容的割给强制下界}
设实际管态射$f:C\to E$已包含正确Tate扭转与取向。
割读出为$f^*\lambda$，割链为$\gamma$，
其像在$E$中生成$R_{f\gamma}$。
\begin{proposition}[割的物理保护]\label{prop:lower}
若$R_{f\gamma}\subset R_\sigma$，则有实际单射
\[
 M([C,\gamma,f^*\lambda])\hookrightarrow M([E,\sigma,\lambda]).
\]
故割最小对象的每个简单组成在物理最小对象中至少有同样重数。
\end{proposition}
\begin{proof}
实际标架关系识别左侧与
$R_{f\gamma}/(R_{f\gamma}\cap I_\lambda)$，
再由包含关系注入$R_\sigma/(R_\sigma\cap I_\lambda)$。
\end{proof}
原链循环生成$E$时条件自动满足。
否则，非零广义割可能来自原链没有激发的方向，只能证明环境几何。
完整物理求和必须先形成总读出，再检验割，不能逐图忽略干涉。

\begin{corollary}[不求完整积分的不可消去证书]\label{cor:prediction}
若实际几何上界与标架已构造，某个割满足上述链条件，
割最小对象含简单因子$S$，且$S\notin\cS_{<G}$，
则物理输出必须含$S$，并在定理\ref{thm:ren}的全部锚定方案中保留。
若上界中$S$的重数等于割下界，则该重数已精确确定。
\end{corollary}
\begin{proof}
依次用命题\ref{prop:lower}及定理\ref{thm:ren}，再夹逼整数重数。
\end{proof}
这是可以贯通本项目的主判据。每个前提都要有证据；
未知的链生成性或regulator非零不能藏在“自动编译”中。
banana指定纤维的相应输入已有独立证明；
一般图可以停在上界或条件下界，但须标明输出强度。

\subsection{由物理变分验证链条件：族与固定纤维分开}
令$S$为连通光滑参数曲线，$\mathcal E$为已经构造的相对上同调
同调对偶变分，具有有理Betti局部系统和代数de Rham连接。
下面在一个有限长度的实际变分类别中取子对象，且其Betti实现为局部系统；
可使用几何产生的可容许混合Hodge变分及其子商。
不据实现结果宣称存在尚未构造的Nori族态射。
固定基点$s_*$、链$\sigma$及单值代数读出$\lambda$。
$R_\sigma^{\rm var}$表示包含$\sigma$的最小子变分，
$I_\lambda^{\rm var}$表示被$\lambda$恒等消去的最大子变分。

\begin{theorem}[物理变分给出的割保护]\label{thm:variation}
设实际管态射$f:\mathcal C\to\mathcal E$和路径$\ell$满足
\[
 (T_\ell-1)\sigma=c\,f_B\gamma,\qquad c\in\Q^\times.
\]
则在变分类别中
\[
 M_{\rm var}(\mathcal C;\gamma,f^*\lambda)
 \hookrightarrow M_{\rm var}(\mathcal E;\sigma,\lambda).
\]
若实际管像包含一个Betti局部系统不可约的子变分$\mathcal A$，
且$f_B\gamma\in\mathcal A_{s_*}$非零、$\lambda|_{\mathcal A}\not\equiv0$，
则$\mathcal A$单射物理最小变分。
\end{theorem}
\begin{proof}
$R_\sigma^{\rm var}$的Betti纤维被所有单值保持，所以同时含$\sigma$
和$T_\ell\sigma$，从而包含$f_B\gamma$。最小性给
$R_{f\gamma}^{\rm var}\subset R_\sigma^{\rm var}$，应用命题\ref{prop:lower}。
第二项中，交$\mathcal A\cap R_\sigma^{\rm var}$的局部系统非零，
不可约性使其等于$\mathcal A$；而
$\mathcal A\cap I_\lambda^{\rm var}$要么为零，要么全部，
后者被非零读出排除。
\end{proof}
这消除了族中链包含的一个实质待证条件，但没有消除实际管比较和交数的责任。
若变分不是单一管周期而是多个割之和，应先对总割求最小对象，
不能把各项分别当作已被保护。

\begin{proposition}[无需完整积分的微分阶数下界]\label{prop:order}
令$F$为物理周期函数，并设其某个变分等于非零割函数$c\,g$。
若$g$的全部解析延拓所张成的解空间$W$不可约且维数为$r$，
则任意以单值亚纯函数为系数、满足$LF=0$的标量齐次算子都有
$\operatorname{ord}L\ge r$。
若$W$由$K[\partial_s]$中的首一$r$阶算子$L_{\rm cut}$给出，
则每个$L\in K[\partial_s]$满足
\[
 L=A\,L_{\rm cut}.
\]
\end{proposition}
\begin{proof}
解析延拓与单值系数算子交换，故$L$杀死$F$的变分和它的所有延拓。
于是它杀死整个$r$维$W$；普通点处$m$阶方程的局部解空间至多$m$维。
对$L$用$L_{\rm cut}$作右除，余项阶数小于$r$且仍杀死$W$，只能为零。
\end{proof}
不可约割局部系统与非零代数读出通常提供这里的$r$：
其平坦向量到周期函数的核被单值保持，不可约性迫使核为零。
如扩大到代数系数，须在相应有限覆盖上仍有不可约性；
连通几何单值群已不可约时满足这一点。

\textbf{固定纤维边界。}
$T_\ell$一般不是$(\mathcal E_{s_*},F^\bullet,W_\bullet)$的Hodge自同态。
所以本定理的$M_{\rm var}$不能替换成任意指定点的$M$。
例如既有banana固定纤维$\operatorname{End}_{\rm Hdg}T=\Q$，
非标量节点反射就不是固定Hodge态射。
指定$s=-64$的秩八最小性仍用原regulator与群证明，
不从本条微分下界重证。
更一般地，某一纤维数值周期偶然为零，也不抹去非零族函数的解空间。

\subsection{一个任意圈数的实际物理入口：二维banana正常阈值}
设$h\ge2$、$n=h-1$、$m_i>0$，$\mu_i=m_i^2$。
采用明确的正参数归一化
\[
 I_h(s)=\int_{\alpha_i>0,\ \alpha_h=1}
 \frac{d\alpha_1\cdots d\alpha_n}
 {(\sum_i\mu_i\alpha_i)U_h-s\prod_i\alpha_i}.
\]
这里不含动量测度的外部$\pi,i$因子；
$h=4$时正取向与本项目欧氏参数周期相同。
在图表中置$y_i=\log(m_i\alpha_i/m_h)$、$y_h=0$，则
\[
 I_h(s)=\int_{\mathbb R^n}\frac{dy_1\cdots dy_n}{f(y)-s},
 \qquad
 f(y)=\Big(\sum_{i=1}^h m_i e^{y_i}\Big)
       \Big(\sum_{i=1}^h m_i e^{-y_i}\Big).
\]
令$M=\sum_i m_i$、$s_0=M^2$。

\begin{theorem}[正物理链的阈值非消失]\label{thm:threshold}
$I_h$在$\mathbb C\setminus[s_0,\infty)$上全纯。
$f$严格凸、恰当，唯一极小点为$y=0$，极小值为$s_0$。
对$t>s_0$，正向边界差为
\[
 I_h(t+i0)-I_h(t-i0)=2\pi i\,\rho_h(t),\qquad
 \rho_h(t)=\int_{\mathbb R^n}\delta(f(y)-t)\,dy>0.
\]
在$\varepsilon\downarrow0$时，
\[
 \rho_h(s_0+\varepsilon)=
 \frac{\pi^{n/2}}
 {\Gamma(n/2)\,M^{(h-2)/2}\sqrt{\prod_i m_i}}\,
 \varepsilon^{(h-3)/2}\bigl(1+O(\varepsilon)\bigr).
 \tag{T}
\]
局部链变分为正常Morse消失球的单次管，取向由上式固定，
交数绝对值为一；尤其物理分支确实激发该割。
\end{theorem}
\begin{proof}
展开为
\[
 f=\sum_i m_i^2+2\sum_{i<j}m_i m_j\cosh(y_i-y_j).
\]
沿任意非零方向$v$（$v_h=0$），Hessian二次型为
$2\sum_{i<j}m_i m_j\cosh(y_i-y_j)(v_i-v_j)^2>0$。
全体差同时为零只在$y=0$，由Cauchy--Schwarz得到极小值。
若$R=\max_i y_i-\min_i y_i$，则
$f\ge m_{\min}^2e^R$且$R\ge\|y\|_\infty$；
故$f$恰当，积分在无穷远有指数可积控制。
对避开割线的任意紧$s$集合，分母有统一相对下界；
可在积分下任意求$s$导数，得全纯性。

严格凸性使每个$t>s_0$的紧水平面均无临界点。
用余面积公式推前正测度，再对$1/(u-t\mp i0)$作一维Cauchy边界值，
得到跳跃式。水平面非空且其密度处处正，故$\rho_h(t)>0$。
这是本参数积分的直接证明，不调用任意维多圈Cutkosky作为前提。

在极小点，令$m=(m_1,\ldots,m_n)^t$，有
\[
 f-s_0=y^tAy+O(\|y\|^4),\quad
 A=M\,\operatorname{diag}(m_1,\ldots,m_n)-mm^t,
 \quad \det A=M^{h-2}\prod_{i=1}^h m_i.
\]
矩阵行列式引理给最后等式。
实解析Morse坐标把$f-s_0$化为$\|u\|^2$，
体积密度在零点为$(\det A)^{-1/2}$。
球面平均消去所有奇次项，故首项为
\[
 (\det A)^{-1/2}
 \int_{\mathbb R^n}\delta(\|u\|^2-\varepsilon)\,du
 =\frac{\pi^{n/2}\varepsilon^{n/2-1}}
 {\Gamma(n/2)\sqrt{\det A}},
\]
下一相对阶为$O(\varepsilon)$，包括$n=1$情形。

最后，取极小点的小球；球外$f-s_0$有正下界，
外部积分在$s_0$邻域全纯，不能贡献这次局部变分。
球内实链与Morse消失胞的对偶局部胞仅横截相交一次；
亦可在径向变量中绕过唯一简单极点，直接得到一次Cauchy管。
其留数正密度恰为上式，因此倍数不是未定整数。
经典局部管与Picard--Lefschetz解释可比较\cite[\S2.3]{muhl}。
\end{proof}
由正$y$图表对应未来向在壳支，以及既有光锥/参数割簇同构，
此处给出了正常阈值处的物理激发。
它没有计算全体异常阈值或所有全局轮廓，也未补完全部动量空间测度比较。

\begin{corollary}[从阈值到不可约几何核]\label{cor:banana-threshold}
若此正常消失球的紧割周期属于一个已证明不可约的$r$维纯局部系统，
则$I_h$的标量齐次微分方程阶数至少为$r$，并含对应割算子作为右因子。
已构造实际管变分时，同时得到定理\ref{thm:variation}的对象保护。
\end{corollary}
\begin{proof}
式(T)保证割函数非零，不可约性使其延拓生成全纯核，
再用命题\ref{prop:order}与定理\ref{thm:variation}。
\end{proof}
$h=3$在不可约椭圆变分上给$r=2$；
$h=4$、$m=(1,2,4,8)$射线利用已证纯单值群给$r=6$。
两者结论是物理函数的下界；完整混合对象和齐次阶数不因此分别等于$2,6$。
$h>4$的同一阈值公式成立，但不能凭阈值幂次就宣布某个CY核不可约。

\section{从边表进入几何：圈星类的结构规则}
\subsection{可终止的图级入口}
原\nolinkurl{conic_star_classification}定义任意圈数的可识别类：
在一个整数基本圈基中，有中心圈与$k$个叶圈，每圈恰有三条私有边，
另有中心分别与各叶共享的$k$条边；$L=k+1$、$E=4k+3$。
枚举生成树并检查基本圈支撑，是对该定义完备的有限识别。
至多检验$\binom{E}{V-1}$棵候选树，每项为有限图运算；
不声称任意图的多项式分类。

每圈三个私有割之差给两条线性方程。
秩二、剩余二次曲线光滑且有两个不同无穷远点时，
经明确有限代数扩域可参数化为$\ell_j=v_j+a_jt_j+b_j/t_j$。
置中心$t_0=z$，共享割给
\[
 F_i=A_i(z)t_i^2+B_i(z)t_i+C_i(z),\quad
 \deg A_i,\deg B_i,\deg C_i\leq2,
\]
\[
 y_i=2A_it_i+B_i,\quad y_i^2=\Delta_i(z)=B_i^2-4A_iC_i,\quad
 \deg\Delta_i\leq4.
\]
这由传播子支撑推导，未按图名预设椭圆答案。
秩、判别式及平方类条件都是明确的精确代数前提。

\subsection{紧割分解和分子选择}
先取平方类独立、四次平方自由且两两不相交的一般参数层。
紧化覆盖$C$的群为$H=(\mathbb Z/2)^k$，
$C_S:y_S^2=\prod_{i\in S}\Delta_i$。

\begin{proposition}[曲线核与标架选择]\label{prop:star}
代数对应给出
\[
 H^1(C)\simeq\bigoplus_{\varnothing\ne S\subset[k]}H^1(C_S),
 \quad g(C_S)=2|S|-1,\quad g(C)=1+2^k(k-1).
\]
单位标量最大割仅选全字符：
\[
 \omega_{\rm sc}=c\,z^{k-1}dz/(y_1\cdots y_k),\quad c\ne0.
\]
其紧宿主亏格为$2k-1$。若$N=\sum_U n_U(z)y_U$，
则$n_Uy_U\omega_{\rm sc}$进入补集字符$[k]\setminus U$。
\end{proposition}
\begin{proof}
群幂等元$e_S=2^{-k}\sum_h\chi_S(h)h$为代数自同构图的有理线性组合。
商映射到$C_S$的拉回与迹复合为覆盖次数；除以该次数，
像即字符$S$部分，平凡字符$H^1(\mathbb P^1)$为零。
Riemann--Hurwitz对$4|S|$个分歧点给亏格，再求和。
共享留数将$dt_i/t_i$除以$F_i/(zt_i)$，得到$\pm z/y_i$；
乘中心$dz/z$即得标量形式。
分子规则来自$y_U/(y_1\cdots y_k)=1/y_{[k]\setminus U}$。
在一般层该标量形式非零且全纯。
\end{proof}
曲线群幂等元分解为标准机制\cite{KR}，
最大割亏格的相关先行工作见\cite{HHMZ}。
本篇吸收原稿的传播子与留数证明，不主张经典公式优先权。

\subsection{宿主不自动等于最小对象}
若选中各$H^1(C_S)$在当前范畴中简单且两两不同构，
最小纯对象恰为链投影与读出同时非零的因子直和。
因为不等型简单直和的子对象按因子分解，链闭包先选非零链分量，
读出再删去零分量。这时才是精确分类，而不仅是亏格上界。

这些条件可在充分泛参数族中通过占优分支变化和单值群证明；
原稿论证使用完整独立离壳外源，不能套用每个指定物理切片。
本篇不凭亏格重证其泛单值定理。
在等源特殊点，不同字符可成为同构对象，发生标架相消：
同一简单对象$A$的两份中，$(v,v)$与$(\lambda,-\lambda)$给零周期。
两份各自非零不足以保证总和可见；同型重数须单独处理。

分支碰撞时，先在代数闭常数域删除平方类依赖，得秩$r$。
每个非零几何类$v$有
\[
 b(v)=\#\{p:\sum_i v_i\operatorname{ord}_p\Delta_i=1\bmod2\},
 \qquad g_v=b(v)/2-1 ,
\]
其中包括无穷远。非零几何类在$\mathbb P^1$至少有两个分支点；
常数非平方仅为算术扭曲，不可算出几何亏格$-1$。
原域可按不可约因子次数加权，同时保留常数域与连通分量。
约去平方因子仍在微分中留下极点，因此相对边界扩张另记。
实际系数为数域时，因子分解与二元线性代数部分有限终止；
这不是指定Jacobian全部等源或一般动机的判定算法。

\subsection{补齐纯核之外的相对边界}
下面给接受曲线核后的实际混合对象，避免把紧亏格直接当作完整周期。
设$C/k$为光滑几何连通射影曲线，$D,B$为互不相交的有限约化点集，
$U=C\setminus D$，$d=\#D(\bar k)$、$b=\#B(\bar k)$。
所有点均带实际代数定义；必要时只作有限扩域。

\begin{proposition}[曲线相对对象的三层结构]\label{prop:curve-boundary}
实际同调对象$E=H_1(U,B)$仅有如下权重分级：
\[
 \begin{split}
 \operatorname{Gr}_{-2}^W E
   &=\operatorname{coker}\{\Q(1)\to H_0(D)(1)\},\\
 \operatorname{Gr}_{-1}^W E&=H_1(C),\\
 \operatorname{Gr}_{0}^W E
   &=\ker\{H_0(B)\to\Q(0)\}.
 \end{split}
\]
空$D$或空$B$相应项记零，其维数为
\[
 \max(d-1,0),\quad 2g(C),\quad\max(b-1,0).
\]
原域上的端点项为Artin及其Tate扭转；点全部分裂后才写为Tate直和。
\end{proposition}
\begin{proof}
光滑曲线除点的局部化及其同调对偶给
\[
 H_2(C)\longrightarrow H_0(D)(1)\longrightarrow
 H_1(U)\longrightarrow H_1(C)\longrightarrow0,
\]
第一映射在几何点基中为对角映射。
相对同调正合段给
\[
 0\longrightarrow H_1(U)\longrightarrow H_1(U,B)
 \longrightarrow H_0(B)\longrightarrow H_0(U).
\]
$H_0(U)=\Q(0)$；各项纯权重及权重严格性给上述分级。
这些先是实际几何对的映射，Hodge实现只用于读取权重；
不从一个抽象分级矩阵推测扩张存在。
\end{proof}
因此在\emph{已构造的这一个曲线模型}内，
开放链和极点不会创造新的紧$H^1$因子，却可能增加
Abel--Jacobi、第三类微分及端点Kummer扩张。
其除子类、支点与路径决定粘合；仅列$b,d,g$不能决定这些扩张。
给定图若还需更高维积分、其他sector或重整化边界，
本命题不能将它们全部压缩成该曲线对象。

\subsection{非齐次源由端点给出：一个精确的微分分类规则}
在参数曲线$S$上考虑上述光滑射影曲线族与互不相交有限平坦$B,D$。
删去有限坏点，再经有限覆盖把点分裂为代数截面；
记函数域为$K$、导数为$\partial$。
令$\lambda\in H^1_{\dR}(U,B)$为实际相对形式，
其在$D$的全部留数为零。
相对类可以含必要的端点分量，移动端点不能在求导时冻结。

令$V\subset H^1_{\dR}(U,B)$为留数核。
由上两条正合列的对偶得到带连接的正合列
\[
 0\longrightarrow H^0_{\dR}(B)/K
 \xrightarrow{j}V\longrightarrow H^1_{\dR}(C)\longrightarrow0.
 \tag{C}
\]
空$B$时左项记零。
把$j(h)$的配对约定为$\sum_i n_i h(B_i)$，
其中$\partial\gamma=\sum_i n_i[B_i]$；在常用映射锥中$j(h)$由$(0,-h)$表示。
令$\bar\lambda$为$\lambda$在$H^1_{\dR}(C)$中的像，
$r$为$\bar\lambda$生成的循环连接模的维数，
$L_{\rm pure}\in K[\partial]$为对应首一$r$阶算子。

\begin{theorem}[端点源与一阶附加界]\label{thm:endpoint}
对任意平坦相对链$\gamma$，$F=\langle\gamma,\lambda\rangle$满足
\[
 L_{\rm pure}F=a(s),\qquad
 a(s)=\sum_i n_i h_i(s)\in K,\qquad
 L_{\rm pure}\lambda=j(h).
\]
若$a=0$，齐次阶数至多$r$；若$a\ne0$，则
\[
 \bigl(\partial-a'/a\bigr)L_{\rm pure}F=0,
 \qquad \operatorname{ord}\le r+1.
 \tag{E}
\]
这里$a$由相对de Rham约化和代数端点求值决定，
不用先求出$F$的完整积分。
\end{theorem}
\begin{proof}
$\bar\lambda$满足$L_{\rm pure}\bar\lambda=0$。
式(C)为带连接的正合列，故$L_{\rm pure}\lambda$落在$j$的像。
点分裂后$H^0(B)$为带平凡连接的$K^b$，
其对平坦零链的配对是$K$中的代数函数。
配对与Gauss--Manin连接相容，所以得到非齐次式。
再对$a$施以$\partial-a'/a$即得(E)。
常数同时加到全部$h_i$不影响结果，因为$\sum_i n_i=0$。
\end{proof}
若用纤维上的恰当形式$df$实现一次约化，其端点贡献为
$f(B_{\rm end})-f(B_{\rm start})$；
一般移动端点还包含相对Cartan修正，均由实际连接确定。
只看紧曲线上的恰当性、把这些端点项置零会错误地降低阶数。
有非零留数时还存在第三类微分的Tate商，
必须先处理其连接；本定理不自动覆盖该情形。

\begin{corollary}[由割和端点精确决定阶数]\label{cor:exact-order}
若$\bar\lambda$所在的$r$维纯连接不可约，
且物理变分已满足命题\ref{prop:order}的$r$维割保护，
则$F$的最小齐次阶数恰为
\[
 \begin{cases}
 r,&a=0,\\
 r+1,&a\ne0.
 \end{cases}
\]
\end{corollary}
\begin{proof}
割保护给$r$阶下界。若$a=0$，由上界得等号。
若$a\ne0$而有$r$阶湮灭算子，则命题\ref{prop:order}
的右整除使它为$L_{\rm pure}$的非零标量倍数，
与$L_{\rm pure}F=a\ne0$矛盾。式(E)给剩余上界。
\end{proof}
对于一般圈星标量割，形式为紧全字符曲线$C_{[k]}$上的非零全纯形式。
若其纯连接已在所用参数曲线上证明不可约，则$r=4k-2$。
实际闭链周期的源为零；实际开链表示且满足割保护时，
齐次阶数恰为$4k-2$或$4k-1$。
这是一条\emph{带实际曲线表示和物理链条件的精确规则}；
尚未构造完整费曼积分到该相对曲线的比较时，
只能将它用于割或已识别的曲线子sector。
同样的表述适用于椭圆$r=2$的闭/开Abel积分。
这些微分阶数不等于任意指定代数点的Nori最小秩。

\section{可核查的图族分类表：从MZV证据到完整物理对象}
\label{sec:familyclassification}

本节以四个物理问题组织主线：为什么某些费曼图的周期是MZV；
哪些实际振幅、局部系数或物理投影已有MZV肯定证明；
哪些需要椭圆、K3等其它周期结构；如何由图和物理输入在完整积分以前预测。
数论与几何为这些判断提供判据和证明，不另立脱离物理分类的主攻目标。
所谓“需要椭圆、K3等”，指由质量与运动学决定的几何模参数和周期结构，
并非新增自由物理参数。本表覆盖已有证据的图族和物理情形，不是任意图的完整分类。

读表时先选定UV留数、振幅、截面或RG系数，再核对维数、质量、分子和减除，
最后确认周期结构在物理求和与投影后是否保留。
数学工作的优先级取决于它能否改进这些判断或覆盖更多物理图族。
完整扩张一栏追踪纯核之外的边界、非齐次源和常数；
给出有用的物理预测，不必等待所有抽象结构都已分类。

\subsection{读表约定：六个不能合并的判断}
\label{subsec:classsemantics}
完整输入记为
\[
 {\cal I}=(G,D,\boldsymbol a,\boldsymbol m,\boldsymbol p,N,
             K,\sigma,\text{解析支},\text{减除方案},\text{物理读出}).
\]
一行中出现的秩可以是完整载体秩、纯核秩、函数最小阶数或固定动机秩；
下表逐一写明对象。除特别说明，质量为正代数数，
从收敛的欧氏正链或明确的正能量链开始延拓。
在函数行中删去几何判别式；不能把这个开集中的结论逐点照抄到边界。

\textbf{MZV肯定证据}指普通实数代数
\[
 {\cal Z}=\Q[\zeta(n_1,\ldots,n_d):n_1\ge2,\ n_j\ge1],\qquad
 \zeta(\boldsymbol n)=\sum_{k_1>\cdots>k_d\ge1}
                       k_1^{-n_1}\cdots k_d^{-n_d}
\]
中的表达证明，或明确标出的形式MZV证明。
测度中的$\pi$、$2\pi i$、代数系数和质量量纲须先声明；
不能把$\overline\Q[{\cal Z},(2\pi i)^{\pm1}]$改称$\cal Z$。
MPL是带变量和字母的函数类，混合Tate是范畴性质，MZV是特定端点的常数类。
本表以MT表示混合Tate，以AT表示混合Artin--Tate；
AT的有限角色经扩域可分裂，但质量对数和端点扩张仍可能存在。

\textbf{非Tate证据}分三层：环境上同调出现非Tate；
指定de Rham形式投到该组成非零；实际双标架的最小对象确实含该组成。
数点、曲面名称、最大割或非零形式各自只证明相应一层。
\textbf{物理可见性}必须另给链、管变分、正性或真实投影的证明。
\textbf{完整扩张}必须说明是在$\mathbb C(s)$的微分范畴，
还是指定数域、指定代数纤维的Nori范畴；两者不能互换。

表中的“数值未判”是一个精确状态：所引证据不判定指定数值是否属于$\cal Z$。
可见非Tate最小动机排除同一范畴内的混合Tate形式表示；
要将它变成具体复数的否定判定，还需要适当共同范畴上的周期映射单射性，
或一项独立的数论证明。本节没有假定该单射性。
相比之下，可见的非虚幺幂单值群已足以给函数层面的MPL障碍。

\subsection{图族总表}
\label{subsec:familytable}
每个编号在后文都有证据卡。M1--M2引自公开定理；N1单列公开非Tate基准。
其余各行也有直接先行工作；表后证据卡与第\ref{sec:literaturematrix}节逐项对照。
本项目的具体推导不自动构成新颖性，内部可审阅性也不等于已通过外部同行评议。
“未知”表示该行没有交付该结论，不表示不存在这种结构。

\begin{landscape}
\begingroup
\renewcommand{\small}{\fontsize{8.8}{11.2}\selectfont}\small
\setlength{\tabcolsep}{1.2mm}\renewcommand{\arraystretch}{1.08}
\begin{longtable}{@{}>{\raggedright\arraybackslash}p{19mm}>{\raggedright\arraybackslash}p{35mm}>{\raggedright\arraybackslash}p{27mm}>{\raggedright\arraybackslash}p{36mm}>{\raggedright\arraybackslash}p{35mm}>{\raggedright\arraybackslash}p{43mm}>{\raggedright\arraybackslash}p{34mm}@{}}
\toprule 图族/编号&输入条件&MZV肯定证据&非Tate结构证据&物理可见性&完整扩张状态&特殊点及例外\\\midrule\endfirsthead
\toprule 图族/编号&输入条件&MZV肯定证据&非Tate结构证据&物理可见性&完整扩张状态&特殊点及例外\\\midrule\endhead
\bottomrule\endfoot
\textbf{M1}\newline 小顶点宽度&无质量、$D=4$本原对数图周期；$N=2h\ge6$；存在边序宽度$\le3$&是；$P_G\in{\cal Z}$，权重至多$N-3$&该周期有MT控制；判据不给非Tate因子&正单纯形上的实际图周期；进入图的UV留数&给肯定归类；不逐图算最小秩及全部扩张参数&有质量、分子及运动学函数不在此结论内；宽度$>3$不是否定证据\\\midrule
\textbf{M2}\newline zig-zag $Z_n$&$n\ge3$；无质量本原$D=4$周期、标准$\Omega/\Psi^2$规范&是；显式有理系数乘$\zeta(2n-3)$&无需非Tate结构来表示该数值&实际图周期公式；可作为局部RG图贡献&数值公式已定；不由公式认定整个图动机或全部扩张&整体振幅的图权相消另查；增加质量后须另立一行\\\midrule
\textbf{M3}\newline 无质量串并联&双端分解证书；仅端点注入；正整数边指数；$D=4-2\varepsilon$，$q^2>0$&Euler规范后全阶系数在单$\zeta$代数；带尺度为$\mathcal Z_1[\log(q^2/\mu^2)]$&该读出有Gamma表示；无需椭圆/K3函数&实际卷积归纳；森林/RG须另核查子图余图及IR条件&所有展开系数已控制；不据此分类原生完整最小动机&减除点系数为MZV；一般尺度对数非自动MZV；质量、非整数指标、内部注入另判\\\midrule
\textbf{M4}\newline 全零质量banana&$D=4-2\varepsilon$，$n\ge2$单位边；非零欧氏两点；指定MS/MOM&全阶单$\zeta$；MS有限项仅有理数与尺度对数&四线谱曲面的正规化有理；完整读出已显式&实际反项、MOM及正谱均已定；不是散射截面&完整正规化函数和全部$\varepsilon$系数已定；原生动机最小化另留&所有质量为零是奇异层；不能套用正质量K3；$D=2$有IR问题\\\midrule
\textbf{A1}\newline 一圈标量族&$D=2d$，整数幂，$\sum a_i>d$；正质量欧氏；同极点集分子&一般无MZV肯定结论；MT不充分&完整相对载体AT；有限扩域后MT&实际参数标架给物理上界；排除椭圆/K3简单因子&三角/盒一般位置完整载体秩$7/8$，分级已定；标架最小化与扩张参数另判&判别式零、零质量IR、非整数维数均须重审\\\midrule
\textbf{E1}\newline 三边banana&$D=2$，正质量，单位标量；正常阈值及其光滑延拓域&无一般MZV肯定结论；特殊数值未判&光滑三次曲线$H^1(E)$；几何群$\mathrm{SL}_2$&正阈值割非零，核完整可见；函数有MPL障碍&在$\mathbb C(s)$上非分裂秩$3$共同单位扩张；固定算术扩张未统一分类&$s=0$及伪阈值/阈值纤维需重算；等质量的光滑族仍可椭圆\\\midrule
\textbf{E2}\newline theta细分族&三条串联路径、零内部注入；路径内异质量；私有阈值；限定分子次数&活动数$r>0$时无本表的MZV肯定证据；$r=0$为零&每个活动质量三元组给一份互异椭圆核&网格值非零当且仅当对应核可见&微分对象完整：秩$2r+1$，群$\mathbb G_a^{2r}\rtimes\mathrm{SL}_2^r$；不是固定绝对群&质量合流见E5上界；阈值重合及高次分子另判；固定圈数不界定$r$\\\midrule
\textbf{E3}\newline 内部注入四边图&$D=2$三点$(2,1,1)$图；正质量；两端面及夹挤满足指定一般位置&没有一般肯定结论；数值未判&绝对极点曲面有理，但相对边界含两个椭圆核&单个标量读出各看见一次；非AT半单秩$4$&纯组成精确；完整粘合和固定纤维最小秩未定&注入为零可降到串联规则；同步退化和端面碰撞另查\\\midrule
\textbf{E4}\newline 三体总率&$D=4$正质量，$m_1\ne m_2$；单位相空间或接触源读出&无一般肯定结论；固定能量未判&$y^2=\lambda(t,m_1^2,m_2^2)\lambda(s,t,m_3^2)$属一&正物理圈激发恰一份$H^1$&函数阶数$2$--$5$；混合扩张未完整确定&阈值、分支/极点碰撞另判；某个等质量极限不能直接套本行\\\midrule
\textbf{E5}\newline theta质量合流&$D=2$正质量；路径内可重复；零内部注入；路径平方分子逐次真分式&无一般MZV肯定；合流本身不保证Tate&至多$R=r_1r_2r_3$个光滑sunrise椭圆载体；只是表示上界&有限喷射给升幂读出；零阶网格值不足；完整求和可见性另判&原积分函数有有限导数表示；原生标架比较及最小扩张未定&$r_a$数不同质量而非边数；曲线退化另判；不套E2的精确秩与群\\\midrule
\textbf{P1}\newline 割点乘积族&接受theta块的实际割点拼接；独立块须跨块私有阈值&无普遍MZV肯定结论；继承各块证据&独立块给纯核外张量积；重复块给对称幂&实际乘法及最小性证明，非仅张量上界&独立时阶数$\prod_i(2r_i+1)$；重复$k$块为$\binom{2r+k}{k}$；微分对象已定&关联块不能当独立块；这些并非全属同一可重整化模型\\\midrule
\textbf{K1}\newline 固定四线banana&$D=2$，$m=(1,2,4,8)$，$s=-64$，$K=\Q$；原正标量标架&数值MZV性未判；形式MT表示被排除&实际秩$6$超越K3，格$U\oplus U(2)\oplus A_2(-2)$&指定正链的最小秩$8$；K3确实可见&原积分标架等式已给；秩$8$继承Picard与非分裂证书；完整群另依赖固定算术输入&不能外推所有代数点；全四参数秩$10$载体不等于此点最小秩\\\midrule
\textbf{K2}\newline 四粒子谱曲面&两对质量互异；$s\ne0$；四个$\lambda(s,a_\pm,b_\pm)\ne0$&无一般肯定结论；数值未判&$(4,4)$分支双覆盖消解为K3；$\rho\ge16$，$\mathrm{rk}\,T\le6$&$D=2$正谱是非零全纯二形式周期；精确简单性另需输入&一般质量只给载体/纯层上界，未给完整相对最小对象&质量相等或分支碰撞重算模型；Picard跳跃可降秩但不自动变Tate\\\midrule
\textbf{K3}\newline 四线能量MOM&$D=4$，$m=(1,2,4,8)$，$s_0<0$，三次减除；指定有限能量代表&无一般肯定结论；$s=s_0$的值严格为$0$&函数中恰一份秩$6$K3无限单值因子&实际球面/谱比较加阈值下界；完整物理复合源读出&完整固定纤维扩张未定；独立森林代表的形式同一另判&$s=s_0$是零值例外；不能称此四维对象秩$8$\\\midrule
\textbf{K4}\newline 多项式四体源&$D=4$，指定质量$(1,2,4,8)$；局部标量积多项式$N$；正平方权&无一般肯定结论；质量壳零类给零谱&无交换分母时非有限单值半单部分恰为一份K3&$\overline N\ne0$当且仅当该K3可见；角矩是有限投影&完整三圈函数继承该运动学核；混合粘合未定&先平均振幅再平方会丢信息；有符号权不适用正平方判据\\\midrule
\textbf{K5}\newline 带交换四体源&同上质量；$f(a)=p(a)+\sum h_i/(a-c_i)$；全源比较取$1<c_i<9$&无一般肯定结论；数值未判&一份K3加每个保留简单极点一份椭圆核&合并极部、实际夹挤及完整谱证明各重数一&完整函数的无限单值半单部分精确；扩张、常数周期未定&零总留数删除；极点碰撞/高阶极部须重审；总阶数只知下界$6+2r$\\\midrule
\textbf{O1}\newline $\phi^4$弹性总截面&同一极点质量、实标量、已声明两圈方案；$g^4$系数及完整角投影&无一般肯定结论；固定能量未判&双Kummer系数投影产生一份椭圆$H^1$&完整带干涉系数确有且仅一份；指定有限耦合变化保持&无限单值简单重数确定；完整扩张/最小总阶数未定&阈值及退化另判；超过四体阈值还须加产生通道\\\midrule
\textbf{O2}\newline $\phi^4$四体产生&等质量，十个真实树图合并平方；$s>16m^2$，阶$g^4$&没有MZV分类结论&整体旋转纤维是Kummer；剩余五变量几何未分类&实际旋转推前保持干涉；总率严格正&旋转纤维已定；全率扩张未知。20点仅角尾界$<6{\cdot}10^{-17}$，总率窄认证待完&阈值归零；不能把Kummer纤维推广为全率MT，也不能套用K1的质量\\\midrule
\textbf{N1}\newline 八圈$G_8$&无质量本原$D=4$、16条边；明确边表见证据卡&没有数值MZV肯定/否定判定&公开K3约化；费曼形式生成非对角$\mathrm{gr}^{13,11}$&de Rham标架非Tate已证；本表不额外声称双标架最小化&完整相对扩张未由所引定理确定&“非Tate形式”不是该数值非MZV的证明；不要与K1的K3混同\\\midrule
\textbf{H1}\newline 高边数banana&任意有限$n\ge5$；正质量；$D=2k$的指定能量MOM&无统一肯定结论&二叉树谱双覆盖是实际载体；高维纯核未全面分类&正支比较成立；正常阈值非零；因子简单性须另证&有实际形式周期比较，未有全部最小纯核及扩张表&双覆盖不等于椭圆；高维载体不自动等于不可约CY核\\\midrule
\textbf{H2}\newline 有质量BPHZ图&四维局部可重整化标量理论、任意有限圈；合并森林；欧氏正质量&无统一肯定结论&输出实际$N,U,F$和相对载体，不预判其Tate性&完整合并链的参数比较已给；不是逐项积分发散项&有限Nori入口及PF上界；最小组成/扩张未分类&零质量IR、一般Minkowski全局链、非整数维数另判\\\midrule
\textbf{H3}\newline 圈星覆盖入口&满足原圈星接受条件的多重二次覆盖；给定分子角色&无统一肯定结论；属零仍需字母端点&各角色曲线属可精确算；相对三层边界可列&先得角色宿主；原完整图与指定链比较必须提交&紧割及相对曲线结构有规则；原图完整扩张非自动&分支碰撞、等源、角色消零须重算；不以局部立方体代K3核\\\midrule
\textbf{H4}\newline 任意图能量支撑&单位标量、泛独立能量根、守恒外能量；指定字符&不是MZV判定器&dual-critical给精确非零字符支撑；未给动机类型&能量积分的代数支撑；空间推前仍有任务&尚不分类空间推前后的完整扩张&根间代数关系与物理特殊化可改变支撑；非零字符不等于非Tate\\
\end{longtable}
\endgroup
\end{landscape}

\subsection{肯定MZV的机制与否定结论的边界}
\label{subsec:classproofs}
\begin{proposition}[零、一字母证书的充分性]\label{prop:classmzv}
设规范化实际周期通过有限个已证明的代数换元、相对Stokes、合法纤维推前和线性关系，
化为有理系数的有限和与积；每项是$[0,1]$上仅含
$\omega_0=dt/t$、$\omega_1=dt/(1-t)$的收敛迭代积分，
或采用标准单位切向基点的shuffle正则化值。
若没有未处理的端点常数、额外代数系数或Tate逆因子，则该数值属于$\cal Z$。
\end{proposition}
\begin{proof}
固定迭代积分约定
\[
 I(a_1\cdots a_w)=\int_{1>t_1>\cdots>t_w>0}
                \omega_{a_1}(t_1)\cdots\omega_{a_w}(t_w).
\]
收敛非空字以$0$开始、以$1$结束，唯一写成
$0^{n_1-1}1\cdots0^{n_d-1}1$，其中$n_1\ge2$、$n_j\ge1$。
展开$1/(1-t)$的几何级数，非负性允许交换逐项积分与求和，
连续积分给严格嵌套求和，结果为$\zeta(n_1,\ldots,n_d)$。
空字给$1$。shuffle恒等式把乘积化为字的有限和。
标准正则化以两个端点对数的常数项定义；
通过shuffle递归消去起首的$1$和末尾的$0$，该常数项是收敛字的有理线性组合。
因此所有终点项均在$\cal Z$中，沿已证明的关系反向传回原周期即可。
这不是声称任意图都存在这样的化简证书。
更广的$M_{0,n}$周期解释见\cite{ClassMzero}；本命题只用上述显式字证书。
\end{proof}

MT对象不自动提供这个证书。一般数域、分歧素数、额外删点、
质量Kummer类和切向基点缩放都可能增加常数。
线性可约性给逐步超对数的计算途径，
还须检查实际端点与分歧；一次约化失败只表示该途径未通过。
它既不是非Tate的充分证据，也不是数值非MZV的证据。

\begin{proposition}[函数MPL障碍与固定数字的区别]\label{prop:classmplobstruction}
在代数参数曲线上，设实际标量函数的最小正则奇异系统含一个简单组成，
其单值表示在任何有限底覆盖后仍非平凡且不可约，维数大于一。
则该函数不能写为有限个对数微分迭代积分的和与积，
配以代数函数系数和任意固定复常数。
该结论包含本表的可见$\mathrm{SL}_2$椭圆标准核和$\mathrm{SO}_6$标准K3核，
但不判定任一固定参数值的MZV性。
\end{proposition}
\begin{proof}
假设有上述有限表达。取有限代数覆盖，使所有代数系数及字母单值，
再删去其极点。总表达只涉及有界长度的迭代积分。
按字长过滤时，续接每个迭代积分只增加更短的字及常数，
故在有限维封闭的字空间上单值可同时上三角化，所有对角元为$1$。
积由shuffle转为同类有限表达；固定复常数不贡献运动学单值。
所以表达所生成的系统以及其全部子商的简单组成都为一维平凡表示。
这与假定的简单组成矛盾。
椭圆及当前K3核在有限覆盖后仍分别具有稠密$\mathrm{SL}_2$和$\mathrm{SO}_6$作用，
满足假设。取一个点的值会丢掉这整个续接表示，故最后一句不需也不能由单值论加强。
\end{proof}

若已经确定\emph{固定}双标架最小Nori对象含非Tate简单组成，
则它不属于混合Tate子范畴，从而该形式周期没有同一范畴中的MT表示；
这使用最小周期对象的表述无关性\cite[\S2.4]{brown}。
本结论适用于K1，而E2的微分群本身不能替代这样的固定算术证明。
CM椭圆或高Picard数K3仍可有非对角Hodge型，
“多了对称性”不能自动改成“变成Tate”。

\subsection{证据卡一：无质量MZV族与非Tate基准}
\label{subsec:classliterature}
\paragraph{M1：纯拓扑的充分判据有严格的物理域。}
给边序$e_1,\ldots,e_N$，令$V_i$为前$i$条边与剩余边都接触的顶点集。
宽度是$\max_i|V_i|$；提交宽度至多三的一条边序就是有限组合证据。
采用
\[
 P_G=\int_{\sigma}\frac{\Omega_{N-1}}{\Psi_G^2},\qquad
 \Psi_G=\sum_{\text{生成树 }T}\prod_{e\notin T}\alpha_e,
 \qquad N=2h_G\ge6,
\]
并要求每个非空真边子图满足$|E_\gamma|>2h_\gamma$，保证本原收敛周期。
Brown的矩阵型与正性定理把该边序送到MZV肯定结论
\cite[Theorems 1--2,118]{ClassWidth}。
这是一条充分规则；它不把含第二Symanzik多项式的有质量族一并分类。
此处没有恢复每个图的全部扩张参数。

\paragraph{M2：显式值是最强的数值肯定证据，但不替代全对象。}
标准zig-zag族按文献图定义，$n\ge3$时
\[
 P_{Z_n}=4\frac{(2n-2)!}{n!(n-1)!}
 \left(1-\frac{1-(-1)^n}{2^{2n-3}}\right)\zeta(2n-3).
\]
这是Brown--Schnetz的Theorem 1.1\cite{ClassZigzag}，
输入是本原无质量图周期及其标准正链。
数值公式与正系数给实际非零值；任何进一步的全图动机最小化须有自己的证明。

\paragraph{N1：作为非Tate证据分级的公开校准。}
$G_8$有顶点$1,\ldots,9$，有序边表为
\[
 (34,14,13,12,27,25,58,78,89,59,49,47,35,36,67,69).
\]
这里$ij$表示无向边$\{i,j\}$；采用\cite[\S6.2,式(30)]{ClassKthree}的版本和边序。
该文的约化得到K3。
Brown--Doryn的Theorem 5进一步证明
\[
 \operatorname{gr}^{13,11}H^{15}_{\dR}(\mathbb P^{15}\setminus X_{G_8})
\]
一维且由费曼形式类生成\cite{ClassFraming}。
这一公开结论强于只对环境数点；所引定理在绝对上同调层。
本表保留完整相对对象、物理双标架最小化与数字非MZV性各自的责任，
不把它们填写为已证。

\subsection{证据卡二：一圈、椭圆族和拼接}
\paragraph{M3--M4：可直接检查的无质量物理肯定族。}
定理\ref{thm:masslesssp}给串并联输入的全阶单$\zeta$证据，
命题\ref{prop:masslesscertificate}给有限边表核查界；
每条物理森林需另满足命题\ref{prop:masslessforestclosure}的闭合与IR条件。
无质量banana则由定理\ref{thm:masslessbanana}实际完成真森林判断，
推论\ref{cor:masslessphysical}给MS、MOM和正谱，
命题\ref{prop:masslessdegeneration}解释与正质量K3同拓扑却异类型的原因。
这里只固定标准单尺度和Euler规范，不把一般尺度上的对数值无条件认定为MZV。

\label{subsec:classelliptic}
\paragraph{A1：二次型及其边界足以给出上界。}
令$U=\sum\alpha_i$，
$F=U\sum\mu_i\alpha_i+\sum_{i<j}\kappa_{ij}\alpha_i\alpha_j$。
实际载体是射影空间去掉$F=0$以及形式必要时的$U=0$，
相对于坐标边界的上同调。
命题\ref{prop:onecarrier}证明完整载体AT；
定理\ref{thm:onegraded}给三角图和盒图的完整分级。
检查有限个受限二次型的判别式就能输入该规则。
7与8为一般位置载体秩，不能填入“最小物理秩”格。
质量比与平方类没有经过MZV端点检验，因此MZV栏保留未知。

\paragraph{E1--E2：给定图与质量，先决定活动椭圆清单。}
theta骨架是两端点之间三条内部不交的路径，
第$a$条路径有$r_a$条边、质量$m_{ai}$，内部二价顶点不注入外动量。
令$\mu_{ai}=m_{ai}^2$，路径内要求两两互异。
每个三元组$t=(i_1,i_2,i_3)$给
\begin{align*}
 E_t(s):\quad &(\mu_{1i_1}X+\mu_{2i_2}Y+\mu_{3i_3}Z)
                         (XY+XZ+YZ)-sXYZ=0,\\
 \mathcal D_t&=\{0,\infty\}\cup
   \{(\pm m_{1i_1}\pm m_{2i_2}\pm m_{3i_3})^2\},\qquad
 \tau_t=(m_{1i_1}+m_{2i_2}+m_{3i_3})^2.
\end{align*}
接受条件为$\tau_t\notin\mathcal D_u$对$t\ne u$。
限定$\deg_{Z_a}P\le r_a-1$，活动集为
\[
 J_P=\{t:P(-\mu_{1i_1},-\mu_{2i_2},-\mu_{3i_3})\ne0\}.
\]
完整相对Stokes下降、私有阈值和非分裂单位扩张分别见
命题\ref{prop:thetastokes}，定理\ref{thm:thetavisible}、
\ref{thm:thetanumerator}、\ref{thm:thetaexactorder}、\ref{thm:thetadgal}。
因此这行是\emph{无限图族的有限输入规则}，不是若干积分值的拟合。
$r_1=r_2=r_3=1$给E1；私有阈值条件此时为空。
固定代数$s$时，微分常数中的质量对数、regulator及绝对群仍需另判。


\paragraph{A1与E1--E2：明确先行对象与本表输出。}
一圈动量空间Nori对象、切割/收缩及图式余作用已有
Mounoud的直接结果\cite{LitMou}；本表的参数域、边指数、分子和特殊化须逐条比较，
不能把AT上界解释成普通MZV肯定结论。
任意质量sunrise的椭圆二重对数\cite{LitABW}、
模空间$\varepsilon$形式\cite{LitBMW}、一般质量余作用与边细分\cite{LitTap}
均为直接近邻。Tapušković的命题3.3/推论3.4保留原边质量；
本表E2的路径内异质量不在该同质量细分操作内，
具体图多项式分界见命题\ref{prop:splitmassidentity}。
其命题4.1/定理4.4是在实现范畴中的秩六载体及余作用；
E1固定质量后的$\mathbb C(s)$秩三不能替代它的全参数算术比较，
常数扩张反例见引理\ref{lem:constantkummerloss}。
本表仍不据这些范围差异宣布新颖性。

\paragraph{E5：质量合流必须保留喷射。}
路径$a$的不同平方质量为$\mu_{aj}$、重数$n_{aj}$，$N_a=\sum_j n_{aj}$。
输入要求$\deg_{Z_a}P<N_a$，所有内部顶点零注入；不涵盖一般张量分子。
定理\ref{thm:confluentcertificate}给有限混合导数系数，
命题\ref{prop:confluentelliptic}把完整欧氏函数写成至多$r_1r_2r_3$种sunrise的质量导数。
每种仍用同一相对对，提高极点阶数不新添极点支撑。
这是椭圆表示上界，不是全部非零系数都幸存的下界。
例如两条同质量边配分子$P=Z_1+\mu$时，零阶值全零，
但约去一个传播子后恰为非零sunrise；一阶喷射检测到它。
原合流图的原生Nori标架比较、全和最小性和退化点仍未由本规则完成。

\paragraph{E3：相对边界不能从表中删除。}
取顶点$A,B,C$，两条$AB$边质量$m,m$，$AC,CB$边质量$a,b$，
各顶点注入守恒外动量；以参数$x,y,u,v$标记这些边。
该行可由图多项式直接识别：
\[
 U=xy+(x+y)(u+v),\quad
 F=[m^2(x+y)+a^2u+b^2v]U-s_Axyu-s_Bxyv-s_Cuv(x+y).
\]
检查$(a^2+b^2-s_C)^2-4a^2b^2\ne0$、两个椭圆端面光滑、
所列普通节点及正常交叉模型，并排除阈值处其它同步退化。
命题\ref{prop:injectedgeometry}给几何，
定理\ref{thm:scalaricevisibility}给单标量物理下界。
此例说明绝对几何不能代替完整相对对象。
DHV定理5.1、7.5及命题7.6已经覆盖超椭圆型Hodge上界及适用域内的相对上界
\cite{DHV}；$(2,1,1),D=2$满足$3D/2=a+c$。
本表的精确物理重数须在这一先行结果上另作比较，不能宣称DHV只研究绝对几何。

\paragraph{E4与P1：改变物理读出或拼接方式必须更换规则。}
E4由实际两体相空间的卷积得到，
定理\ref{thm:threebodyvisible}给精确椭圆简单重数和阶数界。
P1须先满足实际割点动量分配，再检查不同块的私有阈值是否跨块分离；
命题\ref{prop:vertexproduct}、\ref{prop:blockjointgroup}和
定理\ref{thm:blockproductminimal}、\ref{thm:repeatedblock}分别给乘法、独立性和最小性。
独立两个sunrise块阶数为9，重复同一个块为6；
只看两个块的各自标签，无法区分两者。
两顶点平行拼接则是卷积推前，可能从代数根产生椭圆或K3，
必须使用定理\ref{thm:spectralglue}，不能套割点张量公式。

\subsection{证据卡三：K3、完整扩张与实际可观测量}
\label{subsec:classkthree}
\paragraph{K1：本项目目前最完整的固定算术行。}
固定的是质量$m=(1,2,4,8)$，不是质量平方列表；$s=-64$、基域$\Q$。
实际箭头链见第\ref{sec:actualframebridge}节：
引理\ref{lem:bridgefiber}比较单位符号与相对扩张，
定理\ref{thm:bridgephysical}证明链沿$i$提升、形式沿$i^\vee$下降的求值前等式；
推论\ref{cor:bridgeminimal}再用固定Picard与下层非分裂证书给最小秩八。
规范coaction存在、实际矩阵读出及完整群识别不是同一验收。
文献的$15$个维数正规化主积分与这里的$6/10/8$维对象也不能按数字直接等同；
奇异$\varepsilon$基变换的特殊化反例见命题\ref{prop:epsilonlattice}。
六维超越核、质量Kummer一维与顶层单位一维构成活动秩八对象；
分级不是把两个非零扩张拆开的理由。
实际整格、绝对比较、固定纯群和幺幂根基有单独证据，
不由参数单值群替代固定群。
定理\ref{thm:actualdispersionformal}及推论\ref{cor:actualdispersionminimal}
给原周期与谱读出的形式等式$I^{\mathfrak m}=8J^{\mathfrak m}$，
因此完整最小对象及coaction一并转移。
具体核查入口是\texttt{FIXED\_PURE\_GROUP.md}、
\texttt{RELATIVE\_TOP\_AUDIT.md}及主文核心实例的P6/P7/P8证据索引。
它们记录的有限精度和几何输入范围仍保留，表本身不代替证书。
Charles与Elsenhans--Jahnel提供的工具及其条件见\cite{LitCharles,LitEJ}；
不将具体纤维证书称为新的通用Picard或CM排除算法。

\paragraph{K2：K3的几何识别先于周期求值。}
令$a_\pm=(m_1\pm m_2)^2$，$b_\pm=(m_3\pm m_4)^2$，
$A(a)=(a-a_-)(a-a_+)$，$B(b)=(b-b_-)(b-b_+)$，
$\lambda(x,y,z)=x^2+y^2+z^2-2xy-2xz-2yz$。
对表中的接受输入，
\[
 X_s:\quad y^2=A(a)B(b)\lambda(s,a,b)
\]
是$\mathbb P^1\times\mathbb P^1$上的$(4,4)$分支双覆盖。
奇点消解、典范类和上同调给K3而非通过数值周期猜测，
见定理\ref{thm:spectralk3}及命题\ref{prop:k3spectralchain}。
一般质量的$\rho\ge16$不能改成$\rho=16$；
指定质量射线的秩六简单性需要既有单值/格输入。
任意四质量banana系统已有两组直接先行工作\cite{LitDMPS,LitPTWW}；
项目采用的公开纯矩阵归于前者。
特殊质量子族的K3周期可出现椭圆乘积分解\cite{LitDM,LitDMfull}，
故特殊点栏还须区分几何K3、纯核不可约性和整个函数的表示。

\paragraph{K3--K4：保持纯核与完成扩张是两件事。}
偶数维单位谱的纯组成上界、正常阈值下界见
定理\ref{thm:evenk3visible}。
能量MOM的实际比较见定理\ref{thm:actualevenenergy}及
推论\ref{cor:energyfourline}；$D=4,n=4$最少需要三次减除。
多项式分子先在正质量壳约化，再对平方作实际角投影；
定理\ref{thm:angularzero}、\ref{thm:angulark3}、\ref{thm:angularsource}
给正谱及完整局部源函数的充要规则。
未完的部分是相对边界粘合、固定常数周期及对应的最小算术对象，
不能把六维纯下界填成完整阶数六。

\paragraph{K5：交换极点给出可核查的新椭圆行。}
先合并同一$a$极点的所有图，保留非零总留数$h_i$。
在质量$(1,2,4,8)$下，候选曲线明确为
\[
 E_{c,s}:\quad v^2=(c-1)(c-9)(b-16)(b-144)\lambda(s,c,b).
\]
相空间规则允许$0<c<9,c\ne1$；接到表中稳定完整源的范围收紧为$1<c<9$。
定理\ref{thm:exchangeclassification}证明每个保留极点的精确重数；
定理\ref{thm:sourcefullcorrelator}、\ref{thm:dispersionexact}
将它们与K3一并传到完整三圈系数。
“精确”指无限单值半单组成，固定环境可以另含非Tate常数。
极点数不是原始图数，也不是完整扩张秩。

\paragraph{O1--O2：算到物理出口后重新分类。}
O1中的参数$a=x-4$、$x=s/m^2$，椭圆曲线为
$Y^2=v(v+4)(a-v)(a-v+4)$；
推论\ref{cor:elasticone}处理完整双Kummer系数和干涉，
定理\ref{thm:elasticscheme}限制允许的方案变化。
O2的十图完整旋转下降见定理\ref{thm:rotationfullrate}和
命题\ref{prop:rotationkummer}。
一个纤维的Kummer性不能决定它随五个剩余变量推前后的全对象。
四体阈值以上包容截面须合并弹性与产生通道，
两行各自的证据不能自动证明这个物理和的最小对象。
这不影响已有指定能量的包容数值，但也不将它升级为全函数算术分类。


\paragraph{O1--O2的精度与独立验证。}
20点的包容系数已经合并四体贡献，但仍是普通数值。
推论\ref{audit:cor:eta}只给角截断误差；命题\ref{audit:prop:boxes}给剩余积分的有限包围路线，
没有交付窄总率区间。
初次目标是完整正区间$[L,U]$满足$(U-L)/L\le10^{-3}$；
弹性系数、包容传播与独立三圈前向全切割分别验收。
2026年Higgs四体切割积分族已有MPL/椭圆/K3混合先行结果\cite{LitWWW}，
其模型与本$\phi^4$不同，积分族求解也不能代替总截面的完整投影。

\subsection{证据卡四：高维与一般图的有限入口}
\label{subsec:classopen}
\paragraph{H1：载体可构造不等于类型已分类。}
任意$n$的正质量banana可以选满二叉树；每个内部顶点给
$\lambda(h_v,h_L,h_R)$，根变量为$s$的谱变量，叶变量为质量平方。
实际双覆盖$y^2=\prod_v\lambda(h_v,h_L,h_R)$及正链由
定理\ref{thm:actualbinarytree}、\ref{thm:actualevenenergy}给出。
在$D=2k$中能量代表的最少减除次数为$(k-1)(n-1)$；二维不需减除。
变化的是微分、边界和Tate归一化；二次覆盖的次数不限制高维中间上同调的复杂性。
高维CY候选必须另查消解、典范类、实际纯子商与物理链。

\paragraph{H2：一般有质量图已经有实际入口，分类仍在其后。}
将完整BPHZ森林合并为一个收敛有理密度，
保留局部插入、图权与路由，再作有限扇区参数化。
定理\ref{thm:massive-period-entry}、\ref{thm:sectorformalcomparison}
及推论\ref{cor:sectorforestentry}给真实形式周期比较和有限相对模型。
其输出的$U,F$、删点及边界是待分类的对象；
动量纤维的Tate性不迫使参数底空间Tate。
本行不再把该接受域的全圈动量/参数比较误写为完全缺失，
也不声称最小纯核与全部算术扩张已可自动决定。

\paragraph{H3--H4：组合入口必须继续走到双标架。}
圈星的覆盖角色、属、边界三层及端点源见
命题\ref{prop:star}、\ref{prop:curve-boundary}和定理\ref{thm:endpoint}。
任意图独立能量根的dual-critical判据见\texttt{BOUNDED\_GRAPH\_SUPPORT.md}
及原\texttt{energy\_character\_classification}定义。
它决定一个代数能量核的字符是否出现；
后续空间积分仍可能形成新的纯核和扩张。
表中为这些输入保留独立行，正是为了避免把已有精确组合定理夸大成完整周期分类。

\subsection{两个无需完整积分的判别实例}
\label{subsec:classworked}
\begin{proposition}[同一四边拓扑的精确结构选择]\label{prop:classthetaselector}
取E2中路径长度$(2,1,1)$、路径质量$(1,2)$、$(5)$、$(20)$，
零内部注入，$D=2$。设$Z$为第一路径的欧氏动量平方。
分子$P=1$时完整函数的最小齐次阶数为5，含两份互不同构椭圆核；
$P=Z+4$时为3，仅保留质量$(1,5,20)$的sunrise；
$P=Z+1$时为3，仅保留质量$(2,5,20)$的sunrise。
\end{proposition}
\begin{proof}
两个正常阈值为$26^2=676$和$27^2=729$。
对应有限退化集合除$0$外分别为
\[
 \{196,256,576,676\},\qquad\{169,289,529,729\}.
\]
故每个正常阈值只属于自己的集合，满足私有阈值条件。
路径内质量平方为$1,4$，部分分式系数为$1/3,-1/3$。
当$P=1$时两个网格读出皆非零，$r=2$。
当$P=Z+4$时，$P(-1)=3,P(-4)=0$，故完整Stokes关系给
$I_P=S_{(1,5,20)}$；另一分子相同计算给$S_{(2,5,20)}$。
所有输入符合次数界；定理\ref{thm:thetaexactorder}给$2r+1$，得结论。
本证明只用质量代数和已有实际比较，不求这三个完整积分。
\end{proof}

\begin{proposition}[非Tate函数可以有严格MZV零值]\label{prop:classmomzero}
取K3行中的同一四维能量MOM函数$E(s;s_0)$，$s_0<0$为光滑点。
则$E(s_0;s_0)=0\in\cal Z$，且它作为$s$的函数仍含可见秩六K3核。
\end{proposition}
\begin{proof}
实际有限核为
\[
 K_3(x;s,s_0)=\frac{(s-s_0)^3}{(x-s)(x-s_0)^3}.
\]
谱支上$x\ge225>0$；在$s=s_0<0$，分母无零点而分子恒为零，
故实际收敛积分等于零。
邻域内这是三次MOM减除的同一函数，
推论\ref{cor:energyfourline}已给其延拓系统的一份可见K3。
固定一个零值不会抹去这个函数系统。
\end{proof}

这两个实例给表的两条必要设计原则：分子必须写进输入；
特殊点必须写进例外。仅从裸图给一个“MZV/椭圆/K3”标签会遗漏这两类可证明变化。
它们不表示所有非零非Tate周期都能在某点化成MZV。

\subsection{从边表到分类行的执行规则及下一阶段验收}
\label{subsec:classprogram}
先检查图、维数、质量、分子、链与方案，再匹配表中接受条件。
M1允许直接提交边序作为MZV充分证据；
E2允许直接由质量网格和私有阈值确定全部活动核及微分扩张；
K2允许由分支除子和消解确定几何；
其它行则按其现有规则输出上界、下界或明确未知。
若只得到一个$U,F$载体，流程不得自动填入椭圆、K3或CY。

每次扩展分类表，应交付四个具体东西：一项可判定的接受条件；
一个实际对象或受控子商；一个与物理链相容的可见性证明；
以及一份保留扩张和特殊化责任的结论。
允许输出“该行接受失败/需要补充引理”，不把失败解释为相反分类。
已有证明和有限证书可作为输入复用，不通过重复求积增加表的可信度。

接下来的主线依照前四个物理问题连接；第五项数学工作按这些问题的需要展开：
\begin{enumerate}
 \item 扩展\emph{MZV肯定列}：在已有实际参数入口中找出能够证明端点无分歧、
 字母为$0,1$或可化到该情形的图操作族；提交充分证书，随后研究判据的必要边界。
 \item 扩展\emph{非Tate且物理可见列}：从纯核候选走到实际双标架；
 优先补真正增加图族覆盖的推前、收缩和森林规则。
 \item 补\emph{完整扩张列}：尤其是从函数半单分类到固定算术对象，
 需要质量单位、边界一动机、regulator和常数环境，不能仅加微分方程阶数。
 \item 补\emph{特殊点列}：区分分支碰撞、Picard跳跃、CM、质量合流、
 分子零核和MOM零值；每个例外明确哪些结论保留、哪些必须重证。
 \item \emph{必要的数学支撑}：当某个物理分类判断确实需要时，研究分歧、
 Hodge型、代数循环、扩张或周期映射。
 不以不断深化抽象群或几何本身作为主要进度；
 具体数字的MZV否定结论仍须单独承担证明责任。
\end{enumerate}
切割、重整化、RG和临界响应仍属完整总目标。
它们现在通过“同一行在实际操作后如何变到另一行”的方式组织：
例如从图周期进入局部RG前必须先完成森林与局部投影，
从振幅进入总率前必须完成相空间推前及物理求和。
局部性只限制运动学依赖，不自动使剩余常数成为MZV；
有限阶的形式MZV临界响应也不自动认定每个原始图的规范动机提升。
这一分类表是当前理论的可审阅总入口，完整总目标仍未完成。


\section{主定理四：从完整参数模型传递到留数与非齐次源}
\label{sec:pencil}
\subsection{固定环境是必须提交的几何输入}
前述曲线定理要求完整周期已被识别为相对曲线周期。
一般图尚未完成这一步时，可以直接在其高维参数模型上工作。
设$k\subset\mathbb C$为特征零域，$X/k$光滑，
$B\subset X$为固定的简单正常交叉边界。
固定极点集合已从$X$中删去。
令$Y_s\subset X$为代数参数曲线$S$上的光滑除子族，
与$B$每个交叉层横截。
必要时缩小$S$，使所用相对上同调成为带Gauss--Manin连接的局部自由对象。
相对形式和物理链都必须实际属于下列模型：
\[
 J=H^n(X,B),\qquad
 M_s=H^n(X\setminus Y_s,B\setminus Y_s),\qquad
 C_s=H^{n-1}(Y_s,B\cap Y_s)(-1).
\]
这里$M_s$用上同调方向，前文最小物理对象承载于$E_s=M_s^\vee$。
$J$是固定几何的对象；这不表示$J$自动为Tate。
$\omega_s\in(M_s)_{\dR}$是代数相对类，
$\sigma_s\in(M_s)_B^\vee$在一个单连通参数开集上平坦，
$I(s)=\langle\sigma_s,\omega_s\rangle$。
物理边界的移动项已包含于相对Gauss--Manin连接。

\begin{theorem}[留数像与环境背景的精确分离]\label{thm:pencil-separation}
在上述实际模型中有带连接的正合段
\begin{equation}\label{eq:pencil-gysin}
 J\longrightarrow M_s\xrightarrow{\mathrm{res}}C_s
       \longrightarrow H^{n+1}(X,B).
\end{equation}
令$K_s=\ker(\mathrm{res})$、$A_s=\im(\mathrm{res})$。
则$K_s$是$J$的商，$A_s$是$C_s$的子对象，且
\[
 0\longrightarrow K_s\longrightarrow M_s\longrightarrow A_s
 \longrightarrow0.
\]
在每个指定代数纤维，取含$J$全部简单组成的Serre背景$\cS$，有
\[
 \pi M_s\simeq\pi A_s.
\]
因此完整模型在该背景以外的简单组成及其重数，恰等于实际留数像的组成及重数。
物理最小对象在对偶Serre背景$\cS^\vee$以外的重数至多为$A_s^\vee$中的重数；
若实际割和物理标架另给达到此上界的下界，则这些重数被精确确定。
\end{theorem}
\begin{proof}
对$Y_s\subset X$写局部化三角，余维一纯性给
$R\Gamma(Y_s)(-1)[-2]\to R\Gamma(X)\to
R\Gamma(X\setminus Y_s)\to$。
在$B$的每个交叉层上写相同三角。
横截基变换使它们与限制相容，取有限交叉复形及相对映射纤维得到
\[
 R\Gamma(Y_s,B\cap Y_s)(-1)[-2]\longrightarrow
 R\Gamma(X,B)\longrightarrow
 R\Gamma(X\setminus Y_s,B\setminus Y_s)\longrightarrow.
\]
取第$n$上同调即得\eqref{eq:pencil-gysin}；
实际六运算及其Nori实现给逐纤维态射，Gauss--Manin自然性给族连接相容。
这里不额外假设已构造一般“Nori变分类别”。
正合性给$K_s$为$J$的商；Serre商消去$K_s$，得到同构。
对偶后$E_s$是$A_s^\vee$与$K_s^\vee$的扩张，
而最小物理对象是$E_s$的子商，故在$\cS^\vee$外逐简单因子重数不超过相应上界。
达到上界时再用有限长度性。
\end{proof}

这将两个容易混淆的判断分开：$C_s$可能比留数像$A_s$大；
即使$A_s$完全已知，物理标架也可能只看到其一部分。
若$J$是mixed Tate，所有非Tate简单组成都来自\emph{相对}留数对象$C_s$。
但$C_s$自身仍含边界，它不自动等于紧致纯割核。
族中的物理变分可给连接/变分层面的下界，
不能不经比较改称指定纤维的Nori下界。

\subsection{非齐次源不需通过完整积分猜出}
取$K=k(S)$，配一个非常数局部参数的导子$\partial$，
假设常数域为$k$。
在$J_{\dR}$中固定$k$基$e_1,\ldots,e_d$，$d=\dim J_{\dR}$。

\begin{theorem}[实际留数湮灭后的固定环境源]\label{thm:ambient-source}
设$L\in K[\partial]$满足
$L\,\mathrm{res}(\omega)=0$，这里是实际de Rham类的零。
则存在$a_j\in K$，使
\begin{equation}\label{eq:ambient-source}
 L I(s)=\sum_{j=1}^d a_j(s)P_j,\qquad
 P_j=\langle\tau,e_j\rangle,
\end{equation}
其中$\tau$为$\sigma_s$在固定对$(X,B)$中的像，与$s$无关。
每个$P_j$是固定环境$J$的周期。
等式在相应实际形式周期中同样成立，
其中求导由Gauss--Manin连接定义，纤维在$k$或明确的扩域上取值。
特别地，这个结论不使用数值周期映射的单射性。
\end{theorem}
\begin{proof}
留数与连接交换，故$L\omega\in(K_s)_{\dR}$。
由\eqref{eq:pencil-gysin}，泛点处的线性满射
$J_{\dR}\otimes_k K\to(K_s)_{\dR}$可为此向量选一个$K$提升
$\theta=\sum a_je_j$。
$\sigma_s$映入固定对的相对同调，映射水平，故其像$\tau$平坦且为常数。
因此
\[
 LI=\langle\sigma_s,L\omega_s\rangle
   =\langle\tau,\theta(s)\rangle=\sum_j a_j(s)P_j.
\]
提升不唯一，但两个提升之差在映射核中，与$\tau$配对为零。
以上每次等号分别来自连接配对、实际态射函子性及线性关系，
所以也给形式周期等式，而非仅给解析数值相等。
\end{proof}

源的求取现在是有限de Rham约化与线性提升问题，
不要求先积分$I$。
不过，本定理本身没有提供适用于任意图的终止约化算法，
也不把$P_j$自动认作代数数、MZV、低圈周期或局部反项。
哪种固定结构出现，仍由$(X,B)$决定。

\begin{corollary}[不求固定周期值的齐次阶数上界]\label{cor:ambient-order}
设$r=\operatorname{ord}L$，
$q=\dim_k\operatorname{span}_k\{a_1,\ldots,a_d\}$。
存在$L_0\in K[\partial]$、$\operatorname{ord}L_0=q$，
使$L_0LI=0$。故最小齐次阶数满足
\[
 \operatorname{ord}_{\min}(I)\le r+q\le r+d.
\]
若源为零，可直接使用上界$r$。
\end{corollary}
\begin{proof}
$q>0$时选$a_j$张成空间的一组$k$基$b_1,\ldots,b_q$。
微分域的常数恰为$k$，故其Wronskian非零。
用Wronskian的最后一列放入待作用函数，除以最高项系数，
得系数在$K$中的首一$q$阶算子$L_0$，杀死所有$b_i$及$a_j$。
$P_j$对$\partial$为常数，因而$L_0$杀死\eqref{eq:ambient-source}的右端。
$q=0$时取$L_0=1$。
\end{proof}

这是上界，不是最小秩公式，也没有假设各$P_j$独立。
不同提升可改善$q$。
若源是非零$a(s)\in K$，则一阶算子$\partial-a'/a$足够；
若源为$P_1+sP_2$，却未证明$P_1/P_2\in k$，
就不能声称其对数导数在$k(s)$中。
$\partial^2$仍给无需这些常数关系的合法湮灭算子。
扩大常数域到$\mathbb C$与在原算术域$k$上求方程，是不同问题。

\begin{corollary}[有物理保护时的精确两支]\label{cor:ambient-exact}
若$L$同时是被非零物理变分激发的$r$维不可约割的最小算子，
且\eqref{eq:ambient-source}的源属于$K$，
则完整函数的最小齐次阶数为$r$（源为零）
或$r+1$（源非零）。
\end{corollary}
\begin{proof}
前文物理变分定理使任何湮灭$I$的算子都以$L$为右因子。
零源给$r$阶上界。
非零源时，$r$阶算子只能为非零函数乘$L$，不可能杀死$I$；
另一方面$(\partial-a'/a)L$给$r+1$阶上界。
\end{proof}

\subsection{完整相对留数与纯核方程之间的差别}
一个纯核算子$L_{\rm pure}$可能只在$C_s$的某个纯商中杀死留数像。
此时$L_{\rm pure}\mathrm{res}(\omega)$仍可在$C_s$的边界中非零，
定理\ref{thm:ambient-source}不能直接用于$L_{\rm pure}$。
需要先为这个实际边界类找算子$L_{\rm bdry}$，
然后对$L=L_{\rm bdry}L_{\rm pure}$应用定理。
这也解释为何一些完整积分在消去最大割后，仍出现随参数变化的下sector源。

同理，多重动量最大割与参数超曲面的一次留数不是定义上相同。
只有给出实际变量变换、管与相对链比较，才能沿用同一个$C_s$。
本项目banana有专门比较；一般圈星仍须提交其完整sector的实际模型，
不能把紧曲线宿主及其亏格直接当作完整积分的结构表。

\subsection{图多项式铅笔的接受条件}
在一参数切片上，若第二图多项式写为$F_s=A-sB_0$，
可先研究有理映射$A/B_0$及固定边界的消解。
特征零下可在固定空间上消解其不定性，并在稠密开集删去坏纤维，
为上述固定环境模型提供几何候选。
但这一步不自动证明原物理链避开新极点，也不消去例外面的周期。

接受一个图或sector需实际提交四项：
固定$(X,B)$与所有被删除的固定极点；移动除子的横截性及坏参数集合；
原形式的相对de Rham代表；原物理链、边界及其随参数的比较。
完成后，定理\ref{thm:pencil-separation}给结构上界，
定理\ref{thm:ambient-source}给源，
物理管变分给下界，锚定森林定理追踪方案变化。
这是可直接拼接的推导接口。
仅知道$X$有理或具有环面坐标，不能证明删去固定图超曲面后的$J$为Tate。

\subsection{banana：从正割密度恢复完整函数}
沿用前文正质量二维$h$传播子归一化，
$s_0=(\sum_i m_i)^2$、$I_h(s)=\int_{\mathbb R^{h-1}}(f(y)-s)^{-1}dy$。

\begin{proposition}[全物理谱的无减除色散表示]\label{prop:dispersion}
对$s\in\mathbb C\setminus[s_0,\infty)$，有
\[
 I_h(s)=\int_{s_0}^{\infty}\frac{\rho_h(t)}{t-s}\,dt,
 \qquad
 \rho_h(t)=\int_{\mathbb R^{h-1}}\delta(f(y)-t)\,dy .
\]
因此这个完整参数函数由其正物理谱密度确定，无需额外减除常数。
\end{proposition}
\begin{proof}
先取实$s<s_0$。正性允许Tonelli，恰当函数$f$的余面积公式把
$\mathbb R^{h-1}$按水平集分解，直接给等式。
前文指数尾界保证左端收敛，故右端也收敛。
对复平面的任意紧集$K_0\subset\mathbb C\setminus[s_0,\infty)$，
固定实$s_*<s_0$，
$(t-s_*)/|t-s|$在$t\ge s_0,s\in K_0$上一致有界；
右端由已收敛的$\rho_h(t)/(t-s_*)$控制，
从而全纯。左端也全纯，由恒等定理得到结论。
\end{proof}
这里积分遍历整个参数区间上的不同割纤维。
函数重建不表示完整相对动机与单个纯割纤维同构；
沿参数积分正会保留混合扩张。
证明依赖此二维正质量收敛模型，不覆盖四维发散图的减除多项式。
也不以这一解析恒等式代替一般动量与参数形式周期的实际比较；
该问题的范围在Brown的讨论中已被明确区分\cite[\S10.2]{brownamp}。

\subsection{一个已知完整例子在统一规则中的位置}
等质量$m_i=1$的四传播子banana，BKV给出\cite[Theorem2.2.1]{BKV}
\begin{align*}
 L_3={}&s^2(s-4)(s-16)\partial^3
       +6s(s^2-15s+32)\partial^2\\
      &+(7s^2-68s+64)\partial+(s-4),
 &L_3I_4=-24 .
\end{align*}
这里算子和常数源均为已有结果，不能计作本项目首次求得。
该文的模参数描述将纯解空间识别为椭圆秩二表示的对称平方；
在有限模覆盖上，单值包含有限指标模群。
两个相反的非平凡幺幂一参数闭包生成$SL_2$，
其对称平方三维表示不可约，故纯连接绝对不可约。
前文正常阈值证明使$s=16$的割非零；
因此右因子保护和非零常数源给
\[
  \operatorname{ord}_{\mathbb C(s),\min}(I_4)=4,
  \qquad \partial L_3 I_4=0 .
\]
三阶以内的齐次方程不可能湮灭完整物理函数。
这一例子把纯K3核、物理激发、非齐次源和完整函数阶数接在一起。
它不意味着不等质量族只有四阶，也不据函数阶数认定每个代数纤维的动机秩。
本项目指定不等质量秩八的结论仍由其独立相对对象及扩张证明承担。

\section{完整一圈图的接受域：二次曲面、算术角色与物理链}
\label{sec:oneloop}
本节使前述规则在banana之外获得完整积分的实际输入。
一圈图的Tate性质与四维对数、双对数表达是经典结果，
尤其三角图的Hodge维数已见Bloch--Kreimer\cite[\S7--10]{BKone}；
偶数维多边形的二次曲面递归见Spradlin--Volovich\cite{SVone}。
下面重新从实际参数对推导，保留原基域上的二次Artin角色，
并接入本篇的物理最小化、操作规则与失败边界。
这些整合不作未经检索的原创性声明。

\subsection{从物理积分到完整相对对象}
取特征零数域$k\subset\mathbb R$，正质量平方$\mu_i$及欧氏不变量
$\kappa_{ij}=(r_i-r_j)^2$属于$k$。定义
\[
 U=\sum_{i=1}^N\alpha_i,\quad
 F=U\sum_i\mu_i\alpha_i+\sum_{i<j}\kappa_{ij}\alpha_i\alpha_j
   =\alpha^t C\alpha,
\]
其中$C_{ii}=\mu_i$、$2C_{ij}=\mu_i+\mu_j+\kappa_{ij}$。
Minkowski延拓使用$\kappa_{ij}=-s_{ij}$，不另外改变$F$的符号。
设$B$为坐标超平面之并，$Q=V(F)$，$H=V(U)$。

四维单位传播子的正欧氏归一化为
\begin{equation}\label{eq:onephysical}
 I_N=\int_{\mathbb R^4}\frac{d^4\ell}{\pi^2}
           \prod_i\frac1{(\ell+r_i)^2+\mu_i}
     =\Gamma(N-2)\int_{\sigma}
        \frac{U^{N-4}\Omega_{N-1}}{F^{N-2}},\qquad N\ge3.
\end{equation}
$\sigma$为正射影单纯形；取$U=1$时$\Omega_{N-1}$的取向使其成为正体积形式。
特别地，三角图与盒图分别由
\begin{align}\label{eq:onepairs}
 M_3&=H^2\bigl(\mathbb P^2\setminus(Q\cup H),
             B\setminus(Q\cup H)\bigr),\notag\\
 M_4&=H^3\bigl(\mathbb P^3\setminus Q,B\setminus Q\bigr)
\end{align}
承载。它们是完整参数对象，不是只对最大割取的紧宿主。

\begin{proposition}[实际链比较与一般一圈上界]\label{prop:onecarrier}
在正质量欧氏点，式\eqref{eq:onephysical}是收敛积分，
形式与正单纯形给式\eqref{eq:onepairs}中的实际Nori标架。
更一般地，偶数维$D=2d$、整数幂$a_i\ge1$、$A=\sum a_i>d$
的收敛一圈标量积分，由
\[
 H^{N-1}\bigl(\mathbb P^{N-1}\setminus(Q\cup H),
                    B\setminus(Q\cup H)\bigr)
\]
承载；若形式在$H$上无极点可不删除$H$。
这个对象经某个有限代数扩域后为混合Tate，原域上为混合Artin--Tate。
同一极点集合内的分子读出及参数导数不产生新的非Artin--Tate简单组成。
\end{proposition}
\begin{proof}
欧氏Schwinger公式的被积函数非负，故Tonelli允许先完成平方并积分$\ell$。
高斯积分给$U^{-d}\exp(-F/U)$，再写$\alpha_i=t x_i$、$\sum x_i=1$，
尺度积分给$\Gamma(A-d)$和射影形式
\[
 \frac{\Gamma(A-d)}{\prod_i\Gamma(a_i)}
 \frac{\prod_i\alpha_i^{a_i-1}\,U^{A-2d}\Omega_{N-1}}{F^{A-d}}.
\]
在闭单纯形上$F\ge\min_i\mu_i>0$，$U=1$；
单纯形与其全部面避开极点。尺度零端由$A>d$控制，
无穷端由正质量控制，因而没有尚待认证的边界正则化。
同一次高斯/尺度推导给出完整物理积分等于该参数周期。
这里构造参数Nori载体，不宣称动量对象与它已经有全局动机同构。

每个坐标交叉层是射影空间，$Q$在其上仍为二次式，$H$仍为超平面。
选择共同有限扩域使这些有限多个二次型全部分裂。
光滑分裂射影二次曲面有仿射胞腔分解；
退化二次曲面去掉其线性顶点后，是非退化低维二次曲面上的仿射空间丛。
因此二次曲面及其超平面截面均由Tate对象构成。
对$Q\cup H$的闭并交叉列、开补空间局部化以及$B$各交叉层的有限复形，
依次应用正合性，便得相对对象在扩域后混合Tate。
原域上取有限Galois扩域的限制与迹；除以扩域次数给原对象的直和投影。
故其组成是有限Artin表示的Tate扭转，而非仅在Hodge实现上对角。

有理形式的整数幂提高只改变读出；在常秩参数开集上，
Gauss--Manin求导仍在同一相对连接内。
实际子商不会引入环境中没有的简单因子，得到最后结论。
\end{proof}

有限扩域的选择是上界，不是最小根式域。
分裂后的混合Tate性也不等于MZV性：质量比、运动学支点、分支与扩张类仍在。
本命题没有覆盖维数正则化的全部$\varepsilon$系数、任意UV分子或质量零点的IR极限。

\subsection{三角图与盒图的完整算术分级}
为得到精确载体秩，以下取一般位置：$Q$与各坐标层横截，
三角图的$Q$还与$H$及其坐标交叉横截；顶点质量非零。
记$\chi_d$为平方类$d\in k^\times/k^{\times2}$对应的一维二次Artin角色，
$d$为平方时它就是单位角色。对每条坐标边$ij$置
\begin{equation}\label{eq:edgedisc}
 d_{ij}=(\mu_i+\mu_j+\kappa_{ij})^2-4\mu_i\mu_j .
\end{equation}
它是$Q$与该边的二点交的判别式。
三角图另置
\[
 d_\infty=(\kappa_{12}-\kappa_{13}-\kappa_{23})^2
                          -4\kappa_{13}\kappa_{23}.
\]
因为$F|_H=\sum_{i<j}\kappa_{ij}\alpha_i\alpha_j$，
这是$Q\cap H$的二点判别式；等价于$-\det(F|_H)$的平方类。
盒图令$d_Q=\det C$，即四元二次型的有符号判别式。

\begin{theorem}[可直接读出的结构表]\label{thm:onegraded}
上述一般位置中，完整相对对象的权重分级为
\begin{align}\label{eq:trianglegraded}
 \operatorname{Gr}^W_0M_3&=\Q(0),\notag\\
 \operatorname{Gr}^W_2M_3&=\Q(-1)^2\oplus
                 \bigoplus_{1\le i<j\le3}\chi_{d_{ij}}(-1),\notag\\
 \operatorname{Gr}^W_4M_3&=\chi_{d_\infty}(-2),
 \qquad \rk M_3=7;
\end{align}
以及
\begin{align}\label{eq:boxgraded}
 \operatorname{Gr}^W_0M_4&=\Q(0),\notag\\
 \operatorname{Gr}^W_2M_4&=\bigoplus_{1\le i<j\le4}\chi_{d_{ij}}(-1),\notag\\
 \operatorname{Gr}^W_4M_4&=\chi_{d_Q}(-2),
 \qquad \rk M_4=8.
\end{align}
其它权重为零。式中7和8是所列\emph{完整载体}的秩；
尚未证明某组特化标架生成全对象时，它们只是物理最小秩的上界。
\end{theorem}
\begin{proof}
先处理三角图。令$V=\mathbb P^2\setminus(Q\cup H)=\mathbb A^2\setminus Q^\circ$。
光滑射影圆锥曲线$Q$的有理上同调为$\Q(0)\oplus\Q(-1)$，
即使没有$k$有理点亦如此：有限点的零循环可除以其次数。
从$Q$删去$Q\cap H$的二点，留数正合列给
\[
 H^0(Q^\circ)=\Q(0),\qquad H^1(Q^\circ)=\chi_{d_\infty}(-1).
\]
仿射平面的Gysin列因而给
$H^1(V)=\Q(-1)$、$H^2(V)=\chi_{d_\infty}(-2)$。

每个边界分量$B_i^\circ$是仿射直线删去二点，故
\[
 H^1(B_i^\circ)=H^0(Q\cap B_i)(-1)
               =\Q(-1)\oplus\chi_{d_{jk}}(-1).
\]
三分量在三个未删除顶点组成一个三角环。
其规范化交叉复形给$H^1(B^\circ)$的权重零部分$\Q(0)$，
以及权重二部分$\Q(-1)^3\oplus\bigoplus\chi_{d_{ij}}(-1)$；
$H^2(B^\circ)=0$。
限制映射$H^1(V)\to H^1(B^\circ)$由$d\log F$表示，
在每一边的两个根上留数均为$+1$，故在三个平凡角色中是非零对角向量。
它是单射，恰消去一个$\Q(-1)$。相对正合列
\[
 H^1(V)\longrightarrow H^1(B^\circ)\longrightarrow M_3
            \longrightarrow H^2(V)\longrightarrow0
\]
和权重严格性给\eqref{eq:trianglegraded}。

再处理盒图。$V=\mathbb P^3\setminus Q$，$Q$为光滑二次曲面。
扩域后两族直线给$H^2(Q)$的两个生成元；
两族之和为超平面类，差由$\chi_{d_Q}$作用。
Gysin的推前在和上为非零、在差上为零。
所以$H^1(V)=H^2(V)=0$，$H^3(V)=\chi_{d_Q}(-2)$。

边界的四个面各为$\mathbb P^2$删去光滑圆锥曲线，
有理上同调只有$H^0=\Q$；六条交边为$\mathbb P^1$删去二点，
其$H^1=\chi_{d_{ij}}(-1)$；四个三重交为未删除顶点。
交叉复形的$H^0$行是四面体边界的复形，给$H^2$中的一个$\Q(0)$。
六个交边$H^1$位于总次数二，给六个权重二角色。
没有其它项能进出这些位置，故$H^3(B\setminus Q)=0$，且
\[
 \operatorname{Gr}^W H^2(B\setminus Q)
   =\Q(0)\oplus\bigoplus_{i<j}\chi_{d_{ij}}(-1).
\]
相对列$0\to H^2(B\setminus Q)\to M_4\to H^3(V)\to0$
给\eqref{eq:boxgraded}。所有映射均由实际几何诱导，
不是仅由Betti维数猜出的分解；只断言权重分级，不将扩张强行拆开。
\end{proof}

这个结果无需积分，只需二次型与有限个限制判别式。
一般位置失效时，不能保留7、8照抄；须重新计算实际交叉层。
判别式为零首先表示此模型发生退化，并不自动证明所选物理分支发散。
三角图的7维分裂Hodge结果归于\cite[Proposition10.1]{BKone}；
这里显式写出算术下降和与最小物理对象的区别。

\subsection{物理出口、操作闭合与源的范围}
\begin{corollary}[已接受图类中的物理上界]\label{cor:oneoperations}
在定理\ref{thm:onegraded}的常秩参数域内，完整三角图和盒图的物理最小对象
均为混合Artin--Tate；没有椭圆$H^1$、K3超越核或其它非Artin--Tate简单因子。
沿代数参数曲线，其标量齐次微分阶数分别至多7、8。
分裂上述有限Artin角色后，连接具有长度至多三的平凡分级过滤，
几何幺幂作用的连续三次严格降过滤乘积为零。

几何边界、留数/管及同一环境中的分子读出、参数导数保持此组成上界。
若有限方案改变量另有混合Artin--Tate载体，则改变方案后的最小对象仍属此类。
这些结论既不保证完整载体都物理可见，也不保证全部扩张类不随方案变化。
\end{corollary}
\begin{proof}
物理最小对象是$M_N^\vee$的子商，给第一结论。
一个秩$r$连接中的$r+1$个连续导数线性相关，给阶数上界；
平凡分级过滤来自式\eqref{eq:trianglegraded}--\eqref{eq:boxgraded}及有限根式覆盖。
三个严格降阶算子的乘积越过全部权重。
几何操作使用实际正合列和Tate扭转，
读出变化用子商；有限方案用直和载体及本篇加性比较定理。
\end{proof}

\textbf{必须区分的收缩接口。}
$\alpha_e=0$确实给收缩图的二次多项式$F|_{\alpha_e=0}$，
所以参数面上的几何仍为二次型。
但盒图的面载体是$H^2(\mathbb P^2\setminus Q,\ldots)$，
真正四维三角图还需要删除$H$，形式从$\Omega/F^2$改为$\Omega/(UF)$。
后者由前者的限制以及$H$上的局部化补项连接；
该补项也是Artin--Tate，但二者不是未经比较就相同的物理对象。
同样，三角图缩成四维bubble时出现整体UV减除，
须使用实际减除载体，不能把一个边界周期直接当作未经减除振幅。
因此这里证明的是操作的\emph{类别上界闭合}，而非割、收缩和积分的无条件互换。

\textbf{与完整参数源的连接。}
三角图可在$X=\mathbb P^2\setminus H\simeq\mathbb A^2$上应用前节，
其固定边界为三条仿射直线，构成三角环。
它们自身及交点无高上同调，故
\[
 J_3=H^2(X,B\cap X)=\Q(0).
\]
盒图取$X=\mathbb P^3$及四个坐标面，其交叉复形给
$H^2(B)=\Q(-1)\oplus\Q(0)$，
$H^2(X)=\Q(-1)$同构到前一项，因而
\[
 J_4=H^3(X,B)=\Q(0).
\]
两个固定环境的正单纯形比较值均按取向为$1$。
故在这两个\emph{完整}模型中，任何消去实际\emph{相对留数类}的算子$L$
都给代数源$LI_N=a(s)\in k(s)$（在所取参数曲线上）。
这兑现了前节的代数源条件，而不是把它额外假设为“所有图都成立”。
只消去顶权重的秩一角色通常不够：还剩边界中的对数扩张。

\textbf{三角图的函数形式可以直接解释。}
在$U=1$图表写$\alpha_1=x,\alpha_2=y,\alpha_3=1-x-y$，则
\[
 I_3=\int_0^1dx\int_0^{1-x}\frac{dy}{a y^2+b(x)y+c(x)},
 \quad r^2=b(x)^2-4ac(x),\quad\deg_x(r^2)\le2.
\]
一般$a\ne0$时对$y$的原函数为
\[
 \frac1r\log\frac{2ay+b-r}{2ay+b+r}.
\]
代入两个端点，只剩$\int(dx/r)\log R(x,r)$。
对二次根曲线作有限代数扩域后的有理参数化，
它成为有理一形式乘有理函数对数。
部分分式中高阶极点由分部积分降阶；简单极点的原函数由
$\operatorname{Li}_2$、对数乘积及对数组成。
零次分母或$a=0$在一般开集上直接线性积分，或由无夹点极限处理。
由欧氏连续分支起步再延拓，保留所有$2\pi i$分支常数。
这给三角图的实际双对数充分判据；
盒图的经典双对数结论可用\cite[\S8--9]{BKone}，
不能仅从“权重0、2、4”三个数字推断任意变分都有同样的参数化公式。

\subsection{为什么不能从所有子图为Tate推出完整图为Tate}
\begin{proposition}[纤维结构表不对积分推前封闭]\label{prop:fiberwarning}
即使一个代数映射的每个光滑纤维都在有限扩域后为Tate，
它的总空间仍可含非Tate纯核。
因此仅列出子积分在固定外部数据处的结构，再将它们张量相乘，
不能替代完整图的变分及推前。
\end{proposition}
\begin{proof}
取$\lambda\ne0,1$及光滑射影曲线
\[
 E_\lambda:\ y^2=x(x-1)(x-\lambda),\qquad
 \pi:E_\lambda\longrightarrow\mathbb P^1.
\]
四个分支点是$0,1,\lambda,\infty$。
每个非分支纤维只是二点，其动机为$\Q\oplus\chi_x$；
每个分支纤维的约化动机为$\Q$。
然而Riemann--Hurwitz给$2g(E_\lambda)-2=2(-2)+4=0$，故$g=1$。
更精确地，令$j:V=\mathbb P^1\setminus\{0,1,\lambda,\infty\}\hookrightarrow\mathbb P^1$，
在$V$上$\pi_*\Q=\Q\oplus\mathcal L$，
$\mathcal L$在每个分支点的单值为$-1$。
延拓后$\pi_*\Q=\Q\oplus j_*\mathcal L$，所以
\[
 H^1(E_\lambda)=H^1(\mathbb P^1,j_*\mathcal L).
\]
各穿孔无不变向量，故$j_*=j_!$；全局$H^0,H^2$均为零，
Euler特征为$\chi_c(V)=-2$，右边维数为二。
它是权重一的椭圆核，不是固定二点纤维的Tate张量积。
\end{proof}

这里的根式随着下一积分变量变化；使\emph{每个固定纤维}分裂的域
不能当作一个共同\emph{常数}域。
在整个底上进行分裂覆盖，正好产生了$E_\lambda$本身。
这解释了“先积出一个双对数子图”为什么不会自动保住MPL类别：
下一次推前还看见判别式的分支局部系统及其上同调。
对应的结构预测规则应输出
\[
 \text{纤维核}\quad+\quad\text{判别式局部系统}\quad+
       \quad\text{底与边界的实际推前},
\]
而不能只有纤维类型标签。
同样，二维三传播子sunrise的任意真圈子图至多一圈、收缩该圈后的余图亦为一圈，
但完整光滑参数三次曲线有属一核。
这并不违反森林的低圈张量规则：森林规则控制减除项与方案差，
从未断言未减除完整图就是子图和余图的张量积。

本节由此闭合一类实际入口：完整一圈积分的物理链、算术角色、边界源及操作上界
都无需完整积分即可给出；同时指出推广到多圈必须新增的几何信息。
一般图的推前与最小标架非消失仍是下一项理论责任，
而不是通过增加一圈数值精度就能消除的缺口。


\section{串联细分的实际下降、多个椭圆核与分子选择}
\label{sec:serieselliptic}
本节把结构预测推进到一类完整两圈积分，而不先计算其值。
所取图是三传播子sunrise的一条边的串联细分：
其余两条边质量均为$m>0$，细分路径有$r$条边，质量为
$a_1,\ldots,a_r>0$，所有新增二价顶点的外部动量严格为零。
它有$r+2$条边、$r+1$个顶点和两个圈。
$r>1$时图本身已非三边banana，但约化骨架仍是sunrise；
下述结论不覆盖二价顶点有非零动量注入的一般多点图。

sunrise椭圆几何及其非齐次方程是经典结果，见
\cite[\S4--5]{ABWsun}及\cite[\S3--4]{BKVsun}。
本节直接证明适用于当前图族的相对Stokes下降与物理可见性，
不将部分分式或sunrise椭圆性宣称为新发现。

\subsection{图与完整参数周期：串联变量可以实际消去}
置$\mu_j=a_j^2$，$s=-p_E^2$，采用每圈$d^2k/\pi$的欧氏归一化，记
\begin{align}
 K_{\boldsymbol\mu}(s)
 &=\int_{\mathbb R^2\times\mathbb R^2}
 \frac{d^2k\,d^2\ell/\pi^2}
 {(k^2+m^2)(\ell^2+m^2)
       \prod_{j=1}^r((p_E-k-\ell)^2+\mu_j)}, \label{eq:seriesmom}\\
 S_\mu(s)&=K_{(\mu)}(s).
\end{align}
正质量、$s<0$给绝对收敛；以下均从这一实际链延拓。
设两条未细分边的参数为$x,y$，路径参数为$\gamma_j$，$t=\sum_j\gamma_j$。
生成树与二森林直接给
\[
 U=xy+(x+y)t,\qquad
 F_{\boldsymbol\mu}
 =\left(m^2x+m^2y+\sum_j\mu_j\gamma_j\right)U-sxyt.
\]
Schwinger高斯积分与总尺度积分的系数是$\Gamma(r)=(r-1)!$，
所以完整参数形式为
\[
 (r-1)!\,\frac{U^{r-1}\Omega_{r+1}}{F_{\boldsymbol\mu}^{\,r}}.
\]
在正链上令$\gamma_j=t v_j$、$\sum v_j=1$。
这是边界吹起后的代数坐标变换，Jacobian为$t^{r-1}$。
取$x+y+t=1$，令$\sigma_2$为其正三角形，$\Delta_{r-1}$为$v$单纯形，
置$F_\mu=(m^2x+m^2y+\mu t)U-sxyt$、$\bar\mu=\sum_j\mu_jv_j$，得到实际参数周期
\begin{equation}\label{eq:seriessimplex}
 K_{\boldsymbol\mu}
 =(r-1)!\int_{\sigma_2\times\Delta_{r-1}}
 \frac{(tU)^{r-1}\Omega_2\wedge dv_1\wedge\cdots\wedge dv_{r-1}}
 {F_{\bar\mu}^{\,r}}.
\end{equation}
这一步给指定物理积分的完整参数表示。
之后的形式周期恒等式在这个实际参数模型中证明，
不依赖“数值相等则动机周期相等”，也不需要宣称动量动机与参数动机整体同构。

\begin{proposition}[串联消元的实际Stokes恒等式]\label{prop:seriesstokes}
若$\mu_1,\ldots,\mu_r$两两不同，则式\eqref{eq:seriessimplex}所定义的形式周期满足
\begin{equation}\label{eq:seriespartial}
 K_{\boldsymbol\mu}^{\mathfrak m}
 =\sum_{j=1}^r c_jS_{\mu_j}^{\mathfrak m},
 \qquad c_j=\prod_{\ell\ne j}(\mu_\ell-\mu_j)^{-1}.
\end{equation}
右边为完整sunrise相对周期，包括其边界扩张，而不是仅其最大割。
\end{proposition}
\begin{proof}
对$r\ge2$，使用独立坐标$v_1,\ldots,v_{r-1}$，$v_r=1-\sum_{i<r}v_i$。
在$\sigma_2$与$v$空间的乘积上，置
\[
 \eta=\frac{(r-2)!}{\mu_r-\mu_1}\,
 \frac{(tU)^{r-2}}{F_{\bar\mu}^{\,r-1}}\,
 \Omega_2\wedge dv_2\wedge\cdots\wedge dv_{r-1}.
\]
$r=2$时最后的外积为空积。由于底维数为二，沿$x,y$的微分与$\Omega_2$
楔乘为零；其余$v_i$微分亦因重复因子为零。
而$\partial_{v_1}F_{\bar\mu}=tU(\mu_1-\mu_r)$，故
\[
 d\eta=(r-1)!\frac{(tU)^{r-1}}{F_{\bar\mu}^{\,r}}
       \Omega_2\wedge dv_1\wedge\cdots\wedge dv_{r-1}.
\]
只有$v_1=0$与$v_r=0$两个面给非零拉回；它们的诱导取向给
\begin{equation}\label{eq:seriesrec}
 K_{(\mu_1,\ldots,\mu_r)}^{\mathfrak m}
 =\frac{K_{(\mu_1,\ldots,\mu_{r-1})}^{\mathfrak m}
        -K_{(\mu_2,\ldots,\mu_r)}^{\mathfrak m}}
       {\mu_r-\mu_1}.
\end{equation}
需要核对底的边界没有被遗漏。三个坐标边上$\Omega_2$的拉回为零。
吹起$\sigma_2$的三个顶点以分离极点与链：
在$x=1$附近写$(y,t)=\rho(u,v)$，则
$\Omega_2$含$\rho\,d\rho$，$tU$消失二阶，$F_{\bar\mu}$消失一阶；
$\eta$的径向阶数为$r-2\ge0$。
在$y=1$附近相同；在$t=1$附近$tU,F_{\bar\mu}$均消失一阶，
$\eta$的径向阶数为零。所有这些拉回均含$d\rho$，
所以在例外边界上的限制为零。
质量凸组合$\bar\mu>0$保证上述局部系数在实际正链上无新零点。
复代数侧删除极点的真变换，并按需取共同消解；
上述相对形式及链可拉回，边界限制恒零仍成立。

因此这是真实几何对上的Stokes关系，其形式版本由相对上同调边界态射给出。
从$r=1$归纳式\eqref{eq:seriesrec}，或者比较简单极点的留数，
即得式\eqref{eq:seriespartial}。
\end{proof}

特别地，$r=2$时只用一个区间：
\[
 \frac{\partial}{\partial v}\frac1{F_{\mu_2+v(\mu_1-\mu_2)}}
 =(\mu_2-\mu_1)\frac{tU}{F_{\mu_2+v(\mu_1-\mu_2)}^2}.
\]
它明确指出哪两个参数端面产生两个椭圆相对对象。
不是把两个sunrise动机张量相乘，而是将这个特定完整周期降到它们的直和。

\subsection{几何上界：每个端点恰有一个可能的椭圆纯核}
令$E_a(s)\subset\mathbb P^2$为三次曲线$F_{a^2}=0$，在共同光滑参数开集上工作。
吹起三个坐标顶点得到$X$，令$B$为坐标三角形的全变换，$\widetilde E_a$为曲线真变换。
实际载体为
\[
 \mathcal M_a
 =H^2(X\setminus\widetilde E_a,\ B\setminus\widetilde E_a).
\]
$X$是有理曲面，$B$为六个有理分量组成的环。
每个分量被$\widetilde E_a$刺穿，交点随$s$固定；
在必要有限常数扩域后，所有点与边界组成均为Tate。
对$X\setminus\widetilde E_a$的Gysin列给
\[
 \operatorname{coker}\!\left(H^0(\widetilde E_a)(-1)\to H^2(X)\right)
 \longrightarrow H^2(X\setminus\widetilde E_a)
 \longrightarrow H^1(E_a)(-1)\longrightarrow0.
\]
左项为Tate，右端用$H^3(X)=0$。
又因$H^2(B\setminus\widetilde E_a)=0$，
相对长正合列使$\mathcal M_a$满射到中间项，核由边界上同调承载。
所以$\mathcal M_a$的非Artin--Tate组成恰为一次
\[
 \mathcal V_a:=H^1(E_a)(-1).
\]
这是完整相对载体的组成论断，尚未单靠此步证明物理标架取满。
式\eqref{eq:seriespartial}给$K_{\boldsymbol\mu}$的上界
$\bigoplus_j\mathcal V_{a_j}$，每个端点至多贡献一次。
下文在明确条件下证明这些组成全部实际出现。

\subsection{阈值分离把几何上界变成物理下界}
本小节在参数曲线上的变分范畴及其正则奇异微分实现中陈述，
不是将单值变换当作固定纤维的Nori自同态。
记
\[
 \mathcal D_a=\{0,\ a^2,\ (a-2m)^2,\ (a+2m)^2,\ \infty\},
 \qquad \tau_a=(a+2m)^2.
\]
这是几何奇异值的包含集合；重复值只记一次，
标量方程的表观奇点不被加入几何集合。
可直接从三次曲线核对：坐标边上除三个顶点外的交点光滑，
顶点的一阶项非零。环面内以$\alpha_i=e^{y_i}/m_i$，$m_i=(m,m,a)$，
写$F/(\alpha_1\alpha_2\alpha_3)=AB-s$，
$A=\sum_i m_i e^{y_i}$、$B=\sum_i m_i e^{-y_i}$。
当$s\ne0$时$A,B\ne0$，临界方程给
$m_i e^{y_i}B-m_i e^{-y_i}A=0$，
即所有$e^{2y_i}=A/B$；模掉共同尺度后，各坐标仅有符号选择。
因此临界值为$(\pm m\pm m\pm a)^2$，$s=0$另使$F=UM$约化。
在$\tau_a$处只有正质量和对应的一个临界点，正定Hessian给普通节点。
该纤维在边界光滑，故没有第二个奇点。
一个约化三次曲线若分成直线与二次曲线，Bézout迫使两个交点或更高交数奇点，
不能只有这一个普通节点；三直线情形亦不符。
因此正常阈值纤维是不可约的结点三次曲线。

\begin{proposition}[端点椭圆变分的不可约性与区分]\label{prop:seriesirred}
$\mathcal V_a$几何绝对不可约。
若$\tau_a\notin\mathcal D_b$，则$\mathcal V_a$与$\mathcal V_b$
在共同底上非同构，且这一差别在任意有限底覆盖后仍存在。
\end{proposition}
\begin{proof}
不可约结点三次曲线是乘法退化，故$j(E_a(s))$在$\tau_a$有极点；
族非等平凡。它有坐标顶点给出的截面。
取有限覆盖加上级别$n\ge3$的结构，得到到底曲线$Y(n)$的非恒定代数映射。
在去除有限分支值及其逆像后，它是有限覆盖到$Y(n)$的一个开子集；
基本群像在$\Gamma(n)$中有有限指数。
因此$H^1(E_a)$的几何单值群Zariski稠密于$\operatorname{SL}_2$，
标准二维表示绝对不可约，Tate扭转不改变此结论。

$\tau_a$处$\mathcal V_a$的Picard--Lefschetz单值为非平凡幺幂变换，
$\mathcal V_b$则在同一圆盘内光滑延拓，其单值为单位。
若取有限覆盖，前者变为某个正整数次幂，仍非平凡，
后者仍为单位；两表示不可能同构。
\end{proof}

\begin{theorem}[串联两圈图的精确椭圆可见组成]\label{thm:seriesvisible}
设质量为正代数数，$\mu_j$两两不同，且
\begin{equation}\label{eq:seriesseparation}
 \tau_{a_j}\notin\mathcal D_{a_\ell}\quad(j\ne\ell).
\end{equation}
则$K_{\boldsymbol\mu}$的最小物理变分对象，在上同调载体约定下，
其非Artin--Tate半单部分为
\[
 \bigoplus_{j=1}^r H^1(E_{a_j})(-1),
\]
每个互不同构的椭圆组成出现恰好一次。
采用全文的同调标架约定时应取此式的对偶。
任何在$\mathbb C(s)$上消去$K_{\boldsymbol\mu}(s)$的非零齐次标量微分算子
阶数至少为$2r$。
\end{theorem}
\begin{proof}
上界已由真实相对Stokes下降及端点载体给出。
在$\tau_{a_j}$处取一个从共同欧氏基点接入的套索。
其余完整相对对在该值附近光滑，边界交点亦不碰撞，故它们的变分为零。
于是
\[
 (T_j-1)K_{\boldsymbol\mu}
      =c_j(T_j-1)S_{\mu_j}.
\]
前文对所有正质量二维banana的Morse阈值证明适用于这里的三传播子积分：
正链与消失管的交数绝对值为一，
$\operatorname{Disc}S_{\mu_j}=2\pi i\,\rho_j$，
且阈值上方$\rho_j>0$。故右端是非零椭圆留数周期；
$c_j\ne0$不可能消掉它。
命题\ref{prop:seriesirred}使它的延拓轨道生成整个二维纯变分。
所有$j$因此给彼此不同的受保护下界，与每个至多一次的上界相合。

若$LK_{\boldsymbol\mu}=0$，微分算子与解析延拓交换，
所以$L$消去全部这些非零割及其延拓。
每个割的最小算子$L_j$不可约且阶数为二，
从右除法得$L_j$为$L$的右因子。
它们的解局部系统互不同构且简单，和为直和；
最小公共左倍数阶数为$2r$，从而$\operatorname{ord}L\ge2r$。
\end{proof}

这给出精确的纯组成，而不是完整混合扩张的分裂。
常数、质量Kummer和相对链扩张仍须保留，完整齐次阶数也可能大于$2r$。
本定理不推出每个指定代数$s$处的椭圆曲线仍彼此非同源，
更不推出某个积分数值不属于MZV。

\textbf{一个无需积分的明确实例。}
取$r=2$，$a_1=3m,a_2=8m$，则
\[
 \mathcal D_{3m}/m^2=\{0,1,9,25,\infty\},\qquad
 \mathcal D_{8m}/m^2=\{0,36,64,100,\infty\}.
\]
两个正常阈值$25m^2,100m^2$各不属于另一族的奇异集，且
\[
 K_{(9m^2,64m^2)}
   =\frac{S_{9m^2}-S_{64m^2}}{55m^2}.
\]
因此这个四边两圈图的完整物理函数同时需要两个互不同构的椭圆变分，
其非Artin--Tate半单秩恰为四，齐次标量阶数至少为四。
它在欧氏域虽是正积分，表示式却是两个相对周期之差；
两个不同阈值证明这个差不可能将任一椭圆组成抵消。

\subsection{质量碰撞和分子读出：结构怎样实际消失}
\begin{corollary}[合流质量与传播子消去]\label{cor:seriesselect}
对同一细分拓扑，有以下两种不同的结构减少机制。
\begin{enumerate}
\item 全部串联质量平方合流为$\mu>0$时，
\[
 K_{(\mu,\ldots,\mu)}
   =\frac{(-1)^{r-1}}{(r-1)!}\,\partial_\mu^{\,r-1}S_\mu .
\]
该读出在同一个sunrise相对Gauss--Manin对象中，
故非Artin--Tate组成至多为一个$H^1(E_{\sqrt\mu})(-1)$。
\item 在互异且阈值分离的情形，令$z=(p_E-k-\ell)^2$，
取分子$P(z)\in k[z]$，$\deg P\le r-1$。
在实际传播子消去后得到的参数形式周期为
\[
 K_{\boldsymbol\mu,P}^{\mathfrak m}
    =\sum_j P(-\mu_j)c_jS_{\mu_j}^{\mathfrak m}.
\]
其可见椭圆组成恰由$P(-\mu_j)\ne0$的指标给出，
对应齐次阶数下界为这些指标数目的两倍。
\end{enumerate}
\end{corollary}
\begin{proof}
第一式可在绝对收敛欧氏积分中逐次求导，
也可在参数形式中直接用
$\partial_\mu^{r-1}(1/F_\mu)
=(-1)^{r-1}(r-1)!(tU)^{r-1}/F_\mu^r$核对；
同质量时$v$单纯形体积$1/(r-1)!$给相同结果。
形式版本由同一实际相对Gauss--Manin连接及水平正链给出。
求导只改变de Rham读出，不能增加载体的简单组成。
这同时说明不能把式\eqref{eq:seriespartial}中各个发散系数单独取极限。

第二式的输入是有限代数恒等式
\[
 P(z)=\sum_j P(-\mu_j)c_j\prod_{\ell\ne j}(z+\mu_\ell).
\]
每项分子逐个取消对应传播子，留下一个实际sunrise积分；
这一步在参数化前就是同一有理被积形式的等式。
以线性Schwinger参数化定义该插入的实际参数周期，
所余端点与命题\ref{prop:seriesstokes}的相对边界一致。
因此所用关系由有理形式线性、实际传播子消去和相对Stokes组成，
没有诉诸数值周期映射单射，也不声称任意分子的动量/参数动机同构。
在当前次数范围所有留下的积分收敛，无需为这一恒等式补UV反项。
若系数为零，该端点在此表示中消失；
若系数非零，阈值套索和定理\ref{thm:seriesvisible}的证明逐字适用。
\end{proof}

例如两个质量时，$P(z)=z+\mu_1$给$K_{\boldsymbol\mu,P}=S_{\mu_2}$：
第一个椭圆核被分子真正删除，第二个由自身阈值保护。
串联段的“收缩”在本节专指这个实际传播子消去/参数端面；
任意抽象边收缩是否保持物理形式，仍须另证。

若再加入具有Artin--Tate载体的局部项，前文加性Serre定理保证这些
非Artin--Tate重数不变；这是有实际载体输入的方案比较，
不是对所有重整化方案的无条件断言。
本类本身在二维正质量条件下收敛，故不需要借重整化完成上述分类。

\subsection{这一接受规则交付的结构预测}
有限输入只需读出串联路径、零注入条件、质量及分子多项式：
建立各$F_{\mu_j}$与有限集合$\mathcal D_{a_j}$，
核对质量互异和式\eqref{eq:seriesseparation}，
再检查有限个值$P(-\mu_j)$，即可给精确的可见椭圆组成和阶数下界。
不要求先求sunrise的完整特殊函数表达式。
若阈值碰撞或路径上有非零注入，规则明确退出精确分支：
前者仍可用已有直和上界但须重新证明非抵消，
后者的Symanzik多项式已改变，连此下降公式也不能照搬。

这个例子同时落实图数据、实际边界下降、纯核、物理链保护和分子选择。
它解释为何拓扑标签必须带质量与读出，且多个椭圆核应以组成及扩张记录，
而非强行给整图贴一个单独的“椭圆”标签。
一般多圈的判别式推前、非串联图的完整相对比较与RG临界质量接口，
仍是总目标的后续主线。



\section{边界携带的椭圆核：非零动量注入与实际收缩}
\label{sec:boundaryelliptic}
\subsection{从绝对几何到完整相对结构的补充规则}
上一节的差商下降要求内部顶点零外动量。现在允许该顶点注入非零动量，
不再使用串联部分分式。新的入口是边界的权重，而非进一步求解积分。
下面先给适用于一类三维参数空间的规则，再在四边两圈图上落实其几何及物理读出。

在变分层，本文称一个分级组成为Artin--Tate（简称AT），
是指它在某个有限\emph{代数}基变换后成为常值Tate变分。
整个对象允许这些组成间存在扩张。有限单值的代数分支不应误记为椭圆核。
固定数域纤维则采用相应的Artin--Tate对象；这两层不能互相替代。

\begin{theorem}[有理极点曲面的边界结构规则]\label{thm:boundaryweight}
设$X$是光滑射影三维簇，其上同调为AT；$S\subset X$为光滑有理曲面，
也允许$S$为若干光滑有理曲面组成的SNC除子，其双交有理、三交有限。
设$B=\bigcup B_i$为SNC除子，各$B_i$为光滑有理曲面。
允许先作由有限点与有理曲线组成的相容消解。
假设$S+B$为SNC；$B_i\cap S$的正亏格分支正规化记为$C_\alpha$，
其余分支有理；$B_i\cap B_j$的分支均有理，三重交有限。
令
\[
 U=X\setminus S,\quad B^\circ=B\setminus S,\quad
 M=H^3(U,B^\circ;\Q),\quad V_\alpha=H^1(C_\alpha;\Q)(-1).
\]
则$H^*(U)$为AT，且在有理混合Hodge结构中有
\begin{equation}\label{eq:boundaryweight}
 \operatorname{Gr}^W_3 M\simeq\bigoplus_\alpha V_\alpha;
 \qquad\text{其余权重分级均为AT。}
\end{equation}
同一几何正合列在实际有理Nori对象的AT Serre商中给
$M\simeq\bigoplus_\alpha V_\alpha$。
这不是完整混合对象的直和分裂，也不是任意标架都看见全部右端的断言。
在光滑相容参数开集上，相同结论成立于变分层。
\end{theorem}
\begin{proof}
$X$与$S$的Gysin列（多分支时用逐交层的权重谱序列）
表明$H^*(U)$的组成均为AT，故$H^2(U)$与$H^3(U)$均无权重三。
单一光滑有理$S$时，还可具体看出$H^2(U)$纯权重二，
$H^3(U)\hookrightarrow H^2(S)(-1)$纯权重四；
多分支时不需要沿用这两个额外的纯性断言。
对每个有理曲面$B_i$取去除曲线的权重列。
其$H^2(B_i\setminus S)$的权重三部分恰为
$\bigoplus_{C_\alpha\subset B_i}H^1(C_\alpha)(-1)$；
曲线交点以及有理分支只贡献偶权重AT对象。
这一事实可由逐分支Gysin列与交点的Mayer--Vietoris列得到，
因此不要求被去除的曲线并集本身光滑。

在$B^\circ$的覆盖谱序列中，双交是去掉有限点的有理曲线，
三交是有限点。它们所有上同调均为AT，特别没有权重三。
于是从单个面的$H^2$发出的任何可能微分都不能消去权重三部分，
得到$\operatorname{Gr}^W_3H^2(B^\circ)=\bigoplus V_\alpha$。
再使用实际相对列
\begin{equation}\label{eq:boundaryrelative}
 H^2(U)\longrightarrow H^2(B^\circ)
 \overset{\delta}{\longrightarrow}M
 \longrightarrow H^3(U)\longrightarrow H^3(B^\circ).
\end{equation}
两侧环境项没有权重三，权重严格性给出式\eqref{eq:boundaryweight}。
全部其它组成由上述AT项的核、商及扩张构成。

所用Gysin、覆盖与相对列均来自实际几何对；在Nori范畴中先取这些列，
再作AT Serre商，环境项和所有有理交层消失，得到所述同构。
这里没有用Hodge实现的同构反推一个未经构造的动机态射。
族中须保留有理曲面除子类的有限代数单值：
光滑有理曲面的$H^2$由有限个代数除子生成，在泛点有限扩域后可定义，
缩小参数开集使其延伸，便得到所需有限覆盖。故这里的AT断言
具有实际代数除子依据，不是从任意逐纤维Tate性推出全推前。
\end{proof}

\subsection{四边两圈三点图及接受开集}
顶点为$A,B,C$；边$x,y$平行连接$A,B$，质量均为$m>0$；
边$u$连接$A,C$，质量$a>0$；边$v$连接$C,B$，质量$b>0$。
三个入射动量满足$p_A+p_B+p_C=0$，令$s_i=p_i^2$，
并允许$p_C\ne0$。在二维复运动学中，三个$s_i$没有额外多项式Gram关系。
先在真实欧氏区域定义正链积分，再作解析延拓。

生成树及二森林直接给
\begin{align}
 U&=xy+(x+y)(u+v),\\
 M_0&=m^2(x+y)+a^2u+b^2v,\\
 F&=M_0U-s_Axyu-s_Bxyv-s_Cuv(x+y),\\
 I(s_A,s_B,s_C)&=\int_{\sigma_3}\frac{U\Omega_3}{F^2}.
 \label{eq:injectedgraph}
\end{align}
每个二维环动量采用$d^2k/\pi$归一化，故本式的$\Gamma(2)=1$。
正质量欧氏积分UV和IR均收敛。图在文献中称为$(2,1,1)$ ice cream cone；
Doran--Harder--Vanhove已研究其绝对几何及局部系统
\cite[\S7.3, 式(120)--(121), Proposition7.6]{DHV}。
变量对应为其$(y_1,y_2,x_1,z)=(u,v,x,y)$，
$(p_1^2t,p_3^2t,p_2^2t)=(s_A,s_B,s_C)$。
下文重新证明本节所需的曲面与\emph{相对边界}结论，
不把图超曲面的Tate结论称为新的发现。

取以下非空Zariski开集：$a,b,m\ne0$；
\begin{equation}\label{eq:injecteddiscriminant}
 \Delta_C=(a^2+b^2-s_C)^2-4a^2b^2\ne0;
\end{equation}
两个端面三次曲线光滑；下述两个节点普通，曲面在它们之外光滑；
消解后的交层具有上述正常交叉性质。
这是接受条件，特殊参数必须重新分析，不能按连续性自动继承结论。

\begin{proposition}[极点曲面有理，两个端面椭圆]\label{prop:injectedgeometry}
在上述开集上，$S_0=V(F)\subset\mathbf P^3$有两个普通双点，
其光滑消解为有理曲面。存在由有理曲线与有限点的吹起得到的共同模型$X$，
使正链与极点曲面分离，并满足定理\ref{thm:boundaryweight}的条件。
其中仅有两个正亏格边界分支：
\begin{align}
 E_a(s_A):\ &(m^2x+m^2y+a^2u)\{xy+(x+y)u\}-s_Axyu=0,\\
 E_b(s_B):\ &(m^2x+m^2y+b^2v)\{xy+(x+y)v\}-s_Bxyv=0.
 \label{eq:twofaceelliptics}
\end{align}
因此该实际完整相对载体$M$满足
\begin{equation}\label{eq:injectedcarrier}
 \operatorname{Gr}^W_3 M=H^1(E_a)(-1)\oplus H^1(E_b)(-1),
\end{equation}
其它分级为AT，而忘记边界后的$H^3(U)$全部为AT。
\end{proposition}
\begin{proof}
直线$\ell=\{x=y=0\}$包含于$S_0$。其上梯度为
\[
 (F_x,F_y,F_u,F_v)=(d,d,0,0),\quad
 d=a^2u^2+(a^2+b^2-s_C)uv+b^2v^2.
\]
式\eqref{eq:injecteddiscriminant}保证$d$有两个不同根$P$。
在根处取直线局部坐标$w$，二次首项为
\[
 d'(P)w(x+y)+A_Pxy+m^2(u+v)(x+y)^2,
 \qquad A_P=a^2u+b^2v-s_Au-s_Bv.
\]
$d'(P)\ne0$且$A_P\ne0$时此二次型非退化，故为普通双点。
参数族在基点之外由Bertini光滑；其基点集为$\ell$及有限点。
有限点除四个坐标顶点外，位于某坐标边的非零质量线性因子的根，
或$[1:-1:0:0]$。沿该边或质量方向的导数非零，故均光滑；
四顶点也由$F$的非零线性项光滑。
$A_P\ne0$、端面光滑等条件仅去掉真代数闭集，证明接受开集非空。

任选一个节点，并在有限扩域后将其写为$[1:0:0:0]$。
三次方程成为$X_0Q_2(Y)+Q_3(Y)=0$。
从该点投影到$\mathbf P^2_Y$有有理逆
$Y\mapsto[-Q_3(Y):YQ_2(Y)]$，在$Q_2\ne0$上成立。
故曲面及其光滑消解有理。这个论证保留可能的有限算术扩域。

实际正链首先要求处理整条$\ell$，不能只消解两个复节点。
吹起$\ell$，在$x=\rho\xi,y=\rho\eta$图中有
\[
 U=\rho(u+v)(\xi+\eta)+O(\rho^2),\qquad
 F=\rho d(u,v)(\xi+\eta)+O(\rho^2).
\]
$\Omega_3$也含一个$\rho$因子，因此$U\Omega_3/F^2$沿新例外面
没有径向极点。正链方向$\xi,\eta\ge0$满足$\xi+\eta>0$，
且欧氏区域正$u,v$上$d>0$，故链与严格极点曲面分离。
例外面是有理直纹面；其与严格曲面的交为一条截面$\xi+\eta=0$
和$d=0$的两条纤维。在相交点处，严格方程的$\rho$导数
为$A_P\xi\eta\ne0$，所以此吹起也消解了两个节点。

剩余$x,y$坐标顶点吹起时$U,F$的径向阶均为一，
$\Omega_3$阶为二，故拉回形式阶为一，仍无例外极点。
它们的例外曲面与极点曲面交于直线，正方向由正质量首项避开。
为使复杂边界正常交叉所需的进一步吹起，只在极点上的有限点及由此产生的
有理交线进行，不沿整条光滑椭圆曲线吹起。
例如例外面的截面与纤维交点处，局部形式为$S:\rho+dh=0$与$E:\rho=0$；
点吹起后再吹起有理交线即可作嵌入消解。
这些后续中心在极点上，取极点的\emph{约化全逆像}，
则新的补空间与先前$U$同构，不会错误地把新增极点面纳入可积边界。
新增极点分支和其交层均有理，适用定理的多分支版本。
所有新曲面、双交曲线均有理，$X$上同调为AT。

原端面$v=0,u=0$分别给式\eqref{eq:twofaceelliptics}的光滑平面三次曲线。
其它两个面，例如$x=0$，限制为
\[
 F|_{x=0}=y\bigl[(u+v)(m^2y+a^2u+b^2v)-s_Cuv\bigr].
\]
它由直线与二次曲线组成；$y=0$面相同。
上述吹起仅在端面曲线的有限点作修改，保持它们的$H^1$。
应用定理\ref{thm:boundaryweight}即得式\eqref{eq:injectedcarrier}。
\end{proof}

有理曲面并不等于所有定义域上的纯Tate对象；这里保留有限Artin下降。
也不能从绝对载体为AT推出整个标量积分属于MPL，
因为式\eqref{eq:injectedcarrier}中的椭圆层恰由相对边界进入。
这个结论针对几何载体；单个形式的可见性由第\ref{sec:scalarboundary}节另作标架比较。

\subsection{收缩端面的实际态射与联合物理读出}
设$B_v$为$v=0$面的严格变换，$B_{\ne v}$为其余边界并集，
用上标$\circ$表示去掉极点。闭覆盖的相对切除和连接态射给
\begin{equation}\label{eq:facemap}
\begin{split}
 H^2(B_v^\circ,B_v^\circ\cap B_{\ne v}^\circ)
 &\simeq H^2(B^\circ,B_{\ne v}^\circ)\\
 &\longrightarrow H^2(B^\circ)
 \overset\delta\longrightarrow M.
\end{split}
\end{equation}
记复合为$\delta_v$，类似定义$\delta_u$。
这是完整带角点面对象的态射，不是把裸椭圆曲线直接嵌入完整积分。

\begin{theorem}[实际边界读出的精确组成]\label{thm:jointsector}
在式\eqref{eq:injectedgraph}的接受开集上，取联合物理读出
\[
 \mathcal I=(I,\ I[D_v],\ I[D_u]),
\]
其中$D_u,D_v$为相应欧氏传播子分母。则
\begin{equation}\label{eq:cancelledphysical}
 I[D_v]=S_{a^2}(s_A),\qquad I[D_u]=S_{b^2}(s_B).
\end{equation}
两者通过$\delta_v,\delta_u$和同一正相对链实现。
联合读出的最小变分对象，以上同调约定记录其非AT半单组成时，恰为
\[
 H^1(E_a(s_A))(-1)\oplus H^1(E_b(s_B))(-1),
\]
各出现一次，非AT半单秩为四。仅取单个$I$时，本节只证明它的非AT组成
是上述两项的子集；第\ref{sec:scalarboundary}节再以独立的边界链变分证明两项均可见。
\end{theorem}
\begin{proof}
在欧氏动量中置$q=p_A-k-\ell$，可取
\[
 D_u=q^2+a^2,\qquad D_v=(q+p_C)^2+b^2.
\]
分母$D_v$消去后，剩余三条传播子为质量$(m,m,a)$、外动量$p_A$
的sunrise；消去$D_u$后，因$p_A+p_C=-p_B$，得到质量$(m,m,b)$、
外动量$p_B$的sunrise。两积分收敛，同一归一化给式\eqref{eq:cancelledphysical}。
这也由各自的三传播子Schwinger参数化直接给出对应端面形式，
不依赖尚未构造的一般动量/参数动机同构。

连接态射与链边界的配对公式为
\[
 \langle\sigma_3,\delta_v\omega_a\rangle
   =\langle(\partial\sigma_3)_v,\omega_a\rangle,
\]
取$\delta_v$的取向号使右端为正sunrise参数链；$u$面同理。
相对面对象包含其它边和例外边界，故没有遗漏角点或填充项。
定理\ref{thm:boundaryweight}的证明还表明$\delta_v,\delta_u$在权重三上
分别识别$E_a,E_b$的两项。

将$\omega,\delta_v\omega_a,\delta_u\omega_b$作为同一相对载体的三个读出，
联合最小对象是$M$的子商，故由式\eqref{eq:injectedcarrier}得到重数上界一。
两条sunrise物理函数的正常阈值割均非零，
前节的非恒定椭圆变分与不可约单值证明给各自完整$H^1$下界。
在独立参数$s_A,s_B$上，它们互不同构：
沿$s_A=(a+2m)^2$绕行而固定一般$s_B,s_C$，第一个有非平凡幺幂单值，
第二个无单值；交换$A,B$同理。故两项都必须出现，匹配上界。
这里使用的是变分与实际物理面读出，未宣称所有指定代数纤维都互不同源。
\end{proof}

三个读出可进一步作可观测量选择。
只保留$I[D_v]$时，$E_b$方向在该读出中确实消失，只剩$E_a$；
只保留$I[D_u]$则相反。这给出非零内部动量注入下的实际分子筛选与收缩规则，
不需要先求三个函数的特殊函数表达式。
若加入已证由AT对象承载的局部项，加性Serre定理继续保持联合非AT组成。
这不赋予任意局部反项自动为AT的性质。

\subsection{为什么绝对微分方程可能有椭圆边界源}
\begin{corollary}[绝对消去不等于完整消去]\label{cor:boundarysource}
将上述族限制到一条光滑代数参数曲线，令$L$为消去$[\omega]\in H^3_{\dR}(U)$
的有理系数微分算子。则$LI$的实际周期载体可取
\[
 K=\operatorname{coker}\{H^2(U)\longrightarrow H^2(B^\circ)\}.
\]
$K$的非AT组成恰由两个端面椭圆核给出；$LI$的标架可能只看见其中一部分。
特别，绝对载体的有限代数单值，不能单独推出$LI$是代数函数。
\end{corollary}
\begin{proof}
相对列\eqref{eq:boundaryrelative}把$K$识别为
$M\to H^3(U)$的核，且该列与Gauss--Manin连接相容。
因此$L[\omega]$在相对de Rham实现中属于$K$，
正链读出限制到$K$即给$LI$。其组成由定理\ref{thm:boundaryweight}确定。
在形式层若$L\omega=d\eta$，相对锥中的恒等式为
$[(d\eta,0)]=[(0,-\eta|_{B^\circ})]$，明确展示边界源。
该负号采用相对微分$d(\alpha,\beta)=(d\alpha,\alpha|_B-d\beta)$；
改变锥约定只同时改变连接态射和取向号，不改变实际周期。
\end{proof}

上述推论是前文完整\emph{相对}留数定理的补充：
只消去绝对部分时，边界还可能带非Tate源；
消去包含边界的完整相对留数像后，才可使用已证的环境背景结论。
这也说明文献中的绝对Picard--Fuchs算子与完整物理函数的齐次最小算子
不可直接等同。本节不把忘边界的上同调映射等同于尚未比较的全部物理最大割。

若退回$p_C=0,s_A=s_B$，本族落在特殊几何层，前节的实际Stokes公式重新适用；
不能在未检查奇异层及附近循环的情况下，用该特殊层的最小性证明一般参数标架最小性。
本节已闭合的输出为：一般非零注入下的完整载体上界、两个端面核的来源、
实际收缩态射、联合读出的精确组成，以及绝对消去后的源载体。
单个标量$I$的两个椭圆方向由下一节的专属链变分证明均受保护；
不能仅用本节的载体组成替代该证明。


\section{从边界核到单个物理读出：一维变分保护}
\label{sec:scalarboundary}
\subsection{把非零跳跃升级为指定简单因子的判据}
上一节的两个椭圆端面是完整载体的组成。
本节证明：在声明的光滑参数开集上，不必添加收缩分子，
原标量积分$I$本身就看见两者。证明只用局部几何与渐近首项，
不求完整积分、连接矩阵或数值单值。

首先在同调方向表述一个可重复使用的判据。
设$E$是连通参数开集上的有限秩有理局部系统，
带单值不变过滤$W$；相关分级半单。
设$V$是其中恰出现一次的简单局部系统。
令$T$为一个具体闭路的单值，$N=T-1$。
形式$\lambda$为单值代数de Rham读出，$\sigma$为一个平坦Betti芽，
$f=\langle\lambda,\sigma\rangle$。
子商$R/P$承载此标架，是指$\sigma\in R$且$\lambda|_P=0$；
这里的等式均在相应实现及参数开集上理解。

\begin{theorem}[一维变分保护]\label{thm:rankonevisibility}
假设$\rk N=1$是在整个$E$上成立，且$N$在某个
$\operatorname{Gr}^W_kE$中的$V$上诱导非零算子。
若实际函数的跳跃
\[
 \Delta_T f=\langle\lambda,N\sigma\rangle
\]
不是恒零函数，则每个承载此标架的局部系统子商$R/P$都含有$V$。
因而任何相容的实际最小变分对象也保留$V$。
\end{theorem}
\begin{proof}
记$\Gamma=\im N$。由于$N$在$V$上非零，
可以在$W_kE$中选一个提升，使其$N$像在$\operatorname{Gr}_k^W$中非零。
整个像只有一维，故
$\Gamma\subset W_kE$、$\Gamma\cap W_{k-1}E=0$，
而其非零分级像位于$V$。
跳跃非零说明$N\sigma\ne0$，于是
$\Gamma\subset R$且$\Gamma\not\subset P$。
分级半单性与$V$的简单性表明$\operatorname{Gr}_k^WR$包含整个$V$。

反设$R/P$不含$V$。对诱导过滤取分级仍有正合列，
且$V$在$\operatorname{Gr}E$中重数一，
所以$\operatorname{Gr}_k^WP$必须包含该$V$。
选$p\in P\cap W_kE$，使其分级像$v\in V$满足$Nv\ne0$。
单值保持$P$，故$Np\in P$；它在分级上的像非零，
因而$Np$生成$\Gamma$。这与$\Gamma\not\subset P$矛盾。
\end{proof}

判据中的总秩一不可换成“$V$上的变分秩一”。
例如将一个具有非零退化单值的椭圆变分与一个对数扩张取直和，
只读出后者，仍有非零跳跃却完全看不见椭圆层。
同时，本定理不把一个参数闭路当成固定纤维Nori群的元素：
它首先是变分及实际函数的结论。

\subsection{边界普通二重点的总变分秩}
考虑一个光滑三维复流形中的相对对$(U,B^\circ)$，
其中$U=X\setminus S$。设围绕参数$\epsilon=0$只有一处失去横截性，
发生在一个光滑极点分量$S$与一个光滑边界分量$B_v$之间；
其它边界均不经过此点，其余已紧化、解消的对在参数盘上局部平凡。
假设该点的横截局部模型为
\begin{equation}\label{eq:boundarypinchmodel}
 S=\{w=0\},\qquad
 B_v=\{w=z_1^2+z_2^2+\epsilon\}.
\end{equation}
非零单位与参数的可逆解析换元不改变下述结论。

\begin{proposition}[边界夹挤的秩一性]\label{prop:boundaryrankone}
局部闭路对$H_3(U,B^\circ;\Q)$的变分满足$\rk(T-1)\leq1$。
若它在边界曲线的$H^1$组成上是非零Picard--Lefschetz变分，
则总秩恰为一。
\end{proposition}
\begin{proof}
选适配Milnor邻域$D$及固定外圈。
对的单值可取为$D$外恒等，沿外圈也恒等；
把任意相对链细分后，其差分支撑于$D$，
故全局变分像因子化经过
$H_3(U_D,B_D^\circ;\Q)$的全局像。
这一局部化是实际链的切除，不是只对权重分级作估计。

由\eqref{eq:boundarypinchmodel}，$U_D$收缩至一个圆，
所以$H^2(U_D)=H^3(U_D)=0$。
$B_D^\circ$则为复二维Milnor邻域去掉
$C_\epsilon=\{z_1^2+z_2^2+\epsilon=0\}$。
后者的Milnor纤维收缩至一个圆。局部Gysin列与相对正合列依次给出
\[
 H^2(B_D^\circ)\simeq H^1(C_\epsilon)(-1),\qquad
 H^3(U_D,B_D^\circ)\simeq H^2(B_D^\circ).
\]
这里只需两个空间均为一维，不把局部极限的混合Hodge结构
误认为光滑全局纤维的权重分解。
由有理系数的对偶及切除，全局变分秩至多一。
若边界曲线的全局组成已有非零变分，则总变分不为零，秩即为一。
在本节使用的普通椭圆节点情形，分级上的变分平方为零；
总像的一维线在该分级中投影非零，故它与更低过滤的交为零。
于是$N$再次作用于像线必为零，即总算子也满足$N^2=0$。
\end{proof}

局部管、变分与交数的标准背景见\cite[\S2.3]{muhl}。
这里另给了带边界的相对链论证；没有用单圈物理定理替代本两圈问题。
在同调方向，上一节的$V_a=H^1(E_a)(-1)$对应其对偶，
位于权重$-3$；下文陈述几何组成时仍用上同调记号$V_a,V_b$。

\subsection{实际正链的非零阈值首项}
仍采用上一节的四边图、正质量$m,m,a,b$和每圈$d^2k/\pi$归一化。
取正射影链的$x+y=1$截面，写$y=1-x$、$z=x(1-x)$。
在欧氏区域，齐次换元给
\begin{equation}\label{eq:scalaricechart}
 I=\int_0^1dx\int_0^\infty du\int_0^\infty dv\,
       \frac{z+u+v}{F^2}.
\end{equation}
沿已固定的正链取向，$dx\,du\,dv$为正测度。
此式直接来自齐次参数形式，Gamma因子为$\Gamma(2)=1$。
多项式准确分解为
\begin{align}
 F&=F_a+vB(x,u)+b^2v^2,\label{eq:scalariceexpand}\\
 F_a&=(m^2+a^2u)(z+u)-s_Azu,\nonumber\\
 B(x,u)&=m^2+(b^2-s_B)z+(a^2+b^2-s_C)u.\nonumber
\end{align}
固定$s_B,s_C<0$，令$h_a=a+2m$、$\tau_a=h_a^2$。
$F_a$在$s_A=\tau_a$的正域中恰有一个零点：
\[
 x_0=\tfrac12,\qquad u_0=\frac{m}{2a},\qquad v_0=0.
\]
例如在$z>0,u>0$时，
$F_a/(zu)=m^2/u+m^2/z+a^2+a^2u/z-s_A$；
先对$u$、再对$z\leq1/4$取最小值，即得$\tau_a$及唯一等号点。
又由$B>0$与$b^2>0$，全正链上没有其它同时达到此阈值的极点。

设$\epsilon=\tau_a-s_A$、$\xi=x-\tfrac12$、$\eta=u-u_0$。
在$v=0$处有精确展开
\begin{equation}\label{eq:scalaricehessian}
 F_a=a^2\eta^2+\frac{m^2h_a}{a}\xi^2
       +4m(m+a)\eta\xi^2+\epsilon zu.
\end{equation}
记
\begin{align}
 U_a&=\frac{h_a}{4a},&
 c_a&=\frac{m}{8a},&
 A_a&=\frac{m^2h_a}{a},\nonumber\\
 B_a&=m^2+\frac{b^2-s_B}{4}
             +\frac{m}{2a}(a^2+b^2-s_C)>0.\label{eq:scalariceconstants}
\end{align}
由于$\partial_vF=B_a\ne0$，极点曲面在夹挤点本身仍光滑。
非退化的正切Hessian和参数横截项$c_a\epsilon$，
由带参数的解析Morse坐标化给出\eqref{eq:boundarypinchmodel}。

\begin{proposition}[不求完整积分的非零跳跃]\label{prop:scalaricelog}
在上述归一化及从正欧氏链续接的分支上，
\begin{equation}\label{eq:scalaricelog}
 I=-C_a\log\epsilon+O(1),\qquad
 C_a=\frac{\pi U_a}{B_a\sqrt{A_aa^2}}
     =\frac{\pi h_a}{4a B_a m\sqrt{a h_a}}>0
 \quad(\epsilon\longrightarrow0^+).
\end{equation}
局部奇异部分可写成$H_a(\epsilon)\log\epsilon$，
其中$H_a$解析且$H_a(0)=-C_a$。
因此$\epsilon$逆时针一周的实际跳跃非零，极限为$-2\pi iC_a$。
交换$(a,s_A,u)$与$(b,s_B,v)$得到另一正常阈值处的非零跳跃。
\end{proposition}
\begin{proof}
先限制到夹挤点的一个小正半邻域。
利用$\partial_vF\ne0$，以$w=F$作法向坐标，
其下端为$w=F_a$。写解析密度
$J(x,u,w,\epsilon)=(U/\partial_vF)(x,u,v(w),\epsilon)$。
对$J/w^2$作法向积分，常数项产生
$J(x,u,0,\epsilon)/F_a$；
其余项至多产生$\log F_a$或解析项。
在二维切向中，$\log F_a$的积分只有
$\epsilon\log\epsilon$等次领先项，因此不改变首项系数。
同理，$J(x_0,u_0,0,0)=U_a/B_a$以外的密度项不改变该系数。

现在使用\eqref{eq:scalaricehessian}的解析Morse坐标，
切向Jacobian的常数为$1/\sqrt{A_aa^2}$，
临界值为$c_a\epsilon+O(\epsilon^2)$。
径向积分
\[
 \int_{X^2+Y^2<r^2}
       \frac{dX\,dY}{c_a\epsilon+X^2+Y^2}
       =-\pi\log\epsilon+O(1)
\]
即给\eqref{eq:scalaricelog}。
解析密度逐次展开后，角向奇项积分为零，
偶项径向积分为解析项加$\epsilon^j\log\epsilon$；
法向的次领先项具有相同性质。
缩小邻域即可使展开一致收敛，故对数系数是解析函数。
这一步也证明了非零单值，而非仅从实轴上的$O(1)$记号猜测分支。

其余链在先前构造的紧化、吹起模型中不接触极点，
在足够小的参数盘上给解析函数；汇合到无穷远的坐标端
由该模型的收敛性与横截性一并控制。
所以局部首项没有被其它链片抵消。
\end{proof}

式\eqref{eq:scalaricelog}的系数只需要一个局部Hessian和一个旁观传播子的
非零法向系数。它是实际函数的结构证据，不是数值拟合。
本节固定了参数闭路的方向；没有将此符号直接称为所有物理通道的
Cutkosky约定，后者仍须相应Minkowski链比较。

\subsection{单标量积分的精确椭圆组成}
取上一节已接受的连通光滑参数开集。沿两个正常阈值的一般点，
要求解消后的对除了所述端面夹挤外无其它同时退化。
该条件是几何适用条件，不能省去。
在本族中可满足：固定$s_A=\tau_a$后，
$s_B,s_C$仍提供线性系，Bertini排除基集外的新奇点；
在$v=0$的固定端面节点处，
$\partial_vF=B_a$对负的$s_B,s_C$严格非零。
旧基线及坐标顶点的解消照上一节进行，
再排除有限个代数坏条件即可。另一个端面同理。

\begin{theorem}[非零内部注入图的单标量最小组成]\label{thm:scalaricevisibility}
对于此参数族及其正物理链芽，完整标量积分$I$的最小变分对象，
以上同调方向记录时，其非Artin--Tate半单组成恰为
\begin{equation}\label{eq:scalariceexact}
 \bigl(M_{\min}(I)_{\mathrm{nonAT}}\bigr)^{\mathrm{ss}}
       =H^1(E_a(s_A))(-1)\oplus H^1(E_b(s_B))(-1),
\end{equation}
各出现一次，非AT半单秩为四。
这是单个$I$的结论，无需同时加入$I[D_v]$和$I[D_u]$。
\end{theorem}
\begin{proof}
上一节的相对权重列给出载体上界：非AT组成只有$V_a,V_b$，
重数各一。两条sunrise椭圆族均非等常，
其$H^1$几何局部系统绝对不可约，已由前节的椭圆单值论证说明。
沿$s_A=\tau_a$的一般点，$V_a$有普通单节点退化的非零变分，
而$V_b$横截延拓。故两者作为全参数局部系统不相同。

先选$s_B,s_C<0$的一般点，在$s_A$方向绕$\tau_a$。
命题\ref{prop:boundaryrankone}给完整同调载体上的总变分秩一，
它在$V_a^\vee$上非零；命题\ref{prop:scalaricelog}给实际$I$的非零跳跃。
定理\ref{thm:rankonevisibility}因而强制任何最小标架子商保留
$V_a^\vee$。对另一个正常阈值重复论证，得到$V_b^\vee$。
这些基点可在同一光滑参数开集内平行运输到同一芽，
所以两个下界属于同一物理变分对象。
与重数一的上界比较并对偶，即得\eqref{eq:scalariceexact}。
\end{proof}

此处的“秩四”是非AT半单层之秩，不能改写为完整混合对象恰为四维。
Tate层与扩张仍可增加完整最小秩。
同样，本定理不声称任意两个固定代数点的曲线互不同源，
也不以参数变分排除某个孤立数值偶然等于MZV。
与上一节的分子公式合用，得到同一个图中的精确选择表：
标量$I$保留两项；收缩分子$D_v$只保留$V_a$；
$D_u$只保留$V_b$。这给出了实际物理操作对结构的筛选，
而不是仅对载体作标签分类。

\subsection{显式一参数切片的齐次阶数下界}
取$\Lambda>0$、$c<0$，并令
\begin{equation}\label{eq:scalariceray}
 s_A=t,\qquad s_B=-t-\Lambda,\qquad s_C=c.
\end{equation}
在$-\Lambda<t<0$，三个不变量均负，正参数链无极点。
如要求一个实欧氏动量起点，可取$t=-\Lambda/2$、
$-2\Lambda<c<0$；Gram条件允许这样的非共线动量。
其余$t$按该芽解析续接。
选择一般$\Lambda,c$使上述正常阈值处满足单一夹挤条件。
下述有理函数系数域$k(t)$包含固定质量及$\Lambda,c$；
若讨论数域上的算术对象，另取这些常数为代数数。

\begin{corollary}[无需完整积分的微分阶数障碍]\label{cor:scalariceorder}
沿\eqref{eq:scalariceray}，$I(t)$的每个非零齐次有理系数
微分消去算子都至少四阶。有限代数系数扩张后仍有此下界。
\end{corollary}
\begin{proof}
$E_a(t)$的有限退化值属于
\[
 D_a=\{0,a^2,(a-2m)^2,(a+2m)^2\}\subset[0,\infty),
\]
而$E_b(-t-\Lambda)$的有限退化值属于
$-\Lambda-D_b\subset(-\infty,-\Lambda]$，两组不交。
两个仿射拉回椭圆局部系统仍非等常且绝对不可约，
且一个正常阈值处仅其中一个有非零单值。
在$t=\tau_a$和$t=-\Lambda-\tau_b$，
另一通道参数为负，故本节的两个正链跳跃证明均适用。
秩一保护使该一参数标量芽的变分子商也含两个秩二因子。

更明确地，将平坦链$r$映为函数芽$\langle\lambda(t),r(t)\rangle$。
在由$\sigma$的单值轨道生成的有限秩局部系统中，
其恒零函数核是单值不变子系统。
为讨论复函数关系，可将此有限秩局部系统复化；
所用两个椭圆因子绝对不可约，保护论证仍成立。
对这个商可直接用定理\ref{thm:rankonevisibility}，
不需假定所有函数关系都来自动机关系。
若$L$为$r$阶齐次算子，所有续接芽均在其$r$维局部解空间中，
因此$r\geq2+2=4$。
有限代数扩张只将局部幺幂单值换成非零整数次幂，
总秩一与非零跳跃仍保留；椭圆连通几何单值亦不失不可约性。
\end{proof}

本节没有给出四阶上界，因而不能把“至少四阶”写成“恰四阶”。
其作用是用边界类型、普通夹挤与局部正链，
在不计算全函数的前提下把结构表的上界提升为精确简单组成。
加入实际AT背景项后，这些非AT组成还受前文加性比较定理保护。


\section{物理求和的重数与森林背景的阈值判据}\label{sec:physicalsum}
\subsection{从一个图的保护到整个组合的检测}
物理输出通常是多个图、分子及反项的和。同一个椭圆核可在不同图中重复出现；
逐图有非零割，不等于总和保留全部重数。下面将上一节的一维像线论证推广到
同型因子，并明确哪一种跳跃信息足以决定消去。
这里的基本层次是参数局部系统；固定纤维的Hodge/Nori对象另保留其下降条件。

设$E$为连通参数基上的有限秩复局部系统，带单值不变有限过滤$W$，
每个$\operatorname{Gr}^W E$半单。固定一个基点单值$T$，记$N=T-1$。
选定绝对简单局部系统$V$，设$T$在$V$上的变分秩为$d>0$。
以下条件将反复使用：
\begin{equation}\label{eq:strictvariation}
 \begin{gathered}
 \operatorname{Gr}^W N\text{只在权重 }k\text{的同型块 }V^{\oplus m}
       \text{上非零},\\
 \rk N=\rk(\operatorname{Gr}^W N)=md.
 \end{gathered}
\end{equation}
等式要求整个载体没有额外的降权变分；只核对$V$上的Picard--Lefschetz公式不够。
若另一个权重有同一个简单局部系统$V$，$T$在那里也有秩$d$，会违反第一行，
故这里$m$就是$E$中$V$的全部重数。

\begin{theorem}[变分秩在不变子商上给出重数]\label{thm:strictmultiplicity}
在\eqref{eq:strictvariation}下，每个单值不变子商$Q=R/P$均满足
\begin{equation}\label{eq:rankmultiplicity}
 \rk(N|Q)=d[Q:V].
\end{equation}
特别地，$Q$不含$V$当且仅当$T$在$Q$上恒等。
结论包括相同几何核在多个图中重复出现的情形。
\end{theorem}
\begin{proof}
令$\Gamma=\im N$。$N(W_kE)$在$\operatorname{Gr}_k^WE$中的像已经有维数$md$，
等于$\dim\Gamma$，所以
\[
 \Gamma=N(W_kE)\subset W_kE,\qquad
 \Gamma\cap W_{k-1}E=0.
\]
因此$\Gamma$投影同构到$\im\overline N$，其中
$\overline N=\operatorname{Gr}_k^WN$。
对任意单值不变子对象$A\subset E$，$\overline A=\operatorname{Gr}_k^WA$
在半单局部系统$\operatorname{Gr}_k^WE$中有单值不变补空间。于是
\[
 \im\overline N\cap\overline A=\overline N(\overline A).
\]
若$\gamma\in\Gamma\cap A$，其投影可写成$\overline N\bar a$，取
$a\in A\cap W_kE$提升$\bar a$。$\gamma$与$Na$在$\Gamma$中有相同投影，
故$\gamma=Na$。反向包含显然，因而
\[
 NA=\Gamma\cap A,\qquad
 \rk(NA)=\rk(\overline N|\overline A)=d[A:V].
\]
最后$NR\cap P=\Gamma\cap P=NP$，从而
\[
 \rk(N|R/P)=\dim NR-\dim NP
           =d\bigl([R:V]-[P:V]\bigr).
\]
这里使用的是半单分级上的补空间，不是假定原混合局部系统分裂。
\end{proof}

该定理把上一节$m=d=1$的论证提升为可计数版本。
若整个$E$上的秩等式失败，隐藏的对数扩张也会贡献跳跃，
不能再除以$d$解释成几何核重数。反过来，即使某一物理分支恰无跳跃，
也不能据此断言它的全部解析续接分支均无跳跃。

\subsection{有限续接、导数与精确的函数判据}
现把参数基限制为光滑代数曲线$B$，剔除所有连接和读出的奇点。
设$n=\rk E$，$\lambda$为单值的代数de Rham读出，
$\sigma$为基点链，$f(t)=\langle\lambda(t),\sigma(t)\rangle$。
在复解析局部系统范畴中定义
\[
 R_\sigma^{\rm an}
   =\operatorname{span}_{\mathbb C}\{\rho(\gamma)\sigma:\gamma\in\pi_1(B)\},
 \quad
 P_f^{\rm an}
   =\{r\in R_\sigma^{\rm an}:
               \langle\lambda(t),r(t)\rangle\equiv0\},
 \quad Q_f^{\rm an}=R_\sigma^{\rm an}/P_f^{\rm an}.
\]
恒零解析芽的所有续接仍恒零，故$P_f^{\rm an}$不变。
这个商只记录函数芽及其分支，不预设函数关系都来自动机态射。

\begin{proposition}[有限跳跃导数矩阵]\label{prop:jumpjet}
取$\pi_1(B)$的有限对称生成集。所有长度至多$n-1$的词作用于$\sigma$
已张成$R_\sigma^{\rm an}$。从中选出有限生成链$r_1,\ldots,r_\ell$，定义
\[
 h_j(t)=\langle\lambda(t),(T-1)r_j(t)\rangle,\qquad
 \mathcal J_T(f)=\bigl(\partial_t^i h_j(t)\bigr)_{
                              0\leq i\leq n-1,\ 1\leq j\leq\ell}.
\]
这里先在基点作用$T-1$，再平坦延拓；这固定了分支和乘法次序。
在普通点的亚纯芽域$\mathcal M$上有
\begin{equation}\label{eq:jumpjetrank}
 \rk_{\mathcal M}\mathcal J_T(f)=\rk(N|Q_f^{\rm an}).
\end{equation}
\end{proposition}
\begin{proof}
设$R_j$为长度至多$j$的词生成的空间。若$R_j=R_{j+1}$，
每个生成元都保持$R_j$，序列随后稳定。若不稳定，每步维数至少增加一，
从$\dim R_0=1$开始至多$n-1$步到达全轨道空间。
$\sigma=0$时结论平凡。

链到函数的映射在$R_\sigma^{\rm an}$上的核恰为$P_f^{\rm an}$。
因$NP_f^{\rm an}\subset P_f^{\rm an}$，$h_j$的复常数线性包络
同构于$NQ_f^{\rm an}$。设其维数为$r\leq n$。
导数不会破坏常数线性关系，故矩阵秩至多$r$。
反之，取$r$个常数线性独立的亚纯函数$g_1,\ldots,g_r$，
其Wronskian不恒零。对此事实可直接归纳：在$g_1\ne0$处，除以$g_1$，
Wronskian经行列消去等于
\[
 g_1^r W\bigl((g_2/g_1)',\ldots,(g_r/g_1)'\bigr).
\]
右边的$r-1$个函数若有常数线性关系，积分后得到$g_i$的常数线性关系，
与假设矛盾；归纳得非零。亚纯恒等原理使结论不依赖所选开邻域。
前$r$行因而已有非零$r$阶子式，得到反向不等式。
\end{proof}

\begin{corollary}[物理求和后的检测公式]\label{cor:physicalmultiplicity}
在\eqref{eq:strictvariation}下，
\begin{equation}\label{eq:functionmultiplicity}
 [Q_f^{\rm an}:V]=\frac{1}{d}\rk_{\mathcal M}\mathcal J_T(f).
\end{equation}
特别地，某一次实际续接跳跃非零便给至少一份$V$；
所有生成链的跳跃恒零才给函数局部系统中$V$完全消失。
若$M_f$是已实际构造、承载同一函数的Hodge或Nori变分子商，
则在其复局部系统实现中至少有上述数目的$V$。
\end{corollary}
\begin{proof}
对$Q_f^{\rm an}$应用定理\ref{thm:strictmultiplicity}与
命题\ref{prop:jumpjet}。任何承载$f$的实际标架子商，其链轨道到函数芽的映射
仍以这些解析续接为像，因此$Q_f^{\rm an}$是该实现的子商。
Jordan--Hölder重数的单调性给最后一个下界。
\end{proof}

最后一个下界通常不能直接改为固定纤维Nori重数的等号：
解析恒等核是否为实际子对象、以及参数对象如何下降，仍需比较。
当已知几何上界与它相同，则可用上下界夹逼取得相应层次的精确结论，
上一节单积分双椭圆就是这种用法。

本公式是有限证明条件，不是已经解决一般超越函数恒等判定的程序。
证明一个子式在一个普通点非零足以给下界；某点子式为零不证明其恒等为零。
同样，若$h(\varepsilon)=H(\varepsilon)\log\varepsilon+G(\varepsilon)$，
只消去$H(0)$并不消去$H$；$\varepsilon\log\varepsilon$就是反例。
物理分子或图权重是否真正删除几何核，应作用于完整跳跃矩阵，
不能只检查第一个阈值系数。

\subsection{独有退化与有限森林的方案保护}
锚定森林公式已经给出实际差载体$L$：
它由有限个子图参考系数和商图差值的载体取直和、张量、子商得到。
令$\cS$为这些实际载体的张量生成类别再取Serre闭包；
这不是在整个动机范畴中取包含单位的张量理想。
考虑参数空间中一个余维一退化分量$D$及绕它的局部单值$T_D$。

\begin{theorem}[可跨阈值背景的保护判据]\label{thm:forestthreshold}
假设这些背景生成载体的参数局部系统均能跨过$D$延拓；
更一般地，先拉回一个共同的有限分歧覆盖，并在覆盖后的基上陈述下列结论。
若该基上的简单变分因子$V$在$D$处的局部单值非恒等，
则$V$不属于该背景的Serre类别。
因此，在已证明的实际锚定方案比较
$p_b^{\mathfrak m}-p_a^{\mathfrak m}=[L,\tau,\eta]$中，
两边最小对象的参数实现具有相同的$V$重数。
若比较和背景都在指定实际范畴中构造，这也是该范畴Serre商比较的推论。
\end{theorem}
\begin{proof}
延拓使$T_D$在每个背景生成局部系统上恒等，因而在有限直和、
张量、对偶及子商上恒等。这些对象的每个简单组成也如此。
取Serre闭包允许更多扩张，但不会增添新的简单组成，
故其中不可能出现$T_D|V\ne1$的$V$。
实际差载体属于背景，应用定理\ref{thm:addition}的精确商比较，
即得背景外重数保持。也可先取局部系统实现，再在该范畴应用同一个证明。
有限覆盖时以$T_D^e$取代$T_D$，其中$e>0$为分歧指数，
并按覆盖后实际简单组成陈述结论。
\end{proof}

此判据不要求背景全部Tate。背景可以在别处有椭圆退化，
只要在当前$D$平凡而新核在$D$不平凡即可。
对普通椭圆节点，$T_D=1+N$、$N^2=0$且$N\ne0$，
故任何正整数次幂仍非平凡，有限覆盖不能消掉这个检测。

有两处不能省略。首先，只检查背景的某一个物理函数在$D$解析不够；
必须检查承载全部分支的实际局部系统能延拓。
其次，“$T_D$恒等”的对象一般不对扩张封闭：
对数扩张可使总单值幺幂非平凡，其简单组成仍全部局部平凡。
上述证明使用的是组成因子，而没有错误地声称全部Serre背景均无跳跃。
局部多项式反项的系数在外运动学上为常数，
所以其参数局部系统跨阈值延拓；这并不证明那些常数的算术动机为Tate。
若某个收缩图或参考系数在同一个$D$也退化，本判据就不适用，
须回到实际森林比较，不能仅凭圈数减少宣布新核受保护。

\subsection{一个完整交叉对称图组合：三个椭圆通道确实保留}
用上一节的收敛二维四边图，取四条边质量均为$m>0$，
记$I(s_A,s_B;s_C)$为$AB$双边、$AC/CB$各一边的积分。
换元交换$u,v$给$I(s_A,s_B;s_C)=I(s_B,s_A;s_C)$。
构造实际图函数
\begin{equation}\label{eq:cyclicice}
 \mathcal A_{\rm cyc}
 =I(s_A,s_B;s_C)+I(s_B,s_C;s_A)+I(s_C,s_A;s_B).
\end{equation}
它在三个外不变量置换下不变。这里选定的是明确的交叉对称图组合；
没有声称它已经是某个拉氏量的全振幅。
定义$V_A,V_B,V_C$为三个等质量sunrise椭圆局部系统，
分别拉回到$s_A,s_B,s_C$三个独立变量。
接受域是前两节的光滑开集，并在每个正常阈值的一般点排除其它夹挤碰撞。

\begin{proposition}[物理求和后的三通道下界]\label{prop:cyclicice}
在上述接受域，$\mathcal A_{\rm cyc}$的函数局部系统至少含
$V_A,V_B,V_C$各一份；承载它的实际最小参数对象同样有此下界。
由三图直和载体给出的每项上界为两份，因此
\[
 1\leq[Q_{\mathcal A_{\rm cyc}}^{\rm an}:V_i]\leq2\quad(i=A,B,C),
 \qquad
 6\leq\rk\bigl((Q_{\mathcal A_{\rm cyc}}^{\rm an})_{\rm nonAT}^{\rm ss}\bigr)
      \leq12.
\]
同样的组成上下界适用于实际最小参数子商的复局部系统实现。
这是无需计算完整积分的物理组合结论；本命题不决定三个重数究竟为一还是二。
\end{proposition}
\begin{proof}
每个图的非AT载体已由第\ref{sec:boundaryelliptic}节给为两个边界椭圆核各一次。
三图直和因而含$V_A,V_B,V_C$各两次，其余分级为AT。
三个局部系统绝对简单；在$s_A=9m^2$的一般点，$V_A$有普通节点，
而$V_B,V_C$平凡，故它们在三参数基上两两不同构。

固定$s_B,s_C<0$，令$\varepsilon=9m^2-s_A$。
第一、第三项各有上一节证明的一个边界普通夹挤；
第二项在$s_A$方向的一般此点没有夹挤，局部系统延拓。
因而三图直和的整个变分秩为二，恰等于两份$V_A$上的变分秩；
这满足\eqref{eq:strictvariation}，其中同型重数为二、$d=1$。

为说明排除额外夹挤不是空条件，可取$s_B=s_C=-Rm^2$、$R>0$，
并用质量尺度令$m=1$。对第二项，$s_A=9$是旁通道。
置$q=x+y,r=x-y,w=u+v,d=u-v$，其极点三次式为
\[
 4\mathcal F=(q+(1+R)w)(q^2-r^2)+4q^2w-5qw^2+9qd^2.
\]
在$q=1$的图中，用$z=x(1-x)$写为
$\mathcal F=z+w+(1+R)zw-\frac54w^2+\frac94d^2$。
额外驻点必须$d=0$且$x=1/2$或$w=-1/(1+R)$。
驻点处$\mathcal F$分别为
$1/4+(5+R)^2/80$和$-1/(1+R)-5/[4(1+R)^2]$，均不为零。
$q=0,r\ne0$处在曲面上有$w=0$，且
$\partial_w\mathcal F=-(1+R)r^2/4\ne0$。
$q=r=0$仅有前节已消解的两个普通节点$9d^2=5w^2$。
两个椭圆端面在参数$-R$光滑；另外两个端面的剩余二次曲线为
$qw-\frac54w^2+\frac94d^2=0$，非退化。
与坐标层的固定交点按同一消解处理，剩余交叉条件在一般$R$满足。
因此这条旁通道没有强制的$s_A=9$退化；一般点的接受开集非空。

将上一节的正链首项取$a=b=m$，得到
\begin{align*}
 \mathcal A_{\rm cyc}
   &=-D_A\log\varepsilon+O(1),\\
 D_A&=\frac{\sqrt3\,\pi}{m^2}
       \left(\frac1{9m^2-s_B-2s_C}
             +\frac1{9m^2-s_C-2s_B}\right)>0 .
\end{align*}
两个系数相加，未发生首项抵消；其解析对数系数非零，
故逆时针变分极限为$-2\pi iD_A\ne0$。
定理\ref{thm:strictmultiplicity}对函数标架子商给至少一份$V_A$。
同理置换指标给$V_B,V_C$。三图直和给每项至多二，
简单因子秩均为二，得到全部上下界。
\end{proof}

与上一节的单图相比，求和确实产生了同型重数问题，
故本节不以“非零首项”直接宣称取满两份。
若需要精确重数，应在保持这些因子简单且分离的曲线上
应用\eqref{eq:functionmultiplicity}，或在多参数基上用相同的函数包络判据；
曲线限制是否保持简单性和独立性也须证明。
有符号图权重、真实分子和对称性投影可以改变该矩阵的秩。
这给出了总目标中从组合到物理出口的一条严格接口：
先确定实际载体上界，再用完整组合的变分确定下界或重数，
最后由跨阈值的森林背景决定哪些组成对方案变化受保护。
它既不把逐图几何表当作可观测量分类，也不把未计算的矩阵秩宣称为已知。


\section{有限骨架与无限图族：从边表到精确椭圆支撑}
\label{sec:finitecores}
本节将前面的几何条件转成一个可以直接读输入的接受规则。
首先分开两个问题：固定圈数时有多少种约化拓扑，以及这些拓扑的质量与标架
可以产生多少个不同的可见纯核。前者有初等的有限性定理，后者即使两圈也没有
仅依赖圈数的秩上界。随后对三条路径均可任意细分的theta图，
给出不求完整积分的精确支撑公式。图的归约与部分分式本身是经典操作；
这里的责任是证明它们在所选完整相对周期上的合法性，并给出物理非抵消条件。
任意质量sunrise的椭圆几何已有\cite{ABWsun,BKVsun}的研究；
下述多路径论证使用相对形式及正常阈值，不需要其完整特殊函数解。

\subsection{固定圈数与外部标记数的有限骨架}
连通有限多重图允许自环，度数中自环计二。标记所有可能有外部动量注入的顶点，
数目为$n$；即使一个指定标记的净注入偶然为零，也保留该标记。
未标记顶点的动量注入恒等为零。记圈数为$h=E-V+1$。
先反复删去未标记的一价顶点及其边，再压缩未标记的二价顶点，
把两条相邻边的串联路径换成一条边；不得压缩标记顶点。
一个没有标记的纯圈单独保留一个顶点和一个自环。
只剩孤立顶点的情形也单独保留。

\begin{proposition}[有限骨架的顶点与边数界]\label{prop:finitecore}
除上述孤立顶点及无标记纯圈两个特例外，得到的骨架$C$满足
\[
 V(C)\le 2h+2n-2,\qquad E(C)\le 3h+2n-3.
\]
若$C$没有桥，则有更强的界
\[
 V(C)\le 2h+n-2,\qquad E(C)\le 3h+n-3.
\]
所以对固定$h,n$，这样的骨架只有有限种同构类型。
这些是不带边上质量列表、传播子幂次或分子的拓扑类型。
\end{proposition}
\begin{proof}
删叶和二价压缩均不改变$E-V$，故保留$h$；连通性也保持。
非例外终态的未标记顶点度数至少三，标记顶点度数至少一。
握手引理与$E=h+V-1$给
\[
 2h+2V-2=2E\ge3(V-n)+n=3V-2n,
\]
得到第一组界。无桥且不是孤立顶点时，每个顶点度数至少二，
因此标记部分的贡献可换成$2n$，给第二组界。
有限个顶点和有界总边数只有有限种多重邻接数据，自环也受总边数控制。
压缩若产生孤立自环，它所在的连通分量只能是整个图；
无标记时正是已列出的纯圈特例。
\end{proof}

在单位标量被积式中，被删的未标记叶边动量为零，
给一个代数因子$1/\mu_e$；这里假设$\mu_e>0$，且省略整体动量守恒的分布因子。
压缩路径$P$上的每条边动量相同或反向，故其平方都是同一个$Z_P$。
若路径上的质量平方两两不同，则有有理函数恒等式
\begin{equation}\label{eq:corepartial}
 \prod_{j=1}^{r}\frac1{Z_P+\mu_j}
 =\sum_{j=1}^{r}\frac{c_j}{Z_P+\mu_j},\qquad
 c_j=\prod_{\ell\ne j}(\mu_\ell-\mu_j)^{-1}.
\end{equation}
乘回公共分母，在$Z_P=-\mu_j$逐点取值就证明了此式。
质量重复时用Hermite部分分式，得到带幂次的骨架传播子；
并且在可微分积分的域内
\[
 (Z_P+\mu)^{-a}
 =\frac{(-1)^{a-1}}{(a-1)!}\,
       \partial_\mu^{a-1}(Z_P+\mu)^{-1}.
\]

必须保留一个分析与几何关口：式\eqref{eq:corepartial}首先只是被积式等式。
例如二维
$\int d^2k/((k^2+\mu_1)(k^2+\mu_2))$收敛，
右边拆出的两个单传播子tadpole分别发散。
所以有限骨架定理不自动给重整化周期的直和定理，更不自动给任意森林的最小对象。
要把这种归约提升为形式周期，必须另给收敛的相对Stokes关系，
或一个与归约相容的实际减除比较。下面整个theta族满足前一种条件。

\subsection{三条任意长路径的实际相对下降}
令$G_{\mathbf r}$由两个端点$A,B$及三条内部顶点互不相交的$A$到$B$路径构成，
路径长度$r_a\ge1$，$a=1,2,3$。
只有$A,B$承载$p,-p$，内部二价顶点没有外部注入；
压缩这些顶点后恰为三条平行边。记
\[
 N=r_1+r_2+r_3,\quad h=2,\quad
 m_{aj}>0,\quad\mu_{aj}=m_{aj}^2 .
\]
质量是正代数数。取$D=2$、单位传播子、$s=-p_E^2$及每圈$d^2k/\pi$，
三条路径的动量平方为
$Z_1=k^2,\ Z_2=\ell^2,\ Z_3=(p_E-k-\ell)^2$。
欧氏正链定义
\[
 I_{\mathbf r}(s)=
 \int_{\mathbb R^4}\frac{d^2k\,d^2\ell}{\pi^2}
       \prod_{a=1}^3\prod_{j=1}^{r_a}(Z_a+\mu_{aj})^{-1}.
\]
任一圈至少包含两条正质量传播子，整个图至少三条；
相应二维UV次数严格为负，正质量排除IR发散。
逐条取消串联传播子直到每条路径剩一条，所得积分仍收敛。

令$A_a=\sum_j\alpha_{aj}$。生成树和二森林给
\begin{equation}\label{eq:thetaUF}
 U=A_1A_2+A_1A_3+A_2A_3,\qquad
 F=\left(\sum_{a,j}\mu_{aj}\alpha_{aj}\right)U-sA_1A_2A_3 .
\end{equation}
由两个高斯圈积分及尺度积分，完整参数形式为
$\Gamma(N-2)U^{N-3}\Omega_{N-1}/F^{N-2}$。
写$\alpha_{aj}=A_av_{aj}$、$\sum_jv_{aj}=1$、
$\sum_a A_a=1$，置$\bar\mu_a=\sum_j\mu_{aj}v_{aj}$。
在边界吹起后的实际正链上，它变成
\begin{equation}\label{eq:thetaproduct}
 \Gamma(N-2)\,
 \frac{U^{N-3}\prod_a A_a^{r_a-1}}
      {F_{\bar{\boldsymbol\mu}}^{\,N-2}}\,
 \Omega_2\wedge\bigwedge_{a=1}^3 dv_a ,
\end{equation}
其中$dv_a$是$\Delta_{r_a-1}$的定向体积形式，零维时为$1$。
以下形式周期均指这一实际参数模型及其传播子插入的线性参数化。

\begin{proposition}[全部三条路径的完整Stokes下降]\label{prop:thetastokes}
假设每条路径内的$\mu_{aj}$两两不同。对
$t=(i_1,i_2,i_3)$置
\[
 c_t=\prod_{a=1}^3\prod_{j\ne i_a}(\mu_{aj}-\mu_{ai_a})^{-1}.
\]
令$S_t$为三条质量$m_{1i_1},m_{2i_2},m_{3i_3}$的完整二维sunrise相对周期。
则有实际形式周期关系
\begin{equation}\label{eq:thetaformal}
 I_{\mathbf r}^{\mathfrak m}=\sum_{i_1=1}^{r_1}
       \sum_{i_2=1}^{r_2}\sum_{i_3=1}^{r_3}c_tS_t^{\mathfrak m}.
\end{equation}
它保留端点sunrise的边界与混合扩张，不只是割的等式。
\end{proposition}
\begin{proof}
选$r_a\ge2$的一条路径，以$v_{a1},\ldots,v_{a,r_a-1}$为独立坐标。
删去体积形式中的$dv_{a1}$并选取使外微分取正的楔积符号，定义原函数
\[
 \eta=
 \frac{\Gamma(N-3)}{\mu_{a,r_a}-\mu_{a1}}\,
 \frac{U^{N-4}A_a^{r_a-2}\prod_{b\ne a}A_b^{r_b-1}}
      {F_{\bar{\boldsymbol\mu}}^{\,N-3}}\,
 \Omega_2\wedge (dv_1\wedge dv_2\wedge dv_3)_{\widehat{a1}} .
\]
有
$\partial_{v_{a1}}F=A_aU(\mu_{a1}-\mu_{a,r_a})$。
其余纤维微分重复外积，底上的微分与最高次$\Omega_2$重复，故$d\eta$
正好是\eqref{eq:thetaproduct}。
路径$a$的面中仅$v_{a1}=0$和$v_{a,r_a}=0$有非零拉回。
其他路径的整个体积形式在其任何面上拉回为零。
这给递推
\[
 I_{\mathbf r}^{\mathfrak m}
 =\frac{I_{\text{省去路径$a$末质量}}^{\mathfrak m}
          -I_{\text{省去路径$a$首质量}}^{\mathfrak m}}
        {\mu_{a,r_a}-\mu_{a1}} .
\]
这里省去参数是串联边的参数端面，即缩短路径；不是把一条边抽象删除而断开路径。

还须核对底的边界。坐标边的内部$F>0$，$\Omega_2$的限制为零。
在$A_i=1$附近置其余两个$A$之和为$\rho$，
则$U,F$均消失一阶，$\Omega_2$含$\rho\,d\rho$。
若正在缩短路径$i$，$\eta$的径向阶数是
$\sum_{b\ne i}(r_b-1)\ge0$；
若缩短另一路$a$，阶数是
$\sum_{b\ne i}(r_b-1)-1\ge0$，因为$r_a\ge2$。
拉回均含$d\rho$，所以例外边界限制为零。
各$\bar\mu_a>0$使这项检查在闭纤维单纯形上一致；
在复代数侧删去极点真变换并取共同消解，即得到实际相对形式和边界态射。
因此递推属于几何Stokes关系，不能仅用数值等式替代。

迭代递推使三条路径都剩一条，系数由
\eqref{eq:corepartial}的唯一简单部分分式确定，即为$c_t$。
所有出现的参数积分收敛，过程中没有未声明的UV反项。
\end{proof}

\subsection{完全由质量给出的精确接受条件}
每个三元组$t$给椭圆曲线
\[
 E_t(s):\
 (\mu_{1i_1}X+\mu_{2i_2}Y+\mu_{3i_3}Z)(XY+XZ+YZ)-sXYZ=0.
\]
令
\begin{align}
 \mathcal D_t&=\{0,\infty\}\cup
 \{(\epsilon_1m_{1i_1}+\epsilon_2m_{2i_2}
                       +\epsilon_3m_{3i_3})^2:\epsilon_a=\pm1\},\label{eq:thetadiscr}\\
 \tau_t&=(m_{1i_1}+m_{2i_2}+m_{3i_3})^2 .\nonumber
\end{align}
整体翻号不改变平方，因而中间至多四个不同值。
在共同光滑底上取$\mathcal V_t=H^1(E_t)(-1)$。
精确分支的额外接受条件是
\begin{equation}\label{eq:thetaprivate}
             \tau_t\notin\mathcal D_u\quad\text{对所有 }t\ne u .
\end{equation}
这是有限个代数数等式的否定，不需先计算积分或单值矩阵；
它是本规则的充分适用条件，不声称是核独立性的必要条件。

\begin{theorem}[theta细分族的精确椭圆支撑]\label{thm:thetavisible}
在上述正质量、零内部注入、路径内互异及\eqref{eq:thetaprivate}条件下，
完整物理函数的最小变分对象，其非Artin--Tate半单部分在上同调约定下为
\begin{equation}\label{eq:thetass}
       \bigoplus_t\mathcal V_t,\qquad
       \dim(\text{非Artin--Tate半单部分})=2r_1r_2r_3 .
\end{equation}
每项几何绝对不可约，彼此非同构，且恰出现一次。
同调标架约定取对偶。任何在$\mathbb C(s)$上消去该完整函数的
非零齐次标量微分算子，阶数至少为$2r_1r_2r_3$。
\end{theorem}
\begin{proof}
先给无需积分的上界。对每个$E_t$吹起三个坐标顶点，
环境为有理曲面，边界为有理曲线环及有限点。
坐标边上的第三交点只由对应质量比给出，坐标顶点处的切线也与$s$无关。
于是相对边界只有Artin--Tate组成；
Gysin列和相对长正合列与第\ref{sec:serieselliptic}节相同，
使完整sunrise载体的非Artin--Tate组成恰为一次$\mathcal V_t$。
式\eqref{eq:thetaformal}给最小标架对象的直和上界。

式\eqref{eq:thetadiscr}可直接核对。环面坐标
$\alpha_a=e^{y_a}/m_{ai_a}$给
$F/(\alpha_1\alpha_2\alpha_3)=AB-s$，
其中$A=\sum_a m_{ai_a}e^{y_a}$、$B=\sum_a m_{ai_a}e^{-y_a}$。
$s\ne0$时临界方程迫使所有$e^{2y_a}$相等；
模掉共同尺度后留下三个符号，给上述临界值。
$s=0$的约化纤维另行列入。坐标边和吹起顶点的交点不发生新的$s$碰撞。
因此$\mathcal D_t$包含完整相对对的一切有限退化值；
标量PF的表观奇点不在这里冒充几何退化。

在$\tau_t$处全正符号给唯一临界点，正质量使其对数坐标Hessian正定。
三次曲线只有这一个普通节点，因Bézout不能分成直线与二次曲线，
所以是不可约的结点三次曲线。其$j$不恒定，
有限级别覆盖到模曲线的论证遂给$\operatorname{SL}_2$稠密单值，
如命题\ref{prop:seriesirred}的证明，故$\mathcal V_t$绝对不可约。
条件\eqref{eq:thetaprivate}使它在自己的$\tau_t$处非平凡幺幂，
其余$\mathcal V_u$则为恒等，因而各项非同构，有限底覆盖也不消除此区别。

现在从共同欧氏基点取套索。其余完整相对对在$\tau_t$光滑延拓，所以
\[
 (T_t-1)I_{\mathbf r}=c_t(T_t-1)S_t .
\]
正链的正常阈值Morse夹挤给非零消失周期，阈值上方的三体谱密度为正；
此处用的是前文已证明的正质量banana局部管变分。
$c_t\ne0$，故右边非零；其延拓轨道生成整个$\mathcal V_t$。
所有$t$给互不同构的下界，恰与每项一次的上界吻合。
这一参数变分实现的下界也迫使实际变分载体中的对应因子出现，
因为上界已是同一相对参数模型中的实际形式周期表示；
没有用数值周期映射的单射性。

最后，消去$I_{\mathbf r}$的算子也消去每个跳跃及其全部延拓。
这些互不同构的简单二维解局部系统的和为直和，
故其解空间维数至少为$2r_1r_2r_3$，得到阶数下界。
它不是完整混合对象的精确总秩，也不是齐次算子阶数的等号。
\end{proof}

\subsection{分子把结构表变为可选择的支撑表}
令$P(Z_1,Z_2,Z_3)$的系数属于质量所在数域，且
$\deg_{Z_a}P\le r_a-1$。
所选分子允许各路径动量平方的这些多项式，未声称它们都来自同一个
固定可重整化拉氏量。记插入后的完整积分为$I_{\mathbf r,P}$。
定义有限网格读出
\[
 p_t=P(-\mu_{1i_1},-\mu_{2i_2},-\mu_{3i_3}),\qquad
 \mathcal A(P)=\{t:p_t\ne0\}.
\]

\begin{theorem}[物理分子的必要充分支撑判据]\label{thm:thetanumerator}
在定理\ref{thm:thetavisible}的接受域内，
\begin{equation}\label{eq:thetapinsert}
 I_{\mathbf r,P}^{\mathfrak m}=\sum_t p_tc_tS_t^{\mathfrak m},
 \qquad
 \operatorname{ss}_{\mathrm{nonAT}} I_{\mathbf r,P}
       =\bigoplus_{t\in\mathcal A(P)}\mathcal V_t .
\end{equation}
因此$\mathcal V_t$可见当且仅当$p_t\ne0$。
每个三元组子集都可由上述次数范围内的一个显式分子实现。
\end{theorem}
\begin{proof}
矩形次数范围的多变量插值是三个一变量插值的张量积，故有多项式恒等式
\[
 P(\mathbf Z)=\sum_t p_tc_t
       \prod_{a=1}^3\prod_{j\ne i_a}(Z_a+\mu_{aj}).
\]
乘以原传播子乘积，各项分子直接取消串联传播子，留下$c_tp_tS_t$。
这在参数化前是同一有理形式的有限线性等式。
用传播子消去后的线性Schwinger参数化定义插入的实际参数周期，
所有剩余积分仍绝对收敛，端点与命题\ref{prop:thetastokes}一致。
所以该形式关系只用代数线性、实际消去及相对Stokes，
不需要断言任意动量动机与参数动机整体同构。

$p_t=0$时由直和上界排除该因子；$p_t\ne0$时自己的正常阈值
仅有$c_tp_tS_t$贡献，前一定理的非零割论证给整个$\mathcal V_t$。
令
\[
 L_{ai}(Z)=
 \prod_{j\ne i}\frac{Z+\mu_{aj}}{\mu_{aj}-\mu_{ai}} ,
 \qquad
 P_{\mathcal A}(\mathbf Z)=
       \sum_{(i_1,i_2,i_3)\in\mathcal A}\prod_{a=1}^3L_{ai_a}(Z_a).
\]
各$L_{ai}(-\mu_{aj})=\delta_{ij}$，所以$P_{\mathcal A}$在所选网格上取$1$，
其余取$0$，且满足次数限制，证明最后一句。
\end{proof}

一般分子可以有负系数；这里的非抵消来自特有阈值，不依赖总分子正性。
若允许更高的分次次数，部分分式将出现额外接触项或发散积分，
本定理不再原样适用。若质量合流，则合流导数仍属于对应质量族的
Gauss--Manin载体，但不能把发散的$c_t$分别取极限，也不能保留
\eqref{eq:thetass}的泛型秩等号。

\subsection{无需求值的无限实例与圈数分类的限制}
\begin{corollary}[固定两圈下无统一的可见秩上界]\label{cor:thetaunbounded}
任给$r_1,r_2,r_3\ge1$，给全部$N$条边分配互不相同的非负整数$e_{aj}$，
并令
\[
                    m_{aj}=M\,5^{e_{aj}},\qquad M>0
\]
为正代数尺度。则\eqref{eq:thetaprivate}必成立。
在单位分子下，非Artin--Tate半单秩恰为$2r_1r_2r_3$。
因而即使固定$h=2$和两个外部标记点，
允许任意长的零注入细分及任意多个边质量时，
可见非Artin--Tate秩与齐次微分阶数都没有统一有限上界。
\end{corollary}
\begin{proof}
若一个三元组的正质量和平方等于另一三元组的带符号质量和平方，
开平方后再选整体符号，得到
$\sum_e d_e5^e=0$，其中$d_e\in\{-2,-1,0,1,2\}$。
设最大非零指数为$L$，则
\[
 |d_L|5^L\ge5^L>
       2\sum_{e=0}^{L-1}5^e=(5^L-1)/2,
\]
矛盾。因此全部$d_e=0$。由于边指数互不相同，
只有两个三元组使用完全相同的三条边且符号一致才可能等号；
不同三元组被排除。正和亦不为零，故接受条件成立。
路径内质量互异显然成立，应用定理\ref{thm:thetavisible}即得秩公式。
取$r_1=r_2=r_3=r\to\infty$可得$2r^3\to\infty$，
每个图仍恰为两个圈、两个外部标记点和同一个theta骨架。
\end{proof}

这不与固定质量参数数量等附加类型约束下的周期有限性冲突：
本例允许质量列表长度随图增长。它也不宣称出现任意属的曲线；
增长的是互不同构的属一核的数目。
二价质量变化顶点及所选多项式分子属于带物理数据的图输入，
不自动属于一个固定有限场内容拉氏量的全部允许顶点。

本节因而给出一个完整、可拒绝不合条件输入的理论规则：
先由边表和外部标记识别三条路径，核对零注入；
再由正代数质量生成三元组、曲线和有限退化集合；
核对路径内互异及正常阈值分离；
最后由$P$的有限网格值输出精确的可见椭圆支撑、重数和微分阶数下界。
若$R=r_1r_2r_3$，识别路径后至多需$4R(R-1)$次带符号平方值比较，
以及$R$次分子求值；在给定代数数的标准精确表示下这些是可判定操作。
这只是代数操作计数，未把代数数次数与高度的位复杂度隐藏成常数。
计算只涉及有限图操作及代数数比较，无需完整积分、区间运输或数值单值识别。
全部步骤的正确性由上面的几何下降与物理阈值证明承担。

若分离条件失败，输出仍合法的直和上界以及具体碰撞条件，
不得据此宣称椭圆因子不可见；若有非零内部注入，则连
\eqref{eq:thetaUF}也改变，须转用第\ref{sec:boundaryelliptic}节等实际几何。
本节完成无限接受族的纯组成与分子支撑，未分类其所有混合扩张和固定点算术。
有限骨架为一般图提供组织方式，而这些骨架的推前、边界粘合及标架生成性
仍须逐类证明。总目标中的收缩、森林、RG和固定代数点比较不能从骨架有限性跳过。


\section{共同常数扩张与完整函数的精确微分分类}
\label{sec:commonunit}
上一节确定了完整形式周期中的椭圆组成及分子选择。
还需要回答一个不同的问题：选中的纯核怎样接到非齐次源，
完整物理函数的最小齐次微分阶数究竟是多少？
本节在微分域$K=\mathbb C(s)$上回答这一问题。
复常数包括固定质量的对数；这一基域约定使函数分类有精确答案，
但不会抹去原Nori对象中必须保留的算术边界。

\subsection{共享商对象，不能逐项相加边界维数}
先给不依赖费曼积分的范畴命题。
在特征零的中性Tannaka范畴中，取$r\ge1$，设
$A=\bigoplus_{i=1}^r A_i$，其中$A_i$为两两不同构的简单对象，
且均不同构于单位对象$\mathbf1$。
设实际扩张为
\begin{equation}\label{eq:commonextension}
0\longrightarrow A\longrightarrow E
 \overset{\pi}{\longrightarrow}\mathbf1\longrightarrow0,
\qquad
\xi=(\xi_1,\ldots,\xi_r),\quad
\xi_i\in\operatorname{Ext}^1(\mathbf1,A_i).
\end{equation}
有限直和给出这里的扩张群分解。标架仍取同调方向。

\begin{theorem}[共同单位扩张的最小化]\label{thm:commonunit}
若全部$\xi_i\ne0$，则$E$的每个真子对象都包含在$A$中，
且$\operatorname{soc}(E)=A$。
任一满足$\pi(\sigma)\ne0$的纤维向量$\sigma$生成$E$。
设读出$\lambda$在恰好$J\subset\{1,\ldots,r\}$这些$A_i$上非零。
当$J\ne\varnothing$时，
\begin{equation}\label{eq:commonminimal}
 M(E,\sigma,\lambda)
 =E\big/\bigoplus_{i\notin J}A_i,\qquad
 \rk M=1+\sum_{i\in J}\rk A_i .
\end{equation}
这决定实际扩张子商，强于只计算半单化。
\end{theorem}
\begin{proof}
若子对象$R\subset E$在$\mathbf1$中的像非零，该像就是$\mathbf1$。
令$D=R\cap A$。由$A$半单且无重数，$D$是某些$A_i$的直和。
在$E/D$中，$R/D\simeq\mathbf1$给出单位商的截面，
所以推前扩张在$A/D$中的每个分量均为零。
全部$\xi_i\ne0$迫使$D=A$，从而$R=E$。
这证明真子对象的断言，也证明$\sigma$的生成性。
任何不在$A$中的简单子对象都会给出上述截面，故不存在；
$A$本身半单，所以它就是socle。

读出非零时，能被它消去的子对象必为真子对象，
于是包含于$A$，且仍是某些$A_i$的直和。
最大的一个恰为$\bigoplus_{i\notin J}A_i$，得到式
\eqref{eq:commonminimal}。若$J=\varnothing$，
则非零的$\lambda$经单位商因子化，最小对象为$\mathbf1$；
若$\lambda=0$，最小对象为零。
\end{proof}

给定$r$个实际扩张$E_i\to\mathbf1$，共享商的正确装配是
\[
 E=E_1\times_{\mathbf1}\cdots\times_{\mathbf1}E_r
 =(\bigoplus_iE_i)\times_{\mathbf1^r}\mathbf1 ,
\]
最后一箭头为对角嵌入。它的秩是$1+\sum_i\rk A_i$，
扩张类正是$(\xi_i)_i$。
这里要求各个商与\emph{同一个}单位对象的实际识别。
不能只看见$r$个非齐次方程，就在原动机范畴中宣布其边界相同。
下面是在$\mathbb C(s)$上通过明确微分系统实现这一识别。

\subsection{sunrise的有理源与实际微分扩张}
固定正质量$m_1,m_2,m_3$，记$\mu_i=m_i^2$。
采用此前不附加任意质量尺度的归一化，
\[
 S(s)=\int_{\sigma}
 \frac{\Omega_2}{F_s},\qquad
 F_s=M U-sxyz,\quad
 M=\mu_1x+\mu_2y+\mu_3z,\quad U=xy+xz+yz .
\]
椭圆留数的二阶Picard--Fuchs算子写成
$L=\partial_s^2+a(s)\partial_s+b(s)$。
经典公式\cite[\S4,式(37)--(43)]{ABWsun}给
\begin{equation}\label{eq:suninhom}
 LS=f(s),\qquad a,b,f\in\mathbb C(s).
\end{equation}
质量固定后源中的质量对数为复常数。原文的任意尺度因子是常数，
除去它不改变本节结论。该文已经给出齐次三阶上界；
下文补的是不分裂证明、最小性和多核装配，而非重新声称发现该上界。

\begin{proposition}[有理源的共同单位实现]\label{prop:sunextension}
式\eqref{eq:suninhom}由以下秩三微分系统实现：
\begin{equation}\label{eq:sunmatrix}
 \partial_s
 \begin{pmatrix}y\\v\\c\end{pmatrix}
 =
 \begin{pmatrix}0&1&0\\-b&-a&f\\0&0&0\end{pmatrix}
 \begin{pmatrix}y\\v\\c\end{pmatrix}.
\end{equation}
它是纯椭圆系统$V$对单位商的扩张$E$；
物理解为$(S,S',1)$。该扩张在$\mathbb C(s)$上分裂，
当且仅当存在$r(s)\in\mathbb C(s)$满足$Lr=f$。
\end{proposition}
\begin{proof}
$c=0$的子系统就是$Ly=0$，为椭圆周期系统；
商方程为$c'=0$，故为单位对象。
单位商的水平截面必须形如$(r,r',1)$，
且第二行要求$r''+ar'+br=f$，这也充分。
固定质量下，原源为有理函数与固定边界周期常数的有限和，
所以系统定义在$K$上。它可视为实际相对Gauss--Manin系统
在固定边界周期读出后的子商；也可直接按式\eqref{eq:sunmatrix}构造。
前一种描述不把边界的算术周期等同于有理数。
\end{proof}

该系统正则奇异。纯椭圆Gauss--Manin系统正则奇异，
单位系统亦然，而正则奇异微分模对扩张封闭。
在局部变参数积分中，有理源与纯解至多有幂次及对数增长，
亦直接给出这一性质。标量表达引入的表观奇点不改变内在局部系统。

\subsection{无穷远比较给出不分裂性}
\begin{theorem}[正质量sunrise的扩张不分裂]\label{thm:sunnonsplit}
对任意固定正质量，式\eqref{eq:sunmatrix}的扩张类
$\xi\in\operatorname{Ext}^1_K(\mathbf1,V)$非零。
\end{theorem}
\begin{proof}
\textbf{第一步：完整物理积分有平方对数下界。}
取$s=-t$、$t>0$，在$x+y+z=1$的正单纯形积分。
令$a_0=\max_i\mu_i$、$\delta=1/4$。
$t$充分大时，矩形
$a_0/t\le y,z\le\delta$包含于正单纯形，且$x=1-y-z$。
在该矩形上
\[
 M\le a_0,\qquad U\le y+z+yz,\qquad xyz\le yz ,
\]
从而
\[
 F_{-t}\le a_0(y+z+yz)+tyz\le4tyz.
\]
最后一步分别使用$a_0/(ty)\le1$、$a_0/(tz)\le1$及$a_0/t\le1$。
正积分因此满足严格渐近下界
\begin{equation}\label{eq:sunloglower}
 S(-t)\ \ge\ \frac1{4t}
       \log^2\!\left(\frac{\delta t}{a_0}\right).
\end{equation}
另一方面$1/F_{-t}\le1/(MU)$。
右端在单纯形可积：角点附近$U$与两个小坐标之和同阶，
$M$有正下界；其它边界没有更坏奇性。
内部被积函数逐点趋零，控制收敛定理给$S(-t)\to0$。

\textbf{第二步：纯椭圆解只有一次对数。}
这里不需求完整椭圆积分表达式。
\cite[式(38)--(39)]{ABWsun}的多项式首项为
\[
 p_0(s)=3s^7+O(s^6),\quad
 p_1(s)=9s^6+O(s^5),\quad
 p_2(s)=3s^5+O(s^4).
\]
因此$a(s)=3/s+O(s^{-2})$、$b(s)=1/s^2+O(s^{-3})$。
令$u=1/s$，齐次方程成为
\[
 y_{uu}+\left(-u^{-1}+O(1)\right)y_u
       +\left(u^{-2}+O(u^{-1})\right)y=0 .
\]
其正则奇点指标多项式是$(\rho-1)^2$。
Frobenius理论给任一齐次解在固定扇区内
$y=u(h_0(u)+h_1(u)\log u)$，其中$h_0,h_1$全纯。
故在负实方向
\begin{equation}\label{eq:sunpureinfty}
 y(-t)=O(\log t/t).
\end{equation}
以上首项不依赖伪阈值是否重合；质量相等也不改变这一估计。
约去多项式共同因子不改变$a,b$及无穷远指标。

\textbf{第三步：排除有理截面。}
若扩张分裂，命题\ref{prop:sunextension}给有理特解$r$，
且$S-r$为纯齐次解。由$S\to0$和式\eqref{eq:sunpureinfty}，
$r(-t)\to0$。有理函数于是满足$r(-t)=O(1/t)$。
这将推出$S(-t)=O(\log t/t)$，与式\eqref{eq:sunloglower}矛盾。
故扩张不分裂；特别地$f$不恒为零。
\end{proof}

这个论证只使用一个正区域下界及二阶算子的首项。
它不声称已求出完整高能展开，也没有将高精度数值当成非零证明。
正常阈值为节点且$j$函数非恒定，已证明$V$为秩二不可约系统；
故定理\ref{thm:commonunit}还给单个$S$的最小齐次阶数恰为三。

\subsection{分子选择后的精确阶数}
沿用定理\ref{thm:thetanumerator}的所有条件：
二维、三条无内部外动量注入的串联路径、正代数质量、
每条路径内质量平方互异、不同三元组的正常阈值满足
式\eqref{eq:thetaprivate}，以及
$\deg_{Z_a}P\le r_a-1$。
令
\[
 d_t=c_tP(-\mu_{1i_1},-\mu_{2i_2},-\mu_{3i_3}),\quad
 J_P=\{t:d_t\ne0\},\quad r=|J_P|.
\]
完整Stokes下降已给$I_P=\sum_{t\in J_P}d_tS_t$，
不是只有最大割的等式。

\begin{theorem}[接受theta族的完整函数分类]\label{thm:thetaexactorder}
若$r>0$，则$I_P$在$\mathbb C(s)$上的最小齐次线性微分阶数为
\begin{equation}\label{eq:thetaexactorder}
 \boxed{\operatorname{ord}_{\min,\mathbb C(s)}(I_P)=2r+1.}
\end{equation}
其最小微分对象为实际共同单位扩张
\begin{equation}\label{eq:thetamixed}
 0\longrightarrow\bigoplus_{t\in J_P}V_t
 \longrightarrow E_{J_P}\longrightarrow\mathbf1\longrightarrow0 ,
\end{equation}
扩张的每个分量均不分裂。
若$r=0$，$I_P=0$，按零函数的单位算子约定最小阶数为零。
\end{theorem}
\begin{proof}
对每个$t$使用式\eqref{eq:sunmatrix}，将末坐标识别为同一个$c$，
得到明确的$(2r+1)$维系统
\[
 y_t'=v_t,\qquad
 v_t'=-b_t y_t-a_t v_t+f_tc,\qquad c'=0.
\]
这是各$E_t\to\mathbf1$的纤维积。
物理解向量为$(S_t,S_t')_{t\in J_P}\oplus1$；
标量读出为$\lambda=\sum_td_ty_t$。
特有阈值条件已证明$V_t$两两不同构且不可约；
定理\ref{thm:sunnonsplit}给每个扩张分量非零。
在每个$V_t$上，$y_t$读出非零，且$d_t\ne0$。
故定理\ref{thm:commonunit}给完整最小对象就是$E_{J_P}$。

为明确连接到\emph{普通函数}的最小方程，可在Picard--Vessiot域中工作。
物理解向量在单位商的像为$1$，因此它的微分Galois轨道生成全部解空间。
把解向量送到$\lambda$读出的函数，其核是不变子空间。
核与每个$V_t$相交为零；共同扩张定理迫使核为零。
于是$I_P$的Galois轨道张成一个维数$2r+1$的函数空间。
任意湮灭$I_P$的$K$系数算子也湮灭全部轨道，
其阶数至少$2r+1$。反之，对一个$2r+1$维有理系统逐次求导，
$2r+2$个读出行向量线性相关，给至多该阶的齐次算子。
上下界相等，不需要假设动机周期映射单射。
\end{proof}

纯核秩为$2r$，完整函数阶数为$2r+1$：
额外的一维是共同常数商，不能分别计入$r$次。
在原部分分式求和中，各项的最高平方对数系数可能抵消，
但这不使式\eqref{eq:thetamixed}分裂。
分裂与否由每个互不同构简单分量的扩张类决定，
不是由总函数某一个渐近系数决定。

\subsection{完整微分Galois群}
\begin{theorem}[接受theta族的微分群]\label{thm:thetadgal}
上述$r>0$的最小系统在$\mathbb C(s)$上的微分Galois群，
在适当解空间基下为
\begin{equation}\label{eq:thetadgal}
 G_{\partial}=
 \mathbb G_a^{\,2r}\rtimes\prod_{t\in J_P}\mathrm{SL}_2(\mathbb C),
 \qquad \dim G_{\partial}=5r.
\end{equation}
每个$\mathrm{SL}_2$在对应的二维加法分量上作标准作用。
该群连通，幺幂根基为整个$\mathbb G_a^{2r}$。
\end{theorem}
\begin{proof}
先证明纯群是直积，而不从各投影满射直接跳到独立性。
每个非等常椭圆系统的单值闭包为$\mathrm{SL}_2$。
在正常阈值$\tau_t$处，$V_t$有非平凡节点幺幂变换，
其它$V_u$均光滑，单值为恒等。
设纯系统的总单值闭包为$H\subset\prod_t\mathrm{SL}_2$。
则$H\cap\mathrm{SL}_{2,t}$包含该非平凡幺幂元，
并在$H$对第$t$因子的满投影下为正规子群。
$\mathrm{SL}_2$的含非中心幺幂元的闭正规子群是整个$\mathrm{SL}_2$，
所以$H$包含每个因子，因而等于直积$R$。
系统正则奇异，Schlesinger稠密定理
\cite[Proposition1.3.27]{Singer}使这也是纯微分Galois群。

$E_{J_P}$的群$G$保持纯子空间并在单位商上恒等，
故嵌于$A_{\mathrm{add}}\rtimes R$，其中
$A=\bigoplus_t\mathbb C^2_t$，且$G\to R$满射。
核$U$是加法群$A_{\mathrm{add}}$的闭子群；
特征零下它是$R$不变线性子空间。
若$U$投到某个$\mathbb C_t^2$为零，
则该三维扩张表示因子化经约化群$R$。
特征零下$R$表示完全可约，迫使$E_t$分裂，矛盾。
各$\mathbb C^2_t$是两两不同构的简单$R$模，
故所有投影非零迫使$U=A$。
包含全部核且满射到商的$G$就是整个半直积。
这也给连通性及维数$2r+3r=5r$。
\end{proof}

在该表示中，矩阵与余乘法可以写成
\[
 g=\begin{pmatrix}\operatorname{diag}(H_t)&(x_t)_t\\0&1\end{pmatrix},
 \quad \det H_t=1,\qquad
 \Delta H_t=H_t\mathbin{\dot\otimes}H_t,\quad
 \Delta x_t=H_t\mathbin{\dot\otimes}x_t+x_t\otimes1 .
\]
这给完整\emph{微分}群的函数关系和扩张作用。
它不是本项目固定代数纤维的绝对Nori群或绝对动机coaction公式。

\subsection{对总分类目标的具体增益与剩余边界}
本节给一个无需完整积分的闭合推导链：
\[
 \text{图与质量}\ \longrightarrow\
 \text{正常阈值检验}\ \longrightarrow\
 \text{分子网格 }J_P\ \longrightarrow\
 E_{J_P},\ 2|J_P|+1,\ G_{\partial}.
\]
这同时包含纯组成、不可分裂混合、物理读出和精确函数阶数，
比只列“有$r$个椭圆核”更接近完整分类。
若质量导致阈值重合，或分子超出接受次数，须回到实际比较，
本节不输出伪精确答案。

本节的$K=\mathbb C(s)$会把固定质量的所有周期常数作为系数。
例如非零常数函数$\log2$的最小微分阶数是一，
而原算术范畴中质量单位$2$的Kummer扩张并不因此分裂。
所以式\eqref{eq:thetaexactorder}不能宣布原Nori最小秩就是$2r+1$，
也不能自动下降到只含代数常数的系数域。
对于算术分类，边界常数的独立性、Kummer粘合、实际相对链和比较仍须保留。
已有Nori非Artin--Tate重数结论维持原范围；本节新增的是完整函数实现中的精确扩张。
同样，这一无UV发散的二维族不能替代森林/RG主接口，
但可作为后续物理求和与背景扰动的已完全识别的微分输入。



\section{局部减除、完整混合保护与微分选择}
\label{sec:fullprotection}
完整函数分类还须与物理操作接合。
本节先把背景加法定理提升为完整对象的商保护，再给实际的减除与微分规则。
一般范畴结论用于已经构造的形式周期；
theta应用继续在$\mathbb C(s)$的微分模范畴中，
不把函数系数域中的常数与算术动机的单位对象混同。

\subsection{从简单组成保持到完整对象的商}
设$\cS$为有限长度范畴$\cC$中的Serre子范畴，
$p=[E,\sigma,\lambda]$已最小化，即$E=M(p)$。
称$E$满足\emph{背景socle分离条件}，如果
\begin{equation}\label{eq:socleseparation}
 E\text{没有属于}\cS\text{的非零子对象}.
\end{equation}
这等价于$\operatorname{soc}(E)$没有属于$\cS$的简单因子，
不要求$E$的全部组成都在背景外。
例如，上一节的共同单位扩张虽有单位商，
其socle却全部由非平凡椭圆块组成。

\begin{theorem}[完整混合对象的背景保护]\label{thm:fullbackground}
若式\eqref{eq:socleseparation}成立，且
$\ell=[B,\tau,\eta]$由$B\in\cS$承载，则存在实际短正合列
\begin{equation}\label{eq:fullbackground}
 0\longrightarrow B_0\longrightarrow M(p+\ell)
 \longrightarrow E\longrightarrow0,\qquad B_0\in\cS,
\end{equation}
其中$B_0$是$B$的子商。因此
\[
 \rk E\le\rk M(p+\ell)\le\rk E+\rk B.
\]
在中性Tannaka情形，所生成对象的群有忠实平坦商
$G_{M(p+\ell)}\twoheadrightarrow G_E$。
若$M(p+\ell)$也满足式\eqref{eq:socleseparation}，则$B_0=0$，
从而两个完整对象同构。
\end{theorem}
\begin{proof}
令$R\subset E\oplus B$为$(\sigma,\tau)$生成的子对象，
$I\subset R$为总读出$(\lambda,\eta)$消去的最大子对象，
所以$M(p+\ell)=R/I$。
投影$\rho:R\to E$满射，因为其像承载生成向量$\sigma$。
其核$K_0=R\cap(0\oplus B)$是$B$的子对象。

由$E$的最小性，
$I\cap(E\oplus0)=0$；否则该交是被$\lambda$消去的非零子对象。
于是$I\to B$单射，故$I\in\cS$。
$\rho(I)$既为$E$的子对象，又属于$\cS$，
式\eqref{eq:socleseparation}迫使$\rho(I)=0$。
因此$I\subset K_0$，$\rho$下降为$R/I\to E$，
其核$B_0=K_0/I$给式\eqref{eq:fullbackground}及秩界。

若新对象也无背景子对象，其背景核$B_0$只能为零。
群的结论来自实际商对象$E\in\langle M(p+\ell)\rangle_\otimes$：
相应满Tannaka子范畴包含给忠实平坦群商，
参见\cite[Proposition2.21]{DM}.
\end{proof}

此处保留的是完整$E$，包括其已经证明的非分裂扩张，
而不是只有$\operatorname{Gr}E$。
但式\eqref{eq:fullbackground}一般不分裂，
其投影也不必保持原来两个读出的数值。
对象核$B_0$不应直接认作群商的核。
若socle条件失败，完整保护确实可能失败：
取$E=B=\mathbf1$及$p=1,\ell=-1$，新最小对象为零。
只检查某个高权纯商在背景外，不能取代socle条件。

\subsection{锚定森林与函数背景}
\begin{corollary}[满足socle条件的完整方案保护]\label{cor:fullforest}
在式\eqref{eq:forest}已经构造的相容锚定方案范围内，
若$M(p_a)$没有$\cS_{<G}$中的非零子对象，则
\[
 M(p_b)\twoheadrightarrow M(p_a)
\]
是实际商，核属于$\cS_{<G}$。
若两端都满足该条件，则$M(p_a)\simeq M(p_b)$。
\end{corollary}
\begin{proof}
实际森林式给$p_b=p_a+\ell$，其中$\ell$的载体是
严格低圈背景张量积的有限直和，属于$\cS_{<G}$。
应用定理\ref{thm:fullbackground}。
两端条件成立时，第一条商的背景核是新对象的背景子对象，
所以为零。
\end{proof}

该推论加强定理\ref{thm:ren}，但新增了可核查的socle前提。
它不声称所有重整化对象都满足此前提，也不涵盖尚未构造实际比较的幂次森林。
未锚定整图常数仍可能含同圈数算术对象，须另计入背景。

在$\mathbb C(s)$的正则奇异微分模范畴中，
令$\cS_{\mathrm{fin}}$为简单组成均具有有限单值的Serre子范畴。
有限代数覆盖后幺幂的已知MPL系统属于这一类；
这里要求提供实际系统，不能凭函数名称认定。
上一节的$E_J$有socle $\bigoplus_{t\in J}V_t$，
每个$V_t$的单值闭包是$\mathrm{SL}_2$，故不在$\cS_{\mathrm{fin}}$。
因此对任何由$B\in\cS_{\mathrm{fin}}$承载的函数$b$，
完整函数$I_P+b$的最小微分对象仍满射到$E_J$。
特别地，其最小阶数至少$2|J|+1$，微分群满射到
$\mathbb G_a^{2|J|}\rtimes(\mathrm{SL}_2)^{|J|}$。
加入此类背景不能把共同非分裂扩张也消去。
新函数可以额外需要背景方向，所以这里没有无条件阶数等号。

\subsection{有理微分操作的必要充分选择条件}
令$K=\mathbb C(s)$，固定上一节接受theta族及其非空支撑$J$。
写$I=I_P$，纯椭圆算子为
$L_t=\partial_s^2+a_t\partial_s+b_t$。
对任意$D\in K[\partial_s]$作\emph{右除法}
\begin{equation}\label{eq:oreremainder}
 D=Q_tL_t+R_t,\qquad \operatorname{ord}R_t<2,
\end{equation}
零余式也允许。这里系数在左，乘法满足
$\partial_s a=a\partial_s+a'$；左右顺序不可互换。
经典微分算子除法及最小公左倍数见\cite[\S1.2]{Singer}。
定义
\[
 J_D=\{t\in J:R_t\ne0\}.
\]

\begin{theorem}[微分投影的精确支撑与阶数]\label{thm:differentialselection}
对任意$h\in K$，若$J_D\ne\varnothing$，则
\begin{equation}\label{eq:operatororder}
 M_\partial\!\left(D(I+h)\right)
 \simeq E_J\big/\bigoplus_{t\notin J_D}V_t,\qquad
 \operatorname{ord}_{\min}D(I+h)=2|J_D|+1.
\end{equation}
若$J_D=\varnothing$，则$D(I+h)\in K$，
最小阶数为一或零，分别对应非零函数或零函数。
因此，一个纯核被$D$完全消去当且仅当$D$以其纯算子为右因子：
$D\in K[\partial_s]L_t$。
\end{theorem}
\begin{proof}
用上一节系统中的共同坐标$c$写$I+h$的读出为$\lambda+hc$。
沿连接求导，任意$D$给同一对象$E_J$上的有理读出$\lambda_D$，
物理解向量及其生成性不变。
$c$在纯子系统上为零，因此$h$不影响纯块上的新读出。

在$V_t$上，新读出正是原来非零常系数$d_t$乘以$D y_t$。
若$R_t=0$，它在全部纯解上为零。
反之若在全部纯解上为零，右除法后$R_t$也湮灭两个独立纯解。
非零的零阶或一阶算子不可能具有两个独立解，
所以$R_t=0$。这证明支撑条件。

每个未被消去的$V_t$上读出非零。
定理\ref{thm:commonunit}于是给式\eqref{eq:operatororder}的实际商。
上一节的解空间轨道论证把该最小对象的秩识别为普通函数的最小阶数。
若全部纯块被消去，读出经单位商因子化，
所以在物理解$c=1$上为某个$q(s)\in K$。
非零有理函数的最小齐次算子是一阶
$\partial_s-q'/q$；零函数用零阶单位算子湮灭。
\end{proof}

此处的选择判据只用纯二阶算子，不需先求质量对数的数值或完整积分。
它能处理有理系数的有限参数微分组合。
把这样的组合实现为某个特定拉氏量中的分子或规范恒等式，
仍须提供那个理论的实际算子关系；本节不虚构此类关系。

\subsection{构造任意选择，并读出非齐次源}
\begin{proposition}[纯算子的有限选择规则]\label{prop:oreselectors}
对$K_0\subset J$，令$D_{K_0}$为$\{L_t:t\in K_0\}$的
首一最小公左倍数，并取$D_\varnothing=1$。
则
\[
 \operatorname{ord}D_{K_0}=2|K_0|,\qquad
 J_{D_{K_0}}=J\setminus K_0.
\]
若$K_0\ne J$，它保留的完整函数阶数恰为$2|J\setminus K_0|+1$。
若$K_0=J$，对任意$h\in K$都有
\begin{equation}\label{eq:nonzerosource}
 0\ne D_J(I+h)\in K .
\end{equation}
\end{proposition}
\begin{proof}
每个$L_t$的解空间为二维简单表示，且这些表示两两不同构。
把它们作为函数放入共同Picard--Vessiot域，
相加映射的核是不变子空间。
无重数半单性与每个分量上的非零性使该核为零，
故解空间之和是直和，维数$2|K_0|$。

最小公左倍数的解空间正是这个和。
具体地，包含该和的算子阶数至少是其维数；
直和基的Wronskian构造一个该阶首一算子。
其系数在微分Galois群下不变，故属于$K$。
这也证明最小公左倍数不会额外湮灭一个不在$K_0$中的$V_t$，
否则该简单表示会嵌入其余互不同构表示的直和。
得到阶数和精确选择集合。

由定理\ref{thm:differentialselection}，$D_J(I+h)$为有理函数。
若它为零，则$I+h$被阶$2|J|$算子湮灭。
但同一定理取$D=1$给其最小阶数$2|J|+1$，矛盾。
\end{proof}

因此无需先算完整积分，就可以构造
“只保留指定椭圆核”的有限微分读出，或全部消去纯核后得到非零有理源。
这与把非齐次源丢弃完全不同。
在本族中$L_t$的系数属于质量生成的代数数域上的有理函数域；
最小公左倍数可在该域内由有限欧几里得操作形成。
这里给的是理论终止规则，没有启动新的算子生成程序。
有理源的系数仍可含已固定的质量对数，不能据“有理于$s$”称它为有理数。

\subsection{局部Taylor减除及尺度导数}
取无奇性的欧氏点$s_0<0$、整数$d\ge0$，定义
\begin{equation}\label{eq:localtaylor}
 T_{s_0}^{d}I(s)=\sum_{j=0}^d
       \frac{I^{(j)}(s_0)}{j!}(s-s_0)^j,\qquad
 I^{[d]}_{s_0}=I-T_{s_0}^dI .
\end{equation}
在$K=\mathbb C(s)$中，减去的项是多项式；
其系数作为算术周期时则必须另记。

\begin{corollary}[局部减除与有限尺度响应的完整保持]\label{cor:jetstability}
对本节非空支撑$J$，每个$I^{[d]}_{s_0}$以及$I+h$（$h\in K$）
均在同一$E_J$上表示，最小微分对象都恰为$E_J$，
阶数为$2|J|+1$，微分群不变。
对任意$n\ge0$，$\partial_s^n(I+h)$仍有同样的最小对象。
任意非零一阶有理算子作用也保持该对象，特别是
$s\partial_s-\alpha$，$\alpha\in\mathbb C$。
两个减除方案之\emph{差}是多项式，不属于上述非零纯支撑断言。
\end{corollary}
\begin{proof}
减除函数与$I$之差属于$K$，所以定理\ref{thm:differentialselection}
取$D=1$即得同一完整最小对象。
非零一阶算子不能湮灭任何二维纯解空间，故$J_D=J$。
对$\partial_s^n$，若它湮灭某个$V_t$的所有解，
这些解都将是次数小于$n$的多项式。
其单值表示于是平凡，与$\mathrm{SL}_2$纯单值矛盾。
$n=0$是恒等算子情形。
再次使用式\eqref{eq:operatororder}，得到全部断言。
差函数为两个Taylor多项式之差，可为零或有理函数，
所以不声称它也保留$E_J$。
\end{proof}

这里的导数具有实际参数积分解释。
在负实参数的紧集上，完整theta参数式中的
$F=M U-sA_1A_2A_3$满足
$A_1A_2A_3/F$一致有界。
任意有限次$s$导数使
\[
 F^{-(N-2)}
 \longmapsto (N-2)_n
       (A_1A_2A_3)^n F^{-(N-2+n)} .
\]
它仍由原收敛核的常数倍控制，故可以交换导数与积分。
上述原参数式对应单位分子；对本节允许的$P$，
亦可沿已证完整Stokes等式对有限个sunrise项逐项求导。
这给有限导数及其有理组合的真实读出，而非只形式操纵微分方程。
它可能改变传播子幂次及参数分子；是否仍是原理论的某个指定可观测量须另判。

这个无UV发散的二维例子说明局部减除\emph{能做什么}，
并不冒充四维全森林重整化。
在已构造的四维实际方案中，整图局部差可以是外动量多项式，
子图反项插入收缩图则一般不是关于$s$的有理函数；
后一类应使用定理\ref{thm:fullbackground}及其真实背景条件。
局部导数再在移动减除点求值、对质量取导、取极点系数或临界极限，
也不是本节固定质量的一变量微分操作。

\subsection{为什么局部RG结论仍须保留算术数据}
有理微分读出保留参数函数的整个混合对象，
不等于在一个点的周期值具有同样的动机秩。
最直接的例子是
$I^{[d]}_{s_0}(s_0)=0$，
但作为$s$的函数它仍有最小阶$2|J|+1$。
求值并不是正合忠实的微分模纤维函子：
此处求值作用于选定矩阵系数，不能与微分对象的整条纤维混为一谈。

同理，RG系数在外运动学上为常数，
只能说明相关参数变分被局部投影消去，
并不说明留下的常数在算术上为Tate或有理数。
本节提供两个可以接入总目标的规则：
已构造的森林差若满足socle分离，则完整混合对象作为商受保护；
实际一变量有理微分操作则由纯算子的右余式决定核是否消失。
多耦合、临界质量坐标、热插入及所有尚未构造的实际比较仍保持原任务，
不能由这一函数层的闭合结果跳过。



\section{割点拼接的乘法规则：独立核、重复核与完整混合对象}
\label{sec:blockproducts}
此前的规则处理加法、微分和收缩；结构表还必须处理实际图的乘法。
本节沿割点分解一个明确的收敛图类，
把完整参数周期的乘积关系与乘积的最小化分别证明。
后者依赖各块之间的关系，不能仅把逐块秩机械相乘。
图多项式及参数归一化使用经典约定\cite[\S2,\S4]{BWpoly}；
微分模的张量与对称幂背景见\cite[\S1.5]{Singer}。
本节不对这些经典构造作新颖性宣称。

\subsection{割点的实际动量分配与参数尺度}
设连通图$G=G_1\vee_vG_2$由两图只在一个顶点$v$相交得到。
对$G_i$的非$v$顶点保留原外部动量，在$v$放入使该块单独守恒的动量
\[
 p_v^{(i)}=-\sum_{w\in V(G_i)\setminus\{v\}}p_w .
\]
于是$p_v^{(1)}+p_v^{(2)}=p_v$。
这个重新分配是图分解数据的一部分；
不得让每个块在$v$重复携带整个$p_v$。

取偶数维$D=2d$、正整数传播子幂$\nu_e$及正质量，
各环采用欧氏$d^Dk/\pi^{D/2}$归一化，
省略整体动量守恒的分布因子。
写$h_i=h(G_i)$、$\nu_i=\sum_{e\in G_i}\nu_e$、
$a_i=\nu_i-dh_i$。
假设两个块的欧氏积分绝对收敛且$a_i>0$。
参数式为
\[
 I_i=\frac{\Gamma(a_i)}{\prod_{e\in G_i}\Gamma(\nu_e)}
 \int_{\Delta_i}
   \left(\prod_{e\in G_i}x_e^{\nu_e-1}\right)
   \frac{U_i^{a_i-d}}{F_i^{a_i}}\,\Omega_i .
\]
以下首先处理单位标量分子，有限和的分离分子另由线性性处理。

\begin{proposition}[割点拼接的完整周期乘法]\label{prop:vertexproduct}
在上述假设下，具有诱导块运动学的图多项式满足
\[
 U_G=U_1U_2,\qquad F_G=F_1U_2+F_2U_1 .
\]
图积分满足$I_G=I_1I_2$。
在相同代数参数域上的收敛参数形式周期中也有
\[
 I_G^{\mathfrak m}=I_1^{\mathfrak m}I_2^{\mathfrak m}.
\]
因此最小标架对象$M(I_G^{\mathfrak m})$是
$M(I_1^{\mathfrak m})\otimes M(I_2^{\mathfrak m})$的子商；
没有断言整个图的环境相对上同调就是该张量积。
\end{proposition}
\begin{proof}
图的圈空间为两个块的圈空间直和。
生成树恰为各块生成树的并，给$U$乘法；
生成二森林恰在一个块断开、另一个块保持树。
按诱导守恒动量计算断开分量的动量平方，再加质量项，
即得$F$公式。动量路由也能直接取块对角形式，
所以绝对收敛与Fubini已经给数值乘法。

为取得形式周期关系，不以Fubini数值等式代替几何比较。
在参数单纯形内部写
\[
 \alpha_e=t x_e\quad(e\in G_1),\qquad
 \alpha_e=(1-t)y_e\quad(e\in G_2),\qquad 0<t<1 .
\]
选取相容乘积取向，单纯形测度给
$t^{\nu_1-1}(1-t)^{\nu_2-1}$。
齐次性给
\[
 U_G=t^{h_1}(1-t)^{h_2}U_1U_2,\quad
 \frac{F_G}{U_G}=tA+(1-t)B,\qquad A=\frac{F_1}{U_1},\quad B=\frac{F_2}{U_2}.
\]
因此相对尺度方向的全部积分恰为
\begin{equation}\label{eq:vertexbeta}
 \Gamma(a_1+a_2)\int_0^1
 \frac{t^{a_1-1}(1-t)^{a_2-1}\,dt}
      {[tA+(1-t)B]^{a_1+a_2}}
 =\frac{\Gamma(a_1)\Gamma(a_2)}{A^{a_1}B^{a_2}} .
\end{equation}
这不是未计算的特殊函数输入。作有理变量替换
\[
 u=\frac{tA}{tA+(1-t)B},\qquad
 t=\frac{uB}{A(1-u)+uB},
\]
左侧被积式成为
$A^{-a_1}B^{-a_2}u^{a_1-1}(1-u)^{a_2-1}du$。
$a_i$为正整数，最后的多项式有有理原函数，
两端之差为$(a_1-1)!(a_2-1)!/(a_1+a_2-1)!$。

欧氏域中$A,B>0$，有理映射及逆在链内部无极点，并把两端送到$0,1$。
在各底面上，其与顶次底形式的楔积限制为零；
有理映射的边界未定义处可取图的闭包并作共同消解。
收敛假设使该边界表示可用收敛相对链，
尺度方向的等式由上述有理原函数的Stokes关系给出。
故它属于形式周期允许的代数变量替换、线性性、Stokes和乘积关系，
而非假定数值周期映射单射。
消去$t$后剩余的$U_i^{-d}A^{-a_i}$正是两个块的参数形式，
连同式\eqref{eq:vertexbeta}的系数恢复两个$\Gamma(a_i)$。
最后由形式周期乘法的张量标架及最小化得子商结论。
\end{proof}

对任意有限割点树可以归纳应用此命题。
桥上的动量由外部动量唯一确定，其非零传播子是代数有理因子，
单独记录即可。
如果存在UV发散块、质量为零造成IR问题，或分子含跨块环动量收缩，
必须先检查收敛、有限分离展开或实际重整化相容性。
它们不由割点图形自动解决。

\subsection{在共同参数上证明各块独立}
取有限个此前接受的theta块，记其物理函数为$f_i(s)$，
非空活动集合为$J_i$，$r_i=|J_i|$。
继续固定质量，在$K=\mathbb C(s)$上讨论。
假设\emph{全部块合在一起}仍满足私有阈值条件：
每个活动三元组的正常阈值，不在任何其它活动三元组的退化集合中。
仅检查每个块内部的条件不够。

各块最小微分对象为
\[
 0\longrightarrow V_i=\bigoplus_{t\in J_i}V_t
 \longrightarrow E_i\longrightarrow\mathbf1\longrightarrow0,
 \quad
 G_i=V_i\rtimes\prod_{t\in J_i}\mathrm{SL}(V_t).
\]
这里$V_i$作为群表示也记其加性向量群，
每个$V_t$为二维，且$\dim E_i=2r_i+1$。

\begin{proposition}[跨块私有阈值给完整联合群]\label{prop:blockjointgroup}
这些$E_i$生成的联合微分Galois群为
\[
 G_{\mathrm{joint}}=\prod_iG_i.
\]
这既包括纯部分的独立性，也包括所有加法扩张方向。
\end{proposition}
\begin{proof}
先看联合纯系统。
私有阈值的局部单值仅在一个二维因子非平凡幺幂，
其它纯因子恒等；各单因子的闭包是$\mathrm{SL}_2$。
该幺幂元的正规闭包包含整个对应$\mathrm{SL}_2$因子，
逐一应用得到全部纯直积。正则奇异的单值稠密定理
把这个结论识别为纯微分群，而非仅某些回路的子群。

联合群$H$嵌入$\prod_iG_i$，其到纯直积的投影满射。
记核为$N\subset\bigoplus_iV_i$；它是稳定于纯直积的加性子群，
特征零中即为子表示。
所有$V_t$是互不同构简单模，各只出现一次，
所以$N$是其中若干块的直和。

若某个$V_t$未出现于$N$，把对应$E_i$商去其余纯块。
$N$在所得秩三表示上作用平凡，
于是该表示通过约化纯直积，故半单。
这使原来已证非分裂的$V_t$扩张分裂，矛盾。
故$N=\bigoplus_iV_i$。
拥有整个核且满射到纯商的闭子群只能是全部$\prod_iG_i$。
\end{proof}

这一证明不从“质量不同”直接推导独立。
若阈值条件失败，只保留张量载体上界；
相同函数、同源关系或其它关联应另行识别。

\subsection{乘积的精确最小对象，而非只有张量上界}
把各最小标量方程的解空间嵌入共同Picard--Vessiot域$L$，
记为$W_i$。
由此前完整最小性，$W_i$作为$G_i$表示同构于$E_i$：
它由纯解空间$V_i$和一个单位商方向组成。
物理解$f_i$在单位商中非零；归一化后可取商坐标为$1$。
以下证明使用域中的乘法，须排除其隐藏核。

\begin{theorem}[独立块乘积的完整分类]\label{thm:blockproductminimal}
在命题\ref{prop:blockjointgroup}的假设下，
\[
 M_\partial\left(\prod_i f_i\right)\simeq\bigotimes_i E_i,\qquad
 \operatorname{ord}_{\min}\left(\prod_i f_i\right)
       =\prod_i(2r_i+1).
\]
该最小对象的微分群是$\prod_iG_i$，连通且维数$5\sum_i r_i$。
若各块由割点拼接成前述实际图，这就是该图的函数分类。
\end{theorem}
\begin{proof}
乘法给$G_{\mathrm{joint}}$等变线性映射
\[
 \mu:\bigotimes_iW_i\longrightarrow L,\qquad
 \bigotimes_iw_i\longmapsto\prod_iw_i .
\]
不能预先假定$\mu$单射。
每个加法群$V_i$在$W_i$上的不变子空间正是纯子空间$V_i$，
故联合加法群的不变空间为$\bigotimes_iV_i$。
正规幺幂群在任意简单表示上作用平凡；
因此$\bigotimes_iW_i$的socle包含在$\bigotimes_iV_i$中。
后者已为纯直积的半单表示，实际就是socle，且
\[
 \bigotimes_iV_i=
 \bigoplus_{(t_i)\in\prod_iJ_i}\bigotimes_iV_{t_i}.
\]
这些外张量积是互不同构的简单表示。

如果$\ker\mu\ne0$，有限长度性使其包含一个简单子表示，
于是包含上述某个外张量积。
但每一分量上可以选取非零纯解$y_{t_i}$，
它们在域$L$中的乘积非零。
$\mu$在该简单表示上非零，故单射，矛盾。
所以$\mu$没有核。

另一方面，$G_i$的平移把$f_i$变成$f_i+v$，
$v$遍历$V_i$；它们的线性张成是整个$W_i$。
联合群是全直积，故$\bigotimes_if_i$的轨道张成整个张量积。
于是$\prod_i f_i$的Galois轨道张成的函数空间维数正是
$\prod_i(2r_i+1)$。
其Wronskian构成该阶首一算子，群不变性使系数属于$K$；
任何湮灭此函数的$K$算子都必须湮灭整个轨道。
这证明最小阶数及对象等号。

张量表示的核也须核对。
若$\bigotimes_i g_i$为恒等，则每个$g_i$在$E_i$上为标量；
而每个$E_i$的单位商上作用恒等，使该标量只能为$1$。
各$E_i$原表示忠实，故张量表示忠实，得到群等号及维数。
\end{proof}

这给一个不限总圈数的接受族：
先用割点分解，再逐块识别theta质量网格与分子支撑，
最后核对跨块的私有阈值。
这些步骤只处理有限图和代数数比较，不求完整积分或周期初值。
可以取所有原边质量为互不同位的$M5^e$，
沿用最高位证明得到跨块分离。
所有块的外平方取同一$s$在欧氏域可实现：
把各块的一个端点合并，在另一端放相同欧氏向量$p$，
公共顶点放入总守恒动量。
这允许多外腿及较高顶点价，不声称全部来自某个固定场内容的拉氏量。

\subsection{重复同一块：对称幂使乘法秩下降}
独立性失败不等于所有结论消失；重复块可以精确处理。
设$f$为单个接受theta块的物理函数，活动核数$r\ge1$，
其对象$E$的秩为$2r+1$，群为
$G=V\rtimes(\mathrm{SL}_2)^r$，$V=\bigoplus_{t=1}^rV_t$。
相同运动学与相同质量的$k$个图块沿割点拼接，
命题\ref{prop:vertexproduct}给完整函数$f^k$。

\begin{theorem}[重复块的精确对称幂规则]\label{thm:repeatedblock}
对任意整数$k\ge1$，
\[
 M_\partial(f^k)\simeq\operatorname{Sym}^kE,\qquad
 \operatorname{ord}_{\min}(f^k)=\binom{2r+k}{k}.
\]
其微分群仍为$G$，并没有增加到$G^k$。
特别地，独立的两个sunrise块乘积阶数为九，
同一个sunrise的平方阶数为六。
\end{theorem}
\begin{proof}
令$W$为$f$的最小方程解空间。
乘法经对称幂给$\mu_k:\operatorname{Sym}^kW\to L$。
以$z$表示单位商的一个提升，$W=V\oplus\mathbb Cz$作为向量空间；
平移作用为$z\mapsto z+v$、$V$不动。

在多项式模型中，平移的无穷小作用为$v\partial_z$。
多项式环无零因子，所以所有平移不变元恰为$\operatorname{Sym}^kV$。
于是整个表示的socle为
\[
 \operatorname{Sym}^kV=
 \bigoplus_{n_1+\cdots+n_r=k}
       \bigotimes_{t=1}^r\operatorname{Sym}^{n_t}V_t .
\]
每个$\operatorname{Sym}^{n_t}V_t$是$\mathrm{SL}_2$的简单表示，
上述外张量积分量互不同构。
在任意分量选非零纯解$y_t$，函数$\prod_ty_t^{n_t}$非零；
故$\mu_k$在每个简单分量非零。
若$\mu_k$有核，该核含简单子对象，立即矛盾。
所以乘法是单射。

$f$的商坐标非零，平移后归一化为$z+v$。
所有$(z+v)^k$张成$\operatorname{Sym}^kW$：
展开关于$v$各坐标的多项式，有限插值提取每个齐次系数即可。
因此$f^k$的轨道张成整个对称幂，
最小阶数为其维数$\binom{2r+k}{k}$。

最后，$\operatorname{Sym}^k g=1$迫使线性变换$g$是
$k$次单位根标量：它固定全部纯幂的射影直线，
故固定全部$W$的射影直线，再看纯幂确定标量幂。
$W$的单位商把该标量固定为$1$。
因此对称幂表示仍忠实，微分群不变。
\end{proof}

这里比较的是同一底域、同一变量上的微分阶数。
同质量但不同自变量的块不属于“同一$f$的幂”；
不同质量但碰撞的阈值也不能自动归入上述两种精确分支。

\subsection{结构表的张量闭合与操作出口}
两种拼接给出的高权纯因子具有不同来源。
去掉统一Tate扭转后，独立的$k$个椭圆$H^1$张量
具有
\[
 h^{j,k-j}=\binom{k}{j},
\]
同一椭圆的$\operatorname{Sym}^kH^1$则各有
$h^{j,k-j}=1$，$0\le j\le k$。
这是Hodge双分次的乘法与对称幂计算。
例如两个独立因子的张量秩四，
同一因子的对称平方秩三；
它们都可能有$(2,0)$方向，但不能仅凭这一方向
宣称发现一个新的不可约K3图几何，更不能替代其实际代数模型。
“张量生成的高权核”应在结构表中与新的基本几何核分列。

\begin{corollary}[乘积结构的背景保护与coaction]\label{cor:blockoperations}
上述独立乘积和重复幂的最小微分对象，
均没有简单组成为有限单值的非零子对象。
所以加入由$\mathcal S_{\mathrm{fin}}$承载的背景函数后，
原来的完整乘积对象仍作为商保留。
若独立乘积有$k$个非空块，或取一个非空块的$k$次幂，
其物理函数对微分群幺幂根基的深度都恰为$k$。
这里深度定义为幺幂Lie代数的所有$k+1$重作用为零，
但某个$k$重作用非零；不是固定纤维算术扩张长度的断言。

在已构造的形式周期层面，割点乘法与coaction满足
\[
 \Delta I_G^{\mathfrak m}
   =\prod_i\Delta I_i^{\mathfrak m}.
\]
若$i$块单独绕一个仅使它退化的阈值，则实际解析跳跃为
\[
 (T-1)\prod_j f_j=((T_i-1)f_i)\prod_{j\ne i}f_j .
\]
对重复块则是
\[
 (T-1)f^k=\sum_{j=1}^k
             \binom{k}{j}f^{k-j}\big((T-1)f\big)^j ,
\]
不能只保留第一项。
\end{corollary}
\begin{proof}
独立乘积的socle已算为$\bigotimes_iV_i$的无重数简单分解，
每个分量在至少一个$\mathrm{SL}_2$上非平凡，故单值无限。
重复幂的socle分解中$\sum n_t=k>0$，
每个分量也至少在一个$\mathrm{SL}_2$上非平凡。
应用定理\ref{thm:fullbackground}即可。

对独立乘积，在第$i$块选平移生成元$N_i$，
使$N_if_i=y_i\ne0$为纯解。
于是$N_1\cdots N_k\prod_i f_i=\prod_i y_i\ne0$，
而任意$k+1$重作用为零。
重复幂中选$Nf=y\ne0$，有$N^k f^k=k!y^k\ne0$、
$N^{k+1}f^k=0$；其余生成元的更高混合作用也为零，
因为总次数不超过$k$。这证明深度恰为$k$。

形式周期的coaction为代数同态，故对已证明的实际周期乘法给乘积律。
数值函数的跳跃由解析延拓保持乘法，
以及另一些块在该回路上单值恒等而得；
重复幂用二项式展开。
这里不将参数微分群等同固定纤维Nori群。
\end{proof}

这些规则直接扩大可接受的图类：
割点拼接已识别的收敛块，可输出张量载体；
跨块私有阈值给完整秩与群，重复块给对称幂分支，
未判定关联时保留上界而不输出伪等号。
分离分子的有限线性组合仍须经过实际物理求和的最小化，
Ward关系、图对称因子和指定可观测量不由本节虚构。
下一节处理两顶点拼接的正谱接受类；一般无割点骨架、
跨块UV森林和一般算术下降仍是总目标中独立的未完成接口。


\section{两顶点拼接：卷积推前与新几何的产生}\label{sec:twoterminal}
\subsection{增加一个圈的操作}
割点乘法之外，最先需要处理的操作是把两个连通图的两对端点分别识别。
本节只允许两个公共顶点$a,b$，两块无公共边，内部顶点没有外动量注入，
总外动量为$a$处$p$、$b$处$-p$。不删除额外的标记边；
此处的操作是平行拼接$G=G_1\parallel G_2$，并非另一种约定下的图拟阵二和。
圈数满足$h_G=h_1+h_2+1$。新增的一个圈正是跨块传递动量$q$的积分，
所以不能沿用前节的周期乘法。

令$U_i$为余生成树多项式，$V_i$为分离$a,b$的余生成二森林多项式，
$M_i=\sum_{e\in G_i}m_e^2\alpha_e$。
两点第二多项式采用$F_i(s)=M_iU_i-sV_i$，$s=p_{\mathrm{Mink}}^2$。
图多项式的生成树定义及参数规范来自经典文献\cite{BWpoly}；
下面直接证明这一操作的具体公式。

\begin{proposition}[两端拼接的实际图多项式与卷积]\label{prop:parallelpoly}
在上述图与动量条件下，
\begin{equation}\label{eq:parallelpoly}
 U_G=U_1V_2+V_1U_2,\qquad V_G=V_1V_2,\qquad
 F_G=(M_1+M_2)(U_1V_2+V_1U_2)-sV_1V_2 .
\end{equation}
若正质量、单位标量分子、正传播子幂的欧氏积分绝对收敛，
以每圈$d^Dk/\pi^{D/2}$归一化，则
\begin{equation}\label{eq:parallelconv}
 I_G(-p_E^2)=\int_{\mathbb R^D}\frac{d^Dq}{\pi^{D/2}}\,
       I_1(-q^2)I_2(-(p_E-q)^2).
\end{equation}
公式同时给出一个有额外积分变量的实际参数模型，而不是两个固定点对象的张量等式。
\end{proposition}
\begin{proof}
$G$的生成树在每块上的限制都是森林。每个分量必须遇到至少一个端点，
否则在总树中不连通，所以每块至多有两个分量。
两块不可能都是树，否则$a,b$间产生圈；也不能都是分离二森林，
否则总图不连通。恰一块为树、另一块为分离二森林，
按余边单项式求和给$U_G$。
对总分离二森林，两块均须为分离二森林，故$V_G=V_1V_2$。
只有端点注入动量时，其余二森林的动量权为零，
质量项为$(M_1+M_2)U_G$，得到$F_G$。

选每块内部的圈动量和跨块流量$q$为总圈坐标，
边传播子分为两组，整数线性换坐标的Jacobian绝对值为一。
欧氏正性及收敛性给Tonelli/Fubini，从而得\eqref{eq:parallelconv}。
还可直接检查其参数模型：固定各块Schwinger参数，先积内部圈，
设$r_i=V_i/U_i>0$。余下高斯积分为
\[
 \int\frac{d^Dq}{\pi^{D/2}}
  e^{-r_1q^2-r_2(p_E-q)^2}
  =(r_1+r_2)^{-D/2}
   e^{-p_E^2r_1r_2/(r_1+r_2)} .
\]
乘$U_1^{-D/2}U_2^{-D/2}e^{-(M_1+M_2)}$
正好恢复\eqref{eq:parallelpoly}，包括质量及动量项。
高斯步骤在这里是解析参数化检查，未被称为代数态射。
偶数维、整数幂时，径向参数积分后的有理形式则直接由上述$U_G,F_G$
定义实际相对参数对象。对其形式周期的进一步分解仍须相对推前与边界比较，
不能由高斯数值等式宣布完整动机同构。
\end{proof}

若内部有其它外动量、分子跨块耦合或需要UV/IR减除，
标量卷积必须相应扩大；这些情形不在本节自动接受域中。
即使两块的几何核已经知道，$q$依赖及其退化集合仍是必需输入。

\subsection{二维正谱类的精确拼接规则}
给一个可完全推导的接受类。设$D=2$，两个函数具有无减除正Stieltjes表示
\[
 I_i(s)=\int_{(0,\infty)}\frac{d\mu_i(t)}{t-s}.
\]
测度可含原子、可有无限总质量；要求在负实$s$处表示有限，
且拼接卷积有限。单传播子的测度是$\delta_{m^2}$。
正质量banana的无减除表示已在前文由实际正参数链证明，
所以不是在此对任意图假设谱正性。
记
\[
 \lambda(x,y,z)=x^2+y^2+z^2-2xy-2xz-2yz,\qquad
 B_2(s;u,v)=\int_0^1\frac{dx}{ux+v(1-x)-sx(1-x)} .
\]

\begin{theorem}[正谱卷积与精确物理支撑]\label{thm:spectralglue}
在上述接受域内，拼接函数及其谱测度由
\begin{align}
 I_G(s)&=\int d\mu_1(u)d\mu_2(v)\,B_2(s;u,v),\label{eq:spectralglue}\\
 \rho_G(t)&=2\int_{\sqrt u+\sqrt v<\sqrt t}
    \frac{d\mu_1(u)d\mu_2(v)}{\sqrt{\lambda(t,u,v)}},\qquad
 I_G(s)=\int\frac{\rho_G(t)\,dt}{t-s}                       \label{eq:spectraldensity}
\end{align}
给出，密度式在通常的几乎处处意义成立。
若$\inf\operatorname{supp}\mu_i=a_i>0$，且每个$a_i$的右邻域有正测度，
则谱支撑闭包精确为$[(\sqrt{a_1}+\sqrt{a_2})^2,\infty)$。
在密度光滑的阈值上方开区间，
$\operatorname{Disc}I_G(t)=2\pi i\,\rho_G(t)$。
\end{theorem}
\begin{proof}
先在$s=-p_E^2<0$应用\eqref{eq:parallelconv}与Tonelli，
把两个正测度移出$q$积分。一个Feynman参数和
$\int_{\mathbb R^2}d^2q/\{\pi(q^2+\Delta)^2\}=\Delta^{-1}$
给内层$B_2$，证明\eqref{eq:spectralglue}。
这一常数规范由上述初等二维积分固定，不沿用别处相空间的$2\pi$约定。

对$B_2$作有理映射
\[
 t=\frac{ux+v(1-x)}{x(1-x)}=\frac{u}{1-x}+\frac{v}{x}.
\]
它在$(0,1)$的最小值为$(\sqrt u+\sqrt v)^2$，
较大$t$有两个原像。二次方程的判别式为$\lambda(t,u,v)$，
两支分别给
$|dx|/[x(1-x)]=dt/\sqrt{\lambda(t,u,v)}$，故
\[
 B_2(s;u,v)=
 2\int_{(\sqrt u+\sqrt v)^2}^{\infty}
       \frac{dt}{(t-s)\sqrt{\lambda(t,u,v)}} .
\]
再用Tonelli得\eqref{eq:spectraldensity}。
其支撑下界显然；任取$t$严格高于该下界，
两测度在端点附近的一小块乘积有正质量且$\sqrt u+\sqrt v<\sqrt t$，
于是$\rho_G(t)>0$（允许个别点取无穷，整体局部可积）。
Stieltjes积分在割平面上的紧集可用任一负实参考点的积分一致控制，
所以等式全纯延拓；光滑区间上的边值公式给最后一式。
\end{proof}

这是经典相空间递推/色散思想的本规范推导，参见Kreimer\cite{KreimerBanana}
的相关递推；不主张发现该物理机制。此处的用途是保留所有变量、根与链，
使递推能够输出实际几何，而不止输出积分表示。
更一般带符号谱或减除多项式，不得不加检查地套用正测度结论。

\subsection{从代数根到不可约椭圆核}
用$B_n(s;m_1,\ldots,m_n)$表示$n$条正质量边的二维banana，
每圈规范仍如上。单边$B_1=(m_1^2-s)^{-1}$。
把一个$n-1$边块与最后一条边拼接，立即得到
\begin{equation}\label{eq:bananarhoiter}
 \rho_n(s)=2\int_{M_{n-1}^2}^{(\sqrt s-m_n)^2}
 \frac{\rho_{n-1}(t)\,dt}{\sqrt{\lambda(s,t,m_n^2)}},
 \qquad M_j=\sum_{i=1}^j m_i,\quad s>M_n^2 .
\end{equation}
前文收敛与无减除表示保证每一步合法；这是函数等式，不需先求出$B_n$。

\begin{proposition}[三边拼接的椭圆核]\label{prop:glueelliptic}
任意$m_1,m_2,m_3>0$及$s>(m_1+m_2+m_3)^2$，有
\[
 \rho_3(s)=4\int_{(m_1+m_2)^2}^{(\sqrt s-m_3)^2}
 \frac{dt}{\sqrt{\lambda(t,m_1^2,m_2^2)\lambda(s,t,m_3^2)}} .
\]
相应紧曲线
\begin{equation}\label{eq:gluequartic}
 E_s:\quad y^2=\lambda(t,m_1^2,m_2^2)\lambda(s,t,m_3^2)
\end{equation}
为属一曲线。该正割周期非零，其在$s$方向的延拓生成整个二维$H^1$核。
因此两步平行拼接不能用原有代数根的张量标签代替。
\end{proposition}
\begin{proof}
由前一定理$\rho_2(t)=2/\sqrt{\lambda(t,m_1^2,m_2^2)}$，
代入\eqref{eq:bananarhoiter}即得积分。
四个分支点依次为
\[
 (m_1-m_2)^2<(m_1+m_2)^2<(\sqrt s-m_3)^2<(\sqrt s+m_3)^2 .
\]
它们互异；Riemann--Hurwitz给属一，$\omega=dt/y$为非零全纯一形式。
实区间的两张反向薄片闭合成一条整循环$\gamma$，
$\rho_3=2\int_\gamma\omega>0$。
在正常阈值处恰有中间两点碰撞，另两点仍分离，
交比趋于$0,1$或$\infty$，$j$有极点，故族非等模。
在有限级别覆盖后，非恒定代数$j$映射为模曲线的支配泛有限映射；
删去有限分支值后，像基本群在模曲线基本群中为有限指数。
因而$H^1$单值Zariski稠密于$\mathrm{SL}_2$，绝对不可约。
非零物理周期的延拓轨道遂生成整个二维核。
这只陈述函数/变分层的纯核；完整sunrise扩张另由前节的不分裂证明给出。
\end{proof}

两种根可先分别置于四重覆盖上，再取同时翻转两根的商，得到
\eqref{eq:gluequartic}。形式$dt/y$由该有限代数商下降，
链由实根选择给出；这明确指出了推前中出现的新曲线。
它不是“椭圆常数不属于数值MZV”的未经证明断言。

\subsection{两组双线拼接的K3双覆盖}
现在把四条边分为$(m_1,m_2)$、$(m_3,m_4)$两组，
设
\[
 a_\pm=(m_1\pm m_2)^2,\quad b_\pm=(m_3\pm m_4)^2,\quad
 A(a)=(a-a_-)(a-a_+),\quad B(b)=(b-b_-)(b-b_+).
\]
谱拼接给
\begin{equation}\label{eq:rho4surface}
 \rho_4(s)=8\int_{\mathcal R_s}
 \frac{da\,db}{\sqrt{A(a)B(b)\lambda(s,a,b)}},\qquad
 \mathcal R_s=\{a\ge a_+,\,b\ge b_+,\,\sqrt a+\sqrt b\le\sqrt s\}.
\end{equation}
考虑$(a,b)\in\mathbb P^1\times\mathbb P^1$上的双覆盖
\begin{equation}\label{eq:spectralk3}
              X_s:\quad y^2=A(a)B(b)\lambda(s,a,b).
\end{equation}
下面给足够而明确的接受条件：
\begin{equation}\label{eq:k3accept}
 m_i>0,\quad m_1\ne m_2,\quad m_3\ne m_4,\quad s\ne0,\quad
 \lambda(s,a_\epsilon,b_\eta)\ne0\quad(\epsilon,\eta=\pm).
\end{equation}
这些条件对所有实$s>(\sum m_i)^2$成立，前两对质量互异仍须保留。
代数参数时应在对应数域上取模型，几何结论在代数闭包上判断。

\begin{theorem}[不求周期的K3识别与秩上界]\label{thm:spectralk3}
在\eqref{eq:k3accept}下，$X_s$有十二个$A_1$点及一个$A_2$点，
无其它奇点。其最小消解$\widetilde X_s$为K3曲面，且
\[
 \rho(\widetilde X_s)\ge16,\qquad
 \rk T(\widetilde X_s)\le6,\qquad
 \omega_s=\frac{da\wedge db}{y}
\]
延拓为非零全纯二形式。
这里不宣称Picard数必为16，也不由此识别完整整格或旧周期基。
\end{theorem}
\begin{proof}
分支为两条竖直线、两条水平线和曲线$C_s:\lambda(s,a,b)=0$。
分别次数为$(2,0),(0,2),(2,2)$，故总分支次数$(4,4)$，
覆盖线丛为$L=\mathcal O(2,2)$。
$a_\pm,b_\pm$均为非零有限互异值。
竖直与水平线给四个交点；
每条竖直线与$C_s$有两个互异横截点，因为关于$b$的判别式为$16sa_\epsilon$。
水平线同理，共再给八个点。
条件$\lambda(s,a_\epsilon,b_\eta)\ne0$排除三重交会。
这些十二个普通分支节点在双覆盖上为$z^2=uv$型$A_1$点。

仿射$C_s$没有奇点：两个偏导同时为零要求$a-b=s$且$b-a=s$，
与$s\ne0$矛盾。其唯一无穷远点是$(\infty,\infty)$。
置$u=1/a,v=1/b$，该处的齐次方程为
\[
 c(u,v)=(u-v)^2-2suv(u+v)+s^2u^2v^2=0.
\]
沿唯一切线$u=v$的三次项为$-4su^3\ne0$，
故为普通尖点。另四条分支线不经过此点，
它们在此为可逆因子；解析取其平方根后，
覆盖奇点等价于$z^2=t^2+w^3$，即$A_2$。
以上穷尽分支奇点，覆盖在其它点光滑。

有限双覆盖的结构代数为
$\pi_*\mathcal O_{X_s}=\mathcal O\oplus L^{-1}$，
对偶化层公式给
$\omega_{X_s}=\pi^*(\omega_{\mathbb P^1\times\mathbb P^1}\otimes L)
\simeq\mathcal O_{X_s}$。
由Künneth，$H^1(\mathcal O)=H^1(\mathcal O(-2,-2))=0$。
$A_1,A_2$均为有理双点，最小消解保留$H^1(\mathcal O)$及平凡典范，
故$\widetilde X_s$是K3。所用双覆盖与有理双点的标准背景见\cite{HuyK3}；
此处分支与局部型逐一另算。
仿射$\omega_s$正是典范生成元，以上线丛及crepant消解保证其全局延拓。

最后，两个底的ruling类的拉回具有交叉矩阵
$\left(\begin{smallmatrix}0&2\\2&0\end{smallmatrix}\right)$。
十二个$A_1$例外类及$A_2$的两个例外类组成秩十四的负定格，
与两个拉回类正交。因此这十六个类线性无关，
给$\rho\ge16$；K3的$b_2=22$遂给超越秩上界六。
这些类生成的子格未被证明饱和；不能据此输出判别式48。
\end{proof}

\begin{proposition}[实际物理链与K3全纯周期]\label{prop:k3spectralchain}
保留两对质量互异，取$s>(\sum m_i)^2$。
存在$\Gamma_s\in H_2(\widetilde X_s,\mathbb Z)$，使
\begin{equation}\label{eq:k3spectralperiod}
                 \rho_4(s)=4\int_{\Gamma_s}\omega_s>0.
\end{equation}
特别地，谱读出确实有非零超越分量，不能只归入K3的代数除子。
此结论不自动证明整条超越局部系统不可约，亦不等同完整未割对象。
\end{proposition}
\begin{proof}
$\mathcal R_s$是紧曲边三角形，内部三个根均为正。
在其上取$y=\pm\sqrt{AB\lambda}$两片，
把第二片反向定向，沿三条分支边粘合。
由$y$变号，二形式积分变为同号，得到
$2\int_{\mathcal R_s}da\,db/\sqrt{AB\lambda}$。
三角形的三个角位于前述$A_1$点，其余边上覆盖光滑。
因节点的例外集是有理曲线树，$H_1$为零，
上述闭合二循环可提升到消解的整二同调；
两个提升之差由例外二循环给出。
全纯二形式在例外曲线上为零，故积分不依赖这个选择。
局部根奇性可积，在角点是$u^{-1/2}v^{-1/2}$，
没有遗漏非零边界项。结合\eqref{eq:rho4surface}得因子$4$。
内部积分严格为正，故周期非零；
全纯二形式在全部代数除子类上为零，所以$\Gamma_s$的超越投影非零。

还可在代数层记录根的来源：先取
$r_1^2=A(a),r_2^2=B(b),r_3^2=\lambda(s,a,b)$的八重覆盖。
形式$da\wedge db/(r_1r_2r_3)$具有三个符号的乘积角色。
偶数次翻号子群的商正是\eqref{eq:spectralk3}，
$y=r_1r_2r_3$给有限代数映射，拉回形式就是上式。
其迹及实链推前是实际覆盖操作，不依赖数值周期映射单射。
因而这里得到的是被物理谱选中的几何载体，
不是仅凭某个积分长得像曲面周期而命名。
\end{proof}

K3与三圈banana的关系本身已有\cite{BKV}等先行研究；
本节给出的作用是把它放入两端拼接的可检查分支模型中。
旧已认证banana模型与本双覆盖的具体整基、Tate尺度及饱和比较不在此重新认定。
即使两者给出同一物理割函数，也不能未经实际对应把所有整格标记直接搬运。
第\ref{sec:evendimk3}节补出次数四的实际有理对应，
并在指定质量射线上识别有理纯局部系统；该结论仍不搬运整基和绝对周期标记。
两对质量相等时，部分分支交会并改变局部型；
本节的十二加一计数不再适用，不能因此断言K3消失。

\subsection{阈值非零性与操作出口}
\begin{corollary}[任意边数的正常阈值首项]\label{cor:bananathresholdall}
任意$n\ge2$、正质量、单位分子，设$M=\sum_{i=1}^n m_i$。
在本节规范下，$\varepsilon\downarrow0$时
\begin{equation}\label{eq:bananathresholdall}
 \rho_n(M^2+\varepsilon)=
 \frac{\pi^{(n-1)/2}}
 {\Gamma((n-1)/2)\sqrt{\prod_i m_i}\,M^{(n-2)/2}}\,
 \varepsilon^{(n-3)/2}\bigl(1+O(\varepsilon)\bigr).
\end{equation}
因此每一步的物理割非零；任何有限局部多项式减除都不改变这项割。
这不是任意森林反项均为多项式背景的断言。
\end{corollary}
\begin{proof}
$n=2$时$\lambda(M^2+\varepsilon,m_1^2,m_2^2)
=\varepsilon(4m_1m_2+\varepsilon)$，结论直接成立。
设前一步的指数为$\alpha=(n-4)/2$、系数为$C_{n-1}$，
写$M'=M-m_n$、$t=M'^2+u$。
\eqref{eq:bananarhoiter}的上限为
$u_{\max}=(\sqrt{M^2+\varepsilon}-m_n)^2-M'^2
=(M'/M)\varepsilon+O(\varepsilon^2)$。
为控制端点误差，作精确代换$u=u_{\max}v$，$0\le v\le1$。
根式的两个因子分别为
$\varepsilon(1-v)(1+O(\varepsilon))$、
$4M'm_n(1+O(\varepsilon))$，
余项在$v\in[0,1]$一致。
归纳余项也一致，$v^\alpha(1-v)^{-1/2}$可积。
因此
\[
 C_n=\frac{C_{n-1}}{\sqrt{M'm_n}}
       \left(\frac{M'}M\right)^{(n-2)/2}
       \mathrm B\left(\frac{n-2}2,\frac12\right).
\]
用$\mathrm B(x,y)=\Gamma(x)\Gamma(y)/\Gamma(x+y)$约去中间因子，
恰得到\eqref{eq:bananathresholdall}。全部系数为正；
多项式没有割跳跃，故不改变此结论。
\end{proof}

这个首项证明只用了正性与端点局部展开，没有进行数值积分，
也没有求完整高能或完整阈值级数。
对$n=3$它给非零常数首项；对$n=4$给
$2\pi\sqrt{\varepsilon}/(M\sqrt{\prod m_i})$，
与上面消失的实K3二循环相容。

至此，结构表获得一条不同于割点乘法的规则：
先从两端图计算$U,V$与实际卷积；在有已证正谱表示的接受域内，
由谱支持和Källén根给代数推前模型；
再检查分支碰撞、紧化、边界链与全纯读出。
三边给精确二维椭圆核，四边给实际K3载体、秩上界及非零物理投影。
它不把“两个Tate块”映成一个预设标签，新增几何来自必须保留的推前。
下一节证明自能边插入中减除与推前相容的一类实际条件。
完整混合扩张、一般多端骨架、任意森林交换及物理分子可能引起的相消，
仍需各自的态射及非零性证明；本节不是任意图的完整分类。


\section{自能边插入：减除、谱推前与椭圆核的保留}\label{sec:selfinsert}
\subsection{操作的输入、规范与先行结果}
两端拼接之外，一个常用操作是把截肢二点子图$H$插入外层图$C$的一条边。
本节严格保留插入两侧的两个传播子：
原边的$1/(q^2+\mu)$被
$\Sigma_H(-q^2)/(q^2+\mu)^2$替换，其中$\mu>0$是质量平方。
这与把一条边直接换成一个二端网络不是同一个操作。
外层剩余被积式记为$K(k,p)$，$q$是其内动量的仿射线性组合；
每个欧氏圈采用$d^Dk/\pi^{D/2}$，$s=-p^2$。
以下没有包含拉氏量的耦合常数或图对称因子。

自能的色散插入和在壳两项减除是已有方法：
Aleksejevs在\cite[\S II.A, 式(15)]{Aleksejevs}明确写出
减除后平方因子与两侧传播子的消去。
本节不将这一恒等式作为新发现。
新增到本文理论链的是接受条件、共同相对模型中的减除关系、
外层收敛的精确限制，以及一个实际五边图的椭圆可见性证明。
本节的“在壳”只表示在$s=\mu$固定值与一阶导数，
不是宣称任意单图都构成可观测量或满足规范恒等式。

\subsection{谱减除与收敛外图的实际比较}
\begin{proposition}[两项减除后的正谱边]\label{prop:subtractededge}
设$t_0>\mu>0$，$\rho(t)\geq0$，且
\[
 Z_\mu=\int_{t_0}^{\infty}\frac{\rho(t)}{(t-\mu)^2}\,dt<\infty.
\]
假定$\Sigma$在$\mathbb C\setminus[t_0,\infty)$上有一次减除表示
\[
 \Sigma(s)-\Sigma(\mu)
 =(s-\mu)\int_{t_0}^{\infty}
       \frac{\rho(t)}{(t-\mu)(t-s)}\,dt+b(s-\mu),
\]
其中$b$为常数；也允许该式本身定义有限差，而不定义两个发散的单项。
令
\[
 \Sigma_{\rm os}(s)=\Sigma(s)-\Sigma(\mu)-(s-\mu)\Sigma'(\mu).
\]
则
\begin{align}
 \Sigma_{\rm os}(s)
  &=(s-\mu)^2\int_{t_0}^{\infty}
       \frac{\rho(t)}{(t-\mu)^2(t-s)}\,dt,\label{eq:osdisp}\\
 g_\mu(s):=\frac{\Sigma_{\rm os}(s)}{(\mu-s)^2}
  &=\int_{t_0}^{\infty}\frac{d\nu_\mu(t)}{t-s},
  \qquad d\nu_\mu(t)=\frac{\rho(t)\,dt}{(t-\mu)^2}.
 \label{eq:effectiveedge}
\end{align}
$g_\mu$在$s=\mu$无极点，且对于$r\geq0$，
\[
 0\leq g_\mu(-r)\leq\frac{Z_\mu}{r+t_0}
                        \leq\frac{Z_\mu}{r+\mu}.
\]
\end{proposition}
\begin{proof}
在任意避开谱支撑的紧集上，$(t-\mu)/(t-s)$一致有界，
故给定的有限矩控制积分及$\mu$附近的导数。
于是$\Sigma'(\mu)=Z_\mu+b$。
从一次减除式中减去$(s-\mu)(Z_\mu+b)$，通分即得
\eqref{eq:osdisp}。再除平方因子给\eqref{eq:effectiveedge}；
所得积分在$\mu<t_0$处全纯，其欧氏上界由$t\geq t_0$逐点得到。
这里消去的是两个传播子可能产生的极点，
不是把整个非局部自能换成一个局部反项。
\end{proof}

\begin{theorem}[收敛外图的谱插入及共同形式模型]\label{thm:edgepush}
设欧氏标量外图满足$K\geq0$，并且
\[
 C(p;\mu)=\int\prod_{\ell=1}^{h_C}\frac{d^Dk_\ell}{\pi^{D/2}}
             \frac{K(k,p)}{q^2+\mu}<\infty .
\]
在命题\ref{prop:subtractededge}的条件下，实际减除插入为
\begin{align}
 J_\mu(p)&=\int\prod_\ell\frac{d^Dk_\ell}{\pi^{D/2}}
                       K(k,p)g_\mu(-q^2)\notag\\
          &=\int_{t_0}^{\infty}d\nu_\mu(t)\,C(p;t),
 \qquad 0\leq J_\mu(p)\leq Z_\mu C(p;\mu).
 \label{eq:edgepush}
\end{align}
若$K$与$\rho$在正测度集合上严格正，则$J_\mu>0$。

另设$D=2$，$H$为正质量banana，使用其既有实际参数模型
\[
 \Sigma_H(s)=\int_{\sigma_H}\frac{\Omega_H}{A-sB},
 \qquad A=\Bigl(\sum_i m_i^2\alpha_i\Bigr)U_H,\quad
 B=\prod_i\alpha_i .
\]
质量与外参数取代数值，或工作在明确的代数参数域。
在$\mu<t_0=(\sum_i m_i)^2$时，插入的两项减除
在同一个$H$参数空间与外层动量空间的乘积模型中，由有理形式
\begin{equation}\label{eq:hybridform}
 K(k,p)\,\frac{B^2\Omega_H}
              {(A+q^2B)(A-\mu B)^2}
 \prod_\ell\frac{d^2k_\ell}{\pi}
\end{equation}
实际表示。因此在这个指定共同相对模型中，
减除等式是形式周期关系，不只是数值相等。
\end{theorem}
\begin{proof}
Tonelli适用于非负被积式，给出交换次序及可能无穷的等式。
命题\ref{prop:subtractededge}的上界随即证明其有限性。
正性结论也直接来自正测度集合，而不依赖数值求积。

对第二项断言，逐点使用
\[
 \frac1{A-sB}-\frac1{A-\mu B}
 -\frac{(s-\mu)B}{(A-\mu B)^2}
 =\frac{(s-\mu)^2B^2}
       {(A-sB)(A-\mu B)^2}.
\]
令$s=-q^2$，并约掉两侧传播子的平方分母。
banana中$U_H=B\sum_i\alpha_i^{-1}$，
Cauchy不等式给$A/B\geq(\sum_i m_i)^2$；
所以$A-\mu B$在链内部严格正。
已收敛的$\Sigma_H(0),\Sigma_H(\mu),\Sigma_H'(\mu)$控制各项：
额外的$(q^2+\mu)^{-1}$有界于$\mu^{-1}$，
故这些形式的积分分别有限，不存在拆开后发散的非法相消。

将各实动量坐标用有理映射紧化，例如
$x\mapsto(2x-1)/(x(1-x))$把$(0,1)$双射到$\mathbb R$，
再共同消解极点与边界。四个有理形式位于同一相对复形，
上述等式在形式层已经成立；拉回、线性与共同链因此给形式周期关系。
每圈的$\pi^{-1}$统一作为Tate规范保留。
这没有把一般高斯积分宣布为代数态射，
也没有证明任意另选的全动量图对象与全参数图对象同构。
\end{proof}

\subsection{局部子反项为何成为非局部外层背景}
\begin{proposition}[有限局部方案差的外层像]\label{prop:edgescheme}
保持边质量$\mu$及外图固定，设两个子图方案满足
\[
 \Sigma_b(s)-\Sigma_a(s)=a+b(s-\mu),
\]
其中$a,b$不依赖外运动学$p$。在上述收敛条件下，
\begin{equation}\label{eq:edgeschemediff}
 J_b(p)-J_a(p)
       =-a\,\partial_\mu C(p;\mu)-b\,C(p;\mu).
\end{equation}
$\partial_\mu$只作用被标记的边质量，$K$保持固定。
若$a,b$与外图均有已给出的共同形式周期模型，
该关系在相应乘积/共同模型中也成立。
因而在由乘积载体
$E_a\otimes E_{\partial_\mu C}$与$E_b\otimes E_C$
生成的Serre背景商中，
两种插入具有相同的最小对象；背景外简单组成的重数相同。
\end{proposition}
\begin{proof}
在$s=-q^2$处，
\[
 \frac{a+b(s-\mu)}{(\mu-s)^2}
   =\frac{a}{(q^2+\mu)^2}-\frac{b}{q^2+\mu},
 \qquad
 \partial_\mu\frac1{q^2+\mu}=-\frac1{(q^2+\mu)^2}.
\]
正质量与额外有界分母允许在积分下求导，得到公式。
实际形式模型中的相同有理等式与乘积周期给形式关系；
随后应用本文的加性Serre比较定理。
局部的是插入前的$a+b(s-\mu)$，
插入后一般出现依赖$p$的$C$及其质量导数。
此外，$a,b$可能是非Tate常数周期，
不能仅因其不依赖$p$而将整个背景宣布为Tate。
\end{proof}

\subsection{五边图：低阈值消去，而椭圆割保留}
取外图为二传播子bubble：被标记边质量平方$\mu$，
另一边质量为$M>0$；插入子图为质量$m_1,m_2>0$的bubble。
实际五条边为
\[
 (a,c;\sqrt\mu),\quad(c,d;m_1),\quad(c,d;m_2),\quad
 (d,b;\sqrt\mu),\quad(a,b;M),
\]
外动量只在$a,b$注入。图有四顶点、五边、两个圈，
不是二顶点banana。允许二价外部插入顶点，
不宣称它来自某个指定的纯$\phi^4$拉氏量。
记
\[
 a_\pm=(m_1\pm m_2)^2,\quad t_0=a_+,\quad
 A(t)=(t-a_-)(t-a_+),\quad
 \tau=(m_1+m_2+M)^2,
\]
并假设$0<\mu<t_0$。

\begin{proposition}[二维五边图的减除与精确割]\label{prop:fiveedge2}
以$\Sigma_2(s)=B_2(s;m_1^2,m_2^2)$为截肢子图，
其两项减除插入满足
\begin{align}
 J_{2,\rm os}(s)
  &=J_{2,\rm bare}(s)
       +\Sigma_2(\mu)\,\partial_\mu B_2(s;\mu,M^2)
       +\Sigma_2'(\mu)\,B_2(s;\mu,M^2),\label{eq:fiveedge2local}\\
  &=\int_{t_0}^{\infty}
       \frac{2\,B_2(s;t,M^2)}
            {(t-\mu)^2\sqrt{A(t)}}\,dt .
 \label{eq:fiveedge2spectral}
\end{align}
它的谱支撑恰为$[\tau,\infty)$，且$s>\tau$时
\begin{equation}\label{eq:fiveedgecut2}
 \rho_{J,2}(s)=4\int_{t_0}^{(\sqrt s-M)^2}
       \frac{dt}{(t-\mu)^2
          \sqrt{A(t)\lambda(s,t,M^2)}}>0,\qquad
 \operatorname{Disc}J_{2,\rm os}=2\pi i\,\rho_{J,2}.
\end{equation}
正常阈值的首项为
\[
 \rho_{J,2}(\tau+\varepsilon)
 =\frac{\pi}
 {(t_0-\mu)^2\sqrt{m_1m_2M(m_1+m_2+M)}}
       \bigl(1+O(\varepsilon)\bigr),\qquad\varepsilon\downarrow0.
\]
较低点$\tau_{\rm low}=(\sqrt\mu+M)^2$不是减除插入的阈值，
但一般是裸插入的真实分支点。
\end{proposition}
\begin{proof}
二维bubble正谱为$2/\sqrt{A(t)}$，
相应$Z_\mu=\Sigma_2'(\mu)$有限。外层二维bubble绝对收敛，
故定理\ref{thm:edgepush}直接给\eqref{eq:fiveedge2spectral}。
从减除定义或命题\ref{prop:edgescheme}得到
\eqref{eq:fiveedge2local}，包括其中两个正号。
再用前节已独立定标的$\rho_{B_2}(s;t,M^2)=2/\sqrt{\lambda(s,t,M^2)}$，
得到\eqref{eq:fiveedgecut2}。有效区间非空恰为$s>\tau$；
正性和阈值前全纯性分别给精确支撑的两端断言。

当$\varepsilon\to0^+$时，积分区间收缩到$t_0$，
其上$(t-\mu)^{-2}=(t_0-\mu)^{-2}(1+O(\varepsilon))$一致成立。
代入前节三边banana阈值公式，得到所列常数和余项。
又因$\tau_{\rm low}<\tau$，
在两阈值之间，裸插入的割是
\[
 -2\pi i\left[
 \Sigma_2(\mu)\,\partial_\mu\rho_{B_2}(s;\mu,M^2)
 +\Sigma_2'(\mu)\,\rho_{B_2}(s;\mu,M^2)\right].
\]
在$\tau_{\rm low}$上方，
$\rho_{B_2}\asymp(s-\tau_{\rm low})^{-1/2}$，
其质量导数含非零$(s-\tau_{\rm low})^{-3/2}$首项，
因为$\partial_\mu\tau_{\rm low}=1+M/\sqrt\mu>0$。
$\Sigma_2(\mu)>0$，故这一首项不能与后一项抵消。
断言仅针对该单图及其给定减除组合，并非完整散射振幅的阈值分类。
\end{proof}

\subsection{四维反例与继承的整体减除}
\begin{proposition}[有限子图不保证有限插入图]\label{prop:oversubtraction}
对同一五边拓扑改取$D=4$。令
\[
 D_x(s)=m_1^2x+m_2^2(1-x)-s\,x(1-x),\quad v=x(1-x),
\]
并定义有限的一次差和其在壳斜率
\[
 \Sigma_0(s)=-\int_0^1\log\frac{D_x(s)}{D_x(\mu)}\,dx,\qquad
 Z_\mu=\int_0^1\frac{v\,dx}{D_x(\mu)}>0.
\]
设$\Sigma_1(s)=\Sigma_0(s)-(s-\mu)Z_\mu$。
两项减除已使子图有限，但其直接外层插入满足
\begin{equation}\label{eq:oversubdiv}
 \int_{|q|<\Lambda}\frac{d^4q}{\pi^2}\,
 \frac{\Sigma_1(-q^2)}
 {(q^2+\mu)^2((p-q)^2+M^2)}
       =Z_\mu\log\Lambda^2+O(1).
\end{equation}
对比之下，将$\Sigma_1$换成仅减去常数的$\Sigma_0$后，外层积分收敛。
\end{proposition}
\begin{proof}
$\mu<t_0$保证$D_x(\mu)$在$[0,1]$上一致正。
于是$\Sigma_0(-q^2)=O(\log(2+q^2))$，并有精确分解
\[
 \frac{\Sigma_1(-q^2)}{(q^2+\mu)^2}
 =\frac{\Sigma_0(-q^2)}{(q^2+\mu)^2}
       +\frac{Z_\mu}{q^2+\mu}.
\]
第一项乘外层传播子在无穷远为$O(\log|q|/|q|^6)$，
四维径向积分收敛；正质量排除原点和有限点问题。
第二项恰为$Z_\mu$乘四维bubble。
由$\operatorname{vol}(S^3)=2\pi^2$，
其径向主项为$2\int^\Lambda dr/r=\log\Lambda^2$；
固定外动量只改变可积的余项。

这里有限的线性子反项改变了外层的大动量阶数。
若给子图选择了这一更高阶减除，便必须把它引起的外层减除一并处理。
不能以原五传播子图的表观总体收敛度为由丢弃该项，
也不能从一项子图已经有限推出整个森林组合已经完成。
\end{proof}

\begin{theorem}[四维插入的整体MOM与谱推前相容]\label{thm:fourinsertmom}
保持上述子图两项减除，外层在$p=0$作常数减除。则
\begin{align}
 J_{4,R}(s)
 &=\int\frac{d^4q}{\pi^2}\,g_\mu(-q^2)
       \left[\frac1{(p-q)^2+M^2}-\frac1{q^2+M^2}\right]\notag\\
 &=\int_{t_0}^{\infty}
       \frac{\sqrt{A(t)}\,dt}{t(t-\mu)^2}
                  C_{4,R}(s;t,M^2),\qquad s=-p^2,
 \label{eq:fourinsertmom}\\
 C_{4,R}(s;t,M^2)
 &=\int_0^1\log
       \frac{tx+M^2(1-x)}
            {tx+M^2(1-x)-s x(1-x)}\,dx .
\end{align}
所有欧氏积分绝对收敛，两种次序相等，$J_{4,R}(0)=0$。
在$s>\tau$的物理支上，
\begin{equation}\label{eq:fourinsertcut}
 \rho_{J,4}(s)=
 \int_{t_0}^{(\sqrt s-M)^2}
 \frac{\sqrt{A(t)\lambda(s,t,M^2)}}
      {s\,t(t-\mu)^2}\,dt>0,\qquad
 \operatorname{Disc}J_{4,R}=2\pi i\,\rho_{J,4}.
\end{equation}
给定该继承减除，子图有限方案差在整体减除后仍满足
\[
 \delta J_{4,R}
      =-a\,\partial_\mu C_{4,R}(s;\mu,M^2)
                       -b\,C_{4,R}(s;\mu,M^2).
\]
它没有把一般UV森林与任意谱积分的交换作为前提或结论。
\end{theorem}
\begin{proof}
四维bubble有限差的负对数在$t>t_0$出现负实分母区间；
该区间的$x$长度为$\sqrt{A(t)}/t$。
取$s+i0$时负对数的虚部为正$\pi$倍此长度，
故$\rho_{H,4}(t)=\sqrt{A(t)}/t$。
为核对一次减除表示，置
$T(x)=m_1^2/(1-x)+m_2^2/x$；
$\{x:T(x)<t\}$的长度恰为$\rho_{H,4}(t)$。
因此均匀$dx$的推前密度为$\rho_{H,4}'(t)$，并有
\[
 \Sigma_0'(s)=\int_0^1\frac{dx}{T(x)-s}
 =\int_{t_0}^\infty\frac{\rho_{H,4}'(t)}{t-s}\,dt
 =\int_{t_0}^\infty\frac{\rho_{H,4}(t)}{(t-s)^2}\,dt.
\]
最后一步分部积分的两端均为零。
从$\mu$积分至$s$即得命题所需一次减除式，
且$\int\rho_{H,4}(t)(t-\mu)^{-2}\,dt=Z_\mu<\infty$。
因此命题\ref{prop:subtractededge}仍适用。

外层方括号不是非负函数，不能直接援用Tonelli处理它的符号。
不过其绝对值为$O(|q|^{-3})$，
而$g_\mu(-q^2)\leq Z_\mu/(q^2+t_0)$。
乘积在四维无穷远可积，有限处由正质量控制。
对谱积分先应用非负的绝对值版本，所得上界相同，
遂由Fubini得到\eqref{eq:fourinsertmom}。
其中单圈差经Feynman参数合并及平移动量后，
径向积分的有限差恰给所列对数；截断球平移的边界差趋于零。
这一步固定了整体MOM归一化与符号。

在复$s$的避割紧集上，大$t$时$C_{4,R}(s;t,M^2)=O(1/t)$一致成立，
例如利用$x(1-x)/(tx+M^2(1-x))\leq1/t$控制对数。
故谱表示可全纯续接。
在有限物理$s$处，仅$t_0<t<(\sqrt s-M)^2$贡献虚部；
端点平方根可积，单圈虚部为$\pi\sqrt{\lambda(s,t,M^2)}/s$。
于是得到\eqref{eq:fourinsertcut}，支撑仍从$\tau$开始。
局部方案差则在已绝对收敛的整体差中逐点通分，
给与命题\ref{prop:edgescheme}相同的公式。

最后给出实际有理相对模型，而不只留下对数表示。
Taylor积分余项为
\begin{equation}\label{eq:logrationalremainder}
 \frac{\Sigma_1(s)}{(s-\mu)^2}
 =\int_0^1dx\,v^2
       \int_0^1\frac{(1-u)\,du}
              {[D_x(\mu)-u(s-\mu)v]^2}.
\end{equation}
在$s=-q^2$时分母正。右式乘外层有理传播子差，
由上述绝对收敛界形成实际有限周期积分。
把$(x,u,q)$共同紧化消解，给明确相对载体；
若比较不同方案，必须继承同一外层锚点并采用相容共同模型。
此构造没有从谱函数的解析等式直接宣布整个图动机的同构。
\end{proof}

\subsection{椭圆核的精确割重数与方案保护}
\begin{theorem}[插入图割中恰有一个椭圆简单因子]\label{thm:insertelliptic}
取正代数质量$m_1\ne m_2$，$0<\mu<a_+$且$\mu\ne a_-$，
并删除有限个分支或极点碰撞的参数值。
在外参数曲线上令
\[
 E_s:\quad y^2=Q_s(t):=A(t)\lambda(s,t,M^2).
\]
二维$J_{2,\rm os}$和四维$J_{4,R}$的割各由该椭圆族上的实际闭循环表示。
在$\mathbb C(s)$函数范畴，两个割的最小对象的非有限单值简单组成
均恰为一个秩二的$H^1(E_s)$（对偶及整体Tate扭转按所选方向记）。
完整插入函数的最小对象至少包含该因子。

这一因子经由一圈实际局部子反项引起的方案变化仍然保留。
不据此断言完整插入函数的秩、全部扩张、固定纤维Nori群，
或任何指定数值不属于MZV。
\end{theorem}
\begin{proof}
在$s>\tau$，四分支点满足
\[
 a_-<a_+<c_s=(\sqrt s-M)^2<d_s=(\sqrt s+M)^2.
\]
沿$[a_+,c_s]$的正实片走向右端，再从另一片反向返回，
给整循环$\gamma_s$。因$\mu<a_+$且质量非零，
下列微分的极点不在该循环上：
\[
 \omega_2=\frac{dt}{(t-\mu)^2y},\qquad
 \omega_4=\frac{Q_s(t)\,dt}{s\,t(t-\mu)^2y}.
\]
双片取向与$y$的反号相抵，故
\begin{equation}\label{eq:insertcutperiods}
 \rho_{J,2}=2\int_{\gamma_s}\omega_2,\qquad
 \rho_{J,4}=\frac12\int_{\gamma_s}\omega_4.
\end{equation}
两者严格正。注意这些是有极点的微分，
因此不能仅凭正性就把其周期称为非零全纯椭圆周期。
下面补齐这一可见性步骤。

设$D_s$为极点集：二维只需$t=\mu$的两点，
四维还允许$t=0,\infty$上的点。作必要有限代数基变换，
使极点为不交截面；正常阈值附近这些截面避开退化节点。
开曲线的同调正合列为
\begin{equation}\label{eq:openelliptic}
 0\longrightarrow B_s\longrightarrow
 H_1(E_s\setminus D_s,\Q)
 \longrightarrow H_1(E_s,\Q)\longrightarrow0 .
\end{equation}
$B_s$由绕极点小圈生成，唯一总和关系给
$\dim B_s=|D_s|-1$。其单值仅来自点的有限置换，
在Hodge语言中属于Artin--Tate边界。
属一纯局部系统$V=H_1(E_s,\Q)$只出现一次。

当$s\to\tau$时，$a_+$与$c_s$相碰，其余两个根保持不同。
四根交比趋于$0,1$或$\infty$，$j$有极点，
所以椭圆族非等常。
如前节三边割的模曲线论证，在有限基变换及去掉有限值后，
它给非恒定代数映射到带水平结构的模曲线；
再去掉分歧值，此映射为有限覆盖。
其单值为模曲线单值的有限指数子群，
故$V$的Zariski闭包为$\mathrm{SL}_2$，$V$绝对不可约。

以下将局部系统复化，以讨论$\mathbb C(s)$上的最小函数对象。
令$\mathcal R$为$\gamma$在整个参数域中的续接轨道所生成的局部子系统。
实双片消失圈在紧椭圆上的像非零，因而
$\mathcal R\to V$满射。
正常节点的小回路在开曲线上由一个Dehn twist的非零幂实现：
支撑可取在避开所有极点的Milnor圆盘内，
故其$N=T-1$满足
\[
 N^2=0,\qquad \im N=\mathbb C\gamma,\qquad \rk N=1.
\]
非零幂容许分支点碰撞参数所带的有限基变换，
不需要把其整数系数错误固定为$1$。
因为$\mathcal R$满射到$V_{\mathbb C}$，$N\mathcal R=\mathbb C\gamma$。

对任一读出$\omega_i$，令$\mathcal P\subset\mathcal R$
为其周期作为函数恒零的最大局部子系统；
于是割的最小函数对象由$\mathcal R/\mathcal P$承载。
若商中没有$V$，则由\eqref{eq:openelliptic}及$V$只出现一次，
$\mathcal P$必须满射到$V$。
这强制$N\mathcal P=\mathbb C\gamma\subset\mathcal P$，
与\eqref{eq:insertcutperiods}严格非零矛盾。
故下界是一份$V$；开曲线正合列给上界也是一份。
这里用的是$\gamma$的\emph{全局轨道}：
节点自身满足$N\gamma=0$，
并未把割在本阈值的自跳跃误称为非零。

完整$J$的两个解析分支之差为$2\pi i\rho_J$；
所以其分支生成空间包含割的分支生成空间，
完整最小函数对象至少含$V$。
一圈bubble、其有限质量导数和代数减除点的局部系数
由属零二次根的有理化及对数积分承载。
在函数层有限覆盖后的简单因子均为常数，
在相应实际几何模型中属于Artin--Tate背景；
其质量导数通过同一一圈载体表达。
命题\ref{prop:edgescheme}和定理\ref{thm:fourinsertmom}
说明方案差完全由这些一圈背景给出，
不会消去所证非等常椭圆因子。
若允许任意非Tate局部常数，则必须按
命题\ref{prop:edgescheme}重新声明算术背景，
不能把函数层常数直接降为Tate数值。
\end{proof}

\subsection{加入结构表的操作规则}
本节给一条已经连接到真实图操作的有限推导链：
\[
 \text{被标记边与二点子图}
 \ \longrightarrow\ \text{两项局部减除}
 \ \longrightarrow\ \frac{\rho_H(t)\,dt}{(t-\mu)^2}
 \ \longrightarrow\ \text{外图推前与物理割}.
\]
对收敛正外图，有限谱矩给充分且直接可用的收敛控制；
对四维bubble外图，必须继承整体MOM，反例已排除无条件交换。
五边图的割由四根与极点位置预测为带Tate边界的椭圆载体，
全局链轨道和节点秩一变分再把载体上界提升为精确割重数。
这同时说明：局部反项能消去某些图贡献中的阈值，
却未必能消去新生的非等常纯核。
一般带符号谱、分子耦合、多个重叠子图、
任意幂次森林、完整最小混合对象及实际RG出口仍保留原适用条件；
不能把这一完整两步实例宣布为全部一般图分类。


\section{多重局部插入、相容森林与非线性物理读出}
\label{sec:forestobservable}

上一节的一条局部子反项在外图中成为$-a\partial_\mu C-bC$，
而不是整个外动量的局部多项式。
本节把这一关系推进到任意有限串联数及多个边槽，再接入森林与物理图和。
最后处理从振幅到乘积型可观测量时的新问题：
振幅的背景不能原封不动地用作其平方的背景。
所得是有接受条件的操作定理，不是任意幂次森林的全局收敛定理。

\subsection{边槽质量与子图物理质量必须分开}

固定外图的有理动量路由$q_e(k,p)$，选若干标记边槽$e$。
在边$e$上串联$n_e$个已有二点块，编号集为$I_e$，置
\[
 d_e=q_e^2+\mu_e,\qquad
 F_i=\Sigma_i(-q_e^2)\quad(i\in I_e),\qquad I=\coprod_e I_e.
\]
串联的真实分母是$d_e^{\,n_e+1}$。
把某些块作有限局部改变
\begin{equation}\label{eq:multijetlocal}
 F_i\longmapsto F_i+z_i(a_i-b_i d_e).
\end{equation}
这里$a_i,b_i$不依赖该块的当前外动量$q_e$，但可以是质量及参考点的周期。
$z_i$只用来记录同时改变了哪些块，不是新的物理耦合。

令$K(k,p)$含其余传播子、指定分子及测度。
对$S\subseteq I$，令$m_e=|S\cap I_e|$、$p_e=n_e-m_e+1$，定义
\begin{equation}\label{eq:slotbase}
 C_{I\setminus S}(\nu)=
 \int {\mathcal R}\left[
 K\prod_e
 \frac{\prod_{i\in I_e\setminus S}F_i}
 {(q_e^2+\nu_e)^{p_e}}\right].
\end{equation}
辅助量$\nu_e$只替换显示出的边槽分母；
它不进入$F_i,a_i,b_i$，也不进入$\mathcal R$的参考条件。
物理质量仍是$\mu_e$。先作$\partial_{\nu_e}$，最后取$\nu=\mu$。
这样定义的导数是一个明确传播子幂插入，不是对整个图的所有质量依赖求全导数。

本小节的解析接受条件为：$\mathcal R$是固定线性的继承减除操作，
与这些辅助导数交换；公式所用的有限个合并密度和导数在$\nu=\mu$
的邻域有可积支配函数。可取$\mathcal R=1$且验证绝对收敛。
也可取已经证明满足该条件的共同外层森林。
不允许在每一项上改用其各自最小表面发散度的减除。

\begin{theorem}[有限串联插入的显式局部标准式]\label{thm:multijet}
在上述接受条件下，令
\[
 J(z)=\int\mathcal R\left[
 K\prod_e\left\{d_e^{-n_e-1}
       \prod_{i\in I_e}(F_i+z_i(a_i-b_i d_e))\right\}\right].
\]
取$z_S=\prod_{i\in S}z_i$，对$A\subseteq S$置$j_e=|A\cap I_e|$。
则
\begin{equation}\label{eq:multijet}
 [z_S]J(z)=(-1)^{|S|}
 \sum_{A\subseteq S}
 \left(\prod_{i\in A}a_i\prod_{i\in S\setminus A}b_i\right)
 \left.
 \prod_e\frac{\partial_{\nu_e}^{j_e}}{(p_e)_{j_e}}\,
 C_{I\setminus S}(\nu)\right|_{\nu=\mu},
\end{equation}
其中$(p)_j=p(p+1)\cdots(p+j-1)$、$(p)_0=1$。
当全部块都换成局部项时$p_e=1$，故分母变成$j_e!$。
若各块及继承操作已有共同有理相对模型，且各有限系数有实际周期载体，
该等式亦是这些模型中的形式周期等式。
\end{theorem}
\begin{proof}
先在被积函数层取系数。对$A$中的块选$a_i$，对其余$S\setminus A$中的块选$-b_i d_e$。
第$e$边的分母幂因此为
\[
 n_e+1-(m_e-j_e)=p_e+j_e.
\]
另一方面
\[
 \partial_{\nu_e}^{j_e}(q_e^2+\nu_e)^{-p_e}
   =(-1)^{j_e}(p_e)_{j_e}(q_e^2+\nu_e)^{-p_e-j_e}.
\]
两个符号相乘为
$(-1)^{|S|-|A|}(-1)^{|A|}=(-1)^{|S|}$，得式\eqref{eq:multijet}。
这是有限多项式恒等式。线性继承操作、导数交换及所假定的统一可积性
依次允许作$\mathcal R$、积分与有限阶导数，证明解析等式。

形式周期部分不由解析等式反推。
将各块的已给有理参数与外动量同时保留，式中的分母恒等式仍逐点成立。
系数的乘积由乘积链与外积形式实现。
共同有理紧化和极点--边界消解后，各有理形式及相对边界关系落在共同模型中；
有限求和、实际拉回和相对Stokes给同一等式。
这一结论明确要求所述共同模型，不把任意独立提出的参数化自动识别。
\end{proof}

两个块位于同一边、同时换成局部项时，公式是
\begin{equation}\label{eq:twojets}
 \frac{a_1a_2}{2}\partial_\nu^2 C
 +(a_1b_2+a_2b_1)\partial_\nu C+b_1b_2 C.
\end{equation}
这里的$1/2$不能遗漏。若只改变第一个、保留第二个，
基函数的边幂为$2$，得到$-a_1\partial_\nu C_{\{2\}}/2-b_1C_{\{2\}}$，
也不能套用单一空边的$-a\partial_\mu C-bC$。
这些是任意串联长度的闭式规则，不需求解任一完整积分。
若减除参考点随$\nu$移动，则有额外链式项，本定理不适用未经修正的全导数。

\subsection{森林中哪些乘积确实出现}

森林的相容条件是：任意两个成员不交或包含；真重叠成员不能同时出现。
这是经典森林公式的定义\cite[\S2--3.2]{BK}。
对数图的有限character比较已经在前章证明；一般幂次情形须保留局部张量或多项式装饰，
不能把每个反项都换成一个标量。

将一个已给出的有限局部方案比较称为本节的\emph{接受森林比较}，若满足：
每个嵌入及局部装饰有明确标记；每个系数是已经有限的实际形式周期；
每个装饰商图采用规定的继承减除且有实际周期模型；
整个比较由共同相对形式恒等式或已经证明的有限character关系给出。
这四项是输入责任，不以森林符号本身代替收敛及几何证明。
已完成的正质量至多对数MOM图类满足相应零次装饰版本。
定理\ref{thm:multijet}为含二点仿射装饰的更多比较提供具体接受办法。

写任一接受比较的有限正常式为
\begin{equation}\label{eq:decorateddifference}
 p_b-p_a=
 \sum_{\mathfrak f,\boldsymbol\alpha}
 c_{\mathfrak f,\boldsymbol\alpha}\,
 D_{\mathfrak f,\boldsymbol\alpha}^{\mathfrak m}
 (G/\mathfrak f).
\end{equation}
$\mathfrak f$包括嵌入与递归的相容收缩记录；
$c_{\mathfrak f,\boldsymbol\alpha}$是相应局部系数之积，
$D^{\mathfrak m}$是带局部装饰的实际商图周期，
可以是式\eqref{eq:multijet}的边槽导数。
嵌套记录中父块的系数必须是在其已处理子块之后得到的系数；
它一般不是原父图和子图裸系数的乘积。
式\eqref{eq:decorateddifference}是对\emph{已建立}比较的记账，不另假定任意森林都已具备该形式提升。

\begin{proposition}[相容森林的实际背景与重叠排除]\label{prop:forestcarrier}
给式\eqref{eq:decorateddifference}中每个系数因子$c_v$选实际载体$E_{c_v}$，
给相应装饰商图选载体$E_D$。则整个差有实际载体
\begin{equation}\label{eq:forestcarrier}
 L_{ba}(G)=
 \bigoplus_{\mathfrak f,\boldsymbol\alpha}
 \left(\bigotimes_{v}E_{c_v}\right)\otimes E_D.
\end{equation}
故在含这些\emph{整项乘积载体}的Serre背景$\mathcal S$中，
$\pi_\mathcal S M(p_a)\simeq\pi_\mathcal S M(p_b)$。
对真重叠的$\gamma,\eta$，不能额外加入
$c_\gamma c_\eta\,p(G/(\gamma\cup\eta))$作为它们的共同森林项。
若$\gamma\cup\eta$自身可减除，其系数是独立的准备后系数，仍需实际构造。
\end{proposition}
\begin{proof}
有限系数乘积及其与商图的乘积由实际张量积承载，
有限和由直和承载，故式\eqref{eq:forestcarrier}确实给差的标架。
应用定理\ref{thm:addition}得Serre商同构。
森林成员的相容条件排除了所述同时选择；不能从并集的存在推出两个重叠减除算子可同时乘上。
递归只在互不交的最大成员上作同时收缩，在每个成员内部继续处理其自身森林，
因此每条递归记录仍满足相容条件。
并集若许可，只会以自己的节点及自己的系数出现，不能与被排除的乘积混同。
\end{proof}

一个只有两项可选的真重叠减除的森林展开含空项及各自单项，
而不含二者同时项；这不禁止有限系数在后续代数化简中偶然出现同样的数值乘积。
所禁止的是没有嵌入与共同模型依据的图操作。
此外，只将所有单独$E_{c_v}$与$E_D$的简单组成列为背景是不够的：
例如两个椭圆因子的张量可出现新的对称平方或外积组成。
应保留式\eqref{eq:forestcarrier}，或取这些输入生成的张量范畴再作Serre闭包。

\subsection{先带图权求和，再谈物理保护}

固定有限阶物理组合及同一外部标记，令
\begin{equation}\label{eq:weightedobservable}
 A_a^{\mathfrak m}=\sum_{G\in\mathcal G}w_G\,p_a^{\mathfrak m}(G).
\end{equation}
有理图权$w_G$包括已核对的对称因子、符号和固定阶耦合系数；
允许的线性投影须已作用于$p_a(G)$的实际标架，不能混作数值权重。
若权重本身是周期，必须把其实际载体也放入下面的张量积。
这里比较相同权重和相同物理坐标；耦合、场或算符的有限换参
若改变这些条件，须先代回相同坐标并列出新增项。
本式不以逐图归一化条件自动代替一个拉氏量的完整重整化。

\begin{theorem}[带权物理和的实际方案背景]\label{thm:weightedforest}
若每个图均有接受比较及载体$L_{ba}(G)$，则
\[
 A_b^{\mathfrak m}-A_a^{\mathfrak m}
 \quad\hbox{由}\quad L_A=\bigoplus_{G\in\mathcal G}L_{ba}(G)
 \quad\hbox{承载}.
\]
故在含$L_A$的Serre背景外，\emph{整个物理和}的简单组成及重数不随此次方案变化。
若$E=M(A_a^{\mathfrak m})$无背景子对象，则另有
\[
 0\longrightarrow B_0\longrightarrow M(A_b^{\mathfrak m})
   \longrightarrow E\longrightarrow0,\qquad B_0\in\mathcal S.
 \tag{\ref{thm:weightedforest}a}
\]
如果两端都无背景子对象，则两端完整对象同构。
这些结论不等于把各图的组成重数相加。
\end{theorem}
\begin{proof}
用各个$w_G$改变差载体的读出标架，形成直和，便得到差的实际表示；
有周期权重时先张量其载体。加性稳定性给商及重数结论。
无背景子对象的条件正是定理\ref{thm:fullbackground}的socle分离条件，
应用该定理给所示完整商及两端同构。
必须对总读出应用这些结果：同一周期$p$的两份图贡献取权$1,-1$
即可使总周期为零，虽两份最小对象都非零。
所以逐图非消失不能替代总读出的非消失。
\end{proof}

这一表述可与前文的物理割联合使用：
在总读出上先证明某个割非零且其链来自实际物理链，
再用$L_A$的结构证明该因子不属于背景。
只有这两个输入同时成立，才得到该物理组合不可消去的几何结构。
标量理论没有额外的规范Ward恒等式可代为提供这个相消检查。

\subsection{非线性读出需要按张量次数扩大背景}

现固定中性Tannakian环境$\mathcal C$，以及对张量、对偶、子商封闭的背景子范畴
$\mathcal B$，记其在$\mathcal C$中的Serre闭包为$\mathcal S_0$。
选择有限个实际输入载体$E_1,\ldots,E_r$。
在乘积读出中若用到不同边值或对偶，应预先将对应载体也列入这份输入。
定义$\mathcal S_d$为下列有限直和的子商、扩张所生成的Serre类别：
\begin{equation}\label{eq:tensordegree}
 B\otimes E_{i_1}\otimes\cdots\otimes E_{i_j},
 \qquad B\in\mathcal B,\quad 0\le j\le d.
\end{equation}
次数记录使用了多少个指定输入，不是圈数、权重或几何维数。
由张量正合性及过滤的逐层张量，
\[
 \mathcal S_d\otimes\mathcal S_e\subseteq\mathcal S_{d+e}.
\]
这里不是要求$\mathcal S_0$为整个环境的张量理想。

令$p_i$由$E_i$承载，$q_i=p_i+\ell_i$且$\ell_i$由$\mathcal S_0$承载。
多项式$P$的系数有$\mathcal B$内的实际载体，所有乘积均已通过实际外积及所需对角拉回定义。
系数若含额外非背景周期，应把它算入输入次数；不能暗中当有理权。

\begin{theorem}[非线性物理读出的最高张量层保护]\label{thm:nonlinearbackground}
若$P$的总次数至多$d\ge1$，则
\begin{equation}\label{eq:nonlinearbackground}
 P(q_1,\ldots,q_r)-P(p_1,\ldots,p_r)
       \quad\hbox{有}\ \mathcal S_{d-1}\ \hbox{内的实际载体}.
\end{equation}
从而在$\mathcal C/\mathcal S_{d-1}$中，两种读出的最小对象同构；
背景外各简单因子的重数相同。
该结论无需知道所有$p_i$的完整积分。
\end{theorem}
\begin{proof}
把每个单项式中的每个$q_i$替换为$p_i+\ell_i$并展开。
减去原单项式后，每个留下的项至少含一个$\ell_i$，
因而至多含$d-1$个$E_i$输入。
若$\ell_i$仅属于$\mathcal S_0$而不属于$\mathcal B$，
取其有限长度过滤；逐层张量的简单组成仍由背景张量载体给出。
张量正合性因此保证整项在$\mathcal S_{d-1}$。
实际有限和及乘积给式\eqref{eq:nonlinearbackground}，再应用定理\ref{thm:addition}。
证明没有将$P$视为线性物理投影，也没有把Serre商当作带张量结构的动机范畴。
\end{proof}

特别地，对一个已有实际边值模型的二次读出$A_+A_-$，
假定其在物理边界确实等于$|A|^2$。
若两边振幅各加背景$\ell_\pm$，差为
\begin{equation}\label{eq:crossbackground}
 A_+\ell_-+\ell_+A_-+\ell_+\ell_-.
\end{equation}
这些项一般不在原来的$\mathcal S_0$，而在$\mathcal S_1$。
一个最简单的反例是$f^2$与$(f+1)^2$之差$2f+1$：
若$f$有非背景简单因子，加性稳定性保证差中仍有它。
因此振幅背景定理不能直接照抄为截面背景定理。
式\eqref{eq:crossbackground}只是逐点乘积比较；
跨物理割的复共轭必须由已给的两种实际边值及实结构实现，
不能把任意复数共轭直接声明为原域上的Nori态射。

\subsection{一个精确可读的椭圆高层实例}

以下只在$\mathbb C(s)$的微分模及其Picard--Vessiot实现中陈述，
不作固定代数点的算术下降。
设$V$为秩二不可约正则奇异标量方程的解表示，微分Galois群为$\mathrm{SL}_2$；
其解空间作为函数嵌入一个Picard--Vessiot域。
$f$为其中任一非零解。
背景也限于正则奇异对象，其单值简单组成假定有限；取$\mathcal B$为有限单值对象，
$\mathcal S_0$允许其任意有限扩张。

\begin{proposition}[纯椭圆读出幂的不可消去高层]\label{prop:elliptichighlayer}
对$d\ge1$，$f^d$的最小微分对象为$\operatorname{Sym}^d V$，秩$d+1$。
若$h$有$\mathcal S_0$内的载体，且$P(X)$是次数恰为$d$、
首项系数非零的$\mathbb C(s)$多项式，则$P(f+h)$中
$\operatorname{Sym}^d V$出现恰好一次，不能被其低于$d$次的背景改变消去。
因此其最小齐次微分阶数至少为$d+1$。
这里$\mathcal S_{d-1}$按式\eqref{eq:tensordegree}由$E_1=V$生成。
\end{proposition}
\begin{proof}
函数乘法给$\mathrm{SL}_2$等变映射
$\operatorname{Sym}^d V\to L$，其中$L$是相应Picard--Vessiot域。
它非零，因为$f^d\ne0$；$\operatorname{Sym}^d V$不可约，故此映射单射。
非零$f^d$的群轨道生成整个像，所以最小解空间恰有维数$d+1$。
乘非零有理函数只作一维平凡微分对象的规范改变，不改变同构类。

有限单值背景在某个有限指数子群上恒等。
正则奇异性使单值群在微分Galois群中稠密\cite{Singer}；
该有限指数子群在$\mathrm{SL}_2$中的像仍Zariski稠密，
所以任何含至多$d-1$份$V$的张量，其不可约$\mathrm{SL}_2$组成的最高权
都不超过$d-1$。背景扩张不增加新的简单组成。
因此$\operatorname{Sym}^d V\notin\mathcal S_{d-1}$。
将$P(f+h)$与其非零首项$a_df^d$相减，展开后每项至多含$d-1$份$f$；
其余因子均在背景内。定理\ref{thm:addition}于是给出该简单因子重数恰为一。
秩下界即是所述齐次阶数下界。
\end{proof}

例如平方保留一份秩三的$\operatorname{Sym}^2 V$，
但允许新增秩二$V$与背景方向；并非整个平方的最小阶数恒为三。
该命题要求读出确实是纯方程的解。
上一节的穿孔椭圆割以及完整相对振幅还可能含边界和非分裂扩张，
不能只因其含一个椭圆组成就机械套用$f^d$的精确分类。
同样，积分后的总截面另有相空间推前：
被积函数的背景在推前后可能产生新曲线或更高维几何，
须为\eqref{eq:crossbackground}每项建立实际推前载体，
不能把逐点的$\mathcal S_{d-1}$结论不加说明地用于全相空间积分。

\subsection{与总目标的实际接合}

现在有一条明确的有限推导链：
局部二点装饰经式\eqref{eq:multijet}成为有正确组合系数的传播子幂插入；
接受森林保留嵌入、嵌套次序与重叠排除；
整项张量载体给实际方案背景；真实图权先形成总读出；
乘积型可观测量再使用次数降低一层的扩大背景。
这些规则说明哪些结构受保护，也说明保护何处必须重新证明。

本节没有补出所有装饰商图的全局UV收敛、一般物理边值的动机比较、
相空间积分的完整推前或任意幂次森林的共同相对模型。
其无条件实际输入仍来自前文已建立的图类及通过本节接受条件的实例。
经典森林及有限方案理论的优先权属于原文献；
此处逐式给出的标准式和可观测量分层应接受独立数学审阅，
不以内部推导代替系统新颖性检索。


\section{实际相空间推前与可观测量的新几何}\label{sec:phasepush}
上一节的平方和干涉仍是逐运动学点的读出。
本节把最后一次物理积分本身作为几何操作：
先构造有限相空间的实际周期，再证明一个二体接受规则和一个三体椭圆下界。
这补充的是可观测量的结构预测，不能用振幅所属的范畴直接替代。

\subsection{规范、正实链与实际代数载体}
固定四维度规$(+---)$、外态归一化
$\langle p'|p\rangle=(2\pi)^3 2E_p\delta^3(p'-p)$，并定义
\begin{equation}\label{eq:physicalphin}
 d\Phi_n(P)=(2\pi)^4\delta^4(P-\textstyle\sum_i k_i)
       \prod_{i=1}^n\frac{d^3\mathbf k_i}{(2\pi)^3\,2E_i},
 \qquad E_i=\sqrt{\mathbf k_i^2+m_i^2}.
\end{equation}
这里已把$(2\pi)^4$放进$d\Phi_n$。
PDG的式(49.12)将此因子留在率公式外；转换后两者相同
\cite[式(49.8)--(49.13)]{PDGphase}。
以下使用有标记末态相空间，同种末态的阶乘另除，不再重复除入测度。
令$\lambda(a,b,c)=(a-b-c)^2-4bc$，两入射粒子的通量为
$2\sqrt{\lambda(s,M_1^2,M_2^2)}$。

\begin{theorem}[有理加权相空间的实际周期载体]\label{thm:actualphase}
取$n\ge2$、正代数质量$m_i$和代数$s>(\sum_i m_i)^2$。
在含$\sqrt s$的实数域$k$上置$P=(\sqrt s,0,0,0)$，并令
\[
 X_s=\{(k_1,\ldots,k_n): k_i^2=m_i^2,\ \sum_i k_i=P\}
       \subset\mathbb A_k^{4n},\qquad N=3n-4.
\]
其正能量实分支$\Gamma_s$是紧致、无边界的光滑定向$N$维流形，
包含在$X_s$的光滑部分中。
有按约束顺序定向的代数留数$N$形式$\Omega_n$，使
\begin{equation}\label{eq:phasecarrier}
 (2\pi)^{3n-4}\int d\Phi_n\,W
       =\int_{\Gamma_s}W\Omega_n
\end{equation}
对一切在$\Gamma_s$上无极点的$k$-有理权$W$成立。
右端是实际代数上同调周期；删去极点并紧化消解后，
它给明确的Nori载体上界。左端的物理$\pi$因子在允许Tate逆的范畴中处理；
这不是关于其最小对象的等式断言。
\end{theorem}
\begin{proof}
正能量条件给$m_i\le E_i\le\sqrt s$，故所有空间动量有界。
在约束集合中$E_i$不会趋于零，所以正能量分支是闭的，从而紧致。
为检查约束独立性，假设质量壳约束和四个守恒约束的微分有线性关系。
消去公共四维乘子后，每个$k_i$与同一个四向量成比例；
若某比例系数为零，则全部乘子已为零。
因此非平凡关系只能在所有四速度$k_i/m_i$相等时出现。
正能量及总空间动量为零迫使这些速度均为$(1,0,0,0)$，
于是$\sqrt s=\sum_i m_i$，与严格阈值条件矛盾。
Jacobian因而满秩$n+4$。隐函数定理给维数$N$及无边界性。

以$d^{4n}k$对这$n+4$个方程取代数Poincaré留数；
其定义不额外含$(2\pi i)^{-1}$。
在先消去$k_i^0$的实坐标图中，每个正根的Jacobian为$2E_i$；
随后对四个守恒方程取Gelfand--Leray形式。
选择相容取向后，所得$\Omega_n$的实限制正是
$\delta^4(P-\sum k_i)\prod_i d^3\mathbf k_i/(2E_i)$。
由式\eqref{eq:physicalphin}即得\eqref{eq:phasecarrier}。
这个处处非零的实顶形式同时给$\Gamma_s$定向。

取$U_s\subset X_s^{\rm sm}$删去$W$及形式的极点。
$\Gamma_s$是$U_s(\mathbb C)$中的整同调循环，$W\Omega_n$是代数de Rham类。
故右端直接是$H^N(U_s)$与$H_N(U_s)$的比较配对。
必要的紧化、消解和相对边界给等价的实际周期表示，
不以数值近似反推对象存在。紧致性和无极点性保证绝对可积。
将$(2\pi)$与$2\pi i$比较只需代数因子和Tate扭转；
负Tate次幂不能被默认为有效范畴中的对象。
\end{proof}

这一定理也规定了有理树振幅的实际入口：
先求全树图之和$\mathcal A$，再以$W=\mathcal A^2$进入
\eqref{eq:phasecarrier}，保留全部干涉。
实质量与实耦合、无物理极点的树振幅满足此处实结构条件；
代数周期陈述针对剥离耦合后的微扰系数，或明确取代数耦合。
本书独立$\phi^4$四体基准的十个$3|3$树项已有全物理域无极点证明。
因此其领先$\sigma_{24}$具有实际的$N=8$相空间载体，权为
$(\sum_A(Q_A^2-m^2)^{-1})^2$，并另乘通量和$1/4!$。
这个结论不需重作采样；也没有把该八维代数空间叫作K3，
或声称求出了其最小混合对象。

更一般地，若两个振幅边值已经给出在相空间$X$上的实际有理相对积分模型
$Y_\pm\to X$，就取纤维积$Y_+\times_XY_-$及其双链。
只有在链族相容、所有边界已指定、总积分绝对可积时，
才以总空间的有理形式和总链定义乘积后的相空间周期。
Fubini给它与指定迭代物理积分的相等，拉回及链推前给其实际几何含义。
这是一条带输入的构造规则：
任意复共轭、形式函数乘法或未经控制的阈值极限，都不是这些输入的替代品。

\subsection{二体接受规则：有理角权只增加对数结构}
\begin{theorem}[四维二体有理角积分]\label{thm:twobodyrational}
考虑标量$2\to2$，两端质量均为正，且$s$在两端阈值之上。
设全部角依赖为$z=\cos\theta$的有理函数$W(s,z)$，
在$[-1,1]$上无极点；系数在由运动学平方根生成的代数域中。
设末态对称因子为$S_{\rm id}$。则
\begin{equation}\label{eq:twobodyrule}
 \sigma_W(s)=
 \frac{\sqrt{\lambda(s,m_1^2,m_2^2)}}
 {32\pi s\,S_{\rm id}\sqrt{\lambda(s,M_1^2,M_2^2)}}
 \int_{-1}^1 W(s,z)\,dz .
\end{equation}
在分裂角分母的有限代数扩张上，去除显示的通量与Tate因子后，
结果是代数项及代数数（或代数函数）之对数的有限线性组合。
这给Artin--Tate及Kummer的载体上界，不产生新的椭圆角积分因子。
\end{theorem}
\begin{proof}
守恒先消去第二个空间动量；能量$\delta$函数再消去径向变量。
写$p=\sqrt{\lambda(s,m_1^2,m_2^2)}/(2\sqrt s)$，直接得到
\[
 d\Phi_2=\frac{p}{16\pi^2\sqrt s}\,d\Omega
        =\frac{\sqrt{\lambda(s,m_1^2,m_2^2)}}{32\pi^2s}\,d\Omega .
\]
绕入射轴的方位角积分为$2\pi$，再除通量和$S_{\rm id}$，
即得\eqref{eq:twobodyrule}。该推导也与
\cite[式(49.16)--(49.18)]{PDGphase}在上述测度转换后吻合。

在有限分裂域中作部分分式
\[
 W=P(z)+\sum_{\alpha}\sum_{j\ge1}\frac{c_{\alpha,j}}{(z-\alpha)^j}.
\]
多项式与$j\ge2$项给有理原函数的端点差；
$j=1$项给$c_{\alpha,1}\int_{-1}^1dz/(z-\alpha)$，
即由实际区间连续确定分支的对数端点差。
用单个商的对数书写时可能另加$2\pi i$的整数倍，不能任意取主支。
这一过程发生在穿孔$\mathbb P^1$与两个代数端点的相对上同调中，
所以结论是实际载体上界，不只是一种函数记号。
在原域的下降保留极点的有限置换，因此须允许Artin部分。
若质量和$s$变化，则先删去极点碰撞及端点相碰值，
在相应有限覆盖上同样成立。
\end{proof}

系数若本身是非Tate周期，则它们的载体张量必须保留，
定理只控制新增的角积分几何。
有理性、无区间极点和标量角依赖都不可省：
它不覆盖一般圈振幅、任意代数角权、共振正规化或自旋投影。
在本书$\phi^4$树级弹性情形，$W=g^2$、$S_{\rm id}=2$，
式\eqref{eq:twobodyrule}给$\sigma=g^2/(32\pi s)$，与原规范一致。

\subsection{三体相空间：物理链确实看见椭圆因子}
三体相空间的椭圆表达属于经典结果
\cite[式(1.1)--(1.7)]{DDphase}。
下面独立推导本书规范下的积分，并把
“椭圆曲线出现”加强为“指定物理函数的简单因子确实存在”。

\begin{theorem}[三体总相空间的精确椭圆组成]\label{thm:threebodyvisible}
取正代数质量$m_1,m_2,m_3$，假定$m_1\ne m_2$，置
\[
 a_\pm=(m_1\pm m_2)^2,\quad M=m_1+m_2+m_3,\quad \tau=M^2,
 \qquad Q_s(t)=\lambda(t,m_1^2,m_2^2)\lambda(s,t,m_3^2).
\]
删去有限个分支碰撞和以下极点截面碰撞的$s$值。
在$s>\tau$的物理支上，
\begin{equation}\label{eq:threephase}
 \Phi_3(s)=\frac{1}{128\pi^3s}
       \int_{a_+}^{(\sqrt s-m_3)^2}\frac{\sqrt{Q_s(t)}}{t}\,dt
     =\frac{1}{256\pi^3s}\int_{\gamma_s}\frac{y\,dt}{t},
 \qquad E_s:\ y^2=Q_s(t).
\end{equation}
$\gamma_s$是连接中间两根的双片消失圈。
在$\mathbb C(s)$上的正则奇异函数范畴，
$\Phi_3$的最小对象中非有限单值的简单组成恰为一份
$V=H_1(E_s,\mathbb C)$；$V$秩二且几何单值Zariski稠密于
$\mathrm{SL}_2$。
其余简单组成均有有限单值，整个函数阶数在$2$与$5$之间。
这里不声称完整混合扩张已确定，也不作固定点数值非MZV断言。
\end{theorem}
\begin{proof}
插入中间四动量$q=k_1+k_2$及$t=q^2$。
两个正能量质量壳之间的恒等分解把实际链分解为两次二体相空间，
在本测度中为
\[
 \Phi_3(s)=\int_{a_+}^{(\sqrt s-m_3)^2}
        \frac{dt}{2\pi}\,
        \Phi_2(s;t,m_3^2)\Phi_2(t;m_1^2,m_2^2).
\]
二体总相空间为$\Phi_2(s;u,v)=\sqrt{\lambda(s,u,v)}/(8\pi s)$，
故得到\eqref{eq:threephase}第一个等式。
其中正平方根同时固定了真实正链和取向。
经典文献的$I_3$未含任何$2\pi$因子；
两规范满足$\Phi_3=(2\pi)^{-5}I_3$，不是把其$I_3$直接当作$\Phi_3$。

四个根在物理域满足
\[
 a_-<a_+<c_s=(\sqrt s-m_3)^2<d_s=(\sqrt s+m_3)^2 .
\]
光滑紧化$E_s$属一。沿$[a_+,c_s]$的正片前进，再沿负片返回，
得到整圈$\gamma_s$；形式的反号与路径反向相消，
所以闭圈积分是实区间积分的两倍，且严格正。
$\eta=y\,dt/t$一般是有极点的微分，不能仅靠正性宣布其全纯投影非零。

为补上这一步，令$D_s$包含$t=0,\infty$上的全部点，一般共四点。
在有限代数覆盖上这些点为互异截面，均避开正常阈值节点。
开曲线同调给
\begin{equation}\label{eq:phaseopen}
 0\longrightarrow B\longrightarrow
 H_1(E_s\setminus D_s,\mathbb C)\longrightarrow V\longrightarrow0,
 \qquad \rk B\le3.
\end{equation}
边界由小环生成，单值只是有限点置换。
在正常阈值，$a_+$与$c_s$碰撞，其余两根不同。
四根交比的退化使$j(E_s)$有极点，因而$j$非恒定。
取水平结构的有限基变换后，到模曲线的代数映射非恒定；
删去分歧值后它是有限覆盖。
由此全局单值为模曲线单值的有限指数子群，其Zariski闭包为
$\mathrm{SL}_2$，$V$绝对不可约。有限覆盖不会改变该恒等分支。

令$\mathcal R$为$\gamma$的全局续接轨道在
$H_1(E_s\setminus D_s,\mathbb C)$中生成的局部系统。
$\gamma$在紧曲线上非零，故$\mathcal R\to V$满射。
局部退化的Dehn twist支撑避开$D_s$；
允许碰撞参数的非零整数绕数后，$N=T_\tau-1$满足
\[
 N^2=0,\qquad \im N=\mathbb C\gamma,\qquad N\gamma=0.
\]
交数只取决于紧曲线上的像。因此任何满射到$V$的子局部系统，
经$N$的像均为$\mathbb C\gamma$。

设$\mathcal J\subset\mathcal R$是$\eta$周期恒零的最大子局部系统。
实际函数的分支空间为$\mathcal R/\mathcal J$。
若该商没有$V$，由\eqref{eq:phaseopen}中$V$仅出现一次，
$\mathcal J$必须满射到$V$，于是
$\mathbb C\gamma=N\mathcal J\subset\mathcal J$。
这与\eqref{eq:threephase}严格非零矛盾。
故$V$确实保留，重数上下界都是一。
开曲线总秩至多$2+3=5$，Gauss--Manin系统正则奇异，
所以函数的最小齐次阶数亦在$2$与$5$之间。
前因子$1/s$及非零复常数不改变此函数结论。

注意下界使用全局轨道，不是$\gamma$在自身阈值的跳跃：
后者恰有$N\gamma=0$。
同样，有限单值边界的存在允许第三类微分和混合扩张，
本定理并未把它们删去。
\end{proof}

\begin{proposition}[阈值首项与树级接触总率的椭圆下界]\label{prop:contactphase}
上述三体相空间在$\epsilon=s-M^2\downarrow0$满足
\begin{equation}\label{eq:ph3threshold}
 \Phi_3(M^2+\epsilon)=
  \frac{\sqrt{m_1m_2m_3}}{256\pi^2 M^{7/2}}\,
       \epsilon^2(1+O(\epsilon)).
\end{equation}
取五种不同的实标量，有限阶有效相互作用为
$\mathcal L_{\rm int}=-g\,\Phi_1\Phi_2\chi_1\chi_2\chi_3$，
其中$m_i$为三个$\chi_i$的质量，入射质量满足$M_1+M_2<M$。
在树级$\Phi_1\Phi_2\to\chi_1\chi_2\chi_3$，
\begin{equation}\label{eq:contactrate}
 \mathcal M_{\rm tree}=-g,\qquad
 \sigma_{\rm contact}(s)=
       \frac{g^2\Phi_3(s)}{2\sqrt{\lambda(s,M_1^2,M_2^2)}}.
\end{equation}
因此一个无内部边、无圈的常数振幅，
经实际总率投影后保留一份椭圆简单因子。
此表述在分裂通量平方根的有限参数覆盖上成立；
原参数域上应保留相应的代数秩一扭曲。
\end{proposition}
\begin{proof}
记$b=m_1+m_2$，$\Delta=c_s-a_+$。
由$c_s=(\sqrt s-m_3)^2$得
$\Delta=b\epsilon/M+O(\epsilon^2)$。
代入$t=a_++\Delta u$，$0\le u\le1$。
正常阈值处
\[
 a_+-a_-=4m_1m_2,\qquad d_\tau-a_+=4m_3M,\qquad a_+=b^2.
\]
所以一致于闭区间的展开给
\[
 \int_{a_+}^{c_s}\frac{\sqrt{Q_s(t)}}t\,dt
 =\frac{4\sqrt{m_1m_2m_3M}}{b^2}\Delta^2
       \int_0^1\sqrt{u(1-u)}\,du\ (1+O(\epsilon)).
\]
端点因子已显示，余因子光滑且有正下界，故此展开可逐项积分。
以$u=(1+\cos\theta)/2$算得最后积分为$\pi/8$，
再将$\Delta$及$s=M^2+\epsilon$代入\eqref{eq:threephase}，
即得\eqref{eq:ph3threshold}。

五个场互不相同，顶点不带同种场阶乘；
外态均不同，末态$S_{\rm id}=1$。
故通量公式直接给\eqref{eq:contactrate}，不存在被遗漏的图间干涉。
树振幅是非零常数，而相空间的椭圆组成由定理\ref{thm:threebodyvisible}给出。
在通量平方根成为有理函数的有限覆盖上，
乘非零代数通量只作有理规范变换，保持该简单因子；
回到原域则携带有限代数字符，不能悄悄删去。

这里$g$有负质量维数，模型只是明确有限阶的局域有效相互作用，
没有声称是四维可重整化理论。
也没有把它叫作$\phi^4$的$2\to3$过程：单实标量$\phi^4$的
$\mathbb Z/2$对称性禁止该过程。本实例独立于既有$\phi^4$基准。
\end{proof}

这不是新发现三体椭圆相空间；
经典表达见\cite{DDphase}。
它在本分类方案中的作用，是给出一个有实际物理链和最小组成下界的反例：
“振幅只需代数或Tate结构”不能推出“总截面只需同类结构”。
分类输入必须保留可观测量和相空间，裸振幅图的周期表不足以决定物理出口。
固定$s$代数点可能发生算术特化；上述变化结论不自动判断那一个实数的类型。

\subsection{推前以后必须重新定义背景}
上一节的张量次数背景$\cS_d$控制积分之前的多项式读出。
一般推前不正合，也不保持Tate标签。下面明确应补入的边界项。
为避免假定尚未构造的Nori六运算，本段首先在有理可构造层的导出范畴中陈述；
对局部系统取一个使以下全部上同调层局部常值的有限分层即可。

\begin{theorem}[带相邻上同调的推前背景比较]\label{thm:pushbackground}
令$f:X\to S$是有限型复代数态射。
取$\mathsf F=Rf_!$或$\mathsf F=Rf_*$，两者选定其一并始终保持相同的支撑约定。
假设在$X$上已有实际短正合列
\[
 0\to K\to R\to P\to0,\qquad
 0\to I\to R\to Q\to0
\]
（例如共同实际模型中的生成对象及不可见核，而不是只在每个数值点的同构）。
对整数$q$，令$\mathcal S^f_q$为包含
\begin{equation}\label{eq:pushbg}
 H^q\mathsf F K,\quad H^{q+1}\mathsf F K,\quad
 H^q\mathsf F I,\quad H^{q+1}\mathsf F I
\end{equation}
的Serre子范畴。则在该Serre商中
\begin{equation}\label{eq:pushiso}
 \pi H^q\mathsf F P
   \simeq\pi H^q\mathsf F R
   \simeq\pi H^q\mathsf F Q .
\end{equation}
因而背景外的简单重数相同。
若还给出与这两个映射相容的实际物理链与de Rham读出，
并证明积分后两形式周期之差由指定推前背景承载，
则定理\ref{thm:addition}适用于这些\emph{积分后的}最小对象。
单有\eqref{eq:pushiso}并不足以推出这个标架结论。
\end{theorem}
\begin{proof}
对第一条短正合列取导出推前及上同调长正合列：
\[
 H^q\mathsf F K\longrightarrow H^q\mathsf F R
 \longrightarrow H^q\mathsf F P
 \longrightarrow H^{q+1}\mathsf F K .
\]
中间映射的核是第一个对象的商，余核是最后一个对象的子对象。
两者属于$\mathcal S^f_q$，故映射在Serre商中成为同构。
对第二条列重复此论证，即得\eqref{eq:pushiso}。
这只使用实际正合列与导出函子的三角性质，不假定$f_*$或$f_!$正合。
删去有限分层后的每一层上对象为有限秩局部系统，重数结论随之成立。

最后的标架结论明确另有输入。
若总空间的相对链、形式和边界比较已给积分后实际差周期$\ell_f$，
其载体在声明的背景中，则$p_f-q_f=\ell_f$是该实际范畴中的形式关系。
由加性比较定理可得最小对象的Serre商同构；
若有socle分离，还可应用完整商保护定理。
没有这些输入，任取两条不同的读出可能选到不同子商，
环境上同调的同构不能替代物理标架比较。
\end{proof}

若要将此处实现层的比较提升为Nori断言，还须构造相应态射和差载体的Nori提升；
Betti正合列本身不提供这种提升。
特别地，不能把\eqref{eq:pushbg}的$q+1$项删除，
也不能把$K,I$原来所属的标签直接贴在其推前上。
最简单的几何例子是椭圆曲线到一点：
常系数$\Q$推前的$H^1$已非Tate。
定理\ref{thm:threebodyvisible}和命题\ref{prop:contactphase}
把同一机制落实为正物理相空间与总率，而非仅为抽象反例。

于是现在有一条真正接到物理出口的有限规则：
有理树权先给实际壳面周期；
二体有理角积分给Artin--Tate/Kummer上界；
三体正质量相空间给恰一份可见椭圆组成；
方案或投影比较经过同一实际推前后，使用重新生成的背景而不是原背景。
一般圈振幅的边值、非绝对可积IR区域、相空间上的局部反项、
全部混合扩张与固定点算术仍需各自证明。
这部分关闭了明确接受类的相空间入口，
没有把一般散射可观测量的完整动机分类标为完成。



\section{两圈弹性截面的椭圆因子与实际方案背景}\label{sec:elasticelliptic}
上一节证明树权的相空间可以新增几何。
现在回到本书已有的massive $\phi^4$两圈物理基准，
直接分析其完整弹性总截面中的一圈平方和树/两圈干涉。
这里的新增结论不是另一组数值，而是一个此前没有由数值决定的结构：
两个交叉泡的角卷积产生可见椭圆因子，固定极点质量的有限耦合方案变化不能消去它。

\subsection{从完整物理振幅分离唯一的交叉卷积项}
沿用本书已给Schwinger森林与色散比较的规范：
极点质量$m>0$、单位外腿留数、对称点$x_0=4/3$定义实耦合$g$，
$x=s/m^2>4$。函数为
\[
 B(z)=-\int_0^1\log(1-zu(1-u)-i0)\,du,\quad k=B(x_0),\quad b=B-k,
\]
\[
 J(z)=(z-x_0)\int_4^\infty
 \frac{\beta(r)[A(r)-k]}{(r-x_0)(r-z-i0)}\,dr,\quad
 \beta(r)=\sqrt{1-\frac4r},\quad
 A(r)=\frac1{r-4}\int_0^{r-4}B(-v)\,dv .
\]
实际树、一圈、两圈振幅为
\[
 \mathcal M=-g+g^2\mathcal A_1+g^3\mathcal A_2+O(g^4),\qquad
 \mathcal A_1=\frac{\sum_{\xi=s,t,u}b(\xi/m^2)}{32\pi^2},
\]
\[
 \mathcal A_2=\frac{\sum_{\xi=s,t,u}[-b(\xi/m^2)^2-4J(\xi/m^2)]}
                    {1024\pi^4}.
\]
这些输入来自本书独立物理基准的完整通道、反项及LSZ处理，
不是只保留裸泡的截面近似。
以下$+$、$-$为物理实轴上两条已指定解析边值；
延拓到复$x$时分别延拓两个函数芽，不把复共轭当成全局代数态射。

置$a=x-4$，对角函数写
\[
 \langle f\rangle_a=\frac1a\int_0^a f(v)\,dv,\qquad
 F(a)=\int_0^a B(-v)B(-(a-v))\,dv .
\]
同种末态$2!$已除入。总弹性截面定义$c_3,c_4$：
\begin{equation}\label{eq:elasticcoeff}
 \sigma_{\rm el}(s)=\frac1{32\pi s}
       [g^2+c_3(x)g^3+c_4(x)g^4]+O(g^5).
\end{equation}
对$4<x<16$，这也是该阶包容总截面；高于四体阈值时须另加旧基准的四体项。

\begin{proposition}[完整角投影与交叉项系数]\label{prop:crosscoefficient}
令$b_s^\pm=B(x\pm i0)-k$，$\bar b=\langle B(-v)-k\rangle_a$，
$\overline{b^2}=\langle(B(-v)-k)^2\rangle_a$，
$\bar J=\langle J(-v)\rangle_a$，并定义
$\widehat c_4=1024\pi^4c_4$。
则
\begin{align}
 c_3&=-\frac{b_s^++b_s^-+4\bar b}{32\pi^2},\label{eq:c3full}\\
 \widehat c_4
 &= b_s^+b_s^-+(b_s^+)^2+(b_s^-)^2
   +2(b_s^++b_s^-)\bar b+6\overline{b^2}\notag\\
 &\quad+2\langle b(-v)b(-(a-v))\rangle_a
      +4(J_s^++J_s^-)+16\bar J \notag\\
 &=\mathcal R(x)+\frac{2F(a)}a .\label{eq:c4full}
\end{align}
这里$\mathcal R$由其余显示项及
$-4k\langle B(-v)\rangle_a+2k^2$组成。
它的所有简单函数组成均有有限单值；并且其任一解析分支在$x=-4$处正则。
相反，交叉项$F$不能由逐通道平方或单独的角平均替代。
\end{proposition}
\begin{proof}
在物理角上$t/m^2=-v,\ u/m^2=-(a-v)$，
$d z=-2\,dv/a$。截面中的阶$g^4$系数为
$\langle\mathcal A_1^+\mathcal A_1^--\mathcal A_2^+-\mathcal A_2^-\rangle$。
设$h=b(-v)+b(-(a-v))$。第一项的分子为
$b_s^+b_s^-+(b_s^++b_s^-)h+h^2$；
后两项再给$(b_s^+)^2+(b_s^-)^2+2b(-v)^2+2b(-(a-v))^2$
及相应的$J$和。利用$v\leftrightarrow a-v$，
得到\eqref{eq:c3full}--\eqref{eq:c4full}全部系数。
尤其交叉项系数是$2$，不是分别求两次平方后得到的结果。

为证明关于$\mathcal R$的结构而不依赖数值，先有
\[
 B(-v)=2+\frac uv\log r,\qquad
 u^2=v(v+4),\quad r=\frac{v+2-u}{2},\quad v=\frac{(1-r)^2}{r}.
\]
在正实域选$u>0$、$0<r\le1$。故$B(-v)$及其平方均为
穿孔$\mathbb P^1_r$上的对数表达。
有理函数乘有限长度的此类迭代积分再作原函数，
由部分分式、对高阶极点分部积分以及对简单极点添一个字母闭合。
因此$\bar b,\overline{b^2}$没有非有限单值简单因子。

$J$也可直接作同一类型的判定。令
$r_{\rm sp}=4/(1-z^2)$，$h(z)=\operatorname{artanh}z$，
$D_x(z)=4-x+xz^2$，则实际色散式变为
\begin{equation}\label{eq:Jrationalized}
 J(x)=8(x-x_0)\int_0^1
 \frac{(3-k)z^2-(1-z^2)h(z)^2-2zh(z)}
      {D_{x_0}(z)D_x(z)}\,dz .
\end{equation}
$h=\frac12(\log(1+z)-\log(1-z))$，
其余核对$z$有理；移动极点只有
$\pm\sqrt{1-4/x}$及固定$\pm\sqrt{1-4/x_0}$。
部分分式将积分化为这些字母与$0,\pm1$的有限迭代积分。
也可对移动极点微分，积分分部后仅出现极点之间及与端点的差的对数微分；
按长度归纳得到同一结论。
作$\sqrt{1-4/x}$的有限覆盖后，这是一组属零对数系统。
在$x>4$，按已指定$i0$把路径从内部简单极点的两侧绕开，
给$J^\pm$；无需把主值误当作绝对收敛的实积分。

再以$v=(1-r)^2/r$作外角原函数，
$\sqrt{1+4/v}=(1+r)/(1-r)$有理化；
所以$\bar J$仍由属零迭代积分承载。
这些对象的简单局部系统均有有限单值，张量和直和保持该性质。

最后检查实际奇点，而不把支割线上每一点都当作奇点。
$B(x)$的分支只可能在$x=0,4,\infty$发生非平凡局部单值；
式\eqref{eq:Jrationalized}的移动极点碰撞还可能列入$x=x_0$，
但该点的实际$J$由减除定义解析。特别是$x=-4$没有极点碰撞。
角原函数的端点只将单变量奇点移至$a=0,-4,\infty$
（固定减除点的表观碰撞也不等于$a=-8$）。
故在$a=-8$，即$x=-4$，上述函数的每一个续接分支均正则。
这证明关于$\mathcal R$的两项结论。
\end{proof}

\subsection{交叉角积分的实际属一模型与上界}
\begin{proposition}[双Kummer系数下的椭圆载体]\label{prop:crosscarrier}
对$a>0$，$F(a)$是下列属一曲线上的实际相对周期，
系数为两个Kummer扩张的张量：
\begin{equation}\label{eq:crosscurve}
 C_a:\quad u^2=v(v+4),\qquad w^2=(a-v)(a-v+4).
\end{equation}
其纯属一部分与
\[
 E_a:\quad Y^2=v(v+4)(a-v)(a-v+4)
\]
的$H^1$有理同构。
在$\mathbb C(a)$的函数范畴，$F$的非有限单值简单因子
只能来自$H^1(E_a)$，总重数至多四。
这个上界保留Kummer和端点扩张，不将$F$称为一个纯椭圆周期。
\end{proposition}
\begin{proof}
在曲线\eqref{eq:crosscurve}上取实际提升
$0\le v\le a,\ u,w\ge0$，并置
\[
 r_1=\frac{v+2-u}{2},\qquad
 r_2=\frac{a-v+2-w}{2}.
\]
两者沿链属于$(0,1]$。
由上一证明，
\begin{equation}\label{eq:crosslogs}
 F(a)=\int_0^a
  \left(2+\frac uv\log r_1\right)
  \left(2+\frac w{a-v}\log r_2\right)\,dv .
\end{equation}
端点的外观极点可去：例如$v\to0$时
$u/v=O(v^{-1/2})$、$\log r_1=O(v^{1/2})$。
使用曲线上局部坐标$u$，连$ (u/v)\,dv$本身也正则。
故实际链与两端的相对条件完全指定，积分收敛。

每个对数由实际Kummer模型给出：
\[
 \log r_i=\int_0^1
       \frac{r_i-1}{1+\tau_i(r_i-1)}\,d\tau_i .
\]
将此代入\eqref{eq:crosslogs}，
得到$C_a$与一个或两个辅助有理区间的纤维积上的有限个有理形式积分。
物理链上分母严格为正，链紧致，故Fubini成立。
这给一个实际代数周期模型，并非从函数值猜测动机。
必要时去除复极点、作相对紧化和消解。
随$a$变化的比较条件及端点在删除有限判别式后组成同一实际族。

四重覆盖$C_a\to\mathbb P^1_v$的分支为$-4,0,a,a+4$；
各分支点上有两个指数二的点。
Riemann--Hurwitz给$2g(C_a)-2=-8+4\cdot2=0$。
对角对合$(u,w)\mapsto(-u,-w)$没有固定点，
商为$E_a$，映射$(v,u,w)\mapsto(v,uw)$为次数二无分歧覆盖。
拉回与迹相乘为$2$，两边$H^1$均秩二，
故在有理上同调上同构。

去掉实际有理形式及$r_i$所需的有限极点集$D$，
并记端点集为$Z$。两个Kummer局部系统的张量$\mathcal K$秩四，
具有分级秩$1,2,1$的Tate过滤。
带系数的相对上同调$H^1(C_a\setminus D,Z;\mathcal K)$
承载\eqref{eq:crosslogs}。
对$\mathcal K$的过滤逐次使用长正合列，
可见其中简单因子来自四份普通相对曲线上同调及相邻的
$H^0,H^2$，附带Tate扭转。
普通曲线的非有限单值部分仅为一份$H^1(C_a)$；
极点和端点的有限集只给有限单值组成，$H^0,H^2$也不增加属一因子。
因此非有限单值简单因子的类型只有$H^1(E_a)$，重数至多四。
实际读出的最小对象是这个载体的子商，所以继承上界。
在函数范畴中常数Tate扭转不改变局部单值类型；
算术和权重陈述仍须保留这些扭转。
\end{proof}

属一曲线上的对数迭代结构是已有的通用语言
\cite{WeinEll}；此处给的是指定完整角权和端点的实际输入。
上界还不能证明椭圆部分可见，下面给无需计算$F$闭式的下界。

\subsection{两次续接的非零椭圆读出}
置$S=-a=4-x$，定义从$S<0$续接的函数
\[
 H(S)=\int_0^S B(t)B(S-t)\,dt,\qquad F(a)=-H(-a).
\]
以下$\Delta$均取上边值减下边值，并按实际延拓的先后顺序定义。
$S=8$对应$x=-4$，不在正能量散射区；
它是用来检测完整函数组成的续接奇点，不是宣称物理截面在负$s$产生新末态。

\begin{theorem}[完整弹性系数的椭圆下界]\label{thm:elasticelliptic}
先从$4<S<8$的$H$跨越$S=4$取差，
再将这个差函数跨越$S=8$取差，则
\begin{equation}\label{eq:doublecrossdisc}
 \Delta_8\Delta_4H(S)=2(2\pi i)^2L(S),\qquad
 L(S)=\int_4^{S-4}
       \sqrt{1-\frac4t}\sqrt{1-\frac4{S-t}}\,dt>0\quad(S>8).
\end{equation}
$L$的最小函数对象精确为秩二不可约椭圆局部系统
\[
 V_S=H_1\bigl(\mathcal E_S,\mathbb C\bigr),\qquad
 \mathcal E_S:\ y^2=t(t-4)(S-t)(S-t-4).
\]
由式\eqref{eq:c4full}，
\begin{equation}\label{eq:c4disc}
 \Delta_8\Delta_4\widehat c_4(4-S)
       =\frac{4(2\pi i)^2}{S}\,L(S).
\end{equation}
因此完整$c_4$确实含有该椭圆简单因子。
结合命题\ref{prop:crosscarrier}，它的非有限单值简单类型恰为这一种，
重数在一至四之间。此处尚未使用系数角色；
推论\ref{cor:elasticone}将重数确定为一。
完整扩张和最小齐次总阶数仍不由此决定。
\end{theorem}
\begin{proof}
泡的实际跳跃为
\[
 B(t+i0)-B(t-i0)=2\pi i\sqrt{1-\frac4t}\qquad(t>4).
\]
当$4<S<8$时，两个因子的切口区间在积分线上不重叠。
分别对两个因子取差，再以$t\leftrightarrow S-t$对称，得到
\[
 \Delta_4H(S)=
       2(2\pi i)\int_4^S\sqrt{1-\frac4t}\,B(S-t)\,dt .
\]
这一步保留两种位置，不能遗漏前面的$2$。
泡在阈值处连续，有平方根型局部项；
在有限区间上边值取极限及换序均由可积界控制。
将此差函数继续到$S>8$，第一个平方根的分支固定，
第二因子在$4<t<S-4$跨越其切口。
再用泡跳跃公式立即给\eqref{eq:doublecrossdisc}。
这是顺序进行的两个分支差；
不能把$S>8$处未经分离的第一次总跳跃误作同一个表达式。

在$S>8$，四根为$0<4<S-4<S$。
沿中间区间的双片整圈$\gamma_S$有
\[
 2L(S)=\int_{\gamma_S}\eta_S,\qquad
 \eta_S=\frac{y\,dt}{t(S-t)}.
\]
$\eta_S$在四个有限分支点均正则。
在无穷远的两片上
$y=\pm(t^2-St+O(1))$，
所以$\eta_S=\mp(1+O(t^{-2}))dt$；
只有二阶极点，没有留数。
因此$\eta_S$的代数de Rham类属于紧曲线$H^1$，
不需把第三类残数当作纯周期。

当$S\to8$，中间两根相碰，其余分离，$j$有极点。
如前节模曲线论证，$V_S$绝对不可约且单值Zariski稠密于$\mathrm{SL}_2$。
$L>0$证明该de Rham类和周期读出不为零。
因$\gamma_S$的全局轨道生成不可约$V_S$，
周期函数的不可见核只能为零，故$L$的最小函数对象就是$V_S$，
秩精确为二。

由$F(a)/a=H(S)/S$及\eqref{eq:c4full}，
交叉项给\eqref{eq:c4disc}的右端。
其余$\mathcal R(4-S)$的每一分支都在$S=8$正则，
故即使第一次变分非零，第二次变分也为零。
这排除了完整树/两圈干涉与其它一圈平方项对该椭圆读出的相消。
右端属于$c_4$各分支的线性组合，故其非零秩二分支空间
给$c_4$最小对象的椭圆组成下界。
命题\ref{prop:crosscarrier}及$\mathcal R$的有限单值组成给相应上界。

所有结论均为参数函数层。
它不判定某个指定正实或代数$x$处的数值是否可用MZV表达，
也不从这个rank2椭圆核推出整个截面为纯对象。
\end{proof}

\subsection{有限耦合变化后的实际背景与保护}
\begin{theorem}[固定极点质量的弹性方案保护]\label{thm:elasticscheme}
保持$m$、单位留数和同种末态计数不变。
令两个真实有限耦合方案按与运动学无关的实常数$\alpha,\beta$满足
\[
 g'=g-\alpha g^2-\beta g^3+O(g^4).
\]
则同一物理截面的系数满足
\begin{equation}\label{eq:elasticreparam}
 c'_3=c_3+2\alpha,\qquad
 c'_4=c_4+3\alpha c_3+5\alpha^2+2\beta .
\end{equation}
当$\alpha,\beta$给有已指定标架的实际参考周期时，
这两个关系在本节所构造的相空间/色散周期模型中有实际直和张量载体：
阶四差由参考周期载体乘以$c_3$载体，以及相应常数乘积承载。
它不是未经构造地对推前使用正合性。

$\Delta_8\Delta_4$消去全部这些差项，
所以\eqref{eq:c4disc}的非零椭圆组成在上述方案类中保持。
结论不要求把任意参考周期本身叫作算术Tate；
也不允许改变质量坐标后仍省略链式项。
\end{theorem}
\begin{proof}
形式反演给
$g=g'+\alpha g'^2+(2\alpha^2+\beta)g'^3+O(g'^4)$。
代入\eqref{eq:elasticcoeff}：
\[
 g^2=g'^2+2\alpha g'^3+(5\alpha^2+2\beta)g'^4+O(g'^5),
 \quad
 g^3=g'^3+3\alpha g'^4+O(g'^5).
\]
比较系数即得\eqref{eq:elasticreparam}。
对应振幅关系为
$\mathcal A'_1=\mathcal A_1-\alpha$、
$\mathcal A'_2=\mathcal A_2+2\alpha\mathcal A_1-2\alpha^2-\beta$；
它包含树/两圈干涉的补偿，不能只改变一圈泡中的常数。

实际性来自共同模型而非形式展开本身。
对于交叉项，两方案使用同一$C_a$、同一正链及同一
$\tau_1,\tau_2$辅助Kummer链；
减去常数后，未变的$\log r_1\log r_2$项直接相消，
余项是单Kummer项及常数项的实际有理形式。
对于单通道泡、$J$及它们的角原函数，使用
\eqref{eq:Jrationalized}的共同穿孔属零曲线。
在固定物理$x>4$，内部极点可用代数半径的半圆弧绕开，
上下两条轮廓分别给$+$、$-$边值。
轮廓形变及端点对数的可积性证明它们与规定的边值相同；
选定轮廓后是有限的实际周期，不能仅以主值数值充当模型。

参考周期置于独立因子，与这些共同曲线及链作乘积。
由于外角区间紧致，端点可积，且有限轮廓避开极点，
常数加法、乘法与角积分可在这些实际模型上逐项进行。
于是\eqref{eq:elasticreparam}是指定总率模型中的形式周期关系；
旧基准的两种实际耦合规范提供$\alpha,\beta$的一个明确输入。
这里并未声称与任意未经比较的原始动量Nori模型自动同构。

在参数函数层，参考周期是常数；$c_3$的各分支在$S=8$正则。
因此$\Delta_8\Delta_4$消去式\eqref{eq:elasticreparam}全部新增项，
而$c_4$的双变分由\eqref{eq:c4disc}非零。
按实际背景加法定理，也可把该结论表为指定背景以外的组成保持。
常数参考值可能有非Tate算术信息，但它不因此获得本参数的新单值；
函数保护与固定纤维算术比较不可混同。
\end{proof}

这条推导把已经执行的独立物理基准接回统一目标：
完整森林减除振幅$\to$实际上下边值$\to$含干涉的总角投影
$\to$属一载体$\to$不相消的物理函数因子$\to$方案保护。
没有新增求积或把拟合特殊函数当作证明。
本节完成的是该有限物理模型的结构接口，
不代替三圈前向光学验证、四体窄误差认证、一般幂次森林或任意图算法。
对这一具体截面结构的系统优先权检索尚未完成，不宣称文献首创。



\section{带对数系数的角色选择：从代数根到物理投影}\label{sec:eqkummer}
前节的双Kummer模型留下一个一般问题：
若角积分的代数根覆盖有有限对称群，对数系数也随该群变换，
能否在积分以前排除不可能出现的纯核？
只计算覆盖的亏格不够；把每个对数分级都当作平凡角色也会重复计数。
本节将命题\ref{prop:star}的紧割分解延伸到带Kummer系数的相对投影。
其输出先是组成上界；物理可见性仍须实际链证据。

\subsection{接受输入与实际系数系统}
设$S$为光滑参数曲线，$K=\mathbb C(S)$。
在已声明的必要有限覆盖及删除有限坏点后，
有光滑射影曲线族$C\to S$，
\[
 K(C)=K(z)(y_1,\ldots,y_r),\qquad y_i^2=p_i(z).
\]
要求平方类在$\overline K(z)^\times/\overline K(z)^{\times2}$中独立；
故一般几何纤维连通，群为$G=(\mathbb Z/2)^r$。
只来自常数域的非平方不能误计为几何分支。
对$\chi=(\chi_i)\ne0$，记正规化光滑紧模型
\[
 C_\chi:\ y_\chi^2=\prod_i p_i^{\chi_i},\qquad
 b_\chi=\#\{a\in\mathbb P^1_{\overline K}:
             \textstyle\sum_i\chi_i\operatorname{ord}_a p_i\text{ 为奇数}\}.
\]
无穷远、重根约去和共同根均计入赋值。于是
\begin{equation}\label{eq:eqcompact}
 g_\chi=(b_\chi-2)/2,\qquad
 H^1(C)_\chi\simeq H^1(C_\chi),\qquad H^1(C)_0=0.
\end{equation}
这由商映射的拉回、迹和有理群幂等元给出，
所用标准机制见\cite{KR,CarRod}；本节的新增问题是系数角色。

取$G$稳定、互不相交、均非空的有限点集$D,Z$，令$U=C\setminus D$。
端点、极点及其全部$G$轨道必须加入实际模型。
取$U$上的代数单位$f_1,\ldots,f_q$与角色$\psi_j$，要求
\begin{equation}\label{eq:inverseunit}
 g^*f_j=f_j^{\psi_j(g)}\qquad(g\in G)
\end{equation}
作为有理函数恒等式成立；$\psi_j(g)=\pm1$。
不能省略一个未经记录的常数倍。
若群作用混合不同字母，须另给更一般的等变系数模型。

令$\mathbf k=(k_1,\ldots,k_q)$为非负整数。
接受的角权为
\begin{equation}\label{eq:logangle}
 P(s)=\int_{\gamma_s}
       \sum_{0\le j_\nu\le k_\nu}
       \omega_{\mathbf j}(s)\prod_{\nu=1}^q
             \frac{(\log f_\nu)^{j_\nu}}{j_\nu!},
\end{equation}
$\omega_{\mathbf j}$是$U/S$上的代数一形式，
$\gamma_s$带实际提升、取向、端点和对数支。
先要求紧链上$f_\nu>0$、无零极点且端点形式可积；
其它物理边值须另提交避极点轮廓及比较。
落在$D$上的可积端点不能直接算作这里的普通相对端点，
须先给切向端点或实吹起比较。

取单位一动机$[\mathbb Z\to\mathbb G_m],\ 1\mapsto f_\nu$
的对偶Kummer实现，作第$k_\nu$对称幂再张量，记$\mathcal K_{\mathbf k}$。
这个秩二Kummer系统在$f_\nu=1$时仍存在且分裂；
不能用端点碰撞后的普通对$(\mathbb G_m,\{1,f_\nu\})$
的变秩上同调替代它。
可用基$e_{\mathbf j}$和连接写成
\begin{equation}\label{eq:Kconnection}
 \nabla e_{\mathbf j}
 =\sum_{\nu:j_\nu>0}e_{\mathbf j-\mathbf e_\nu}\,d\log f_\nu,\qquad
 \operatorname{ev}_{\boldsymbol\ell}(e_{\mathbf j})
 =\prod_\nu\frac{\ell_\nu^{j_\nu}}{j_\nu!},\quad \ell_\nu=\log f_\nu .
\end{equation}
降阶算子交换且$d(d\log f_\nu)=0$，故连接平坦；
$\operatorname{ev}_{\boldsymbol\ell}$是沿链的水平对偶系数。
因此这是带系数相对周期，不是常系数代数微分积分。

\begin{proposition}[Kummer过滤及总角色]\label{prop:eqkummer}
令$N=\prod_\nu(k_\nu+1)$。
$\mathcal K_{\mathbf k}$的总次数过滤是$G$稳定的平坦过滤，
每个分级项为常局部系统，群角色为
$\psi^{\mathbf j}=\prod_\nu\psi_\nu^{j_\nu}$。
在Kummer相对上同调约定中还带$\Q(-|\mathbf j|)$；
下文仅在参数函数的单值公式中省略常Tate扭转。

给定$\alpha\in G^\vee$，若每个非零项满足
\begin{equation}\label{eq:equivform}
 g^*\omega_{\mathbf j}
   =\alpha(g)\psi^{\mathbf j}(g)\,\omega_{\mathbf j},
\end{equation}
则$\Omega=\sum\omega_{\mathbf j}e_{\mathbf j}$位于
$H^1_{\mathrm{dR}}(U,Z;\mathcal K_{\mathbf k})_\alpha$。
实际周期由该角色的相对上同调承载，
不要求原始物理链本身为$G$不变链。
\end{proposition}
\begin{proof}
式\eqref{eq:inverseunit}给
$d\log(g^*f_\nu)=\psi_\nu(g)d\log f_\nu$。
定义$g^*e_{\mathbf j}=\psi^{\mathbf j}(g)e_{\mathbf j}$。
代入\eqref{eq:Kconnection}，每项均乘$\psi^{\mathbf j}(g)$，
故与连接相容。次数过滤保留、分级连接为零，
给出所列角色与Kummer扭转。
单位之间若还有乘法关系，张量系统仍给上界；
不能声称其每个分级都进入最小对象。

实际性可由
\[
 \log f_\nu=\int_0^1
       \frac{f_\nu-1}{1+\tau(f_\nu-1)}\,d\tau
\]
直接看出。幂对应有限个独立区间的乘积及对称投影。
正单位保证辅助分母不消失；紧链与所要求的端点可积性给Fubini。
因此系数系统、边界和形式有实际周期来源。

式\eqref{eq:equivform}与$e_{\mathbf j}$的角色相乘，
给$\alpha(g)(\psi^{\mathbf j}(g))^2=\alpha(g)$。
有理幂等元$e_\alpha=2^{-r}\sum_g\alpha(g)g$是实际对应，且
\[
 \langle e_\alpha\Omega,\gamma\rangle
       =\langle\Omega,e_\alpha\gamma\rangle .
\]
链的系数也须按群作用运输。故可作带系数的角色平均而不改变周期；
分支点处的固定点不妨碍有理平均。
不能只平均几何链而忘记水平系数。
\end{proof}

\subsection{积分以前可读出的纯核上界}
令$\mathcal F$为参数局部系统中由有限单值简单对象生成的Serre类别。
它容许非平凡扩张，因此属于$\mathcal F$不等于整个单值群有限。
下式只比较不属于$\mathcal F$的Jordan--Hölder重数，
不假定存在规范的非有限直和子对象。

\begin{theorem}[带系数角色选择上界]\label{thm:eqselection}
定义
\begin{equation}\label{eq:nchar}
 n_\rho=\#\{\mathbf j:0\le j_\nu\le k_\nu,\ \psi^{\mathbf j}=\rho\}.
\end{equation}
若实际读出满足命题\ref{prop:eqkummer}的$\alpha$角色条件，
则每个非有限单值简单局部系统$V$满足
\begin{equation}\label{eq:eqbound}
 [Q_P:V]\le\sum_{\chi\ne0}
                  n_{\alpha\chi}[H^1(C_\chi):V].
\end{equation}
潜在非有限纯核只来自$n_{\alpha\chi}>0$的$C_\chi$；
$g_\chi=0$者不贡献紧核。
不同$C_\chi$若具有同构简单因子，右端必须将重数相加。

记一般几何纤维的$g=g(C)$、$d=\#D$、$z_0=\#Z$。
此外有有限微分阶数上界
\begin{equation}\label{eq:logrankbound}
 \operatorname{ord}_{\min,\mathbb C(S)}P
             \le N(2g+d+z_0-2).
\end{equation}
这些上界不承诺允许核或边界扩张确实可见。
\end{theorem}
\begin{proof}
系数短正合序列给
\[
 H^1(U,Z;\mathcal K')\longrightarrow
 H^1(U,Z;\mathcal K)\longrightarrow H^1(U,Z;\mathcal K'').
\]
中间项的简单重数不超过两侧之和。
特征零下$e_\alpha$精确，故可以逐角色、逐过滤使用。
角色为$\rho=\psi^{\mathbf j}$的常分级贡献
$H^1(U,Z)_{\alpha\rho}$，附带相应Tate扭转。

普通相对曲线的点集项来自$H^0(Z)$和$H^0(D)$及其子商，
随参数变化只有有限置换单值。紧$H^0,H^2$也只有常项。
局部化及相对长正合列表明，非有限简单因子仅来自
$H^1(C)_{\alpha\rho}$，再用\eqref{eq:eqcompact}即得\eqref{eq:eqbound}。
系数扩张可以连接相邻次数，所以整个带系数对象不能直接写成这些紧核的直和。

因为$D\ne\varnothing$，$U$有一维CW型；
因为$Z\ne\varnothing$且$U$连通，在$Z$为零的水平截面恒零。
故$H^0(U,Z;\mathcal K)=H^2(U,Z;\mathcal K)=0$。
带系数Euler示性数为$N(2-2g-d-z_0)$，
从而
\[
 \dim H^1(U,Z;\mathcal K_{\mathbf k})=N(2g+d+z_0-2).
\]
实际周期及其参数导数由此Gauss--Manin系统承载，
超过该秩的de Rham迭代必线性相关，给\eqref{eq:logrankbound}。
闭合链可加入一个不影响原积分的有限轨道作为相对点集，仅使用所得上界。
\end{proof}

这里有两种不同的禁止：
$n_{\alpha\chi}=0$是系数角色禁止，$g_\chi=0$是几何禁止。
即使二者都不禁止，形式仍可能是带系数全微分，或链落入不可见核。
因此角色支撑不是物理支撑的无条件等号。

\subsection{顺序变分、精确化及真实方案差}
取共同基点和已指定的续接路径，
令$W$为闭路作用与恒等作用的有限线性组合，
例如$W=(T_b-1)(T_a-1)$。
必须按实际先后顺序定义枝差，不能把不同片上的虚部直接相乘。

\begin{theorem}[角色上界、变分下界与方案保护]\label{thm:eqsandwich}
设实际总读出$p$有共同相对载体，另有允许的Serre背景$\mathcal B$。
对非有限简单对象$V$，作三项检查：
\begin{enumerate}
 \item 角色公式及其它实际相加载体给$[Q_p:V]\le1$；
 \item 实际链变形给$Wp=c(s)\ell(s)$，
       $c(s)\in\mathbb C(S)^\times$，
       $\ell$是不可约$V$的非零周期函数；
 \item 两个真实方案满足$p'-p=b$，
       且$b$的实际载体属于不含$V$的$\mathcal B$。
\end{enumerate}
则$[Q_p:V]=[Q_{p'}:V]=1$。
若还证明$Wb=0$，同一非零变分读出对两方案相同。
没有第三项时，$Wb=0$本身不能排除背景中的全部$V$。

一个有限可审阅的第三项充分条件是：
实际差只用一个向下封闭的次数集合$A$，
其系数仍满足\eqref{eq:equivform}；
把\eqref{eq:nchar}限制到$\mathbf j\in A$后，
\eqref{eq:eqbound}对$V$的右端为零。
向下封闭指$\mathbf j\in A$且$j_\nu>0$时，
$\mathbf j-\mathbf e_\nu\in A$。
\end{theorem}
\begin{proof}
$Wp$在$p$的有限枝空间中，其全部续接生成其中的子对象。
有理非零因子$c(s)$不改变简单微分模类型。
不可约$V$的非零周期没有非零不可见子模，故至少有一份$V$；
第一项上界使重数恰一。

第三项是实际周期关系，故定理\ref{thm:addition}给
$Q_p$与$Q_{p'}$在背景Serre商中同构，背景外重数保持。
这不依赖数值周期映射的单射性，也不把两个环境动机认成同一个对象。
$Wb=0$时线性性给$Wp'=Wp$。
反之一个特定$W$可以杀死某份$V$里的特定向量，
所以该等式不是一般的背景排除证明。

向下封闭保证式\eqref{eq:Kconnection}保留$A$所张成的实际子系统。
其过滤和角色继承，重复上界定理可知推前载体不含$V$。
必须先证明实际减除差落入这个子系统，
不能只凭局部性删去最高对数项。
参考常数可以来自指定的非Tate周期；
当前函数范畴中它只改变常系数，不据此抹去固定纤维的算术信息。
\end{proof}

这把代数上界、实际变分和真实方案差接成有限的理论证据链。
它不要求先求完整$p$，也不将检查下界的数学工作伪装成万能算法。

\subsection{完整弹性截面的椭圆简单重数恰为一}
\begin{corollary}[双Kummer角投影的精确椭圆重数]\label{cor:elasticone}
在第\ref{sec:elasticelliptic}节的同一质量、链、振幅和边值约定下，
$F(a)$及完整$c_4(x)$的非有限单值半单部分
精确为一份$H^1(E_a)$，
\[
 E_a:Y^2=v(v+4)(a-v)(a-v+4).
\]
有限耦合方案类中的$c'_4$亦如此。
这将前节一至四的粗界收紧为一；
不确定完整混合扩张、全部有限单值组成或最小齐次总阶数。
\end{corollary}
\begin{proof}
记$\chi_1,\chi_2$为$u,w$的翻号角色。前节单位满足
\[
 r_1=(v+2-u)/2,\quad r_2=(a-v+2-w)/2,\qquad
 r_1(u)r_1(-u)=r_2(w)r_2(-w)=1.
\]
故$\psi_1=\chi_1,\psi_2=\chi_2$，$k_1=k_2=1$。
四个Kummer分级的角色各出现一次：
$0,\chi_1,\chi_2,\chi_1\chi_2$。
两个单平方根商均为属零；唯一非零紧角色是$\chi_1\chi_2$，
对应$H^1(E_a)$。

实际带系数形式为
\[
 \Omega=\left(4e_{00}+\frac{2u}{v}e_{10}
       +\frac{2w}{a-v}e_{01}
       +\frac{uw}{v(a-v)}e_{11}\right)dv .
\]
代数系数的角色依次也是$0,\chi_1,\chi_2,\chi_1\chi_2$，
因此每项总角色平凡。
取$\alpha=0$，式\eqref{eq:eqbound}只允许一份$H^1(E_a)$，
不能把四个对数分级各算成一份。

端点$v=0,a$处，采用相应平方根局部坐标，
$(u/v)dv$和$(w/(a-v))dv$正则，$r_i$非零。
无穷远及其它极点全轨道组成$D$，端点全轨道组成$Z$；
二者不交，链合法。
原链取实对数，翻号后取$\log(r_i^{-1})=-\log r_i$，
作带系数平均保持实际归一化。

定理\ref{thm:elasticelliptic}给非零不可约椭圆下界，
所以$F$的重数恰一。
完整$c_4$与$2F/a$之差$\mathcal R$属于$\mathcal F$，
背景加法保留这一唯一非有限组成。
有限耦合方案差$3\alpha c_3+5\alpha^2+2\beta$也属于$\mathcal F$，
从而$c'_4$同样恰一。
这只是参数函数的简单重数，不是完整混合对象的纯性声明。
\end{proof}

\subsection{有限结构规则与输出边界}
\begin{proposition}[单变量对数投影的有限输出]\label{prop:anglecompiler}
对已提交的实际模型，分支奇偶、系数角色、$n_\rho$及
式\eqref{eq:eqbound}的曲线载体上界可由有限代数步骤得到，无需求积分。
平方类、因子分解和单位作用已核对后，
若分支并集有$B$个几何点，角色表的记账运算数为
\[
 O\!\left(2^r\left(rB+\sum_\nu(k_\nu+1)\right)\right).
\]
此界不包含输入多项式因子分解的位复杂度。
没有简单性、非同构性和链变分证据时，输出仍为曲线核上界。
\end{proposition}
\begin{proof}
先将各$p_i$零极点的赋值约至模二，以二元线性代数检查几何独立性，
再对各$\chi$求$b_\chi,g_\chi$。常数域下降另记。
系数计数是在有限群环中计算
\[
 \sum_\rho n_\rho[\rho]
     =\prod_{\nu=1}^q\left(\sum_{j=0}^{k_\nu}[\psi_\nu]^j\right).
\]
逐因子更新$2^r$个整数，给显示运算界。
单位作用和形式角色恒等式可在平方根单项式基中清分母核对。
没有实际恒等式时，规则返回不适用而非猜测角色。

点集轨道与亏格给微分阶数上界；
链合法及可积性仍须已提交的证明，不能由纸面分母猜测。
输出$C_\chi$及$n_{\alpha\chi}$即给可靠载体上界。
Jacobian进一步分解、族的同源碰撞与物理不可见核不由分支数决定。
若另有定理\ref{thm:eqsandwich}的链证据，则升级相应条目的重数；
否则保留上界标签。
\end{proof}

本节接通的规则是：
已减除振幅或割的单变量代数根与对数字母
$\to$实际等变Kummer相对系统
$\to$角色加权的几何核表
$\to$变分精确化与方案保护。
从任意边表恢复这种表达仍未解决；
属$g>1$也不自动意味着不可约高秩核或Calabi--Yau。
曲线与Tate系数的上界不能判断具体特殊值非MZV，
更高维剩余推前、一般森林和全部RG出口仍按原目标继续。



\section{从森林周期方向到临界谱：可见性、方案与共振}\label{sec:criticalspectral}

本节把结构分类接到局部RG的实际输出。
前文的背景保护回答一个几何组成能否被指定操作消去；
这里回答另一个问题：已有合法周期方向经过完整森林和图权后，
究竟能否改变临界指数。
两者不是同一个核。尤其不能把单耦合MS型的反演结论直接搬到多耦合谱。

\subsection{接受输入：周期方向必须先经过真实森林}

固定有限圈阶、拉氏量、外腿及算符基、方案和归一化条件。
记$A$为承载所需有限Laurent系数及有限匹配系数的交换$\Q$代数。
若要作动机陈述，$A$中的这些元素必须已有实际形式周期提升；
任意给数值取一个MZV表达不提供这种比较。
算术方向是保持已知关系的导子$\delta:A\to A$，
并固定维数展开参数、外部物理参数和有理图权。
等价地，使用平方零提升
\[
 a\longmapsto a+\tau\delta a,\qquad \tau^2=0.
\]
它不是允许独立任意改变每一个图的周期值。
没有提供这样的方向时，以下判据仍给响应公式，但不给非零算术见证。

以已接受的连通分级交换Hopf森林模型说明实际入口。
令$P$为声明方案中权$-1$的Rota--Baxter减除投影，$\Phi$为正则化角色，
并假定$\delta P=P\delta$。对连通正次数$x$写
\[
 \overline\Phi(x)=\Phi(x)+\sum_{(x)} C(x')\Phi(x''),\qquad
 C(x)=-P\overline\Phi(x),\qquad
 R(x)=(1-P)\overline\Phi(x).
\]
一般MOM的运动学装饰须进入相应模型，不能仅凭此符号假定
任意Taylor方案都由同一个标量Rota--Baxter投影实现。
下式是经典Birkhoff递归的平方零形式\cite{GZcrit}。

\begin{proposition}[带一个方向标记的完整森林]\label{prop:markedrgforest}
在上述接受模型中，
\[
 \delta\overline\Phi(x)=\delta\Phi(x)+
 \sum_{(x)}
 \bigl(\delta C(x')\Phi(x'')+C(x')\delta\Phi(x'')\bigr),
 \quad
 \delta C=-P\delta\overline\Phi,\quad
 \delta R=(1-P)\delta\overline\Phi.
 \tag{CF1}\label{eq:markedforest}
\]
因此，在固定有限阶和足够的Laurent预算内，
实际RG映射$\mathcal F_L$的方向
$b=\delta\beta=D\mathcal F_L(\delta\Phi)$
由有限个森林项给出；每个乘积项恰有一个因子被标记。
图权、场归一化、耦合反演和实际算符混合须在$\mathcal F_L$中全部保留。
\end{proposition}
\begin{proof}
把角色$\Phi$替换为$\Phi+\tau\delta\Phi$。
导子的Leibniz律保证这仍是到平方零代数的角色；
$\delta\Phi$本身通常不是角色。
对连通分级递归取$\tau$系数即得\eqref{eq:markedforest}，
因$x',x''$的次数下降而终止。
展开任意有限乘积时，两个标记相乘含$\tau^2$，故只留一个标记；
原余积的图权及嵌套次序不变。
有限阶的单位级数反演同样是有限多项式递归，因而也与该提升相容。
若一个目标Laurent系数为$k$，相乘因子的最低次数为$v_j$，
则第$i$个因子只需保留至$k-\sum_{j\ne i}v_j$；
对有限森林逐项取最大值即有有限预算。
这一步解释为何不能先丢掉所有有限项再求插入的极点。
最后按已给出的RG定义和实际物理权重组合即可。
\end{proof}

若某个方向在这一步已给$b=0$，则稳定谱自然看不见它。
但一项图非零、一种纯核出现或一个未减除极点非零，
都不能代替完整$b$的非零性。
一般图的实际动机减除仍由前文接受条件承担，
本命题并未补造缺失的几何态射。

\subsection{固定点移动必须先于稳定矩阵的求导}

令$\beta(g)$有$n$个已指定的独立耦合，固定RG时间为$t=\log Q$，
并给定固定点分支$p$：
\[
 \beta(p)=0,\qquad J=D_g\beta(p),\qquad \det J\in A^\times.
\]
必要时在声明的系数扩张或局部化中工作，
并要求对$p$的代入有定义。
若必须倒置维数参数，所得结论只在那个局部化中成立。
非冗余相关方向的指数常取$-\lambda_i(J)$；
本文先统一记录$J$的特征值，符号由物理约定另换。

记$b(g)=\delta\beta(g)$为固定坐标$g$时的方向，
$D_gJ[v]$表示矩阵函数$D_g\beta(g)$沿$v$的导数在$p$处取值。

\begin{theorem}[临界谱的协变响应与精确一阶核]\label{thm:criticaltrace}
固定点及稳定矩阵的总方向分别为
\[
 \delta p=-J^{-1}b(p),\qquad
 H:=\delta\!\left(D_g\beta(p)\right)
   =D_gb(p)-D_gJ[J^{-1}b(p)].
 \tag{CF2}\label{eq:criticalH}
\]
对任意$n$，所有特征多项式系数的一阶响应为零，当且仅当
\[
 \operatorname{tr}(J^kH)=0,\qquad k=0,\ldots,n-1.
 \tag{CF3}\label{eq:criticaltraces}
\]
若在一个忠实平坦的分裂代数内，$J$有互异特征值$\lambda_i$，
且$\lambda_i-\lambda_j$均可逆，则这又等价于
\[
 \delta\lambda_i=\operatorname{tr}(\Pi_iH)=0\quad(1\le i\le n),
 \qquad
 H=[X,J]\ \text{对某个矩阵 }X,
 \tag{CF4}\label{eq:commkernel}
\]
其中
$\Pi_i=\prod_{j\ne i}(J-\lambda_jI)/(\lambda_i-\lambda_j)$。
最后的交换子等式在该分裂代数内理解。

在无显式$t$依赖的可逆方案坐标$h=f(g)$下，
$\widetilde J=SJS^{-1}$，$S=D_gf(p)$。
若$f$也依赖周期系数，则
\[
 \widetilde H=SHS^{-1}
       +[(\delta S)S^{-1},\widetilde J],
 \tag{CF5}\label{eq:criticalcov}
\]
这里$\delta S$包含$p$的移动。
所以\eqref{eq:criticaltraces}及所有有标号的简单特征值响应与方案坐标无关。
\end{theorem}
\begin{proof}
对$\beta(p)=0$取方向得$b(p)+J\delta p=0$。
再对$D_g\beta(p)$用链式法则即得\eqref{eq:criticalH}。
令$s_m=\operatorname{tr}(J^m)$，迹的循环性给
$\delta s_m=m\operatorname{tr}(J^{m-1}H)$。
在$\Q$代数上，Newton恒等式在前$n$个幂和与特征多项式系数之间
给可逆三角递归，证明\eqref{eq:criticaltraces}，不需单谱。

在单谱分裂基中，$\delta\lambda_i=H_{ii}$。
这也由$\Pi_i$投影后对特征方程求方向得到。
此时\eqref{eq:criticaltraces}的系数矩阵是
$(\lambda_i^k)_{0\le k<n,1\le i\le n}$，Vandermonde行列式为单位，
所以等价于全部$H_{ii}=0$。
令$X_{ii}=0$、$X_{ij}=H_{ij}/(\lambda_j-\lambda_i)$，
便有$[X,J]=H$；反向由迹循环性或对角项立即得到。
忠实平坦性保证迹的零性没有因基变换而丢失。

由$\widetilde\beta(f(g))=D_gf(g)\beta(g)$求一次耦合导数，
在$p$处所有乘$\beta(p)$的项消失，给$\widetilde J=SJS^{-1}$。
对这个等式作总方向求导并用
$\delta(S^{-1})=-S^{-1}(\delta S)S^{-1}$即得\eqref{eq:criticalcov}。
任何$\operatorname{tr}(\widetilde J^k[X,\widetilde J])$为零，
故所列响应不变。
\end{proof}

这不是求一个新的高精度矩阵的要求。
它把给定方向的物理检验压缩为有限个精确标量：
\[
 \delta\Phi
 \ \xrightarrow{\ D\mathcal F_L\ }\ b
 \ \xrightarrow{\ \eqref{eq:criticalH}\ }\ H
 \ \xrightarrow{\ \operatorname{tr}(J^k\,\cdot)\ }\ 
 \text{临界谱响应}.
 \tag{CF6}\label{eq:criticalpipeline}
\]
第一箭头负责真正的图间相消，第二箭头负责固定点漂移，
第三箭头负责只观察谱造成的信息损失。
在声明的单谱范围内，某个迹非零是至少一个临界指数含该方向的充要条件。
一般非零$H$不是充分条件。

\subsection{算符混合与物理指数需要自己的读出}

设一个\emph{实际闭合}的有限算符空间满足
$d\mathcal O/dt=-\Gamma(g)\mathcal O$。
它须已处理所需的冗余算符、运动方程和混合；
未经证明的有限截块不能代替闭合空间。
经典维数若在该块相同，只给固定标量平移；
若不同，则须把完整尺度矩阵作为$\Gamma$来处理。

\begin{proposition}[复合算符临界响应的相同检验]\label{prop:operatortrace}
写$c(g)=\delta\Gamma(g)$，则固定点处的总方向为
\[
 K=c(p)-D_g\Gamma(p)[J^{-1}b(p)].
 \tag{CF7}\label{eq:criticalK}
\]
对一个$m$维简单谱闭合块，全部算符尺度响应为零当且仅当
$\operatorname{tr}(\Gamma(p)^kK)=0$，$0\le k<m$。
有限基变换$\mathcal O'=Z(g)\mathcal O$满足
\[
 \Gamma'=Z\Gamma Z^{-1}-(\beta\cdot\partial_g Z)Z^{-1};
 \qquad \Gamma'(p')=Z(p)\Gamma(p)Z(p)^{-1}.
 \tag{CF8}
\]
因此这些迹检验同样与基及耦合方案相容。
\end{proposition}
\begin{proof}
第一式直接将$\delta p=-J^{-1}b(p)$代入$\delta(\Gamma(p))$。
对$\mathcal O'=Z\mathcal O$沿RG流求导，得到\eqref{eq:criticalK}
之后的变换式；在\emph{移动后的}固定点上$\beta$恒零，
所以对该固定点等式求方向只多一个交换子。
再对$\Gamma(p)$应用定理\ref{thm:criticaltrace}的矩阵证明即可。
\end{proof}

固定点处的相似变换是经典事实；文献\cite{DScrit}还说明，
有evanescent混合时，随意截取有限物理块会出错。
本节不将有限维假设扩张到任意算符塔。
若$\nu$来自相关质量方向、$\eta$来自场或算符归一化，
应使用实际指定的这两种读出，
不能仅凭耦合稳定矩阵的谱声称已检测所有临界量。

\subsection{何时谱不可见真的只是换坐标}

矩阵交换子只描述固定点处的一阶线性化。
整个RG函数芽还含非线性项。
以下命题给出将二者等同所需的额外条件，
其同调方程是经典Poincaré--Dulac机制\cite{CMVcrit}。

\begin{theorem}[形式非共振范围内的方案方向判据]\label{thm:rgresonance}
在特征零域$K$上考虑$\beta$在固定点的形式函数芽。
设$J$可在$K$上对角化，特征值非零、互异，且
\[
 \alpha\cdot\lambda-\lambda_i\ne0
 \quad\text{对所有 }|\alpha|\ge2\text{ 和 }1\le i\le n.
 \tag{CF9}\label{eq:nonresonance}
\]
则一个一阶变形$b$的全部稳定特征值响应为零，
当且仅当存在形式向量场$v$使
\[
 b=Dv\,\beta-D\beta\,v.
 \tag{CF10}\label{eq:rggauge}
\]
它表示$g\mapsto g+\tau v(g)$的无穷小坐标变化。
结论只涉及形式芽，不保证坐标变换收敛、满足特定物理归一化，
也不保证关于额外参数无极点。
\end{theorem}
\begin{proof}
先移固定点到原点。对角线性场$Jx$上的同调算子为
\[
 \mathcal L(v)=Dv\,Jx-Jv,\qquad
 \mathcal L(x^\alpha e_i)
   =(\alpha\cdot\lambda-\lambda_i)x^\alpha e_i.
 \tag{CF11}\label{eq:homological}
\]
对$\beta$从次数二起逐次作齐次坐标变换。
条件\eqref{eq:nonresonance}允许在每个次数解同调方程消去该次项；
每个固定次数只受有限步影响，故在$x$进完备环中得到形式线性化。
这是纯代数的逐次构造，不使用解析小除数估计。

在此固定的线性化坐标下，常数方向由
$\mathcal L(v_0)=-Jv_0$消去，等价于固定点漂移；
线性方向的非对角项由$\lambda_j-\lambda_i$消去。
剩余对角项正是$\delta\lambda_i$。
次数二及以上再由\eqref{eq:homological}和非共振条件逐次解出$v$。
若这些对角项为零，则得到\eqref{eq:rggauge}并拉回原坐标。
反之，任何这样的坐标变化在固定点处只改变共轭，
由定理\ref{thm:criticaltrace}知全部特征值响应为零。
\end{proof}

\paragraph{共振反例。}
取$K$中$\lambda\ne0$和
\[
 \beta(x,y)=(\lambda x,\,2\lambda y),\qquad
 b(x,y)=(0,\,a x^2),\quad a\ne0.
\]
固定点及其全部稳定指数完全不动，但$x^2e_2$的同调特征值
$2\lambda-2\lambda$为零，故$b$不在\eqref{eq:rggauge}的像中。
不只是一阶现象：含参数$a$的完整场的解为
\[
 x(t)=x_0e^{\lambda t},\qquad
 y(t)=e^{2\lambda t}\bigl(y_0+a x_0^2t\bigr).
\]
非线性响应保留一个$t=\log Q$系数，线性临界谱却不含$a$。
这是形式RG数据的反例，未声称已由某个指定拉氏量实现。

\paragraph{参数正则性的反例。}
在$K[[\epsilon]]$上令
$\beta=(\epsilon x,2\epsilon y)$、$b=(a x^2,0)$，
$a\in K^\times$。
其稳定谱方向仍为零；在$K((\epsilon))$中
$v=(a x^2/\epsilon,0)$可消去$b$，
而$K[[\epsilon]]$中的正则$v$不可能做到：
所需单项式系数必须满足$\epsilon v_{20,1}=a$。
所以在局部化中把一个方向判为方案方向，
不能自动判为允许的$\epsilon$正则有限重整化。
这正是临界点与Gaussian点合流时必须保留的算术/形式边界。

\paragraph{重谱的反例。}
若$J=\lambda I_2$，$H=\operatorname{diag}(a,-a)$，
则\eqref{eq:criticaltraces}全为零，却出现
$\lambda+\tau a$、$\lambda-\tau a$两个相反的谱分裂；
$[X,J]$对所有$X$都为零。
一般半单重谱应先在每个重根空间对$\Pi H\Pi$作退化微扰，
非半单情形还可能需要分数幂分支。
因此“特征多项式一阶不变”与“每个指数一阶不变”只能在单谱条件下等同。

\subsection{单耦合原结果的接合及有限阶责任}

前述局部RG附篇已经证明单耦合MS型截面上的完整形式反演。
这里记录它在本节接口中的位置，不将它当作新的圈积分。
令
\[
 \beta=-\epsilon g+gB(g),\quad B=b_1g+b_2g^2+\cdots,\quad
 B(G)=\epsilon,\quad b_1\in A^\times,
\]
并取
$\omega=GB'(G)$、
$\nu=(2+\gamma_m(G))^{-1}$、
$\eta=2\gamma_\phi(G)$，
其中$\gamma_m=c_1g+\cdots$、$\gamma_\phi=d_2g^2+\cdots$。
固定$b_1$且方向消去低于$L\ge2$圈的全部系数，记
$a=\delta b_L,c=\delta c_L,d=\delta d_L$。

\begin{corollary}[单耦合首次完整局部方向不会全体消失]\label{cor:criticalfirst}
上述方向满足
\[
 \begin{pmatrix}
 [\epsilon^L]\delta\omega\\
 [\epsilon^L]\delta\nu\\
 [\epsilon^L]\delta\eta
 \end{pmatrix}
 =
 \begin{pmatrix}
 (L-1)b_1^{-L}&0&0\\
 c_1/(4b_1^{L+1})&-1/(4b_1^L)&0\\
 0&0&2b_1^{-L}
 \end{pmatrix}
 \begin{pmatrix}a\\c\\d\end{pmatrix}.
 \tag{CF12}\label{eq:criticalfirst}
\]
故非零的完整局部三方向至少被一个临界响应检测到。
这与多耦合的交换子核及共振反例不矛盾。
\end{corollary}
\begin{proof}
$B(G)=\epsilon$给
$\delta G=-a\epsilon^L/b_1^{L+1}+O(\epsilon^{L+1})$。
对$\omega$求方向时，直接项贡献
$La\epsilon^L/b_1^L$，固定点漂移贡献
$-a\epsilon^L/b_1^L$；对$\nu$用常数项$1/2$，
对$\eta$用$\gamma_\phi$没有一次项，得到\eqref{eq:criticalfirst}。
矩阵行列式为$-(L-1)/(2b_1^{3L})$，是单位。
\end{proof}

因此，在\emph{该截面}内，几何给出的合法算术方向若通过实际森林后仍非零，
完整临界三元组必有响应。
原六圈$\delta_{35}\omega=-5776/81$只是这一已证明规则的一个已核对实例；
不据此声称某一个图贡献独自受保护，
也不将motivic MZV导子的非零改写为数值超越性。

有限阶使用本节必须带误差理想：
先明确已知$\beta,\Gamma$到什么阶，再检查$J^{-1}$、
特征值差和代入$p$可能降低多少$\epsilon$阶。
若逆矩阵含$\epsilon^{-r}$，不能在取逆前按最终目标阶过早截断。
对于单耦合上式，$G=O(\epsilon)$保证
$b_{L+1}g^{L+2}$只影响$\omega$的$\epsilon^{L+1}$及以上，
这是明确的阶数预算；一般多耦合须重新证明相应预算。
一个数值截断固定点上的精确矩阵恒等式，
不因此变成完整QFT所有圈的结论。

本节完成的是有实际输入的“森林方向---固定点---临界读出”接口。
它既给有限可见性检验，也说明哪些信息会被谱丢弃；
仍不生成任意图的局部动机，不解决临界点存在性，
不替代有限质量MOM到正则临界质量坐标的实际比较。
后续应提交一个确有闭合混合空间的多耦合物理模型或普遍森林几何输入，
而不是把形式反例当作新的物理基准。




\section{乘积临界系统：实际交叉耦合、周期响应与读出重建}\label{sec:productcritical}

上一节给一般已接受RG输入的临界谱判据。
本节用一个真实可重整化模型族确定其中的箭头，
并得到无需新多圈积分的全阶组合规则。
对象是解耦固定点及其线性扰动，不是耦合后未知固定点的完整解。
能量乘积的尺度律是已有的多临界点理论，
见Calabrese--Pelissetto--Vicari的原PDF式(4.1)\cite{CPVprod}；
这里将其证明、有限方案匹配与本项目的周期方向及有限读出接口接合。

\subsection{拉氏量、封闭空间与规范}

取$k$个实向量场$\phi_i\in\mathbb R^{N_i}$，$N_i\geq1$，
在$d=4-\epsilon$使用无质量MS及RG时间$\log\mu$。
保留各个质量源，最后调到共同临界面。
用重整化耦合写相互作用部分及其反项，规范为
\[
 {\cal L}=\frac12\sum_i\bigl((\partial\phi_i)^2+r_i\phi_i^2\bigr)
 +\frac{16\pi^2\mu^\epsilon}{4!}
 \left\{\sum_i u_i(\phi_i^2)^2+
             6\sum_{i<j}w_{ij}\phi_i^2\phi_j^2\right\}
 +{\cal L}_{\rm ct}.                                      \tag{PC1}\label{eq:productlag}
\]
因此$N_i=1$时$u_i$与前文单标量的$u$一致；
本节交叉系数是$6w_{ij}$，不沿用文献中$2w$的同名字母规范。
两者的比较须先匹配整个裸势，再匹配各自重整化约定。

保留全部$k(k-1)/2$个允许的交叉耦合。
二次不变量恰为$E_i=\phi_i^2/2$；
四次不变量恰为$(\phi_i^2)^2$与$\phi_i^2\phi_j^2$。
在四维附近按规范维数作无质量MS重整化时，
二次插入除真空/接触项外构成$k$维封闭块。
对积分的四次扰动还须商去总导数及运动方程冗余，
并包含场重整化；这给上述完整四次耦合空间。
此处没有费米子evanescent算符塔。
这不是随意截取一张无限混合矩阵。

为固定符号与因子，可直接计算一次UV极点。
若$\lambda_{ABCD}$是除去$16\pi^2\mu^\epsilon$的对称四阶张量，
三条一圈通道给
\[
 \beta_{ABCD}=-\epsilon\lambda_{ABCD}
 +\lambda_{ABMN}\lambda_{MNCD}
 +\lambda_{ACMN}\lambda_{MNBD}
 +\lambda_{ADMN}\lambda_{MNBC}+O(\lambda^3).
                                                               \tag{PC2}
\]
确切地，非零欧氏$p$下的一圈泡为
\[
 \int\frac{d^dq}{(2\pi)^d q^2(q+p)^2}
 =\frac{(p^2)^{-\epsilon/2}}{(4\pi)^{d/2}}
   \frac{\Gamma(\epsilon/2)\Gamma(1-\epsilon/2)^2}
        {\Gamma(2-\epsilon)}
 =\frac{2}{16\pi^2\epsilon}+O(1).
\]
每条通道的对称因子$1/2$取消这里的$2$。
同一场块内
$\lambda_{abcd}=(u_i/3)(\delta_{ab}\delta_{cd}
+\delta_{ac}\delta_{bd}+\delta_{ad}\delta_{bc})$，
两块间$\lambda_{ab\,cd}=w_{ij}\delta_{ab}\delta_{cd}$。
指标收缩于是给
\begin{align}
 \beta_{u_i}&=-\epsilon u_i+\frac{N_i+8}{3}u_i^2
                 +3\sum_{j\ne i}N_j w_{ij}^2+O(g^3), \notag\\
 \beta_{w_{ij}}&=-\epsilon w_{ij}
  +\left(\frac{N_i+2}{3}u_i+\frac{N_j+2}{3}u_j\right)w_{ij}
  +4w_{ij}^2+\sum_{\ell\ne i,j}N_\ell w_{i\ell}w_{j\ell}
  +O(g^3).                                               \tag{PC3}\label{eq:productoneloop}
\end{align}
例如第一个收缩在第三块给$N_\ell w_{i\ell}w_{j\ell}$；
另两个收缩各给$2w_{ij}^2$。
这是规范核对和闭合性论证，不是新的多圈计算。

尤其，若仅保留路径上的$w_{12},w_{23}$并把$w_{13}$永久设为零，
式\eqref{eq:productoneloop}已经生成$N_2w_{12}w_{23}$。
下文的观察图只指定\emph{读取哪些线性指数}，
不能用它删去重整化所需的耦合。
观察图也不是某一费曼图的边表。

\subsection{因子化给全阶交叉指数}

在$w=0$，归一化泛函积分为各场块的乘积。
取各单块已指定的Wilson--Fisher形式分支
\[
 u_i^*=\frac{3\epsilon}{N_i+8}+O(\epsilon^2),\qquad
 t_i=\nu_{N_i}^{-1}=2+O(\epsilon).
\]
逐阶形式存在性可将$u_i=\epsilon z_i$代入$\beta_{u_i}=0$后
对$z_i$使用隐函数递归；$N_i+8$为单位。
本文不由此证明非整数维数的非微扰场论存在。

\begin{theorem}[实际能量乘积及其交叉尺度]\label{thm:productcrossover}
在上述乘积固定点，去除单位/接触项的重整化能量插入满足
\[
 \Delta(E_i)=d-t_i,\qquad
 \Delta(E_iE_j)=2d-t_i-t_j\quad(i\ne j).
\]
交叉耦合的长度尺度指数和$\mu$方向稳定特征值分别为
\[
 \boxed{y_{ij}=t_i+t_j-d,\qquad
        \omega_{ij}=d-t_i-t_j=-y_{ij}.}              \tag{PC4}\label{eq:producty}
\]
这些是形式$\epsilon$展开的逐阶恒等式。
二次插入的物理尺度块为
$\operatorname{diag}(d-t_1,\ldots,d-t_k)$；
在解耦坐标的切空间中，完整四次稳定矩阵可写成
\[
 J_*=
 \begin{pmatrix}
  \operatorname{diag}(\omega_{N_i})& *\\
  0&\operatorname{diag}(\omega_{ij})_{i<j}
 \end{pmatrix}.                                        \tag{PC5}\label{eq:productJ}
\]
相同特征值不妨碍\eqref{eq:producty}，因为张量因子已经指定相应算符。
\end{theorem}
\begin{proof}
在$w=0$，无传播子或顶点连接不同场块。
各块的质量源导数定义$E_i$，其RG长度指数为$t_i$，
故共轭插入的维数为$d-t_i$。
不同块的算符在同一点相乘没有跨块Wick收缩；
剔除各自单位混合后，
其任意分离点关联函数按两个块严格因子化。
尺度变换的生成元在张量积上为两个生成元之和，
所以乘积维数相加。
积分扰动$w_{ij}\int E_iE_j$的长度指数为$d-\Delta(E_iE_j)$，
而$\log\mu$的线性特征值取相反号，得到\eqref{eq:producty}。

这一论证也可完全在微扰重整化中进行：
保留恰一个$w_{ij}$插入时，各块的带能量插入图独立，
插入重整化常数在乘积规范中为两个常数的乘积。
对它取RG对数导数得到两个反常维数之和。
两块之外不可能产生线性混合；不同交叉插入间的耦合至少含两个$w$。
在解耦流形上各$\beta_{u_i}$只依赖$u_i$；
交叉行对所有$u$的导数为零。
因而稳定矩阵具有\eqref{eq:productJ}。
可能出现的上右块不改变其特征多项式，
无需假定该块在每一种有限方案中都为零。
对重整化常数的上述乘积关系取有限阶展开，
每个系数只涉及已接受单块的有限森林与乘积；
若这些系数已有实际形式周期提升，乘积使用相应张量载体。
没有用未经比较的数值表达式替代实际提升。
\end{proof}

由一圈式可独立核对
\[
 y_{ij}
 =\left[1-\frac{N_i+2}{N_i+8}
          -\frac{N_j+2}{N_j+8}\right]\epsilon+O(\epsilon^2).
                                                               \tag{PC6}
\]
两个Ising块给$y_{12}=\epsilon/3+O(\epsilon^2)$，
说明该解耦点沿交叉方向在红外不稳定。
不稳定固定点仍有定义良好的交叉指数；
本节没有把它称为通常两标量模型的最终稳定临界点，
也不由未重求和的级数判断$d=3$的全局相图。

\subsection{实际有限重定义及共同算术方向}

这里可以给出具体有限方案，而不只诉诸抽象相似性。
对两个标量取$a,b,c\in\Q$，定义
\[
 \widetilde u=u+au^2,\qquad
 \widetilde v=v+bv^2,\qquad
 \widetilde w=wF(u,v),\quad F=1+c(u+v).                \tag{PC7}\label{eq:productscheme}
\]
固定裸拉氏量，用这个形式可逆映射重写其重整化参数和反项。
逆映射在原点正则，仅添加允许四次算符的有限局部反项。
它是一个实际有限重整化约定；
不是未计算的某个指定有限动量MOM匹配。

\begin{proposition}[有限方案与周期方向的相容]\label{prop:productscheme}
式\eqref{eq:productscheme}保持解耦流形和交叉指数。
更一般地，无显式RG时间的正则可逆有限换参与封闭插入空间的
可逆有限基变换保持\eqref{eq:producty}的物理读出。
若$\delta$是固定$\epsilon,N_i$的合法共同周期导子，则
\[
 \delta y_{ij}=\delta t_i+\delta t_j
       =-\nu_i^{-2}\delta\nu_i-\nu_j^{-2}\delta\nu_j .
                                                               \tag{PC8}\label{eq:productdirection}
\]
右端是单块固定点移动后指数的总导数。
\end{proposition}
\begin{proof}
链式法则给
$\widetilde\beta_w=F\beta_w+w(\beta_u\partial_uF+\beta_v\partial_vF)$。
在移动后的解耦固定点处$\beta_u=\beta_v=0$；
对$\widetilde w=Fw$求线性导数，第二项消失，
第一项的$F$与$1/F$取消，故$\widetilde\omega_{12}=\omega_{12}$。
完整耦合换参给上一节的$J_*\mapsto SJ_*S^{-1}$。

若$\mu\,dO/d\mu=-\Gamma O$且$O'=ZO$，
直接求导得
$\Gamma'=Z\Gamma Z^{-1}-(\beta\cdot\partial Z)Z^{-1}$。
在同一固定点后项为零；二次块共同的规范维数不改变相似性。
例如$Z=\left(\begin{smallmatrix}1&\kappa u\\0&1\end{smallmatrix}\right)$，
$\kappa\in\Q$，是两个偶标量插入的实际有限基变化。
指数仍由谱给出，不能以非对角条目本身判断新增物理方向。
最后对已证明的\eqref{eq:producty}取$\delta$并用倒数法则即得
\eqref{eq:productdirection}。
因为$t_i$已经在$u_i^*$求值，固定点漂移包含在$\delta t_i$中。
\end{proof}

若两个块完全相同，它们的周期系数是同一代数中的同一个元素，
故$\delta t_1=\delta t_2$。
给两个同值常数随意规定相反变化，不定义实际周期代数的导子。
因子化的张量记账不能取消这种对角关系。

\subsection{从交叉指数重建热指数：观察图的精确核}

设$\mathcal O=(\{1,\ldots,k\},E)$为简单观察图，
选取实际模型中的一组$y_{ij}$作为读出。
令$B_{\mathcal O}$为无符号边顶点关联矩阵：
边$ij$对应的行在$i,j$两列为$1$，其余为$0$。
在任意交换$\Q$代数中，
\[
 z_{ij}:=y_{ij}+d=t_i+t_j,\qquad
 z=B_{\mathcal O}t,\qquad \delta y=B_{\mathcal O}\delta t.
                                                               \tag{PC9}\label{eq:observationmatrix}
\]

\begin{theorem}[交叉读出的核与恢复条件]\label{thm:oddcritical}
对每个连通分量：
若它含奇圈，则$B_{\mathcal O}$在该分量上单射且有有理左逆；
若它是二分图，则核恰为
\[
 \{x_i=\sigma_i a:\ a\in A\},\qquad
 \sigma_i=+1\text{或}-1
\]
其中$\sigma$为二分着色。孤立点包括在第二种情形中。
在域上$\rk B_{\mathcal O}=k-b$，$b$为二分连通分量数。

因此，全部分量含奇圈时，所有$t_i$可由观测交叉指数恢复。
三角形给最短公式
\[
 t_i=\frac{y_{ij}+y_{i\ell}-y_{j\ell}+d}{2}.          \tag{PC10}\label{eq:trianglethermal}
\]
对任意固定阶$L$，若底代数含有理规范及$d=4-\epsilon$，
各$t_i$的$\epsilon^{0,\ldots,L}$系数生成的代数，
等于所观察$y_{ij}$相应系数生成的代数；
由于$t_i(0)=2$，也等于各$\nu_i$相应系数生成的代数。
\end{theorem}
\begin{proof}
核条件沿每条边给$x_j=-x_i$。
由一条根路径，连通分量内所有$x$均由根值及路径奇偶决定。
奇圈使根值等于其负数；因$2$可逆，根值为零。
二分图中所有路径的奇偶相容，恰留下一个自由根值。
这也包括无边的单点分量。

若有长度$2r+1$的奇圈$v_0,\ldots,v_{2r},v_{2r+1}=v_0$，
则逐项消去得
\[
 2t_{v_0}=\sum_{a=0}^{2r}(-1)^a z_{v_av_{a+1}}.
\]
再沿生成树用$t_j=z_{ij}-t_i$恢复全部顶点，构成有理左逆。
其每个系数为整数或半整数，不需要任何特征值差的除法。
三角形为$r=1$的特例。
有限阶上正、反映射都为有理线性式；
$t_i^{-1}$和$\nu_i^{-1}$的形式反演因常数项为单位而只用有限多项式递归，
所以所述系数代数相同。
\end{proof}

二分分量若再读取一个$t_i$，其代数核也被消除。
但这些核是允许的形式热数据空间的核。
真实周期方向的不可见集合还必须与
$\{\delta t:\delta\text{保持实际系数关系}\}$相交；
不能从一维代数核便断言物理理论中确有该变化。

此结论是能量乘积尺度律与初等关联矩阵定理的结合，
不声称发现新的图论秩公式。
它给有限的理论规则：在已有单块输入上，选择交叉扰动、
找奇圈与生成树即能确定读出损失。
若只研究线性判定，寻找二分分量与奇圈需$O(k+|E|)$次图操作；
恢复与一致性检查再用同阶次标量运算。
输入指数本身的计算与算术关系判定不包含在这个复杂度内。

\subsection{六圈共同周期方向：所有正整数分量均保留}

在前文已声明的形式MZV提升中，令
\[
 Z_{35}^{\mathfrak m}=
   \left(\sum_{0<n<m}n^{-3}m^{-5}\right)^{\mathfrak m},\qquad
 \delta_{35}=\partial_5\partial_3-\partial_3\partial_5,\qquad
 \delta_{35}Z_{35}^{\mathfrak m}=5.                         \tag{PC11}
\]
沿用同一Brown导子约定\cite{BrownMZVprod}；
这里的$Z_{35}$在递减求和记法中为$\zeta(5,3)$。
这个交换子是导子，并消去单zeta生成的乘法子代数。
已发表六圈图系数及本项目的单耦合反演给\cite{KPprod}
\[
 [\epsilon^6]\delta_{35}\nu_N=A(N)
 =\frac{27(N+2)(42N^3+2279N^2+21282N+50512)}{(N+8)^7},
                                                               \tag{PC12}\label{eq:ANcritical}
\]
较低$\epsilon$阶响应为零。维数约定为
$\epsilon=2\epsilon_{\rm KP}$；这里重用已核对的稀疏系数，
没有重新计算六圈图。

\begin{corollary}[实际多耦合出口中的非零算术方向]\label{cor:productzeta}
对任意两个正整数$N_i,N_j$，
\[
 [\epsilon^6]\delta_{35}y_{ij}
       =-4\bigl(A(N_i)+A(N_j)\bigr)<0.                \tag{PC13}\label{eq:productzetadirection}
\]
因而该交叉指数的指定六阶形式MZV提升不属于单zeta生成的乘法子代数。
两个Ising块具体给
\[
 [\epsilon^6]\delta_{35}y_{12}=-\frac{2440}{243},
 \qquad
 [\epsilon^6]\delta_{35}\omega_{12}=\frac{2440}{243}. \tag{PC14}
\]
\end{corollary}
\begin{proof}
$\nu_N(0)=1/2$且较低阶$\delta_{35}\nu_N$为零，
故在\eqref{eq:productdirection}取六阶系数时，
$\nu_N^{-2}$只有常数项$4$能贡献。
$A(N)$分子各系数为正，且$N\geq1$时分母为正，
于是\eqref{eq:productzetadirection}严格非零，不存在两块之间的抵消。
由旧单块结果$A(1)=305/243$得到\eqref{eq:productzetadirection}的具体数。
单zeta乘法子代数被$\delta_{35}$消去，
所以非零响应给所述形式算术障碍。
\end{proof}

这个正性只是显式有理响应系数的正性，
不是把动机周期赋予一种未定义的次序。
结论仍属mixed Tate的深度信息，不产生非Tate因子；
它不证明相应实数不可能有另一种单zeta表达，
也不等于实验上连续调节zeta常数。
若要将该MZV表达进一步认作每张原始图Nori对象的规范提升，
仍须提交原图比较；此处保留前章的形式提升范围。

\subsection{对结构分类与物理出口的确切贡献}

本节闭合的推导链为
\[
 \text{已接受单块局部周期}
 \longrightarrow (t_i)
 \longrightarrow \bigl(\Delta(E_iE_j),y_{ij}\bigr)
 \longrightarrow \text{指定观察图的精确核}.
\]
在有限阶、已有形式周期提升的范围内，
中间第二箭头只作线性组合，因此其载体可取各$t_i$系数载体与单位的直和；
它不凭空引入新的椭圆或K3核。
若观察图每个分量含奇圈，所读取的系数代数又能恢复全部热系数，
给这个多耦合接受族的物理出口无信息损失结论。
它没有恢复单块$\beta$的全部方案数据、场指数$\eta$、
所有非线性交叉RG系数或原图的有限运动学部分。

特别需要区分三种图：原始费曼图决定积分；
拉氏量的跨块耦合必须在重整化下闭合；
观察图仅表示哪些交叉指数被读取。
任意删除耦合、把相同块的周期独立化、或者离开解耦固定点后继续套用
\eqref{eq:producty}，都会越过本节证明。
本节提供一条全阶结构组合规则和一个六圈算术出口，
不是任意相互作用固定点的普遍分类。



\section{有质量森林的实际周期入口与有限局部插入}
\label{sec:massiveforestperiods}

前面的零点参数构造覆盖至多对数的图，sunset篇又直接处理了一个二次发散实例。
这里将存在性入口扩展到四维有质量标量理论的任意有限圈图：
先在动量空间保留完整BPHZ森林总和，再取其实际相对周期。
这一步使用经典收敛定理，不重新宣称证明重整化定理；
新增的接合在于把有限局部喷射、收缩插入和方案比较放进同一周期代数。
它提供几何载体和微分方程的有限上界，并不直接给最小几何类型。

\subsection{接受域、森林顺序与测度}

固定带实嵌入的数域$k$和普通、局部、可重整化的四维标量$\phi^4$理论。
允许有限种标量及严格正的质量平方$\mu_e$；顶点及外部标记均保留。
本节只处理有限圈、截肢1PI图及其有限局部反项插入。
真空项单独归一化，不讨论零质量极限或非交换相位。
对标量分量选定有理动量路由，剥离耦合、对称因子及每圈$(2\pi)^{-4}$后，裸密度为
\begin{equation}
 I_G(\ell,p,\mu)=
 \frac{N_G(\ell,p)}{\prod_{e\in E(G)}(q_e(\ell,p)^2+\mu_e)^{a_e}},
 \qquad \ell\in\mathbb R^{4h_G}.
 \label{eq:massive-rational-bare}
\end{equation}
原始传播子$a_e=1$，有限次求导和局部插入可增大其幂次或改变多项式分子。
反项顶点的动量次数必须计入幂次计数；本节不允许任意高次分子仍沿用$4-N_{\rm ext}$。
标准质量、动能和四点反项的典范维数不超过四。

对一个重整化部分$\gamma$，在其\emph{独立外动量}中定义
\begin{equation}
 t^{\omega_\gamma}_{\gamma,0}f
 =\sum_{|\alpha|\leq\omega_\gamma}
 \frac{(p^\gamma)^\alpha}{\alpha!}
 (\partial_{p^\gamma}^\alpha f)(0),
 \qquad t^{\omega_\gamma}_{\gamma,0}=0\quad(\omega_\gamma<0).
 \label{eq:massive-taylor}
\end{equation}
对无导数原始$\phi^4$图，二点部分取$\omega=2$，四点部分取$\omega=0$。
在子图内进行Taylor运算后，再代回原图的外动量与外层圈动量；
不能把子图的外层圈动量永久设成零。
嵌套操作从内到外，强不重叠的部分按标准动量空间森林约定组成森林。
所有收缩保留外部结构和嵌入重数。简记
\begin{equation}
 R_0 I_G=(1-t^{\omega_G}_{G,0})
 \sum_{F\in\mathcal F_{\mathrm{prop}}(G)}
 \prod_{\gamma\in F}^{\text{由内到外}}
 (-t^{\omega_\gamma}_{\gamma,0})I_G .
 \label{eq:massive-full-forest}
\end{equation}
其中代回与局部顶点替换算子已包含在乘积记号中。
森林公式和其绝对收敛输入采用Zimmermann定理\cite{ZimMass}；
运算顺序与收缩约定可与\cite{BGHMass}第二节式(31)、(33)、(37)--(38)核对。
本节仅使用该文的\emph{交换}理论部分。

\subsection{合并森林的实际相对对象}

\begin{theorem}[有质量BPHZ周期入口]
\label{thm:massive-period-entry}
在上述接受域，对任意固定代数欧氏外动量和严格正的代数质量平方，
\begin{equation}
 A_0(G;p,\mu)=\int_{\mathbb R^{4h_G}}
 R_0I_G(\ell,p,\mu)\,d^{4h_G}\ell
 \label{eq:massive-ren-period}
\end{equation}
绝对收敛，并有一个从该\emph{完整有理密度及其链}构造的有效相对Nori载体。
可选取光滑$Y$、除子$P,E$和标架，使
\begin{equation}
 M_{G,0}=H^{4h_G}(Y\setminus P,E\setminus P),\qquad
 A_0(G)=\operatorname{per}[M_{G,0},[\eta_G],[\Sigma_G]]^{\mathfrak m}.
 \label{eq:massive-relative-carrier}
\end{equation}
必要时使用有限直和及数域扩大后向$k$的限制标量。
恢复每圈测度因子使用$(2\pi i)^{-4h_G}$，因而在允许Tate逆的周期代数中进行；
不据此声称物理归一化后的数仍是\emph{有效}周期。
\end{theorem}

\begin{proof}
每个Taylor操作只取有限次有理函数求导与代入，故合并密度仍为$k$上的有理函数。
所有实际欧氏分母严格为正；UV无穷远是这里需处理的收敛问题。
经典有质量BPHZ收敛定理适用于完整森林和上述相容减除次数，
给出\eqref{eq:massive-ren-period}的绝对收敛，包括嵌套和重叠子发散。
这是本证明的外部分析输入，不能替换为对每个森林项分别计数。

将各实坐标写成$\ell_j=x_j/(1-x_j^2)$，$-1<x_j<1$。
该变换在实区间严格递增且覆盖$\mathbb R$，Jacobian为
$(1+x_j^2)/(1-x_j^2)^2$。于是得到开立方体上的绝对可积有理顶次形式$\eta_G$。
此有理映射在复代数意义下不必双有理；使用的是它在指定实链上的一次覆盖。

现在对立方体边界及\emph{合并形式}的零极点做正规交叉消解，
应用绝对可积有理积分到正规交叉相对周期的比较定理
\cite{HMSMass}，Lemma12.2.4。
该定理的构造保留积分链：在消解空间中取链的严格变换，
删去拉回形式的真正极除子$P$，其边界位于剩余除子$E\setminus P$。
局部单项式化后，绝对可积性排除了链所邻接实象限上的负整数极阶，
所以拉回形式在链上正则。
相对上同调、形式和链给出\eqref{eq:massive-relative-carrier}，不是仅凭数字相等猜一个对象。
若合并密度恒零，可取零对象。
最后$(2\pi)^{4h_G}=(2\pi i)^{4h_G}$，测度的恢复由Tate对象实现。
\end{proof}

这一结论是经典收敛及周期比较定理的接合，不是新的“所有费曼周期均属某一类型”定理。
不同消解可由共同细化比较其\emph{指定形式周期}；
大载体本身可以含有额外上同调，不能称为唯一最小对象。
它也没有自动证明这个动量空间载体与已有参数空间对象在Nori范畴中同构。

最简单的必要性检验是一圈泡图：
\begin{equation}
 \frac1{(\ell^2+\mu)((\ell+p)^2+\mu)}
 -\frac1{(\ell^2+\mu)^2}.
 \label{eq:massive-bubble-cancel}
\end{equation}
固定$p$时此差为$O(|\ell|^{-5})$，故四维绝对可积；
两项各自均有对数发散。
因此，将\eqref{eq:massive-full-forest}的每项各自赋一个普通周期再相加，是一个已经在一圈失败的操作。
同理，“子反项局部”只表示外动量多项式性，并不表示其系数必须是有理数或Tate周期。

\subsection{有限喷射和参数族：无需先求完整积分}

\begin{proposition}[继承减除的有限喷射]
\label{prop:massive-jets}
固定一个有限的外动量、正质量及代数参考数据的导数列表。
对完整减除振幅取这些导数，在任意接受的代数欧氏点得到的数均为实际周期
（恢复测度后允许Tate逆）。每个系数可由相应的\emph{已合并、已求导}有理密度构造载体。
\end{proposition}

\begin{proof}
有质量BPHZ的参数局部收敛估计可在外动量紧集和$\mu_e\geq c>0$上取一致。
这里可保持一个小复参数邻域：若内部实线性组合为$q$，外参数产生的虚部为$\xi$，则
\[
 \left|(q+i\xi)^2+\mu\right|
 \geq q^2+\mu-|\xi|^2\geq\tfrac12(q^2+\mu)
 \quad\text{当 }|\xi|^2\leq\mu/2.
\]
有限多种森林路由允许同取这样的邻域；质量复扰动再缩小邻域即可。
UV森林余项的幂次估计在该紧参数邻域不变。
由其一致可积主控及Cauchy估计，任意预先固定的有限阶参数导数可移入\emph{完整}积分。
若参考点随参数移动，导数须同时作用于参考点和减除算子。
所得函数仍有理且绝对可积，应用定理\ref{thm:massive-period-entry}的构造即可。
这证明的是继承原森林规则的导数；它不把求导后的图重新按另一个最低减除次数重整化。
\end{proof}

\begin{proposition}[代数一参数族的有限微分上界]
\label{prop:massive-pf}
设外动量分量及质量平方为$k(s)$中的有理函数，
在一个非空实开区间内满足上述欧氏、正质量条件。
则剔除有限个参数值并局部选定物理支后，完整$A_0(G;s)$由一个有限秩相对Gauss--Manin系统承载。
若该系统秩为$b$，存在不全为零的$a_j(s)\in k(s)$使
\begin{equation}
 \sum_{j=0}^{r}a_j(s)\partial_s^j A_0(G;s)=0,
 \qquad r\leq b.
 \label{eq:massive-pf-upper}
\end{equation}
如几何构造使用有限代数扩域，先限制标量，并相应扩大$b$。
此为完整相对对象的齐次上界，不是最小阶或纯核的分类。
\end{proposition}

\begin{proof}
在$k(s)$上对合并形式、立方体边界和零极点作消解，再将有限的构造数据铺到参数开集。
删去坏纤维，使紧化及除子各层光滑，所得相对上同调有限维，带Gauss--Manin连接。
在不跨越上述坏值的足够小实区间内，立方体严格变换给水平的相对链类。
原有绝对收敛保证此链避免真正极除子，前命题保证微分与周期配对相容。
因此$A_0(s)$是水平链与代数de Rham类$v(s)$的配对。
有限维向量空间内$v,\nabla_s v,\ldots,\nabla_s^b v$线性相关；
与该链配对即得\eqref{eq:massive-pf-upper}。
所需的消解、铺展和上同调都来自有理输入，不需要求出积分值。
但本文没有实施该通用符号过程，也没有给它多项式时间界。
特殊参数上的阶数下降、链消失或代数关系，须另行判断。
\end{proof}

这个入口也适用于有限物理图和：先保留各图的真实权重，再取载体直和及总标架。
总和的最小微分阶可以因实际抵消而下降。
因此“可以有限生成一个系统”与“可以有效识别全部物理组成”仍是不同的任务。

\subsection{幂次局部项的有限方案递归}

为每个标记的发散1PI图规定局部多项式空间$L_G$及投影$T_{G,\rho}$，
满足$T_{G,\rho}|_{L_G}=1$。例如在代数欧氏参考点$\rho_G$取不超过$\omega_G$阶的Cartesian Taylor多项式。
更一般地，可用有限个代数欧氏喷射定义一个非退化的代数插值投影。
不发散图取$L_G=0$。同一标记图重复出现时采用同一投影。
任意独立分配投影并不自动保持洛伦兹、置换或内部对称；
若要得到同一拉氏量的物理方案，必须在图和、张量结构及反项类型上施加这些相容条件。

记$W$为一个非空的真子图\emph{收缩族（spinney）}：其部分同时可作局部收缩，彼此不重叠且不嵌套。
嵌套信息由下面的递归包含，而非在同一$W$中重复插入。
给局部多项式$f(\gamma)$选一个有限单项式基，定义
\[
 \mathcal I^0_{G,W}\bigl(\{f(\gamma)\};p\bigr)
 :=\text{将这些多项式作为顶点插入$G/W$后所得的零点重整化振幅}.
\]
此记号对各插入多线性。二点的动能项保留实际内边动量，不能与常数项混同。
展开该记号时，系数与实际带装饰收缩图相乘；后者同属上述可重整化接受域。

\begin{theorem}[由有限周期喷射构造局部方案比较]
\label{thm:massive-finite-scheme}
在任意固定圈数截断内，选择定理\ref{thm:massive-period-entry}及
命题\ref{prop:massive-jets}给出的实际周期提升后，以下递归定义实际有限周期系数的局部多项式：
\begin{align}
 B_\rho^{\mathfrak m}(G;p)
 &=A_0^{\mathfrak m}(G;p)
   +\sum_{\varnothing\ne W\subsetneq G}
    \mathcal I^{0,\mathfrak m}_{G,W}
       \bigl(\{f_\rho^{\mathfrak m}(\gamma)\};p\bigr),
       \label{eq:massive-finite-preparation}\\
 f_\rho^{\mathfrak m}(G;p)&=-T_{G,\rho}B_\rho^{\mathfrak m}(G;p),
 \qquad
 A_\rho^{\mathfrak m}(G;p)=B_\rho^{\mathfrak m}(G;p)+f_\rho^{\mathfrak m}(G;p).
 \label{eq:massive-finite-recursion}
\end{align}
喷射的提升用其已合并有理密度定义，不能把一次数值微分当作形式周期导子。
这些数值周期满足$T_{G,\rho}A_\rho=0$，并等于该局部减除条件下的BPHZ振幅。
在这些\emph{递归构造的提升}之间，有限方案转换、其逆及连续转换的合成都在同一周期代数中成立。
它们由局部插入组合实现，不要求将外动量Taylor投影误认为一个统一的标量Rota--Baxter投影。
\end{theorem}

\begin{proof}
按圈数归纳。真重整化部分至少含一圈，故其局部系数及$G/W$均只涉及较低圈数。
假设此前的$f_\rho(\gamma)$系数已是实际周期。
将各局部多项式在单项式基展开，\eqref{eq:massive-finite-preparation}变成有限个
零点带装饰振幅与已知周期系数的乘积。
应用$T_{G,\rho}$只需有限个代数欧氏喷射；前命题给出每个喷射的实际载体。
有限和、张量积与代数系数运算遂给$f_\rho^{\mathfrak m}(G)$的载体。
$T_{G,\rho}^2=T_{G,\rho}$立即给归一化条件。
注意$f_\rho(G)$可以含有同圈图在参考点的周期；局部性不意味着严格低圈。

物理比较可在共同UV正规化下进行：令$C_{0,a},C_{\rho,a}$为递归准备后的反项，
保留所有外部多项式装饰，用局部插入组合写
\[
 f_{\rho,a}=C_{\rho,a}\star C_{0,a}^{-1},
 \qquad A_{\rho,a}=f_{\rho,a}\star A_{0,a}.
\]
这些是有限正规化下的森林重排恒等式。将真子插入合并后，归一化条件唯一决定整图局部项，
正好给\eqref{eq:massive-finite-recursion}的正规化版本。
采用经典BPHZ允许的共同UV截断，零点振幅及所需有限喷射在$a\to0$时收敛。
按圈数归纳，上式中的\emph{有限方案差}收敛到上述$f_\rho$，
而无需$C_{0,a}$或$C_{\rho,a}$各自存在有限极限。
因此$A_\rho$就是按所声明条件进行的局部重整化振幅。

局部插入组合具有单位，严格正圈数部分三角，故在固定圈数截断内逐圈存在唯一逆。
对两方案$a,b$写$g_{ba}=f_b\star f_a^{-1}$，则
$A_b=g_{ba}\star A_a$且$g_{cb}\star g_{ba}=g_{ca}$。
每项都是同一组实际周期的有限代数表达，故这些恒等式在所构造的形式提升中成立。
这不依赖周期猜想；也不证明一个\emph{独立选取}的$\rho$方案消解载体
与该递归提升具有相同的形式周期或同构的大对象。
后一个比较仍需实际几何图表。
\end{proof}

\subsection{物理读出的背景：先去整体局部自由度}

固定只作用于原图外动量的代数线性读出$\mathcal D_G$，
假设$\mathcal D_G L_G=0$。例如四点对数图可取两个运动学点之差；
二点图的完整局部空间次数为二，需要消去整个二次多项式，
不能一般地只减去一个锚点。
选定Taylor锚定余项$1-T_{G,*}$或总阶大于二的外动量导数均可。

\begin{corollary}[实际收缩背景与形式商保护]
\label{cor:massive-scheme-background}
对前定理的两方案，$\mathcal D_G(A_b^{\mathfrak m}-A_a^{\mathfrak m})$
可由有限个\emph{完整张量积}
\begin{equation}
 M_{g(\gamma_1),\alpha_1}\otimes\cdots\otimes
 M_{g(\gamma_r),\alpha_r}\otimes
 M_{\mathcal D_G A_a(G/W;\alpha_1,\ldots,\alpha_r)}
 \label{eq:massive-background-tensor}
\end{equation}
的直和承载，其中$W$非空且真，每个部分和收缩图都严格低于原图圈数。
令$\mathcal B_{G,\mathcal D}$为这些实际对象生成的Serre张量背景。
则两读出在形式周期向量空间模$\mathcal B_{G,\mathcal D}$的像相同。
若前文的标架/割下界在该商中证明一个非背景组成，则有限方案变化不能消去这个商类。
\end{corollary}

\begin{proof}
在$A_b=g_{ba}\star A_a$中，空子图给$A_a(G)$，整图局部项为$g_{ba}(G)\in L_G$。
读出消去后者。余下各真收缩族贡献逐项展开为有限周期系数与带装饰收缩振幅的乘积。
读出可以作用于插入多项式中的原图外动量，因此先保留实际插入再微分；
Leibniz展开仍是有限个相同类型的带装饰项。
它们的载体正是\eqref{eq:massive-background-tensor}，包括该展开产生的装饰。
积分中依赖圈动量的局部因子不被取到积分号外，而包含在带装饰图中。
周期的乘法由张量积给出，求和由直和给出，故差落入该背景的形式周期子空间。
取商及使用已经证明的非背景下界即得结论。
这里不声称“给同一数的任意两个动机在Serre商中同构”；
使用的是已构造的形式周期恒等式及相容标架。
\end{proof}

三个限制对总目标尤其重要。
第一，二次局部项$a+b p^2$在单锚点减除后仍剩$b(p^2-p_*^2)$，
故仅在对数图成立的单锚点保护不能原样推广。
第二，嵌套系数是对子图完成准备后的$f(\gamma)$，重叠部分不能同时当作独立局部顶点收缩。
第三，背景的算术是\emph{系数与收缩图的张量积}；
若只记收缩图类型，就会漏掉局部系数中的椭圆、K3或混合扩张信息。

由此，任意有限圈的有质量欧氏$\phi^4$图已有一个有适用域的入口：
边表和局部物理数据给有理森林密度，完整减除给实际相对周期，
有限喷射给实际局部方案系数，去整体局部自由度后的读出进入严格低圈收缩背景比较。
参数族同时获得有限微分上界。
仍未解决的是通用载体的\emph{最小}几何/算术分解、独立参数模型的全面几何比较、
质量趋零时的红外控制、MS到该有质量方案的实际匹配及所有Lorentz物理边界值。
因此这一节扩大了可证明的入口和操作范围，未将任意图精确分类宣布完成。


\section{合并森林的参数化、Tate纤维与边界比较障碍}
\label{sec:forestparametric}

第\ref{sec:massiveforestperiods}节已经把收敛的完整森林密度变成实际周期。
本节向参数几何推进：先给保持形式周期的有理关联空间，再消去动量变量，
最后区分光滑纤维比较与退化边界上的尚缺箭头。
Feynman参数化和二次型积分是经典方法；
Bloch--Esnault--Kreimer的第6、10节已经分别讨论解析变换和几何对应\cite{BEKforest}。
这里的责任是把\emph{合并后的有分子、重整化密度}及绝对可积性纳入同一陈述，
而不把经典Schwinger技巧重新称为发现，也不把解析相等冒充全部动机同构。

\subsection{先通分，再作保持正单纯形的代数换元}

固定上一节的一个代数欧氏参数点，且$h\geq1$。
将完整$R_0I_G$通分、合并分子，写成
\begin{equation}
 R(\ell)=\frac{N(\ell)}{\prod_{j=1}^{r}D_j(\ell)^{a_j}},\qquad
 D_j=\left|\sum_{u=1}^h v_{ju}\ell_u+p_j\right|^2+\mu_j,
 \quad a_j\in\mathbb Z_{>0},\quad \mu_j>0 .
 \label{eq:forestcommonfraction}
\end{equation}
这里$\ell_u\in\mathbb R^4$，$v_{ju}\in k$，$p_j\in k^4$，$\mu_j\in k$，
而$N$是实际合并的多项式。来自Taylor运算的同一路由不同外动量分母应分别保留。
常数正分母可先吸收入$N$；$N=0$另作零周期处理。
以下设$N\ne0$，$R\in L^1(\mathbb R^{4h})$，记
$A=\sum a_j$、$d=\deg_\ell N$、$B_a=\prod\Gamma(a_j)/\Gamma(A)\in\mathbb Q_{>0}$。
$\Delta_{r-1}$是$\alpha_j\geq0$、$\sum\alpha_j=1$的有序实单纯形。
$d\alpha=d\alpha_1\cdots d\alpha_{r-1}$取其正体积取向。

\begin{theorem}[实际关联空间与绝对Fubini]
\label{thm:forestincidence}
设$Q_\alpha(\ell)=\sum\alpha_jD_j(\ell)$、
$\alpha^{a-1}=\prod\alpha_j^{a_j-1}$。则
\begin{equation}
 \int R(\ell)\,d^{4h}\ell
 =B_a^{-1}\int_{\Delta_{r-1}}\!\int_{\mathbb R^{4h}}
 \frac{N(\ell)\alpha^{a-1}}{Q_\alpha(\ell)^A}
 \,d^{4h}\ell\,d\alpha .
 \label{eq:forestincidenceintegral}
\end{equation}
右端双重积分绝对收敛。更强地，原动量形式周期与右端\emph{尚未消去动量}
的关联空间形式周期有实际代数换元和K\"unneth比较，故在相容选取的Nori形式周期代数中相等。
此处不需要周期映射单射。
\end{theorem}

\begin{proof}
对每个固定实$\ell$，令
\begin{equation}
 \beta_j=\frac{\alpha_jD_j(\ell)}{Q_\alpha(\ell)},\qquad
 \alpha_j=\frac{\beta_j/D_j(\ell)}{S},\qquad
 S=\sum_i\beta_i/D_i(\ell).
 \label{eq:forestbetamap}
\end{equation}
这是保持每个坐标面的正单纯形微分同胚，两个方向都是有理映射。
在$\ell$固定的相对顶次形式中，Jacobian为$\prod D_j^{-1}S^{-r}$；
又有$Q_\alpha=S^{-1}$，所以
\[
 \frac{\alpha^{a-1}}{Q_\alpha^A}\,d\alpha
 =\frac{\beta^{a-1}}{\prod D_j^{a_j}}\,d\beta.
\]
总空间中含$d\ell$的交叉微分与$d^{4h}\ell$楔乘后为零。
多项式单纯形积分$\int_\Delta\beta^{a-1}d\beta=B_a$由逐次初等积分得到。
正核给精确等式
\[
 B_a^{-1}\int_\Delta\!\int_{\mathbb R^{4h}}
 \frac{|N(\ell)|\alpha^{a-1}}{Q_\alpha^A}
 \,d^{4h}\ell\,d\alpha=\int_{\mathbb R^{4h}}|R(\ell)|\,d^{4h}\ell<\infty,
\]
先用Tonelli，再对带符号分子用Fubini，即得解析陈述。

形式陈述使用同一换元而非上述数值结论：原周期与多项式单纯形周期作外积，
后者的形式周期就是有理数$B_a$，其归约由实际边界映射和多项式原函数给出。
在$D_jQ_\alpha S\ne0$的适当开集上，\eqref{eq:forestbetamap}是代数同构，
保持链内部及全部单纯形面。把动量无穷远也纳入两侧共同的图闭包，
再对图闭包、极点和边界作相容分解；实际链是同一个开链的严格变换。
其可积性保证按上一节的正则化链构造得到合法相对类。
外积、代数拉回及链推前于是给同一形式周期。
这只比较所标记的周期，不要求两侧所有无关上同调同构。
\end{proof}

这个步骤可用于完整森林，恰因为通分发生在参数化之前。
如果分别参数化仍发散的各个森林项，证明中的$L^1$等式便没有可用的右端。

\subsection{实际分子的有限二次型消去公式}

定义$h\times h$矩阵、四向量列和标量
\begin{align}
 M(\alpha)&=\sum_j\alpha_j v_jv_j^t,&
 b_u(\alpha)&=\sum_j\alpha_jv_{ju}p_j,\nonumber\\
 c(\alpha)&=\sum_j\alpha_j(|p_j|^2+\mu_j),&
 U&=\det M,\nonumber\\
 F&=Uc-\sum_{u,v}\operatorname{adj}(M)_{uv}\,b_u\cdot b_v,&
 V&=F/U.\label{eq:forestaugmentedUF}
\end{align}
在$U\ne0$处令$\ell_*=-M^{-1}b$及
\[
 \mathcal D_M=\sum_{u,v}(M^{-1})_{uv}
 \sum_{\lambda=1}^4\frac{\partial^2}
 {\partial\ell_u^\lambda\partial\ell_v^\lambda}.
\]
$U,F$来自\emph{合并分母列表}，不预设等于原图的两个Symanzik多项式。

\begin{theorem}[有限有理参数形式]
\label{thm:forestmoments}
在上述非零绝对可积条件下，$v_j$张成$k^h$，且$A>2h+d/2$。
正单纯形内部有$M>0$、$V\geq\sum\alpha_j\mu_j>0$。
记$K=\lfloor d/2\rfloor$，则
\begin{equation}
 \int_{\mathbb R^{4h}}R\,d^{4h}\ell
 =\pi^{2h}\int_{\Delta_{r-1}}\Psi_N(\alpha)\,d\alpha,
 \label{eq:forestparametricfinal}
\end{equation}
其中下式必须先相加，随后才作参数积分：
\begin{equation}
 \Psi_N=
 \frac{\alpha^{a-1}}{U^2\prod_j\Gamma(a_j)}
 \sum_{k=0}^{K}
 \frac{\Gamma(A-2h-k)}{4^k k!}
 \left(\frac FU\right)^{2h+k-A}
 (\mathcal D_M^kN)(\ell_*).
 \label{eq:forestmomentformula}
\end{equation}
$\Psi_N\in k(\alpha_1,\ldots,\alpha_{r-1})$，在仿射单纯形图中的极点仅可能落在$UF=0$。
合并后的$\Psi_N$绝对可积，故按实际相对周期构造给一个参数空间Nori周期。
此结论首先识别其\emph{数值}为左端；跨越退化参数边界的形式周期比较另见下文。
\end{theorem}

\begin{proof}
若$v_j$不张成，所有分母沿某个非零圈动量方向平移不变。
非零多项式$N$沿该实直线不可能绝对可积；对其他变量的一个开集应用Fubini产生矛盾。
因此$M$在正内部正定。补平方给
$Q_\alpha=(\ell-\ell_*)^t(M\otimes I_4)(\ell-\ell_*)+V$。
其最小值至少为正质量项的凸组合，得到$V$下界。

取球面上一个开集，使$N$的最高齐次项及各非零路由的最高二次项均非零。
该开集存在，因为只排除了有限个真代数集。
沿此锥域，$|R(t\theta)|$与$t^{d-2A}$有同阶正上下界。
径向绝对收敛迫使$d-2A+4h<0$。
这也保证固定内部$\alpha$时，分子的每个有限矩都收敛。

现仅为计算\emph{解析纤维积分}使用Laplace辅助参数：
$Q^{-A}=\Gamma(A)^{-1}\int_0^\infty t^{A-1}e^{-tQ}dt$。
它不被当作代数动机态射。有限多项式Gaussian矩的生成公式为
\[
 \int N(\ell)e^{-t(\ell-\ell_*)^t(M\otimes I_4)(\ell-\ell_*)}d^{4h}\ell
 =\frac{\pi^{2h}}{t^{2h}U^2}
 \sum_{k=0}^{K}\frac{(\mathcal D_M^kN)(\ell_*)}{4^k k!t^k}.
\]
它由一维偶矩和有限Taylor展开逐项证明；奇矩为零，因子$1/(4^kk!)$
包含所有配对重数。再积分$t$，$\Gamma(A-2h-k)$均为正整数阶乘。
与定理\ref{thm:forestincidence}的$B_a^{-1}$合并即得公式。

在每个$\alpha$处，纤维积分的绝对值不超过绝对值的纤维积分；
上一证明的总绝对积分有限，因而完整$\Psi_N$可积。
由于$M^{-1}=\operatorname{adj}(M)/U$，公式只有$U,F$分母，且全部Gamma值有理。
对这个\emph{已合并}的有理参数形式和实际单纯形再次应用HMS的链构造，
得到所述参数Nori对象。此最后一步不单靠积分相等推导独立对象同构。
\end{proof}

一个简单反例说明为什么不能在所有步骤任意拆开分子。
取$x,y\in\mathbb R^4$，
\[
 R(x,y)=\frac{|x-y|^2}{(1+|x-y|^2)^4(1+|x|^2)^3}.
\]
以$z=y-x$换元，完整积分是两个收敛正积分的乘积。
但把分子拆为$|x|^2+|y|^2-2x\cdot y$后，第一项的$x$积分对数发散。
有限矩公式在每个内部参数处成立，不因此允许将其各部分分别积分到所有参数面。

裸测度规范在\eqref{eq:forestparametricfinal}中已经固定。
恢复每圈$(2\pi)^{-4}$后，整体系数为$(16\pi^2)^{-h}$；
在标准Tate约定中$\pi^{2h}=(-1)^h2^{-2h}(2\pi i)^{2h}$。
这里不沿用文献中未匹配实链、射影体积或测度的常数。

\subsection{为何光滑动量纤维只给一个Tate因子}

在参数开集$S=\{U F\ne0\}$上，令$X_\alpha\subset\mathbb P^{4h}$为
$Q_\alpha$的齐次化零集，$H_\infty$为无穷远超平面，
并设
\[
 Y_\alpha=\mathbb P^{4h}\setminus X_\alpha,\qquad
 E_\alpha=H_\infty\setminus(X_\alpha\cap H_\infty).
\]
$X_\alpha$及其无穷远截面均光滑。

\begin{proposition}[实际二次型对的Tate上同调]
\label{prop:forestquadricpair}
对$S$中的代数参数点，在有理系数下
\[
 H^j(Y_\alpha,E_\alpha)=0\quad(j\ne4h),\qquad
 H^{4h}(Y_\alpha,E_\alpha)\simeq\mathbb Q(-2h).
\]
这是实际二次型对的Nori结论；在$S$上其Betti局部系统及Hodge变化是常Tate变化。
在实正分支选择通常的$\mathbb R^{4h}$相对链后，
其dR标架读出就是公式\eqref{eq:forestmomentformula}的纤维系数。
因而对边界也完全处于$S$内的代数参数对，
相对Leray比较把总空间积分降低$4h$个上同调次数并附加$\mathbb Q(-2h)$。
\end{proposition}

\begin{proof}
写$n=4h$。光滑奇数维射影二次曲面$X_\alpha$的有理上同调全由超平面类生成。
其Gysin映射在相应偶数次数为非零同维映射，故
$H^j(Y_\alpha)=0$对$j>0$，$H^0(Y_\alpha)=\mathbb Q(0)$。
偶数维二次曲面$X_\alpha\cap H_\infty$另有两族最大线性空间之差的中间原始类。
Gysin序列于是给$H^0(E_\alpha)=\mathbb Q(0)$及唯一其他项
$H^{n-1}(E_\alpha)=\mathbb Q(-n/2)$，后者可能携带的二阶特征由二次型有符号判别式决定。
此处无穷远二次型为$M\otimes I_4$，故判别式平方类
$(-1)^{n/2}\det(M\otimes I_4)=U^4$平凡。
两族不被Galois或参数单值交换。
在有限分裂扩域上，线性空间周期给代数Tate类；取迹并除以扩域次数使原始类下降。
这用有理系数，不断言整系数二次型分裂。

上述Gysin关系均来自实际代数嵌入与代数循环。
Nori Betti实现的忠实精确性将这些满秩Tate态射识别为相应同构，
而非从Hodge维数猜测一个动机同构。
对$(Y_\alpha,E_\alpha)$的长正合列，$H^0(Y_\alpha)\to H^0(E_\alpha)$为同构，
剩余项给所述唯一的$H^n$。
相同Gysin与循环构造在光滑参数域给常Tate变化；
其Leray谱序列仅有一条非零相对纤维行，因此对边界仍在$S$内的代数参数对退化。

实正二次型的极点不遇实链，$\mathbb R^n$在射影对中的相对基本链保留标准取向。
其边界在$\mathbb RP^{n-1}$有相应的二重覆盖重数，不可无故换成单位闭球面。
在$M=I$、$b=0$、$V=1$处，以
$(2h)!d^n\ell/(1+|\ell|^2)^{2h+1}$为dR标架，正链周期为$\pi^{2h}$。
这固定了上一公式的Tate常数。其余有限矩为同一个一维相对纤维中的类，
所以解析矩公式只是在一个已由代数循环识别的Tate行上确定其有理系数。
\end{proof}

重要的是这是一条\emph{光滑参数域}上的实际比较。
典型多圈图的整个闭单纯形不满足其边界也在$S$内，
因而单靠此命题不能宣称\eqref{eq:forestparametricfinal}已完成全部边界上的形式比较。下节通过共同扇区模型补齐这一接受域内的箭头。

\subsection{把剩余障碍定位到明确的拟阵边界}

\begin{proposition}[正参数链上的退化面恰由路由秩给出]
\label{prop:forestrankfaces}
令$J$是一个非空参数面之内部的支撑，即$\alpha_j>0$恰在$j\in J$成立。
则
\[
 U(\alpha)>0\ \Longleftrightarrow\
 \operatorname{span}\{v_j:j\in J\}=k^h.
\]
在这些满秩面上$F>0$；在非满秩面上$U=F=0$。
因此实际正参数链接触的退化集合恰为不含路由向量基的坐标面之并。
这是一个有限组合判据，无需求完整积分；它不排除复参数空间中其他$UF=0$分量。
\end{proposition}

\begin{proof}
Cauchy--Binet给
\[
 U(\alpha)=\sum_{|I|=h}\det(v_i:i\in I)^2\prod_{i\in I}\alpha_i.
\]
每个实系数非负，故在指定面内为正当且仅当该支撑含一个基。
满秩时补平方及正质量下界给$F=UV>0$。
非满秩时$U=0$；每个分量$b^\lambda=\sum\alpha_jv_jp_j^\lambda$
属于$M$的像，因而$\operatorname{adj}(M)b^\lambda=0$，这在秩小于$h-1$时也成立。
代入\eqref{eq:forestaugmentedUF}即得$F=0$。
\end{proof}

\begin{corollary}[一圈接受域的完整标架比较]
\label{cor:forestoneloopcomparison}
若$h=1$，则在吸收常数正分母之后，整个闭参数单纯形都在$S$内。
定理\ref{thm:forestincidence}与命题\ref{prop:forestquadricpair}因此给完整的实际形式周期关系
\[
 A_{\rm mom}^{\mathfrak m}
 =-\frac{(\mathbb L^{\mathfrak m})^2}{4}\,
 A_{\rm par}^{\mathfrak m},\qquad
 \operatorname{per}(\mathbb L^{\mathfrak m})=2\pi i,
\]
其中$A_{\rm par}$用公式\eqref{eq:forestmomentformula}的完整$\Psi_N$及正单纯形定义。
对标准每圈$(2\pi)^{-4}$测度，系数相应为
$-\tfrac14(\mathbb L^{\mathfrak m})^{-2}$。
该结论包括接受域内完成整体局部减除的一圈密度及其有限继承喷射，
不以数值周期映射单射为前提。
\end{corollary}

\begin{proof}
此时每个剩余$v_j$是非零标量，所以闭单纯形上
$U=\sum\alpha_jv_j^2\geq\min_jv_j^2>0$。
严格正质量再给$F>0$。取$S$中的坐标面为参数相对边界，
完整实参数链的闭包不接触删去的判别式。
关联空间的形式在只删去$Q_\alpha=0$的较大空间中仍正则；
沿开嵌入推前选定链不会改变其形式周期。
于是可对整个链使用单行相对Leray及Tate纤维比较。
裸Tate常数$\pi^2=-\mathbb L^2/4$与恢复测度后的逆Tate常数均由前述归一化得到。
没有退化参数面的极限需要额外交换。
\end{proof}

对于一般多圈，本节将问题定位为一个明确的相对边界问题：
对这些秩亏坐标面及其交叠取共同分解，延伸实际关联空间的光滑Tate比较，
计算其退化纤维和边界连接对\emph{指定形式与链}的贡献。
这与BEK第10节采用秩分层而不直接宣称对应为同构的做法一致。
共同分母中的路由可能来自不同森林，故其面拟阵不能未经证明替换为原图的子图格。

目前已经证明三件不同的事：原动量周期与关联空间的\emph{形式}比较；
全闭单纯形上绝对收敛的\emph{有理参数表达}及实际参数周期存在；
光滑域上的\emph{Tate纤维}比较。
一圈时已由推论\ref{cor:forestoneloopcomparison}闭合；
一般多圈所需的上段额外工作，在第\ref{sec:sectorboundary}节由贪心扇区和同步动量缩放完成。
同样，所有大载体的同构远强于指定标架的比较，既非本节结论，也非物理最小分类的必要前提。
本节为总目标提供了有限的$N,U,F$结构输入与精确边界位置，
尚未判定其中哪些椭圆、K3或混合组成被最终物理读出保留。


\section{贪心扇区的共同边界模型与全圈标架比较}
\label{sec:sectorboundary}

本节处理上一节留下的多圈秩亏边界。方法是将参数的Hepp分区与圈动量的基变换同步进行：
参数趋零时，相应动量坐标按参数平方根重标度。
经过这一变换，所有嵌套面同时落在一个光滑二次型族中；
完整分子的绝对收敛保证变换后的总形式没有参数负幂。
结果是一条含实际边界的形式周期比较，而非任意大载体的同构。

Hepp分区、贪心基和UV幂次分析均有经典来源。
Panzer的Lemma2.8及Proposition2.9说明贪心基如何控制有序参数扇区\cite{PanzerSector}。
本文保留路由配置的行列式平方系数，不能将它们直接换成任意拟阵的单位系数基多项式。
此外，欧氏绝对动量收敛推出绝对参数收敛的正核论证，
已见de la Cruz--Kosower--Novichkov的Appendix B\cite{CKNSector}；
上一节的这部分分析不是新发现。
以下另给合并分子、共同边界、指定形式与链的完整推导，不宣称一般优先权。

\subsection{沿路由平坦的精确UV条件}

仍用式\eqref{eq:forestcommonfraction}的严格正质量欧氏密度，
吸收常数分母且$N\ne0$。此处先不预设绝对收敛。
称由若干$v_j$张成的子空间$W\subset(k^h)^*$为一个路由平坦；
允许$W=0$。重复的张成空间只记一次。
令$K_W$为$W$的零化子，$k_W=h-\dim W$，并定义
\[
 a_{\mathrm{act}}(W)=\sum_{v_j\notin W}a_j,\qquad
 \delta_W(N)=\deg_{K_W\otimes k^4}N.
\]
后一个次数是固定其余坐标、沿$K_W$方向同时缩放四个动量分量时的最高次数。
它取整个多项式的次数，不是在一个特殊外层动量处取值后的次数；
选择线性补空间不改变它。

\begin{lemma}[绝对收敛的路由平坦必要条件]
\label{lem:sectoruv}
若$R\in L^1(\mathbb R^{4h})$，则路由张成$(k^h)^*$，并对每个真路由平坦有
\begin{equation}
 2a_{\mathrm{act}}(W)-4k_W-\delta_W(N)>0.
 \label{eq:sectoruvcriterion}
\end{equation}
特别地，$W=0$给$2A-4h-d>0$。
\end{lemma}

\begin{proof}
不张成时存在所有分母共同不依赖的实方向，非零多项式沿此方向的绝对积分发散，
故首先必须张成。对真$W$选择$k^h=K_W\oplus L$，写动量为$(z,y)$。
把$y$限制在一个有正体积的紧小邻域，使$N$关于$z$的最高齐次项不恒零。
在$z$的单位球面上再取一个紧小片，使这个首项与每个活跃路由的限制均不为零。
这些条件只排除有限个真代数集，故可以同时实现，并可缩小两个小邻域获得一致界。
当$z=t\theta$且$t$充分大时，$|N|$与$t^{\delta_W}$同阶，
$v_j\notin W$的分母与$t^2$同阶；$v_j\in W$的分母只依赖$y$且有正上下界。
于是该锥域上的径向绝对积分与
\[
 \int^{\infty}t^{4k_W-1+\delta_W-2a_{\mathrm{act}}(W)}dt
\]
同阶。收敛迫使式\eqref{eq:sectoruvcriterion}严格成立。
改变补空间只将$z$加上依赖$y$的线性项，不改变关于$z$的最高次数。
\end{proof}

这些条件检测的是实际合并密度的绝对收敛。
它们不把角平均后的相消、维数正则化中的消失系数或解析延拓的有限部分算作$L^1$收敛。

\subsection{每个扇区中的同步参数与动量坐标}

给定一个参数排序，重排指标后在对应扇区取射影坐标
\[
 x_1=1\geq x_2\geq\cdots\geq x_r\geq0,
 \qquad \alpha_j=x_j/\sum_i x_i.
\]
选取正链诱导的取向；不额外附加排列符号。
令$A_k=\sum_{j=k}^r a_j$，并作平方比例换元
\begin{equation}
 x_j=s_j^2,\qquad s_1=1,\qquad
 s_j=\prod_{k=2}^j t_k,\qquad 0\leq t_k\leq1.
 \label{eq:sectorsquarecoords}
\end{equation}
按顺序贪心选择路由基，其枢轴位置为
$b_1=1<b_2<\cdots<b_h$。
以$v_{b_i}$为行组成可逆矩阵$C$，写
\[
 v_j=\sum_i c_{ji}v_{b_i},\qquad
 c_{ji}=0\ (b_i>j),\qquad c_{b_i l}=\delta_{il}.
\]
令$y_i=v_{b_i}\ell$、$z_i=s_{b_i}y_i$。
每个$y_i,z_i$都是四维向量；四份相同的线性换元使取向因子为$(\det C)^4>0$。
记$W_k=\operatorname{span}(v_1,\ldots,v_{k-1})$及$k_k=h-\dim W_k$。

\begin{proposition}[全闭扇区的非退化二次型模型]
\label{prop:sectorgreedy}
在上述坐标中，$Q_x=\sum x_jD_j$变为多项式二次型
\begin{equation}
 \widetilde Q(t,z)=
 \sum_{j=1}^r\left|
   \sum_{b_i\leq j}c_{ji}\Bigl(\prod_{k=b_i+1}^j t_k\Bigr)z_i+s_jp_j
 \right|^2+\sum_{j=1}^r x_j\mu_j .
 \label{eq:sectornormalquadric}
\end{equation}
在整个实闭立方体上，其圈指标矩阵$\widetilde M\geq I_h$。
补平方后的最小值$\widetilde V\geq\mu_1>0$；
$\widetilde U=\det\widetilde M$及$\widetilde F=\widetilde U\widetilde V$
均不在该闭立方体上为零。
此外，存在统一常数$c>0$使
$\widetilde Q(t,z)\geq c(1+|z|^2)$。

若式\eqref{eq:sectoruvcriterion}对所有真路由平坦成立，则
\begin{equation}
 P(t,z)=\left(\prod_{k=2}^r t_k^{2A_k-1-4k_k}\right)
 N\left(C^{-1}\operatorname{diag}(s_{b_i}^{-1})z\right)
 \label{eq:sectorpolynomialnumerator}
\end{equation}
实际上属于$k[t,z]$，且$\deg_zP\leq d$。
完整关联空间形式在此扇区的拉回为
\begin{equation}
 \frac{2^{r-1}}{B_a(\det C)^4},
 \frac{P(t,z)}{\widetilde Q(t,z)^A},d^{4h}z,dt_2\cdots dt_r.
 \label{eq:sectorregularform}
\end{equation}
它在所有$t_k=0,1$面及其交叠处都没有参数极点。
\end{proposition}

\begin{proof}
按射影齐次次数，$\alpha^{a-1}Q_\alpha^{-A}d\alpha$
在$x_1=1$图中等于$\prod_{j=2}^r x_j^{a_j-1}Q_x^{-A}dx_2\cdots dx_r$。
$\sum x_i$的全部因子恰好消去。三角Jacobian再给
\[
 \prod_{j=2}^r x_j^{a_j-1}dx_2\cdots dx_r
 =2^{r-1}\prod_{k=2}^r t_k^{2A_k-1}dt_2\cdots dt_r.
\]
若$c_{ji}\ne0$，则$b_i\leq j$，从而$s_j/s_{b_i}$是所示乘积多项式，
直接代入便得式\eqref{eq:sectornormalquadric}。
每个枢轴行$j=b_i$在和式中贡献$|z_i+s_{b_i}p_{b_i}|^2$，
因此$\widetilde M\geq I_h$。
质量项含$x_1\mu_1=\mu_1$，补平方后的最小值仍至少为$\mu_1$。
最小点$z_*(t)$连续且有界；由$|z-z_*|^2+\mu_1$的下界得到统一的$c$。

动量Jacobian为
$(\det C)^{-4}\prod_i s_{b_i}^{-4}
=(\det C)^{-4}\prod_{k=2}^r t_k^{-4k_k}$。
在$y=C\ell$坐标中，$K_{W_k}$恰由那些$b_i\geq k$的$y_i$坐标张成。
所以替换$y_i=z_i/s_{b_i}$后，分子关于$t_k$可能出现的最大负次数
恰为$\delta_{W_k}(N)$。
若$k_k>0$，所有不在$W_k$中的路由均属于尾部$j\geq k$，因而
$A_k\geq a_{\mathrm{act}}(W_k)$。
利用整数幂及严格UV条件，得到
\[
 2A_k-1-4k_k-\delta_{W_k}(N)\geq0.
\]
若$k_k=0$，则不存在该负次数，且$2A_k-1\geq1$。
逐一看每个单项式的每个$t_k$次数，均非负，故$P$确为多项式。
这同时控制所有角点，不依赖依次逼近不同参数面的次序。
混合的$dt$项与顶次参数体积楔乘后为零，故上述Jacobian给完整形式。
\end{proof}

这一模型还将路由退化的公共因子显式化。
用$U_x,F_x$表示式\eqref{eq:forestaugmentedUF}在非归一化$x$参数中的版本，则
\begin{equation}
 U_x=(\det C)^2\Bigl(\prod_i x_{b_i}\Bigr)\widetilde U,
 \qquad
 F_x=(\det C)^2\Bigl(\prod_i x_{b_i}\Bigr)\widetilde F.
 \label{eq:sectorUFfactor}
\end{equation}
第一式由矩阵合同变换取行列式得到；第二式再用补平方最小值不变。
这里$\prod_i x_{b_i}=\prod_{k=2}^r t_k^{2k_k}$。
剥去共同单项式后，两个剩余多项式在正闭链上都是单位。
不需要声称它们在整个复空间中没有零点，亦未将$(\det C)^2$换成$1$。

\subsection{一个有限且双向的收敛判据}

\begin{theorem}[正质量欧氏密度的有限接受判据]
\label{thm:sectorconvergence}
对本节给定的非零多项式分子与正质量二次分母，
$R\in L^1(\mathbb R^{4h})$当且仅当路由张成$(k^h)^*$且
式\eqref{eq:sectoruvcriterion}对每个真路由平坦成立。
此判据只需有限线性代数和多项式次数，不需求积分。
\end{theorem}

\begin{proof}
必要性已证。反过来，$W=0$先保证$2A>4h+d$。
对$|N|$使用正Dirichlet核的Tonelli等式，此时允许两侧取$+\infty$，
所以论证并未预设收敛。将参数链分为有限个有序扇区，并逐个采用前述坐标。
由于$P$在闭立方体上系数有界且$\deg_zP\leq d$，有
\[
 \frac{|P(t,z)|}{\widetilde Q(t,z)^A}
 \leq C\frac{1+|z|^d}{(1+|z|^2)^A}.
\]
右端在$\mathbb R^{4h}$可积，并且与$t$无关。
每个扇区的绝对积分有限，有限求和及Tonelli等式证明原密度绝对可积。
路由平坦至多由$2^r$个子集生成；各方向次数可在一个有理基中展开后读取。
因此测试是有限且精确的，但这里不宣称关于$r$为多项式复杂度。
\end{proof}

作为纯理论构造，这也说明为何此前必须先合并分子。
完整$N$的$\delta_W$可以因真实多项式相消而降低。
在原动量坐标中随意拆成单项式，可能破坏式\eqref{eq:sectoruvcriterion}；
在式\eqref{eq:sectorpolynomialnumerator}成为真正多项式之后再拆分，
各项才获得上述共同可积界。
本结论是指定正质量欧氏域的UV判据，
不解决一般Minkowski或无质量分子空间的参数多面体猜想\cite{CKNSector}。

\subsection{包含全部边界的实际形式周期关系}

\begin{theorem}[接受域内任意有限圈数的标架比较]
\label{thm:sectorformalcomparison}
设式\eqref{eq:forestcommonfraction}满足上述等价接受条件，$h\geq1$任意有限。
令$A_{\rm mom}^{\mathfrak m}$为其裸动量测度的实际形式周期，
$A_{\rm par}^{\mathfrak m}$为式\eqref{eq:forestmomentformula}
完整有理形式$\Psi_N$及正单纯形所定义的参数形式周期。
在相容的实际Nori形式周期代数中有
\begin{equation}
 A_{\rm mom}^{\mathfrak m}
 =(-1)^h2^{-2h}(\mathbb L^{\mathfrak m})^{2h}
 A_{\rm par}^{\mathfrak m}.
 \label{eq:sectorformalidentity}
\end{equation}
若每圈动量测度含$(2\pi)^{-4}$，对应关系为
\begin{equation}
 A_{\rm norm}^{\mathfrak m}
 =(-1)^h2^{-2h}(\mathbb L^{\mathfrak m})^{-2h}
 A_{\rm par}^{\mathfrak m}.
 \label{eq:sectornormalizedidentity}
\end{equation}
这里$\operatorname{per}(\mathbb L^{\mathfrak m})=2\pi i$。
比较包含所有正参数边界及嵌套角点，不假设周期映射单射，
不宣称任意选取的大上同调载体同构。
\end{theorem}

\begin{proof}
\emph{第一步：实际链的分割与提升。}
定理\ref{thm:forestincidence}已经通过有理换元和K\"unneth给出动量周期到关联空间的形式比较。
在关联空间中沿$\alpha_i=\alpha_j$分割实际参数链，
将这些代数超平面加入暂时的相对边界，得到各有序扇区周期之和。
相邻扇区的公共内部面具有相反边界取向；求和恢复原相对链。
这一步是相对链加法与边界映射，不是根据积分数值相等作的推断。

在每个扇区采用式\eqref{eq:sectorsquarecoords}及$\ell=C^{-1}\operatorname{diag}(s_{b_i}^{-1})z$。
前者是代数有限覆盖在正实分支上的提升，后者在所有$t_k\ne0$上为有理同构。
正开立方体映到原扇区内部恰一次；不能因复覆盖的次数再除以$2^{r-1}$。
该因子已经是式\eqref{eq:sectorregularform}中的Jacobian。
采用两个模型的共同图闭包与相容分解，链提升、形式拉回和链推前相匹配。
在$t_k=0$处原动量坐标可能去向无穷远，正是此共同图闭包所记录的边界。

\emph{第二步：边界处没有遗漏的极点项。}
命题\ref{prop:sectorgreedy}证明拉回的\emph{整个顶次形式}的分子为$P\in k[t,z]$。
闭立方体上$\widetilde Q$没有实零点，故在所有有限$z$及所有参数角点均正则。
再把$z$紧化为$\mathbb P^{4h}$。
若一个分子单项式关于$z$的次数为$e\leq d$，
其形式在一个无穷远仿射图中的法向次数为
\[
 2A-e-(4h+1)\geq0.
\]
不等式来自$2A>4h+d$及整数性。
因此只需删去二次型的复极点及其无穷远截面；
实际正链的紧化边界不落在其中。
共同图模型中的实际形式/链由这些正则拉回给出，
而不是把原来发散的单项式分别赋予周期。
HMS的实际链实现及代数换元的函子性\cite{HMSMass}于是适用于这些相对模型。
全证明不使用参数截断后取极限来代替相对态射。

\emph{第三步：在全闭扇区上消去实际Tate纤维。}
令$S_\sigma=\{\widetilde U\widetilde F\ne0\}\subset\mathbb A^{r-1}$，
参数边界为其中$t_k=0,1$的并。
整个闭立方体都在$S_\sigma$中；
齐次化的$\widetilde Q=0$及其无穷远截面均光滑。
应用命题\ref{prop:forestquadricpair}的实际Gysin/Tate比较：
圈指标矩阵仍出现四次，故其有符号判别式为$\widetilde U^4$，
两族线性空间的原始类没有二阶交换特征。
有理系数下的代数循环、有限分裂域取迹及相对Leray
将指定总空间类降到参数相对类并附加$\mathbb Q(-2h)$。
与上一节不同，现在这个比较覆盖\emph{全部实际边界}。
在正定标准纤维中，通常实动量链固定其周期为$\pi^{2h}$，
即$(-1)^h2^{-2h}\mathbb L^{2h}$。

\emph{第四步：识别推前后的指定参数形式。}
将式\eqref{eq:sectorregularform}的多项式$P$作有限矩消去。
所有矩由统一可积界控制，结果是$S_\sigma$上的有理形式。
在$t_k>0$内部，补平方、线性坐标变换及有限Gaussian矩公式
表明它恰为$\Psi_Nd\alpha$的扇区拉回。
这是有理系数函数的恒等式：两边表示同一一维Tate纤维的dR类；
矩公式只是计算这个类的有理系数，而非以一般数值周期关系代替动机关系。
它因而延伸到整个正则扇区模型。
逆用同一代数参数换元并将扇区相对链相加，
就恢复$A_{\rm par}^{\mathfrak m}$，得到式\eqref{eq:sectorformalidentity}。

最后，每圈$(2\pi)^{-4}$在标准Tate规范中为$(\mathbb L^{\mathfrak m})^{-4}$。
乘上$h$份即得式\eqref{eq:sectornormalizedidentity}。
改变扇区细分或路由基可在共同细分/图模型中比较，故不改变这个指定周期关系。
结论并不要求所有辅助模型的无关上同调相同。
\end{proof}

\subsection{对总目标的作用与下一处真正缺口}

\begin{corollary}[完整森林的有限参数几何入口]
\label{cor:sectorforestentry}
在第\ref{sec:massiveforestperiods}节的有质量欧氏BPHZ接受域内，
任意有限圈完整森林、有限继承喷射及已证明收敛的局部插入，
均可先通分再通过有限代数步骤输出：合并$N,U,F$、
覆盖实际参数链的正则扇区模型、指定参数形式和链，以及其与动量形式周期的Tate比较常数。
构造不需求先算出完整积分；它给实际几何入口，尚不给最小纯核分类。
\end{corollary}

\begin{proof}
前述BPHZ入口给非零完整密度的$L^1$条件与代数二次分母。
因此引理\ref{lem:sectoruv}保证所有方向次数条件。
至多$r!$个有序扇区中，贪心基、式\eqref{eq:sectornormalquadric}与
式\eqref{eq:sectorpolynomialnumerator}只用有限线性代数和多项式运算。
矩公式再输出有理参数形式，定理\ref{thm:sectorformalcomparison}提供实际标架比较。
若完整分子为零，直接输出零周期，不需任何扇区。
不为各个未经整体减除的发散森林项单独作此断言。
\end{proof}

这里没有因“Tate纤维”得到“所有图都是混合Tate”的结论。
$S_\sigma$只是在\emph{实际正链}上避开$\widetilde U\widetilde F=0$；
这些复代数判别式、参数相对边界及其粘合仍可承载非Tate组成。
消去动量后的几何复杂性留在参数底空间，正是总分类目标下一步需要研究的对象。

这也给出清楚的执行次序：先建立上述实际比较，随后从参数几何给组成上界，
再由物理链的管变分、非齐次源及最小标架检验确定哪些组成保留。
接受域中的多圈秩亏边界比较现已由本节处理，不能继续把它写成该域的未完成项。
真正未完成的范围包括无质量/IR边界、一般Minkowski物理链的全局比较、
维数正则化的非整数幂以及一般参数载体的最小几何/算术组成。
有限接受算法和指定标架比较没有自动解决这些问题，也不证明具体周期数字超越于MZV代数。


\section{偶数维四粒子谱：同一K3核与物理读出的保持}\label{sec:evendimk3}
前两节把完整减除的正质量欧氏密度送入实际参数几何。
接下来需要判断几何上界是否被物理读出达到。
本节给一个跨维数的具体结论：在指定非等质量射线上，
单位四粒子相空间在每个正偶数维都保留同一个六维K3简单因子。
维数改变积分微分及混合边界，不能直接由曲面名称推出这一结论。
证明依次构造实际代数对应、控制极点背景，再用正常阈值的局部变分排除消去。
所有最小阶数及简单重数均在参数函数/正则奇异微分系统层陈述，
不把它们当成某个固定数值的超越性结论。

\subsection{从banana环面到谱曲面的四重有理对应}
沿用\eqref{eq:spectralk3}的$A,B$及$\lambda$，记
\[
 Q_s(a,b)=A(a)B(b)\lambda(s,a,b),\qquad \mu_i=m_i^2.
\]
banana的环面模型及裸留数为
\begin{align}
 Z_s:&\quad f(\alpha)=\Bigl(\sum_i\mu_i\alpha_i\Bigr)
                 \Bigl(\sum_i\alpha_i^{-1}\Bigr)=s,\label{eq:evenztorus}\\
 \Omega_s&=\operatorname{Res}_{f=s}
 \frac{d\log(\alpha_1/\alpha_4)\wedge d\log(\alpha_2/\alpha_4)
                  \wedge d\log(\alpha_3/\alpha_4)}{f-s}.
 \label{eq:evenzres}
\end{align}
这里留数采用$d(f-s)\wedge\operatorname{Res}(\cdot)$恢复分子的约定。
以下$T_Z,T_X$均指一般纤维代数除子局部系统的正交补；
在一般纤维上它就是超越上同调。
特殊Noether--Lefschetz纤维可能多出代数类，本节不宣称每个光滑纤维的Picard数都相同。
令$u_i=\mu_i\alpha_i$，取共同尺度使$u_4=1$，并置
\[
 r=\frac{u_1}{u_2},\quad t=\frac{u_3}{u_4},\quad
 z=\frac{u_1+u_2}{u_3+u_4},\quad
 a=(r+1)(\mu_1/r+\mu_2),\quad b=(t+1)(\mu_3/t+\mu_4).
\]
这些是稠密开集上的有理坐标，不要求在全部边界上可逆。

\begin{theorem}[实际商对应及有理纯层比较]\label{thm:evenquotient}
在光滑一般参数上，公式
\begin{equation}\label{eq:evenquotient}
 (r,t,z)\longmapsto (a,b,y),\qquad
 y=(\mu_2r-\mu_1/r)(\mu_4t-\mu_3/t)(bz-a/z)
\end{equation}
给$Z_s\dashrightarrow X_s$的次数四有理映射$\pi$，且
\begin{equation}\label{eq:evenformpull}
                  \Omega_s=-\pi^*\omega_s.
\end{equation}
在指定质量$m=(1,2,4,8)$的光滑参数射线上，
它给超越局部系统的有理同构$T_X\simeq T_Z$，秩为六。
此处采用本项目已经证明的该射线$T_Z$单值闭包$\mathrm O(Q_6)$及绝对不可约性。
拉回将杯积乘以四；结论不是整格同构，不转移旧整周期基标记。
\end{theorem}
\begin{proof}
由定义直接得到
\[
 f=(z+1)(a/z+b),\qquad
 \begin{cases}
 \mu_2r^2+(\mu_1+\mu_2-a)r+\mu_1=0,\\
 \mu_4t^2+(\mu_3+\mu_4-b)t+\mu_3=0,\\
 bz^2+(a+b-s)z+a=0.
 \end{cases}
\]
三个判别式分别为$A(a),B(b),\lambda(s,a,b)$。
各自有另一因子不含的分支素除子，因而它们的平方类独立。
故$k(Z_s)=k(a,b)(\sqrt A,\sqrt B,\sqrt\lambda)$为次数八扩张。
三个根可用\eqref{eq:evenquotient}中的三个括号表示。
偶数次翻号的四元子群固定其乘积$y$，固定域正是$k(X_s)$。
反向恢复为$u_3=t,u_2=z(t+1)/(r+1),u_1=ru_2$，
证明这是原环面的函数域，而非另设的一张根式覆盖。

再有$da=(\mu_2r-\mu_1/r)d\log r$，$db$有同样公式，
$z\partial_z f=bz-a/z$，以及
\[
 d\log u_1\wedge d\log u_2\wedge d\log u_3
       =-d\log r\wedge d\log t\wedge d\log z.
\]
取留数即得\eqref{eq:evenformpull}，包括其负号。

对有理映射取图闭包的光滑消解，得到实际代数对应。
在超越上同调上，推拉复合为$\pi_*\pi^*=4$，
消解例外类的修正为代数类；所以$T_X\hookrightarrow T_Z$且配对乘四。
全纯二形式拉回非零，因此这是非零子局部系统。
指定射线的$T_Z$绝对不可约、秩六，像必为整个$T_Z$。
结合定理\ref{thm:spectralk3}的秩上界，得到有理同构。
这一步使用既有纯单值定理，未进行新运输或点数。
它不依赖周期映射单射，也没有从同一数值积分猜测动机相等。
\end{proof}

\subsection{升维微分的极点只增加边界背景}
在$\widetilde X_s$上，把$a=0,\infty$或$b=0,\infty$的逆像
连同其所经例外曲线作成正常交叉除子$D_s$；必要时再作有限次点爆破。
记$U_s=\widetilde X_s\setminus D_s$。本节所有这些符号均指适当公共光滑开基上的族。
对整数$k\ge1$定义
\begin{equation}\label{eq:eveneta}
 \eta_{k,s}=\left(\frac{Q_s(a,b)}{sab}\right)^{k-1}\omega_s.
\end{equation}

\begin{proposition}[一个K3因子与Artin--Tate边界]\label{prop:evenboundary}
$\eta_{k,s}$在$U_s$上正则。
沿指定质量射线，$H^2(U_s)$的组成因子中，$T_X$恰出现一次，
其它因子均为Artin--Tate型；这里允许非平凡边界扩张。
在正常阈值$\tau=(\sum_i m_i)^2=225$附近，$D_s$没有新增退化。
\end{proposition}
\begin{proof}
在$a,b$非零有限的区域，$Q_s/(sab)$为正则函数，
故乘以全局全纯$\omega_s$不会在分支或有限$A_1$消解处产生极点。
其它极点均已列入$D_s$。
其归一化分量都是有理曲线。例如在$a=0$上，
\[
 y^2=A(0)B(b)(s-b)^2,
\]
去掉平方因子后是只有两个分支点的双覆盖；在$a=\infty$上，
适当齐次化后得到$v^2=B(b)$。另两个坐标边界同理。
所有新增点爆破及有理双点消解只加入有理曲线。

正常交叉开曲面权谱序列的各项来自紧曲面、边界曲线和交点。
有理边界曲线没有$H^1$，其$H^0,H^2$及交点为Artin--Tate对象。
紧曲面的一般代数除子局部系统亦是Artin--Tate型：在有限代数扩域上选取有限组除子即可。
紧曲面的$T_X$不可能被只含代数类的Gysin像消去。
定理\ref{thm:evenquotient}给其秩六与不可约性，故重数恰为一。
在参数局部系统层，有限代数扩展对应的置换因子可有有限单值，
但不能等同于具有$\mathrm{SO}_6$连通单值的$T_X$。

最后检查边界随$s$的变化。$a=0$上的两支在$b=s$相交，
此点撞上$B$的分支点仅当$s=b_\pm$；$b=0$同理仅当$s=a_\pm$。
两零坐标的交点只在$s=0$合并。无穷远归一化方程与$s$无关，
在$(\infty,\infty)$的$A_2$局部式中$s$为非零单位，
因而其标记消解及边界交点在$s\ne0$的小圆盘上同型。
这些坐标图也覆盖各零、无穷远边界的交会；再爆破其移动节点得到相对正常交叉除子。
因此边界可能新增碰撞的值包含在$\{0,a_\pm,b_\pm\}$内。
正常阈值严格大于$a_+,b_+$，不在这个集合中。
这只是一处阈值附近的结论，没有宣称整个参数平面不存在边界单值。
\end{proof}

\subsection{实际偶数维相空间与严格正周期}
取有标号的四个末态粒子，使用
\[
 d\Phi_4^{(D)}=(2\pi)^D\delta^{(D)}(P-\textstyle\sum_i p_i)
       \prod_{i=1}^4\frac{d^{D-1}\mathbf p_i}{(2\pi)^{D-1}2E_i},
 \qquad P^2=s,\quad E_i>0.
\]
不在这里除以$4!$；有相同粒子的另一个物理过程必须自行加入其对称因子。
写$\Phi_4^{(D)}(s)=\int d\Phi_4^{(D)}$，并定义
\[
 c_D=\frac{\operatorname{vol}(S^{D-2})}{2^{D-1}(2\pi)^{D-2}}.
\]
偶数维包括$D=2$，此时$\operatorname{vol}(S^0)=2$。

\begin{proposition}[升维只改变指定代数微分]\label{prop:evenphase}
对$D=2k\ge2$、$s>\tau$，前述物理双三角链$\Gamma_s$满足
\begin{equation}\label{eq:evenphase}
 \Phi_4^{(2k)}(s)=\frac{c_{2k}^3}{8\pi^2}
                         \int_{\Gamma_s}\eta_{k,s}>0.
\end{equation}
尤其在四维
\begin{equation}\label{eq:fourphasek3}
 \Phi_4^{(4)}(s)=\frac1{2048\pi^5}
       \int_{\mathcal R_s}\frac{\sqrt{Q_s(a,b)}}{sab}\,da\,db
       =\frac1{4096\pi^5}\int_{\Gamma_s}\eta_{2,s}.
\end{equation}
设$M=\sum_i m_i$，则它的四维正常阈值首项为
\begin{equation}\label{eq:fourphasethreshold}
 \Phi_4^{(4)}(M^2+\varepsilon)
 =\frac{\sqrt{\prod_i m_i}}{13440\pi^4M^5}\,
              \varepsilon^{7/2}\bigl(1+O(\varepsilon)\bigr),\quad\varepsilon\downarrow0.
\end{equation}
\end{proposition}
\begin{proof}
在总质心系积分二体能量$\delta$函数，径向导数为
$p(E_1+E_2)/(E_1E_2)$，$p=\sqrt{\lambda(s,a,b)}/(2\sqrt s)$。
连同角球面积及原始$2\pi$测度，得到
\[
 \Phi_2^{(D)}(s;a,b)=c_D
            \frac{\lambda(s,a,b)^{(D-3)/2}}{s^{(D-2)/2}}.
\]
两对分组的实际相空间分解给
\[
 \Phi_4^{(D)}=
 \int_{\mathcal R_s}\frac{da\,db}{(2\pi)^2}\,
 \Phi_2^{(D)}(s;a,b)\Phi_2^{(D)}(a;\mu_1,\mu_2)
                         \Phi_2^{(D)}(b;\mu_3,\mu_4).
\]
这是经典相空间递归\cite{KreimerBanana,PDGphase}，此处从所声明的测度另行固定常数。
当$D=2k$时，三个根的积为$Q_s^{k-3/2}$，分母为$(sab)^{k-1}$。
两片定向粘合使$\int_{\Gamma}\eta_k$等于该平面正积分的两倍，
得到\eqref{eq:evenphase}；$c_4=1/(8\pi)$给\eqref{eq:fourphasek3}。
链远离$D_s$，其三个角的例外曲线上，正则二形式积分为零。
所有内点的平面密度严格为正，根端点可积。

为求阈值常数，令$E=\sqrt s-M$，用动量守恒消去第四个三动量。
非相对论能量二次型在每一空间分量上的矩阵为
\[
 A_0=\operatorname{diag}(m_1^{-1},m_2^{-1},m_3^{-1})
                         +m_4^{-1}\mathbf1\mathbf1^t,
 \qquad\det A_0=\frac{M}{\prod_i m_i}.
\]
九维橢球的能量壳积分为
$(2\pi)^{9/2}E^{7/2}/(\Gamma(9/2)(\det A_0)^{3/2})$。
乘原测度$(2\pi)^{-8}/(16\prod_i m_i)$，
再用$E=\varepsilon/(2M)+O(\varepsilon^2)$，恰给\eqref{eq:fourphasethreshold}。
对动量作$\sqrt E$缩放后，正定球面上余项一致为$O(E)$，故所写相对误差成立。
同一缩放在$D$维给正首项和指数$(3D-5)/2$，不需要完整积分。
\end{proof}

\subsection{排除“非零但只在边界”的可能}
以下引理说明为何正性还必须配合实际局部变分。
令$H$为一个有限秩同调局部系统，$H\twoheadrightarrow T$为到简单局部系统的商，
$\ell$为代数de Rham读出。对一个循环的全局平行运输轨道取张成$R$，
再商去所有使$\ell$的周期函数恒零的平坦向量$P$；$R/P$是这个读出的最小函数局部系统。

\begin{lemma}[秩一局部变分迫使简单因子可见]\label{lem:evenvisibility}
设局部单值$T_*$满足$\im(T_*-1)=\mathbb C\delta$，
$\delta$在简单商$T$上的投影非零，$(T_*-1)|_T\ne0$，
且$\ell(\delta)$不是恒零函数。
则由$\delta$生成的物理周期的最小函数局部系统含有$T$。
若整个$H$只含一份$T$，这个重数也恰为一。
\end{lemma}
\begin{proof}
$\delta$的全局轨道$R$投到$T$的非零子系统，所以满射到$T$。
记投影为$p$，周期恒零子系统为$P\subset R$。
若$R/P$不含$T$，则$R/(P+\ker p)$不可能是非零的简单商$T$；
故$p(P)=T$。由假设，$T_*-1$在$T$上非零。
于是$(T_*-1)P$非零，只能等于$\mathbb C\delta$，从而$\delta\in P$，
与读出非零矛盾。本节的反射情形由$T_*\delta=-\delta$明确满足非平凡性假设。
重数上界随子商保持，得到最后断言。
\end{proof}

\begin{proposition}[谱曲面正常阈值的实际秩一变分]\label{prop:evenlocal}
在$m=(1,2,4,8)$上，绕$s=\tau=225$的局部单值对$H_2(U_s)$满足
\[
 \rk(T_\tau-1)=1,\qquad T_\tau\delta=-\delta,
 \qquad \int_\delta\eta_{k,s}\not\equiv0\quad(k\ge1),
\]
其中$\delta$为一个有理局部循环，其紧曲面投影位于$T_X$且非零。
这不是把奇异谱曲面的阈值直接称为普通节点。
\end{proposition}
\begin{proof}
正常阈值使三条分支$a=a_+,b=b_+,\lambda=0$在同一点相交，
其三条切线互异。双覆盖局部型为$D_4$，主项为
$y^2=uv(u+v)$，而非$A_1$。
其Jacobian代数的基可取$1,u,v,uv$，Milnor数为四；
标准有理双点的局部交叉格是非退化负定$D_4$格\cite{HuyK3,WahlD4}。
在邻近$s\ne\tau$，三对分支相交给三个分离的$A_1$。
在小参数圆盘的穿孔部分，它们逐个有标签：
固定竖横交点、竖线与$\lambda$的近交点、横线与$\lambda$的近交点。
它们没有绕阈值交换。消解后的三个例外类$E_1,E_2,E_3$
被单值逐个固定，交叉矩阵为$-2I_3$。

取包围三点的小Milnor邻域。把各$A_1$的小邻域换成其消解后，
局部二同调仍为秩四的$D_4$空间：有理双点的同时消解定理
将其局部变形与$D_4$配置的变形识别\cite{WahlD4}；
$A_1$的平滑和消解也是同一$-2$圆盘丛的相应定向实填充，边界均为同一链接。
同时消解所用的有限基变换只用来识别这一个局部交叉空间，
下面的单值仍沿原$s$平面的圆周计算，不以基变换杀掉反射。
边界除子$D_s$远离此邻域，且由命题\ref{prop:evenboundary}在阈值附近同型。
因此可取外部相对边界的平凡化，使单值变分像支撑于这一个局部空间。
对任意开曲面二循环$\gamma$，交叉配对不变性给
\[
 ((T_\tau-1)\gamma,E_j)
          =(\gamma,(T_\tau^{-1}-1)E_j)=0.
\]
秩四空间中三个非退化例外类的正交补只有一维，故变分秩至多一。

物理双三角形连同三个消解填充全部可放在这个局部邻域内。
写其类为$\Gamma_s=c\delta+\sum_j c_jE_j$，
其中$\delta$生成上述一维正交补，$c,c_j$为有理数。
所有$\eta_k$在这里正则，限制到复曲线$E_j$为零。
特别地，$\int_\Gamma\omega>0$证明$c\ne0$。
$D=2$的阈值缩放给$\int_\Gamma\omega=\varepsilon^{1/2}h(\varepsilon)$，
$h$在零点解析且非零；解析性由固定缩放区域上的正定解析密度及可积端点给出。
所以该一维正交方向单值为$-1$，秩确为一。
再由\eqref{eq:evenphase}的正性，$\int_\delta\eta_k=c^{-1}\int_\Gamma\eta_k\ne0$。

紧曲面的十六个已构造除子类在此小圆盘内均有标签并被固定；
它们由两个ruling、十二个$A_1$和两个$A_2$类构成，张成一般纤维的$\mathrm{NS}_\mathbb Q$。
变分像与它们正交，因而$\delta$投到$T_X$。
其全纯周期非零又排除投影为零。
这里使用的是实际局部循环及其配对，不是把高维相空间的正数直接解释成纯K3类。
\end{proof}

\begin{theorem}[任意正偶数维的单位四粒子谱保留一份六维K3核]\label{thm:evenk3visible}
固定$m=(1,2,4,8)$，令$D=2k\ge2$。
把$\Phi_4^{(2k)}(s)$乘以\eqref{eq:evenphase}中的非零常数的倒数后，
其最小正则奇异微分系统恰含一份六维$T_X\simeq T_Z$简单因子。
其最小齐次有理微分方程阶数至少为六。
维数增长可能改变其它Artin--Tate组成及扩张，本定理不宣称完整系统恰为六维。
\end{theorem}
\begin{proof}
实际相空间等式给$H^2(U_s)$中的形式$\eta_k$与实际链$\Gamma_s$。
命题\ref{prop:evenboundary}给至多一份$T_X$的上界。
正常阈值邻域内，例外曲线的全部$\eta_k$周期为零，
而$\Gamma_s$与$\delta$的周期函数相差非零有理常数；
解析延拓保持这个函数等式。
命题\ref{prop:evenlocal}给秩一非平凡变分和非零读出，
定理\ref{thm:evenquotient}给简单纯商$T_X$。
从开曲面到紧曲面的二同调像是对边界除子交叉为零的子空间；
边界均为代数类，故再投到超越部分是满射。
对$H_2(U_s)$到紧曲面超越部分的投影应用引理\ref{lem:evenvisibility}，
得到至少一份$T_X$，上下界相等。
代数Gauss--Manin系统的正则奇异性及微分子商对应将结论转为
$\mathbb C(s)$上最小微分模的组成重数；其秩至少为六。
有限代数换基不消去$\mathrm{SO}_6$的不可约作用，
但本证明没有用有限覆盖之后仍有同一局部反射这一错误前提。
\end{proof}

对其它质量，\eqref{eq:evenphase}及四重对应仍在声明的接受域成立。
若另有该质量射线纯部分不可约及相应边界/阈值条件的证据，同一论证可以再应用。
本节未将指定质量的秩六结论直接推广到等质量、Picard跃迁点或任意其它代数质量。
维数变化的结果是一条带假设的结构保持规则，而非“维数与周期类型无关”。

\subsection{四维未割读出与局部减除的出口}
四种自由实标量场的局部复合算符$\mathcal O=\phi_1\phi_2\phi_3\phi_4$
具有单位四粒子树形矩阵元。其二点谱密度正比于上述有标号相空间。
为固定归一化，直接取$\rho(s)=\Phi_4^{(4)}(s)$，令$s_0<0$。

\begin{corollary}[三次减除的四维谱读出及K3保护]\label{cor:evendispersion}
下式绝对收敛，并给一个完整的三次减除二点读出：
\begin{equation}\label{eq:evendispersion}
 F_R(s)=(s-s_0)^3\int_{\tau}^{\infty}
              \frac{\rho(x)\,dx}{(x-s)(x-s_0)^3},\qquad
              \operatorname{Disc}F_R(s)=2\pi i\rho(s).
\end{equation}
在指定质量射线上，$F_R$的最小微分系统含有六维K3因子。
固定质量及复合算符，只改变这个二点读出的有限局部接触项减除条件时，
允许的二次多项式差不能消去该因子。
\end{corollary}
\begin{proof}
由\eqref{eq:fourphasek3}，$\sqrt{A(a)}/a\le1$、$\sqrt{B(b)}/b\le1$，
且$0\le\sqrt{\lambda(x,a,b)}/x\le1$。
积分区域包含在$0\le a,b\le x$中，故$\rho(x)=O(x^2)$。
三次减除使色散核在无穷远至多为$O(x^{-2})$；
阈值处\eqref{eq:fourphasethreshold}保证可积。
因此在割外局部一致绝对收敛，Sokhotski公式给所写跳跃及其正号。
这与经典banana运动学减除色散式的$D=4,n=4$情形一致\cite{KreimerBanana}；
这里的$\rho$是跳跃除以$2\pi i$，不混同该文的虚部归一化。

相空间和最后一个色散变量给实际代数积分载体与正链；
在无穷远的收敛减除允许同样的相对周期紧化。
在$x=s$处取留数，其管链变分就是$2\pi i$倍的相空间链，
故这是实际变分态射，而不仅是两个数值函数碰巧相等。
函数层的跳跃是$F_R$延拓轨道中的非零线性组合，
由定理\ref{thm:evenk3visible}知$F_R$必含该K3因子。
次数至多二的多项式没有这个跳跃，因而任何这样的有限局部差都不能改变下界。
这没有断言局部多项式的常数系数只能是Tate周期，
也没有把不同理论的任意子图森林差化为一个多项式。
若连同复合算符本身的有限混合一起改变，须另算新矩阵元与其它粒子道，不能套用这里的接触项结论。
\end{proof}

本节的物理对象是单位四粒子谱及其二点色散读出。
已有$\phi^4$的$2\to4$截面还带十张树图的平方与干涉权，
需要另外追踪那些有理分母和标架；不能把本节当作该完整加权出口的证明。
同样，本节没有完成一般图分类、任意固定点算术或混合对象的全部coaction。
它具体补上了一个主线箭头：已知纯K3核经实际商对应、升维读出及局部减除后，
仍然进入一个四维物理谱和完整二点函数，而无需先计算这些函数的全部积分。


\section{完整干涉权的有限判据：物理K3保持与新增椭圆边界}\label{sec:weightedk3}
上一节的单位四粒子谱还没有包含物理树振幅的平方。
本节处理一个明确的无限输入类：振幅只依赖总不变质量$s$及两对不变质量$a,b$，
传播子极点严格位于对应两粒子阈值以下。
这涵盖任意阶的成对局部导数算符及相应的非共振交换图，
并给出从完整振幅到可见K3因子的有限充要判据。
不依赖这两个不变量的角权、其它配对通道和圈振幅不被自动纳入。

\subsection{完整权、接受条件与几何上界}
固定质量$m=(1,2,4,8)$；一般正质量的局部阈值公式会另行注明。
继续使用$A(a),B(b),Q_s=A B\lambda$及$\eta_{k,s}$。
记$M_A=m_1+m_2,M_B=m_3+m_4,M=M_A+M_B$，所以
$a_+=M_A^2,b_+=M_B^2,\tau=M^2$。
所有系数取一个给定实嵌入下的实代数数，物理耦合也在这一实分支上取值。
输入的完整树振幅写成
\begin{equation}\label{eq:weightedamplitude}
 F(s,a,b)=\sum_\nu F_\nu(s,a,b)
       =\frac{P(s,a,b)}{q(a)r(b)},\qquad P\in K[s,a,b],
\end{equation}
其中已经先通分、合并干涉振幅并约去公因子，而不是先把各图平方。
允许$q=r=1$；一般地
\begin{align}
 q(a)&=\prod_{i=1}^{u}(a-c_i)^{e_i},&
 r(b)&=\prod_{j=1}^{v}(b-d_j)^{f_j},\label{eq:weightedpoles}\\
 0<c_i<a_+,\quad &c_i\ne a_-,&
 0<d_j<b_+,\quad &d_j\ne b_-,
\end{align}
各列极点互异，$e_i,f_j$为正整数。
实物理区域内这些分母一致非零。定义有标号末态的完整正读出
\begin{equation}\label{eq:weightedrate}
 R_{F,D}(s)=\int d\Phi_4^{(D)}\,|F(s,a,b)|^2,\qquad D=2k.
\end{equation}
实系数下$|F|^2=F^2$，其代数延拓使用同一个有理函数$F^2$。
末态角变量已经严格积分，因为$F$与这些角变量无关；这不是角平均的近似。

\begin{proposition}[极点位置决定额外几何，极点阶数不复制纯因子]\label{prop:weightedupper}
在适当光滑开基上，$R_{F,2k}$由实际微分$F^2\eta_{k,s}$及物理链$\Gamma_s$承载，
归一化为
\[
 R_{F,2k}=\frac{c_{2k}^3}{8\pi^2}\int_{\Gamma_s}F^2\eta_{k,s}.
\]
除既有坐标边界外，新增极点除子的非有理归一化分量为
\begin{align}
 E_{c_i,s}:&\quad v^2=A(c_i)B(b)\lambda(s,c_i,b),\label{eq:weightedvertical}\\
 E'_{d_j,s}:&\quad w^2=B(d_j)A(a)\lambda(s,a,d_j).\label{eq:weightedhorizontal}
\end{align}
它们在一般参数上是属一曲线。开曲面载体的半单组成在Artin--Tate背景以外的上界由
\begin{equation}\label{eq:weightedupper}
 T_X\ \oplus\ \bigoplus_i H^1(E_{c_i,s})(-1)\
       \oplus\ \bigoplus_j H^1(E'_{d_j,s})(-1)
\end{equation}
给出，总秩不超过$6+2(u+v)$；$T_X$在整个载体中恰出现一次。
这个上界不宣称每一个椭圆候选都进入指定函数，也不是完整微分阶数的上界。
所有新增边界在正常阈值$\tau$附近同型，远离该处的局部消失区。
\end{proposition}
\begin{proof}
相空间分解与上一节相同，只需将有理标量$F^2$乘入微分；
链没有穿过极点。这是同一个代数曲面上的实际形式，未引入数值等式替代。
在$a=c_i$上，$A(c_i)\ne0$为常数，
归一化后得到\eqref{eq:weightedvertical}。
它在$b$轴上的四个分支点为
$b_-,b_+,(\sqrt{s}-\sqrt{c_i})^2,(\sqrt{s}+\sqrt{c_i})^2$。
一般参数下互异，Riemann--Hurwitz给亏格一。水平极点同理。
公式保留$A(c_i),B(d_j)$，不在基域上漏掉常数二次扭曲；
在复数上去掉这些常数不改变此处的秩界，但并非基域上同构的证明。
正常交叉消解只再加入有理例外曲线和有限交点。

开曲面权谱序列的紧曲面项包含一份$T_X$；
新增曲线的$H^1(-1)$出现在权三，其余边界项为Artin--Tate型。
这一谱序列及其退化为经典输入\cite{DeligneWeight}。
具体几何上界则来自本节逐个识别的极点除子。
Gysin像只能由除子代数类生成，不能杀掉紧曲面的$T_X$。
各$H^1$秩二，故其简单因子不可能复制绝对不可约的秩六$T_X$。
极点重数仅改变同一开曲面中的微分，不增加除子分量，
所以$e_i,f_j$的增大不复制上述纯几何列表。

竖直曲线的分支碰撞只能发生在
$s=0$或$s=(\sqrt{c_i}\pm m_3\pm m_4)^2$，
后者均不大于$(\sqrt{c_i}+M_B)^2<\tau$。
水平曲线的相应值也严格小于$\tau$。
竖横极点相交处，额外分支条件是
$\lambda(s,c_i,d_j)=0$，其两值也小于$\tau$。
它们与零坐标边界的碰撞值包含$s=c_i,d_j$；
与无穷远边界的局部式为非零单位的解析变化。
结合上一节已审阅的原边界，所有可能碰撞都避开$\tau$。
有限个曲线及交点因而可在这个小圆盘上同时正常交叉化。
又因$c_i<a_+,d_j<b_+$，它们与$(a_+,b_+)$的局部消失区分离。
\end{proof}

\subsection{不求完整积分的有限首项判据}
本小节允许一般正质量，并仍要求$F$在$(\tau,a_+,b_+)$没有极点。
先设$P\ne0$，把分子在阈值点展开为
\begin{equation}\label{eq:weightedjet}
 P(\tau+\varepsilon,a_++u,b_++v)
       =P_d(\varepsilon,u,v)+P_{d+1}+\cdots,\qquad P_d\ne0,
\end{equation}
其中$P_j$齐次次数为$j$；$d$是有限整数且不超过$\deg P$。
记$L_0=q(a_+)r(b_+)\ne0$，定义
\[
 h_d(x,y)=L_0^{-1}P_d\left(1,\frac{M_A}{M}x,\frac{M_B}{M}y\right),
 \quad \alpha=\frac{D-1}{2},\quad
 \mathcal T=\{x,y\ge0,\ x+y\le1\}.
 \]
在$\mathcal T$上使用概率测度
\[
 d\nu_\alpha=
 \frac{\Gamma(3\alpha)}{\Gamma(\alpha)^3}
 x^{\alpha-1}y^{\alpha-1}(1-x-y)^{\alpha-1}\,dx\,dy.
\]

\begin{theorem}[正谱首项由有限代数喷射精确决定]\label{thm:weightedjet}
对任意$D=2k\ge2$，有
\begin{equation}\label{eq:weightedthreshold}
 R_{F,D}(\tau+\varepsilon)
 =C_D\,J_\alpha(h_d)\,
 \varepsilon^{(3D-5)/2+2d}\bigl(1+O(\varepsilon)\bigr),
 \qquad J_\alpha(h_d)=\int_{\mathcal T}h_d^2\,d\nu_\alpha>0,
\end{equation}
其中
\begin{equation}\label{eq:weightedCD}
 C_D=
 \frac{2^{1-3D}\pi^{(5-3D)/2}}{\Gamma((3D-3)/2)}
 \frac{(\prod_i m_i)^{(D-3)/2}}{M^{2D-3}}.
\end{equation}
若$h_d=\sum h_{ij}x^iy^j$，则这个严格正系数由有限和给出：
\begin{equation}\label{eq:weightedmoments}
 J_\alpha(h_d)=
 \sum_{i,j,p,q}h_{ij}h_{pq}
 \frac{(\alpha)_{i+p}(\alpha)_{j+q}}{(3\alpha)_{i+j+p+q}}.
\end{equation}
因此无需数值求积或完整周期求值，就能判定完整干涉后的首个非零阈值阶数。
\end{theorem}
\begin{proof}
作精确变量代换
\[
 E=\sqrt{s}-M,\qquad
 \sqrt a=M_A+Ex,\quad \sqrt b=M_B+Ey.
\]
物理区域现在恰为固定三角形$\mathcal T$。
$A(a)=Ex$乘一个严格正解析单位，$B(b)=Ey$乘一个严格正解析单位。
而
\[
 \lambda(s,a,b)=(\sqrt s-\sqrt a-\sqrt b)(\sqrt s+\sqrt a+\sqrt b)
                (\sqrt s-\sqrt a+\sqrt b)(\sqrt s+\sqrt a-\sqrt b)
\]
为$E(1-x-y)$乘一个在闭三角形上一致严格正的解析单位。
包括$da\,db$后，无权密度因此等于
$E^{(3D-5)/2}[xy(1-x-y)]^{\alpha-1}$乘一个正解析单位，
其零阶单位与$x,y$无关。
又因$E=\varepsilon/(2M)+O(\varepsilon^2)$，
\[
 F=\varepsilon^d\bigl(h_d(x,y)+O(\varepsilon)\bigr)
\]
在闭区域上一致成立。端点权可积，故可逐项取首项，并得到相对$O(\varepsilon)$余项。

一般$D$下的非相对论壳有$3(D-1)$个独立实变量。
每个空间分量的能量矩阵行列式仍为$M/\prod_i m_i$。
其椭球壳积分是
\[
 \frac{(2\pi)^{3(D-1)/2}}{\Gamma(3(D-1)/2)}
 \left(\frac{\prod_i m_i}{M}\right)^{(D-1)/2}
 E^{(3D-5)/2}.
\]
乘原测度$(2\pi)^{D-4(D-1)}/(16\prod_i m_i)$，
并把$E$换成$\varepsilon$，得到\eqref{eq:weightedCD}。
归一化角变量之后，剩余三角形概率正是$d\nu_\alpha$。
这同时核对$D=4$的系数为$\sqrt{\prod_i m_i}/(13440\pi^4M^5)$。

因为$P_d$齐次且非零，将第一变量置为$1$后不可能成为零多项式；
后两个非零比例也不改变这一点。因此$h_d$在三角形内部一个开集上非零，
其平方的正权积分严格大于零。
最后对各单项式用二维beta积分，
$\int x^ry^t\,d\nu_\alpha=(\alpha)_r(\alpha)_t/(3\alpha)_{r+t}$，
即得有限式\eqref{eq:weightedmoments}。
正性来自真实平方和积分，并不要求展开后的每个交叉系数为正。
\end{proof}

\subsection{完整树权的K3充要判据}
\begin{theorem}[合并振幅非零当且仅当物理谱保留K3]\label{thm:weightediff}
在\eqref{eq:weightedamplitude}--\eqref{eq:weightedpoles}的接受域、
质量$m=(1,2,4,8)$及任意正偶数维下，以下三件事等价：
\[
 P\not\equiv0;\qquad
 R_{F,2k}\not\equiv0;\qquad
 [M_{\mathrm{fun}}(R_{F,2k}):T_X]=1.
\]
非零情形的最小齐次有理微分阶数至少为六。
在无额外极点的局部多项式振幅类，除该K3因子外仅有Artin--Tate组成。
有极点时则保留命题\ref{prop:weightedupper}的椭圆上界，不据此断言其全部可见。
\end{theorem}
\begin{proof}
$P=0$显然给零函数。$P\ne0$时，
定理\ref{thm:weightedjet}给严格正、有限阶的正常阈值首项。
新增极点除子在该阈值附近同型并避开局部消失区；
因此上一节的实际$D_4$局部变分证明原样适用于这个更小的开曲面：
变分像仍是一维，三个$A_1$例外类仍被固定，其正交方向反号。
这里依赖的是相对边界在邻域外同型，而不是仅检查实积分没有极点。

形式$F^2\eta_k$在消失区正则，故在三个例外曲线上积分为零。
物理链$\Gamma=c\delta+\sum c_jE_j$的函数读出仍是非零常数乘$\delta$的读出，
而\eqref{eq:weightedthreshold}证明后者非零。
从开曲面同调到紧曲面纯部分的投影满射：
新增除子是代数类，其交叉条件不会排除超越部分。
于是引理\ref{lem:evenvisibility}强制物理最小函数系统含有$T_X$。
命题\ref{prop:weightedupper}给它在载体中的重数恰一，因此下界也是精确重数。
余下的断言由同一谱序列上界及子商关系成立。
\end{proof}

这个判据是有限的。给出代数系数、已分解的极点表及指定实嵌入后，
先做有限符号比较检查接受条件，再合并有理振幅、检查$P=0$；
若非零，平移多项式并找到次数不超过$N=\deg P$的第一个齐次项即可。
稠密三变量多项式的直接平移至多用$O(N^6)$次域算术，
二维首项平方及矩公式至多用$O(N^4)$次域算术；
这里未把代数数比较和多项式通分的位复杂度隐藏成常数。
这是上述有限表示输入的规则，不是从任意边表自动得到这种表示的算法。

\begin{proposition}[任意阶成对导数算符的实际物理实现]\label{prop:weightedoperator}
对$q=r=1$及任意实多项式$P(s,a,b)$，
存在局部复合算符，其四个不同自由标量粒子的矩阵元恰为$P(s,a,b)$。
因此任意非零这样的完整算符读出都保留一份六维K3因子；
任意有限实线性组合的消去判据，就是合并后的矩阵元多项式恒零。
\end{proposition}
\begin{proof}
令$O_{12}=:\!\phi_1\phi_2\!:$、$O_{34}=:\!\phi_3\phi_4\!:$，
对单项式$s^ia^jb^\ell$取
\[
 (-\Box)^i\left[(-\Box)^j O_{12}\,(-\Box)^\ell O_{34}\right].
\]
在平面波矩阵元中，每个$-\Box$分别给对应总动量平方，
故这就是$s^ia^jb^\ell$。按$P$的系数取有限和即可。
各物种不同、正规序及所定单位单粒子归一化避免额外置换系数。
其谱密度是完整矩阵元的绝对平方，因此上一条定理直接适用。
若从更一般的导数表达式输入，须按实际在壳数据及适用的算符恒等式取得该矩阵元；
尤其局部插入的整体导数仍给总动量因子，不能以作用量的分部积分约定直接删去。
不能将形式上不同的算符项误作独立物理读出。
本命题未覆盖依赖其它不变量或自由角变量的全部导数算符。
\end{proof}

\subsection{接触项与交换图：首项相消不会移除K3}
下面给真正含树图干涉的部分衰变率。
以$m_1$为质量单位，四种不同末态质量为$1,2,4,8$；
引入质量为$2,5$的标量$A_0,B_0$及一个标量源$\chi$。
相互作用取
\[
 \mathcal L_{\mathrm{int}}
 =-g_0\chi\phi_1\phi_2\phi_3\phi_4
  -\lambda_0\chi A_0B_0-g_{12}A_0\phi_1\phi_2-g_{34}B_0\phi_3\phi_4.
\]
接触项是明确的局部有效相互作用，不冒称五价耦合在四维可重整化。
在固定的最低振幅阶数，令$h=\lambda_0g_{12}g_{34}$，可同时保留
$g_0$和$h$；例如用一个记账参数使两者都为三阶。

\begin{proposition}[包含完整干涉的显式正物理出口]\label{prop:weightedinterference}
上述四粒子道的完整最低阶振幅及部分宽度为
\begin{equation}\label{eq:weightedexample}
 \mathcal M=-\left(g_0+\frac{h}{(a-4)(b-25)}\right),\qquad
 2\sqrt{s}\,\Gamma_{\chi\to1234}=R_{F,4}(s).
\end{equation}
只要$(g_0,h)\ne(0,0)$，这个读出保留一份六维K3因子。
其非Artin--Tate上界另包含至多两份属一边界，分别来自
\[
 v^2=-15(b-16)(b-144)\lambda(s,4,b),\qquad
 w^2=-1071(a-1)(a-9)\lambda(s,a,25).
\]
在$h=-595g_0$、$g_0\ne0$时，阈值处两张图的振幅首项完全相消，
但完整部分率仍满足
\begin{equation}\label{eq:weightedcancellation}
 R_{F,4}(225+\varepsilon)
 =\frac{8}{13440\pi^4\,15^5}\,
   g_0^2\frac{87085}{33\cdot2975^2}\,
   \varepsilon^{11/2}\bigl(1+O(\varepsilon)\bigr)>0.
\end{equation}
因此这项首项相消不消去K3。
\end{proposition}
\begin{proof}
接触顶点给$i\mathcal M=-ig_0$。
交换树的三个$-i$顶点和两个$i$传播子给
$i\mathcal M=-ih/((a-4)(b-25))$；末态物种指定了唯一的两对分组。
相加后再平方得到\eqref{eq:weightedexample}，包括交叉项。
物理区$a\ge9,b\ge144$，所以分母严格正且无共振。
$A_0,B_0$都在相应两体阈值以下，符合\eqref{eq:weightedpoles}。
分子
\[
 P=g_0(a-4)(b-25)+h
\]
只有在$g_0=h=0$时才是零多项式，故定理\ref{thm:weightediff}适用。
宽度的$1/(2\sqrt s)$是标准通量因子；作为$s$上的微分模它给代数秩一扭曲，
在加入$\sqrt s$后与上述K3因子相同，不能把该因子遗漏。

阈值分母为$(9-4)(144-25)=595$。
一般耦合时首项系数是无权系数乘$(g_0+h/595)^2$。
在所列相消条件下，置$u=a-9,v=b-144$，
\[
 P=g_0(119u+5v+uv),\qquad d=1,\qquad
 h_1(x,y)=\frac{g_0}{2975}(119x+20y).
\]
四维$\alpha=3/2$的矩为
$\langle x^2\rangle=\langle y^2\rangle=5/33$、$\langle xy\rangle=1/11$。
故
\[
 J_{3/2}(h_1)=\frac{g_0^2}{2975^2}
       \left[\frac{5(119^2+20^2)}{33}+\frac{2\cdot119\cdot20}{11}\right]
       =g_0^2\frac{87085}{33\cdot2975^2}.
\]
代入\eqref{eq:weightedthreshold}给出结果；全程未求完整相空间积分。
两条属一曲线只是载体边界上界，未在本命题中证明它们都进入完整部分率。
本过程还有$A_0B_0$等其它可开放道，本命题只陈述所指定的四粒子部分宽度。
它也不替代旧等质量$\phi^4$十张$3|3$树分割的角依赖截面。
\end{proof}

\subsection{非零分子仍可消失：一个精确反例}
\begin{proposition}[有符号权不能使用平方振幅的充要判据]\label{prop:weightedsigned}
对同一谱曲面及任意$D=2k\ge2$，非零多项式权
\begin{equation}\label{eq:weightedsigned}
 W_k(s,a,b)=(2k-1)(sab)^{k-1}\,\partial_bQ_s(a,b)
\end{equation}
满足$\int_{\Gamma_s}W_k\eta_{k,s}=0$，从而其有符号相空间读出恒零。
这说明不能把定理\ref{thm:weightediff}的完整平方权替换成任意非零分子。
\end{proposition}
\begin{proof}
在$y^2=Q_s$上直接微分：
\[
 W_k\eta_k=(2k-1)Q_s^{k-1}\partial_bQ_s\,\frac{da\wedge db}{y}
              =-2\,d\!\left(yQ_s^{k-1}da\right).
\]
右边的一形式在物理闭链附近正则，故闭链积分为零。
也可在平面正支上验证：
$W_k$乘相空间密度对$b$为$2\partial_b(Q_s^{k-1/2})$，
而$b=b_+$及$b=(\sqrt s-\sqrt a)^2$两个端点都有$Q_s=0$。
由于$k\ge1$，端点原函数均为零。
$\partial_bQ_s$不是零多项式，故这是非平凡有符号权的实际消去，
而非把振幅本身设为零。
正概率型读出与单独的有符号干涉项在此有实质区别。
\end{proof}

\subsection{加权谱的有限减除与范围}
记$p_a=\deg q,p_b=\deg r$。对非零分子
$P=\sum p_{ij\ell}s^ia^jb^\ell$，定义一个可直接读取的增长上界
\[
 n_\infty=\max_{p_{ij\ell}\ne0}
       \left(i+\max(j-p_a,0)+\max(\ell-p_b,0)\right).
\]

\begin{corollary}[由振幅次数给有限色散入口]\label{cor:weighteddispersion}
四维下$R_{F,4}(x)=O(x^{2+2n_\infty})$。
取$s_0<0$及$J=2n_\infty+3$，则
\begin{equation}\label{eq:weighteddispersion}
 G_{F,R}(s)=(s-s_0)^J\int_\tau^\infty
       \frac{R_{F,4}(x)\,dx}{(x-s)(x-s_0)^J}
\end{equation}
绝对收敛，跳跃为$2\pi iR_{F,4}$。
非零振幅时，该指定谱分量的未割读出保留K3；
固定谱及算符时，有限局部接触项不改变这个下界。
对\eqref{eq:weightedexample}，$n_\infty=0$，三次减除已经足够。
\end{corollary}
\begin{proof}
在$a\ge a_+$上，极点严格低于端点给
$|a^j/q(a)|\le C(1+a)^{\max(j-p_a,0)}$；对$b$同理。
又有$a,b\le x$，所以$|F(x,a,b)|\le C(1+x)^{n_\infty}$。
无权相空间为$O(x^2)$，得到增长上界。
所选$J$使无穷远密度至多为$O(x^{-2})$，阈值处由定理\ref{thm:weightedjet}可积。
于是色散式、实际管链变分和K3下界与上一节一致。
这里的$J$是保证收敛的有限上界，不宣称所有分子都需要这么多次减除。
若是局部复合算符的四粒子谱，这给它的相应二点读出；
交换模型存在其它粒子道时，只重建指定四粒子谱分量，
不将其冒称完整的$\chi$自能。
\end{proof}

由此获得一条完整且有限的接受类规则：
合并物理振幅、检查实际极点、读取有限阈值喷射、确定K3重数和椭圆候选上界，
最后按增长次数接入减除色散。
它处理真实分子与干涉，同时用\eqref{eq:weightedsigned}展示规则的必要边界。
一般角依赖树权、任意图推前、全部椭圆候选的可见性及固定数值算术仍需进一步证明。


\section{交换极点的精确分类：从留数到实际椭圆出口}\label{sec:exchangeelliptic}
上一节保留了一条实质缺口：极点曲线出现在载体中，不保证进入物理率。
本节在一个无限输入类中补齐这条箭头。
固定四维、末态质量$(1,2,4,8)$及有标号测度，考虑完整实振幅
\begin{equation}\label{eq:exchangeclass}
 f(a)=p(a)+\sum_{i=1}^{r}\frac{h_i}{a-c_i},
 \qquad p\in K[a],\quad h_i\in K^\times,\quad
 0<c_i<9,\quad c_i\ne1 .
\end{equation}
$K$带指定实嵌入；$c_i$为互异实代数数。
同一极点处的所有树图先合并，零总留数直接删除。
因此$r$计数实际保留的极点，不是原图数。
先处理简单交换极点、与$s,b$无关的留数，避免把下面的首项论证无条件用于高阶极点。
物理率记为$R_f(s)=\int d\Phi_4^{(4)}f(a)^2$。

\subsection{私有退化值与候选因子的区分}
对$c\in(0,9)\setminus\{1\}$，记$u=\sqrt c$，定义
\[
 E_{c,s}:\quad v^2=A(c)(b-16)(b-144)\lambda(s,c,b),
 \qquad \sigma_c=(u+12)^2.
\]
保留$A(c)=(c-1)(c-9)$这一基域上的二次扭曲。

\begin{proposition}[全部交换曲线具有互异的私有阈值]\label{prop:exchangeprivate}
$H^1(E_c)$在共同光滑底上绝对不可约，几何单值闭包为$\mathrm{SL}_2$。
各$E_c$的正常退化值$\sigma_c\in(144,225)\setminus\{169\}$
不在其它交换曲线、谱K3或原坐标边界的退化集合中。
不同$c_i$给互不同构的秩二局部系统，有限底覆盖也不能消除这一差别。
\end{proposition}
\begin{proof}
四个分支点是$16,144,(\sqrt s-u)^2,(\sqrt s+u)^2$。
它们的有限碰撞值包含在
\[
 0,\quad(12+u)^2,\quad(12-u)^2,\quad(4+u)^2,\quad(4-u)^2 .
\]
因为$0<u<3$，后三值严格小于$144$；唯一大于$144$的值是$\sigma_c$。
函数$c\mapsto\sigma_c$严格递增，故不同极点的这一值不同。
谱K3的节点值为$1,9,25,49,81,121,169,225$，原坐标边界的新增值
包含在$0,1,9,16,144$中。所选开区间中只有$169$可能相交，
恰对应$c=1$，已经排除。

在$\sigma_c$处只有$(\sqrt s-u)^2$与$144$碰撞，其余分支点分离。
交比退化给$j(E_c)$一个极点，因而族非等平凡。
取有限覆盖以选择截面及水平结构，到相应模曲线的非恒定代数映射
在去掉有限分支值后给有限覆盖。
故单值在$\mathrm{SL}_2(\mathbb Z)$中有有限指数，
Zariski闭包为$\mathrm{SL}_2$；这也可直接应用经典结果
\cite{LvovskiEll}，不是本节新证明的普遍椭圆定理。
碰撞处的Picard--Lefschetz变换为非平凡幺幂变换。
这里允许非平凡transvection的整数次幂，不把分支点的标记参数当成通用平滑参数。
其它曲线在该点光滑，单值为单位；有限覆盖只是将前者取正整数次幂，
仍非平凡。因此两个系统不可能同构。
\end{proof}

\subsection{双极点的留数及完整物理链}
令
\[
 C_4=\frac1{2048\pi^5},\qquad
 \alpha(s)=(\sqrt s-12)^2,\qquad
 \beta(s,a)=(\sqrt s-\sqrt a)^2 ,
\]
并在$s>225$的物理支上写
\begin{align}
 H(a,s)&=\frac{\sqrt{A(a)}}{sa}
       \int_{144}^{\beta(s,a)}
       \frac{\sqrt{(b-16)(b-144)\lambda(s,a,b)}}{b}\,db,
       \label{eq:exchangeinner}\\
 R_f(s)&=C_4\int_9^{\alpha(s)} f(a)^2H(a,s)\,da.
       \label{eq:exchangeouter}
\end{align}
这两式来自实际相空间分解，包含完整平方和交叉项。
内积分是椭圆闭循环周期的一半；此处保留半周期以匹配平面积分归一化。
其平方根分支沿下面指定路径续接，不在非物理区强行取正根。

\begin{proposition}[物理链的极点夹挤给非零椭圆跳跃]\label{prop:exchangepinch}
固定一个极点$c=c_i$，令
\[
 h=h_i,\qquad f_i^{\mathrm{reg}}(a)=
 p(a)+\sum_{j\ne i}\frac{h_j}{a-c_j}.
\]
从物理支沿避开其它退化值的路径续接，可取得一个绕$\sigma_c$的局部变分，
使
\begin{equation}\label{eq:exchangedisc}
 \operatorname{Var}_{\sigma_c}R_f
 =2\pi i C_4\,\mathcal R_c(s),\qquad
 \mathcal R_c(s)=h^2\partial_aH(c,s)
              +2h f_i^{\mathrm{reg}}(c)H(c,s).
\end{equation}
圆周及管取向按此式固定；反向绕行改变整体符号。
这是真实链变分，不是仅由形式留数猜测物理跳跃。
设$\delta=s-\sigma_c$，则
\begin{equation}\label{eq:exchangereslead}
 \mathcal R_c(s)=-2h^2\sigma'(c)K(c)\,\delta+O(\delta^2),
 \qquad \sigma'(c)=\frac{\sqrt c+12}{\sqrt c},\qquad K(c)\ne0 .
\end{equation}
尤其接触项和其它交换图不能消去此首项。
\end{proposition}
\begin{proof}
先只分析内积分的局部退化。将$a$限制在$c$的一个小复邻域，选定$\sqrt a$、
$\sqrt{A(a)}$的解析支，令$\sigma(a)=(\sqrt a+12)^2$。
置$z=b-144$、$L=\beta(s,a)-144$，并令
$\beta_+=(\sqrt s+\sqrt a)^2$。
积分的局部被积式为
\[
 \frac{\sqrt{(128+z)z(L-z)(\beta_+-144-z)}}{144+z}\,dz .
\]
在$s=\sigma(a)$处，$\beta_+-144=4\sqrt a(\sqrt a+12)$非零。
用$z=Lx$把链变成$0\le x\le1$，得到$L^2$乘一个解析非零函数。
由于$L=12(s-\sigma(a))/(\sqrt a+12)+O((s-\sigma(a))^2)$，
实际半周期满足
\begin{equation}\label{eq:exchangeHlead}
 H(a,s)=K(a)(s-\sigma(a))^2+O((s-\sigma(a))^3),
 \qquad
 K(a)=\frac{\sqrt{A(a)}}{\sigma(a)a}
       \frac{\pi\sqrt{32\sqrt a}}{2(\sqrt a+12)^{3/2}} .
\end{equation}
系数使用$\int_0^1\sqrt{x(1-x)}\,dx=\pi/8$。
该展开在$(a,s)$的小邻域中解析，$K(c)\ne0$。
常数可能为复数，但其非零性不依赖所选平方根支。

现在把外层实际积分链续接。
从$s>225$出发，可先越过正常阈值的穿孔邻域，再令$\alpha(s)$沿复路径从$9$趋近$c$。
绕开有限个其它极点、分支值及其对应退化值。
内层链的末端始终为$\beta(s,a)$，在$a=\alpha(s)$处收缩到$b=144$。
因此在外层末端的小邻域，内层必为上面从$144$至$\beta$的消失半周期；
它的局部系数由原物理双积分为一，连续运输不会变成零。
若路径作了额外绕行，可出现其它管链项；选择不作这些额外绕行的路径即可。
在$\sigma_c$附近，命题\ref{prop:exchangeprivate}保证没有其它全局退化。
故外层链远离$c$的部分可取局部解析平凡化，全部非单值来自末端撞上$c$。

在这一小双圆盘内$H(a,s)$已经解析，而
\[
 f(a)^2=\frac{h^2}{(a-c)^2}
          +\frac{2h f_i^{\mathrm{reg}}(c)}{a-c}+O(1).
\]
Laurent展开中的简单留数正是$\mathcal R_c(s)$。
原函数包含$\mathcal R_c(s)\log(a-c)$、一个有理极部及单值解析部分。
末端$\alpha(s)-c$在$\sigma_c$处有简单零，逆时针绕行增添$2\pi i$；
固定远端的部分不变，于是得到\eqref{eq:exchangedisc}。
等价地，这一差链是围绕$a=c$的小圆与内层循环的管链。
高阶极点的留数通过局部解析展开取得，没有把二阶极点错误地当作简单Poincaré留数。

最后对\eqref{eq:exchangeHlead}在固定$s$下求$a$导数：
\[
 \partial_aH(c,s)=-2\sigma'(c)K(c)\delta+O(\delta^2),
 \qquad H(c,s)=K(c)\delta^2+O(\delta^3).
\]
平方极部给$\delta$项，全部干涉的简单极部只从$\delta^2$开始。
这就证明了\eqref{eq:exchangereslead}。
\end{proof}

在上式中，$a$导数是椭圆族的实际Gauss--Manin导数。
在平面积分中也可直接求导：末端的被积式为零，导数只有可积的平方根端点。
因此$\mathcal R_c$确是$E_c$上一个代数meromorphic一形式的周期，
需要时删去$b=0,\infty$及其极点，而不是未知的额外特殊函数。

\subsection{精确半单组成、联合群与有限规则}
\begin{theorem}[保留的交换极点数决定精确椭圆重数]\label{thm:exchangeclassification}
设\eqref{eq:exchangeclass}的完整振幅非零。按上同调载体的命名，
$R_f$的最小函数系统在Artin--Tate背景以外的半单组成为
\begin{equation}\label{eq:exchangeclassification}
 T_X\ \oplus\ \bigoplus_{i=1}^{r}H^1(E_{c_i})(-1).
\end{equation}
每个组成恰一份，因而任何消去$R_f$的非零齐次有理微分算子阶数至少$6+2r$。
按同调标架表述时取相应对偶和Tate约定。
本结论没有计算完整混合扩张，也不宣称总阶数恰为$6+2r$。
\end{theorem}
\begin{proof}
上一节的实际开曲面与权谱序列给至多一份$T_X$及每条$E_c$至多一份的上界。
完整正平方的阈值判据已经给出一份$T_X$。
余下要证明每条属一曲线不仅是候选。

取命题\ref{prop:exchangepinch}给出的实际跳跃$\mathcal R_c$。
它是穿孔$E_c$上代数一形式沿正常退化消失循环$\gamma_c$的周期。
在$\sigma_c$附近，穿孔点远离$b=144$处的节点；$b=0,\infty$及其它极点无碰撞。
局部Picard--Lefschetz变分在穿孔曲线同调上仍为秩一，
像为$\mathbb C\gamma_c$；到紧椭圆$H_1(E_c)$的投影满射，
且$\gamma_c$的紧曲线投影非零。
命题\ref{prop:exchangeprivate}给这个秩二商的绝对不可约性；
\eqref{eq:exchangereslead}给指定消失周期读出非零。
因此可应用引理\ref{lem:evenvisibility}，证明$\mathcal R_c$的最小函数系统含完整$H^1(E_c)$。
这里仅凭“留数非零”尚不够，秩一变分和简单商排除了只有穿孔Tate部分可见的可能。

$\mathcal R_c$又是$R_f$的两个实际续接之差，所以其最小系统为$R_f$延拓轨道的子商。
各$c$互不同构，不能把多个私有阈值所迫使的因子算成同一份。
逐个取得下界后恰达到上述上界，得到\eqref{eq:exchangeclassification}及阶数下界。
这些陈述在$\mathbb C(s)$的正则奇异微分模层成立；
不是对任意代数点的数值超越性断言。
\end{proof}

\begin{corollary}[非Artin--Tate半单部分的完整几何群]\label{cor:exchangegroup}
采用上述直接和表示，其几何单值Zariski闭包为
\[
 \mathrm O(q_X)\times\prod_{i=1}^{r}\mathrm{SL}_2 ,
 \qquad
 G^\circ=\mathrm{SO}(q_X)\times\prod_{i=1}^{r}\mathrm{SL}_2 .
\]
这是函数族指定半单部分的群，既不识别整个混合群，也不识别固定纤维的绝对Nori群。
\end{corollary}
\begin{proof}
各投影的闭包已经分别为已证的$\mathrm O(q_X)$与$\mathrm{SL}_2$。
在$\sigma_{c_i}$处，联合表示的单值只在第$i$个椭圆因子中非平凡。
其幺幂幂次的Zariski闭包给该单独因子内的一参数子群。
在联合群中取正规闭包；因为第$i$个投影稠密于$\mathrm{SL}_2$，
这个正规闭包包含整个第$i$个$\mathrm{SL}_2$。
对所有$i$重复，联合群包含这些完整直积因子。
把任意联合元素的椭圆分量用这一已包含的直积消去，
剩下的纯K3投影稠密于$\mathrm O(q_X)$，因此整个群等于所写直积。
这一步使用私有阈值的实际单值，不把若干抽象不同构表示直接当作独立。
\end{proof}

规则具有有限输入输出。将同一$c$的留数先相加，删去总留数零的项，
检查代数不等式和$c\ne1$，即得到精确的$r$及曲线方程。
不需求相空间积分，也不求最小微分算子才能给出半单组成。
在已经提供这类部分分式的输入下，合并只需有限代数比较和加法；
可用$O(n^2)$次比较的直接方法处理$n$个原始项。
若输入是未分解有理式，还须真实执行分解、重根与实嵌入检查，不能把它们视为免费。
未通过这一表示或极点条件的输入输出“规则不适用”，不是输出“无椭圆因子”。

\subsection{物理实现与首项相消的双层解释}
\begin{proposition}[稳定交换粒子的部分率实现及一个精确分支表]\label{prop:exchangephysical}
在更强的物理窗口$1<c_i<9$内，
式\eqref{eq:exchangeclass}可由局部源相互作用及稳定粒子实际实现。
本节的代数周期分类仍允许$0<c_i<9,\ c_i\ne1$；
$c_i<1$时质量为$2$的外腿会衰变，不能称稳定四粒子部分宽度，
详见命题\ref{prop:sourcewindow}。
引入质量$1<\sqrt{c_i}<m_1+m_2$的不同标量$A_i$，取
\[
 -\lambda_i\chi A_i\phi_3\phi_4-g_i A_i\phi_1\phi_2
\]
并加入具有矩阵元$p(a)$的局部$\chi\phi_1\phi_2\phi_3\phi_4$导数接触项。
在所声明的最低振幅阶，$h_i=\lambda_i g_i$且$\mathcal M=-f(a)$，
有$2\sqrt s\,\Gamma_{\chi\to1234}=R_f(s)$。
特别地，对$f(a)=g+h/(a-4)$，精确分类为
\[
 \begin{array}{c|c}
 (g,h) & \text{非Artin--Tate半单组成}\\ \hline
 (0,0)&0\\
 h=0,\ g\ne0&T_X\\
 h\ne0&T_X\oplus H^1(E_4)(-1)
 \end{array}
 \qquad
 E_4:\ v^2=-15(b-16)(b-144)\lambda(s,4,b).
\]
在$h=-5g\ne0$时，真实四粒子阈值$225$的首项相消，
但两种几何因子都保留；其中椭圆因子由续接的私有阈值$196$检测。
\end{proposition}
\begin{proof}
四价与三价顶点及一个传播子给
$(-i\lambda_i)\,i/(a-c_i)\,(-ig_i)=-i\lambda_i g_i/(a-c_i)$，
即$i\mathcal M_i=-i\lambda_i g_i/(a-c_i)$。
与接触项相加给$\mathcal M=-f(a)$。
用共同记账参数令各$\lambda_i,g_i$为一阶、接触矩阵元为二阶，即可固定此最低阶比较。
导数接触项按上一节的成对算符构造实现多项式$p$。
有标号四个物种没有$4!$，初态通量给$1/(2\sqrt s)$。
该代数通量只是秩一扭曲，非零物理部分宽度的几何因子不会因此消失。

对$c=4$，$\sigma_c=(2+12)^2=196$。
只要$h\ne0$，命题\ref{prop:exchangepinch}的首项正比$h^2$且非零，
所以任何$g$都不能移除这份椭圆因子；$h=0$时新增除子不存在。
若$h=-5g$，则$f=g(a-9)/(a-4)$，上一节的阈值喷射为
$h_1(x,y)=gx/25$。于是
\[
 R_f(225+\varepsilon)
 =\frac{8}{13440\pi^4\,15^5}\frac{g^2}{4125}
      \varepsilon^{11/2}\bigl(1+O(\varepsilon)\bigr).
\]
因为$J_{3/2}(gx/25)=g^2(5/33)/625=g^2/4125$，系数严格正。
这是$225$处的K3检测；$196$处的椭圆检测是同一函数的另一个解析支变分。
四粒子部分率在实轴$s<225$仍为零，不能把后者解释为四粒子道提前物理开通。
该理论还具有$A_i\phi_3\phi_4$等其它实过程；
本命题没有把所指定部分宽度与全部包容宽度等同。
\end{proof}

\begin{corollary}[三圈二点图的指定割与局部减除保护]\label{cor:exchangethreeloop}
上述完整四粒子部分谱由实际三圈$\chi$二点图的四粒子切割给出。
在这个指定谱分量的减除色散重建中，
$T_X$及全部$H^1(E_{c_i})(-1)$都是未割最小函数系统的组成下界。
固定该谱，只改变二点局部接触项的方案不能消去这些因子。
\end{corollary}
\begin{proof}
将振幅与其共轭的四条有标号末态线逐一粘合，保留两条外部$\chi$线。
接触--接触项有$(I,V)=(4,2)$，
接触--交换项有$(I,V)=(5,3)$，
交换--交换项有$(I,V)=(6,4)$。
它们都是连通图，均满足$L=I-V+1=3$。
导数顶点只改变分子，不改变这个计数。
切割这四条线并取正能量支，实际树乘积及相空间恰为已定义的完整$R_f$；
相同中间质量先合并的规则也保留全部交叉项。

按上一节的次数界选择足够多的减除，
所得代数色散积分有$\operatorname{Disc}G=2\pi iR_f$。
其管链变分给原函数轨道中的子商，因而定理\ref{thm:exchangeclassification}的
全部因子都进入$G$作为下界。有限局部多项式无此跳跃，不能改变这些下界。
这里证明的是固定四粒子谱的色散分量；真实二点图可能另有三粒子等切割。
本证明不把一个切割的正确树乘积当作全部切割的光学验证，
也不把$G$冒称完整三圈自能或完成的全局Cutkosky比较。
\end{proof}

本节由完整交换留数的合并取得可见椭圆因子的精确数目，
展示固定四粒子质量及K3主核下仍可有无界多的不同椭圆组成。
增加的是传播子质量及粒子内容，且这里允许相应有限图和；
不是一个固定带物理数据的单图无端获得更多因子。
结合上一节的有限减除色散，指定谱分量的未割函数也保留这些因子作为下界；
未割推前可能再带来扩张或其它组成，不在此声称完整分类。
同时含$a,b$交换极点、一般角依赖树权、高阶极点以及随$s$变化的留数，
仍须另行分析其留数周期与实际链，不从本节简单极点类自动外推。


\section{从指定四粒子割到完整三圈源关联函数}\label{sec:fullsourcebridge}
上一节分类了一个真实割函数，却没有把它冒称完整自能。
本节在同一交换类上构造局部源模型，由全部低阶中间态证明这个比较。
结果保留所规定的微扰阶、粒子内容和源归一化；
一般角依赖散射图及任意圈Cutkosky比较仍须另证。

\subsection{稳定粒子的物理接受域}
\begin{proposition}[交换质量窗口的双向稳定性检查]\label{prop:sourcewindow}
上一节的代数周期结论仍以$0<c_i<9,\ c_i\ne1$为域。
若要求$\phi_1,\phi_2,A_i$均为稳定粒子，并保留非零顶点
$-g_i A_i\phi_1\phi_2$，本节须采用更强条件$1<c_i<9$。
在$0<c_i<1$时，$A_i\to\phi_1\phi_2$虽被禁止，
但$\phi_2\to A_i\phi_1$已经允许。
\end{proposition}
\begin{proof}
三条质量不等式为$\sqrt{c_i}<1+2$、$2<1+\sqrt{c_i}$和
$1<2+\sqrt{c_i}$，严格交集为$1<\sqrt{c_i}<3$。
若$0<c_i<1$且$g_i\ne0$，则$\lambda(4,1,c_i)>0$，
衰变宽度最低项$g_i^2\sqrt{\lambda(4,1,c_i)}/(16\pi\,2^3)$严格为正。
这修正物理接受条件，不改变前节在更宽区间内的代数积分论证。

赋予$\phi_1,\phi_2,A_i$依次为$(1,0),(0,1),(1,1)$的$\mathbb Z_2^2$荷。
保留这些荷并固定物理质量时，质量为$1$或$2$的粒子没有同荷的更轻多粒子末态；
$A_i$的同荷两粒子最低质量为$3$，含另一$A_j$的更高粒子道质量又大于$3$。
一粒子混合以在壳对角化处理。本节的$\phi_3,\phi_4$保持自由，
外部源不在圈内传播。这是有限阶微扰粒子的稳定性检查，不断言非微扰束缚态谱。
\end{proof}

\subsection{局部源及全部有限极部}
令$\epsilon$仅为耦合计阶变量，取
\begin{align}
 {\cal L}_{\rm int}&=-\epsilon\sum_i g_i A_i\phi_1\phi_2,
 \qquad m(A_i)^2=c_i, \label{eq:sourcebulk}\\
 {\cal O}_\epsilon&=\epsilon\sum_i\lambda_i A_i+
       \epsilon^2p(-\Box):\!\phi_1\phi_2\!:,\qquad
 {\cal J}_\epsilon={\cal O}_\epsilon\phi_3\phi_4 .
 \label{eq:sourceoperator}
\end{align}
$-\Box$仅作用于$\phi_1\phi_2$的乘积，给成对不变质量$a$。
仍记$h_i=\lambda_i g_i$、$f(a)=p(a)+\sum_i h_i/(a-c_i)$。
通过$-\chi{\cal J}_\epsilon$耦合外部源$\chi$。
这是有限阶局部复合算符问题，不要求导数源耦合具有幂次可重整化性。
需要有下界的经典势时，可加只含内层场、保持上述荷的正四次项，
系数取$\epsilon^2$阶并足够大。由Young不等式可用该四次项和正二次项控制三次项。
它在本节阶数只增加可由在壳条件去掉的一粒子质量修正；
新的多粒子矩阵元平方从更高阶开始。以下仅作有限阶微扰。

固定极点质量和单位留数，要求不同$A_i$无在壳混合，并选源反项使
\begin{equation}\label{eq:sourceonshell}
 \langle0|{\cal O}_\epsilon(0)|A_i\rangle
       =\epsilon\lambda_i+O(\epsilon^5).
\end{equation}
它消去$\epsilon^3$的一粒子振幅修正，不能仅靠质量反项代替。
零总留数先合并删除；无耦合自由源项不参与所研究的$\epsilon^4$系数。

\begin{proposition}[局部反项实现有限极部消去]\label{prop:sourceprincipal}
记$z$为内层动量平方，$z_0<0$。标准物理二点规范的泡函数为
\begin{equation}\label{eq:sourcebubble}
 B(z)=\frac{z-z_0}{2\pi}\int_9^\infty
 \frac{\Phi_2^{(4)}(x;1,2)}{(x-z)(x-z_0)}\,dx,\qquad
 \Phi_2^{(4)}(x;1,2)=\frac{\sqrt{(x-1)(x-9)}}{8\pi x}.
\end{equation}
因此$\operatorname{Disc}B=i\Phi_2^{(4)}$，且$B$在所有$c_i$处解析。
内层完整$\epsilon^4$二点系数可由局部反项选为
\begin{align}
 C_4(z)&=F(z)-Q(z),\label{eq:sourceprincipal}\\
 F(z)&=f(z)^2B(z)-
 \sum_i\left\{\frac{d_i}{(z-c_i)^2}+\frac{e_i}{z-c_i}\right\},\nonumber\\
 d_i&=h_i^2 B(c_i),\qquad
 e_i=h_i^2B'(c_i)+2h_i f_i^{\rm reg}(c_i)B(c_i).\nonumber
\end{align}
$Q$是由参考条件固定的源--源局部多项式。
$C_4$在全部$c_i$处正则，连续谱恰为$f^2\Phi_2^{(4)}$。
\end{proposition}
\begin{proof}
两次源之间的一个$\phi_1\phi_2$泡给
$p^2B+2p\sum_i h_iB/(z-c_i)+
\sum_{ij}h_i h_jB/((z-c_i)(z-c_j))=f^2B$。
这保留不同交换粒子的全部交叉项。有标号泡的态完备性没有$2!$。
在共同UV正规化下定义
\[
 L_{ij}B(z)=
 \frac{(z-c_j)B(c_i)-(z-c_i)B(c_j)}{c_i-c_j}\ (i\ne j),
 \qquad L_{ii}B(z)=B(c_i)+(z-c_i)B'(c_i).
\]
实对称的矩阵反项$g_i g_jL_{ij}B$是局部质量/动能多项式。
$B-L_{ij}B$在$c_i,c_j$同时为零，$i=j$时为二阶零，
实现固定极点、单位留数和无在壳混合。
源混合反项再把$p(z)B(z)/(z-c_i)$换为
$[p(z)B(z)-p(c_i)B(c_i)]/(z-c_i)$，实现\eqref{eq:sourceonshell}。
这是局部$A_i$源插入，不是非局部传播子反项。
所得非局部表达式为
\[
 p^2B+
 2\sum_i h_i\frac{p(z)B(z)-p(c_i)B(c_i)}{z-c_i}
 +\sum_{ij}h_i h_j
       \frac{B(z)-L_{ij}B(z)}{(z-c_i)(z-c_j)} .
\]
部分分式分解恰为$F$。$d_i$来自质量反项，
$e_i$合并了留数、场混合和源归一化；并非人为删除物理道。
泡的UV常数在差中消去，余下多项式由$Q$去除。
所有$c_i$均低于两粒子阈值，故反项实，有限方案差存在。
标准在壳子圈的双零消去已有文献\cite[式(15)]{Aleksejevs}；
本节另行保留源混合、多质量交叉项及物理读出。
\end{proof}

令$d=\max(0,\deg p)$，$p=0$时取$d=0$。
取$N=2d+1$，将$Q$选为$F$在$z_0$的$N-1$阶Taylor多项式，则
\begin{equation}\label{eq:sourceinnerdisp}
 C_4(z)=\frac{(z-z_0)^N}{2\pi}
 \int_9^\infty\frac{f(x)^2\Phi_2^{(4)}(x;1,2)}
 {(x-z)(x-z_0)^N}\,dx .
\end{equation}
这是减除Cauchy公式：$F$已无有限极点，唯一物理割为$[9,\infty)$，
在无穷远至多为$z^{2d}\log z$，减除核使大圆消失并给绝对收敛。
$N$是充分值，不称最小。去掉极部之后才能写成仅有连续谱的形式。

\subsection{完整中间态求和}
\begin{theorem}[该阶全部非局部谱恰为四粒子谱]\label{thm:sourcefullspectrum}
在上述模型、固定极点质量及\eqref{eq:sourceonshell}下，
$ {\cal J}_\epsilon$完整Wightman两点函数的$\epsilon^4$阶谱密度为
\begin{equation}\label{eq:sourcefullrho}
 w_{{\cal J},4}(s)=\int d\Phi_4^{(4)}\,f(a)^2=R_f(s).
\end{equation}
这是该系数的全部中间态求和。
完整关联函数在$\epsilon^2$阶仍有$A_i\phi_3\phi_4$三粒子谱。
\end{theorem}
\begin{proof}
取标准相对论态规范，写
$\widehat W_{\cal O}(P)=\theta(P^0)w_{\cal O}(P^2)$。
单位一粒子矩阵元给$2\pi\delta(a-c)$，
两个不同粒子给$\Phi_2^{(4)}(a)$。最低阶矩阵元为
\[
 \langle0|{\cal O}_\epsilon|A_i\rangle=\epsilon\lambda_i+O(\epsilon^5),
 \qquad
 \langle0|{\cal O}_\epsilon|\phi_1\phi_2\rangle
       =\epsilon^2 f(a)+O(\epsilon^4).
\]
第二式的符号也由$(c_i-a)A_i=-\epsilon g_i\phi_1\phi_2$检查。
固定极点位置和第一式使一粒子谱在$\epsilon^4$阶
既没有权重修正，也没有质量位移的$\delta'$。
其它同荷多粒子矩阵元至少为$\epsilon^3$阶，其平方从$\epsilon^6$开始。
因此
\[
 w_{\cal O}(a)=2\pi\epsilon^2\sum_i\lambda_i^2\delta(a-c_i)
 +\epsilon^4 f(a)^2\Phi_2^{(4)}(a;1,2)+O(\epsilon^6).
\]
这是一项有限阶状态计数，不假设全阶谱已经求出。

旁观场没有体相互作用，分离点严格因子化为
$W_{\cal J}(x)=W_{\cal O}(x)\Delta_+(x;16)\Delta_+(x;64)$。
代入上式，用三粒子态完备性得到
\[
 w_{{\cal J},4}(s)=\int_9^{(\sqrt s-12)^2}\frac{da}{2\pi}\,
 f(a)^2\Phi_2^{(4)}(a;1,2)\Phi_3^{(4)}(s;\sqrt a,4,8).
\]
实际递推$d\Phi_4=(da/(2\pi))\,d\Phi_3\,d\Phi_2$
证明\eqref{eq:sourcefullrho}，并恢复前节双积分常数$1/(2048\pi^5)$。
$\epsilon^2$阶则为$\sum_i\lambda_i^2\Phi_3^{(4)}(s;\sqrt{c_i},4,8)$；
三粒子道仅在所考察的高阶系数中被在壳条件去掉。
配置空间传播子乘积与多体谱的联系属于经典方法\cite{GKPsource}。
此处新增的接口是指定源、计阶下的全部态分类和在壳抵消，
没有用任意图的Cutkosky口号替代这些证明义务。
\end{proof}

\subsection{完整时间序函数与三圈图}
明确固定
\[
 \Pi_{\cal J}(s)=i\int d^4x\,e^{iP\cdot x}
 \langle0|T{\cal J}_\epsilon(x){\cal J}_\epsilon(0)|0\rangle,
 \qquad {\cal G}_{\cal J}=2\pi\Pi_{\cal J},\quad P^2=s.
\]
于是$\operatorname{Disc}\Pi_{\cal J}=iw_{\cal J}$，
而$\operatorname{Disc}{\cal G}_{\cal J}=2\pi iw_{\cal J}$。
额外$2\pi$使其与前节色散规范一致，不能默默省去。

\begin{theorem}[完整三圈系数的实际色散比较]\label{thm:sourcefullcorrelator}
令$ {\cal G}_4=[\epsilon^4]{\cal G}_{\cal J}$，
采用固定的Lorentz不变局部时间序延拓。
取$s_0<0$、$J=2d+3$，则
\begin{equation}\label{eq:sourcefullG}
 {\cal G}_4(s)=P_{J-1}(s)+(s-s_0)^J
 \int_{225}^{\infty}\frac{R_f(x)}{(x-s)(x-s_0)^J}\,dx ,
\end{equation}
$P_{J-1}$为至多$J-1$次局部源--源多项式。
在$s_0$令前$J$个Taylor系数为零时，$P_{J-1}=0$。
此非局部系数由前节三类实际三圈二点图及其局部反项组成。
\end{theorem}
\begin{proof}
实际Wightman函数在非重合点决定时间序函数。
前定理给了$\epsilon^4$阶全部谱，故非局部部分由其Källén--Lehmann核决定。
具体地，先将谱截于有限上限，逐个自由质量的时间序Fourier核为$(x-s)^{-1}$。
减去$s_0$处$J$个Taylor项后，该核变为
$(s-s_0)^J/((x-s)(x-s_0)^J)$。
因$R_f(x)=O(x^{2+2d})$，减除尾部为$O(x^{-2})$，可去掉谱截断；
阈值处正半整数幂同样可积。

两个保持所声明缩放次数的局部延拓之差支撑于$x=0$，
因此是有限个$\delta(x)$导数。Lorentz不变性使其Fourier变换为
至多$2d+2=J-1$次的$s$多项式。
全部余项恰为$P_{J-1}$，不是任意无割函数。
这是从实际源得到完整谱再作比较，不是先任选谱而赋予一个整图名称。

分离点的Wick展开与Feynman图一致：每个内层泡外接两条旁观传播子。
接触--接触、接触--交换、交换--交换分别有
$(I,V)=(4,2),(5,3),(6,4)$，均为三圈。
命题\ref{prop:sourceprincipal}的质量、动能和源反项是相应局部子图替换；
最后的源--源延拓处理整体局部自由度。
普通森林实现须保持这些完整相容条件，不能把分别发散的森林项当成收敛周期。
允许的稳定四次项在此阶只产生上述已消去的质量插入。

本证明给已标记微扰函数、实际图和色散读出的相等。
欧氏值可进入前文收敛载体构造，但没有额外断言
两个独立Nori大载体同构或其全部混合扩张已经识别。
\end{proof}

\subsection{局部子反项的椭圆背景与正常割保护}
令$S_c$为单位源$A\phi_3\phi_4$的减除两点函数，
仍用$ {\cal G}=2\pi\Pi$，故
$\operatorname{Disc}S_c=2\pi i\Phi_3^{(4)}(s;\sqrt c,4,8)$。
这是实际两圈sunset函数，其属一曲线是$E_c$；
常数二次扭曲不改变几何单值。

\begin{proposition}[显式收缩背景与阈值投影]\label{prop:sourceforestbackground}
在本节阶数，改变有限质量/动能反项和高一阶源混合，
但保持最低阶$f$、粒子列表及自由旁观条件，与在壳规范的差具有形式
\begin{equation}\label{eq:sourcebackground}
 \widetilde{\cal G}_4-{\cal G}_4
 =P(s)+\sum_i\bigl(a_i S_{c_i}(s)+b_i\partial_{c_i}S_{c_i}(s)\bigr).
\end{equation}
系数与$s$无关，允许为有限周期；所有这些项在$s=225$的局部单值为恒等。
同一物理质量及源矩阵元已匹配时，$a_i,b_i$被匹配条件去掉，余下只有$P$。
\end{proposition}
\begin{proof}
内层局部项给多项式；质量/动能插入给一次式乘
$((z-c_i)(z-c_j))^{-1}$，源混合给多项式乘$(z-c_i)^{-1}$。
部分分式后，非多项式部分仅为$(z-c_i)^{-1},(z-c_i)^{-2}$。
外接旁观线，分别给$S_{c_i}$及其质量导数，符号吸收到系数。
内层多项式支撑在重合点，外接后只改变整体局部延拓，贡献$P$。
所以一个内层有限局部质量反项会产生非局部的椭圆函数$\partial_cS_c$：
“子反项局部”不能替换为“整个大图只差多项式”。
系数算术须与收缩图共同保留；仅在函数局部系统层，系数是常数。

$S_c$的有限奇异值位于
$0,(12\pm\sqrt c)^2,(4\pm\sqrt c)^2$，在$1<c<9$时全部小于$225$。
即使实轴处于它的物理割上，每个提升到解析曲面的局部支在$225$仍正则。
参数导数不会在共同正则邻域制造新奇点，因此绕$225$杀掉全部差项。
如果比较的是同一物理理论，须同时匹配质量、场和源；
未匹配的极点位置/权重差不属于同一物理两点函数的自由度。
\end{proof}

\begin{corollary}[完整系数中全部可见核的保护]\label{cor:sourcevisible}
若$f\ne0$，则$ {\cal G}_4$和\eqref{eq:sourcebackground}中的
$\widetilde{\cal G}_4$，在$\mathbb C(s)$上的最小函数系统都至少含
\[
              T_X\oplus\bigoplus_{i=1}^rH^1(E_{c_i})(-1).
\]
非零齐次有理消去算子阶数至少$6+2r$。
这是完整指定三圈系数的下界，不声称未割系统仅有这些组成。
\end{corollary}
\begin{proof}
\eqref{eq:sourcefullG}给$\operatorname{Disc}{\cal G}_4=2\pi iR_f$，
是同一实际时间序函数两条续接之差。故$R_f$的延拓轨道给
$ {\cal G}_4$最小函数系统的子商，前节精确分类给下界。

方案保护也可直接由正常阈值$\tau=225$读出。
完整正谱满足$R_f(\tau+t)=t^{7/2+2d_0}u(t)$，
$u$解析且$u(0)\ne0$，$d_0$为实际首个非零阈值喷射阶。
把色散积分分为阈值邻域及解析的远端；
前者的减除Cauchy核绕端点续接，由Plemelj关系得到
$(T_\tau-1){\cal G}_4=\kappa R_f$，
$\kappa=\pm2\pi i$由起始支及闭路取向固定。
半整数幂保证该常数非零。
前命题的全部背景在$\tau$局部单值平凡，
所以同一变分对$\widetilde{\cal G}_4$仍给$\kappa R_f$。
继续这个差函数便得到前节各私有椭圆变分。
因此即使方案背景含同种椭圆，它也不能取消正常割所生成的下界。
\end{proof}

本节闭合一条有明确源、稳定窗口及有限阶归一化的接口：
树图合并、完整中间态谱、实际三圈关联函数、森林背景下的可见组成。
自由旁观因子化、一粒子源条件、局部反项的真实插入及正常阈值分离，
各自承担不可省略的证明义务。
最小混合扩张、一般角依赖模型、所有图的物理链比较、
固定值算术和RG出口仍保留在总目标中。


\section{完整色散周期的固定载体与精确运动学组成}\label{sec:dispersioncarrier}
上一节已把完整三圈源系数与指定正谱比较。
现在补上上界：色散变量是固定代数空间的一个坐标，外运动学仅移动其极点除子。
因此完整函数的新运动学组成受同一个谱纤维控制。
这应用第\ref{sec:pencil}节的固定环境定理，但须实际提交三维空间、
积分链、无穷远和留数，不能只给一个形式Cauchy核。

\subsection{三维空间、真实链与减除无穷远}
采用第\ref{sec:fullsourcebridge}节的物理接受域和相同源归一化。
取带指定实嵌入的数域$k$，包含所用代数质量、耦合及一个代数参考点$s_0<0$，令
\[
 A(a)=(a-1)(a-9),\quad B(b)=(b-16)(b-144),\quad
 Q_x(a,b)=A(a)B(b)\lambda(x,a,b).
\]
取三维双覆盖的一个固定光滑模型
\[
 {\cal Y}:\ y^2=Q_x(a,b),\qquad \pi:{\cal Y}\longrightarrow\mathbb P^1_x .
\]
这里先取射影闭包和正规化，再消解固定边界及$\pi$的不定性。
后续的$x=s$是该固定模型的纤维，不是重新选择一个依赖$s$的环境。
记
\[
 \eta_x^f=f(a)^2\frac{y\,da\wedge db}{xab},\qquad
 \Omega_s=\frac{(s-s_0)^J\,dx\wedge\eta_x^f}
                   {(x-s)(x-s_0)^J},\quad J=2d+3 .
\]
在原仿射实片上，固定三链为
\[
 \Xi=\{x\ge225,\ a\ge9,\ b\ge144,\ \sqrt a+\sqrt b\le\sqrt x,\
              y=+\sqrt{Q_x(a,b)}\},
\]
取$dx\wedge da\wedge db$的方向。所有根式不等式都是半代数条件；
链与$s$无关，初始$s<0$避开其支撑。令$C_4=1/(2048\pi^5)$。

\begin{proposition}[固定三维相对对承载完整减除系数]\label{prop:dispersionpair}
可选固定光滑代数对$(U,B_U)$、有限集$\Sigma\subset\mathbb P^1(k^{\rm alg})$，
使$s\in S=\mathbb P^1\setminus\Sigma$时的移动除子$D_s=\pi^{-1}(s)\cap U$
光滑并与$B_U$各层横截，且
\[
 M_s=H^3(U\setminus D_s,B_U\setminus D_s),\qquad
 {\cal G}_4(s)-P(s)=C_4\langle\Xi_s,\Omega_s\rangle .
 \tag{DP}\label{eq:dispersionpair}
\]
右侧是实际相对周期。可在有限公共消解上实现；$C_4$为固定代数Tate因子。
$D_s$与谱K3纤维双有理，其相对边界仅添加有理曲线和有限点，
新增的非有理极点曲线仍仅为$a=c_i$上的$E_{c_i,s}$。
\end{proposition}
\begin{proof}
在$s<0$处，$\Omega_s$的积分由上一节的收敛界和相空间分解绝对收敛。
直接代入$R_f(x)=C_4\int_{\Xi_x}\eta_x^f$，Tonelli对正核给(DP)，
一般初始复$s$再由一致收敛和解析性得到。这里没有先把相同数值解释为同构。

固定的有限极点位于$a=0,b=0,x=0,x=s_0,a=c_i$；
它们都不穿过物理链。有限物理边界位于$y=0$的分支面；
$x=225$处只是这些面的共同收缩端点。射影无穷远也并入固定边界。
为检查减除后的无穷远，改写
\[
 \Omega_s=
 \frac{(s-s_0)^J(x-s_0)}{x-s}\,
 \frac{dx\wedge\eta_x^f}{(x-s_0)^{J+1}} .
\]
第一因子在$x=\infty$处为正则单位乘$(s-s_0)^J$。
第二个固定顶次形式在物理链上绝对可积：
先对截面积分，尾部由$x^{2+2d}/x^{J+1}=O(x^{-2})$控制。
因此选择消解时不需要把$s$引入无穷远中心。

对固定分子、分母及链的不等式边界作共同嵌入消解和实扇区细分。
局部顶次形式成为单项式乘解析单位。
凡链有全维开扇区趋近的边界面，绝对可积性迫使该面的整数极阶指数大于$-1$，
即至少为零。若若干面只在角点接触，继续对该固定角点主理想化后作同一检查。
删除不接触闭链的真实极点面，余下闭链的边界放入$B_U$。
这是消解收敛代数积分为良好相对周期的标准步骤；
必要的有限细分保持有理链系数，并以公共模型推送，不能把坐标覆盖度当成额外物理倍数。

对这个固定模型再去掉$\pi$及边界各层的有限临界值，
便得到$S$上的光滑横截除子。多余的固定竖直纤维不与一般$D_s$相交。
水平消解中心在曲面纤维上至多给点吹起，新增例外曲线有理。
原分支边界$A=0,B=0,\lambda=0$的归一化也有理：
前两类是直线；后者是光滑平面二次曲线的一部分。
坐标无穷远与零边界已由命题\ref{prop:weightedupper}识别。
唯一新增非有理除子仍是各$a=c_i$，归一化为$E_{c_i,s}$。
相对加入上述有理边界时，其$H^1$仅由有限交点图和Tate方向贡献，
不会新添无限单值的简单因子。

最后，取代数参数的纤维，相对上同调给实际Nori对象；
全族至少给实际混合Hodge结构变分和正则奇连接。
这一构造不假设一般Nori变分类别的存在。
规范常数可以外提：$C_4=i/(64(2\pi i)^5)$，必要时在$k(i)$上写出，
再保留下降。它不改变下文的复几何单值分类。

本节选择该色散对作为完整函数的一个实际几何提升。
若另行采用原动量或森林图载体，上一节的解析函数比较尚不自动成为
两种形式周期的相等；这里的形式留数和管关系只在已明确指定的色散提升中断言。
\end{proof}

\subsection{环境常数与运动学结构必须分开}
令$\mathcal F$为复正则奇局部系统中，由有限单值的简单对象生成的Serre子范畴。
其对象允许幺幂扩张；不要求整个对象本身具有有限单值。
“$\mathcal F$外的半单组成”只指无限单值的简单Jordan--Hölder因子。
固定K3的常值局部系统属于这个背景，但其固定Hodge或动机对象通常非Tate。
因此本节不把这个几何商称为算术mixed Tate商。

\begin{theorem}[完整载体的非平凡运动学来自实际谱纤维]\label{thm:dispersiongysin}
上述模型有实际留数正合段
\[
 H^3(U,B_U)\longrightarrow M_s
 \xrightarrow{\mathrm{res}} C_s
 \longrightarrow H^4(U,B_U),\qquad
 C_s=H^2(D_s,B_U\cap D_s)(-1).
 \tag{DG}\label{eq:dispersiongysin}
\]
在复几何局部系统的Serre商中，
\[
 \pi_{\mathcal F}M\simeq\pi_{\mathcal F}C.
 \tag{DQ}\label{eq:dispersionquotient}
\]
特别地，完整周期最小函数系统在$\mathcal F$外的上界为一份$T_X$
及每个不同极点至多一份$H^1(E_{c_i})$。
此句忽略Tate扭转；上同调载体中的对应项为
$T_X(-1)$和$H^1(E_{c_i})(-2)$。
\end{theorem}
\begin{proof}
固定对已满足定理\ref{thm:pencil-separation}的全部几何条件。
在$B_U$各交叉层上对$D_s$使用余维一纯性，并取相对映射纤维，得到(DG)。
令$K=\ker\mathrm{res}$、$A=\operatorname{im}\mathrm{res}$。
$K$是固定$H^3(U,B_U)$的商，而$C/A$是固定$H^4(U,B_U)$的子对象。
它们在复几何局部系统层均为常值，故均在$\mathcal F$中。
对$0\to K\to M\to A\to0$及$0\to A\to C\to C/A\to0$
取精确Serre商，得到(DQ)。

命题\ref{prop:weightedupper}及前命题识别的相对有理边界，给$C_s$所述上界。
固定环境即使含很复杂的算术周期，也不能在$s$方向生成新的无限单值简单因子。
物理最小函数系统为$M^\vee$的子商，因而继承上界；
几何极化识别所列纯局部系统及其对偶时，Tate扭转不影响这一复分类。
Poincaré留数和权谱序列是经典工具\cite[\S3.1.5、\S3.2]{DeligneWeight}；
这里实际需要补交的是命题\ref{prop:dispersionpair}的完整积分模型。
\end{proof}

\subsection{管链关系确实是形式周期关系}
\begin{proposition}[谱到完整色散对象的实际留数与管链]\label{prop:dispersiontube}
在以上相对对中，代数留数严格满足
\[
 \mathrm{res}_{x=s}\Omega_s=\eta_s^f .
 \tag{DR}\label{eq:dispersionresidue}
\]
对相容的相对Betti二循环$\delta$，取正向法圆周管链$\operatorname{Tub}\delta$，有
\[
 \int_{\operatorname{Tub}\delta}\Omega_s
       =2\pi i\int_\delta\eta_s^f .
 \tag{DT}\label{eq:dispersiontube}
\]
在代数参数纤维，(DT)同时是实际Gysin态射诱导的形式周期等式；
两侧的$2\pi i$须替换为Tate形式周期$\mathbb L^{\mathfrak m}$。
上、下物理边值之差由这个管链构造给出。
\end{proposition}
\begin{proof}
局部用$t=x-s$，核的系数为
$(s-s_0)^J/(s+t-s_0)^J=1+O(t)$，故(DR)成立。
在$\delta$各单形上取小法圆，沿边界的管面落在$B_U$的相应部分；
内部面按相反方向抵消。积分$dt/t$为$2\pi i$，其它项随管半径收缩为零，
给(DT)。这正是(DG)留数与Betti管映射的比较，
所以按态射函子性给形式周期关系，无需周期映射单射。

在实$s>225$且避开其它退化值时，从欧氏点走上下两条路径。
固定三链的$x$轮廓分别从$x=s$两侧绕开；两轮廓之差是一次正或负法圆。
其纤维链恰为物理截面$\Xi_s$，分支边界均相容地落在$B_U$。
因而链差为$\pm\operatorname{Tub}\Xi_s$，与上一节的谱边值规范一致。
$x=\infty$不贡献额外轮廓，因减除后尾部绝对收敛。
这给当前模型的实际链比较，不援引任意图的全局切割定理。

这里的Nori态射是留数及其对偶管映射。
没有把一个绕参数奇点的单值变换宣称为指定纤维的Nori自同态；
从一个物理链生成的固定纤维最小对象仍须另做算术标架分析。
\end{proof}

\subsection{从下界升为完整函数的精确半单分类}
\begin{theorem}[匹配物理源的完整三圈系数恰保留这些运动学核]\label{thm:dispersionexact}
设$1<c_i<9$、$h_i\ne0$且互异极点已经合并，$f\ne0$。
令$\mathcal V({\cal G}_4)$为完整源系数在$\mathbb C(s)$上的最小正则奇函数系统。
则其$\mathcal F$外的半单组成为
\[
 \operatorname{ss}_{\notin\mathcal F}\mathcal V({\cal G}_4)
       = T_X\oplus\bigoplus_{i=1}^rH^1(E_{c_i}) ,
 \tag{DC}\label{eq:dispersionexact}
\]
每项恰出现一次。该半单部分的几何单值闭包为
\[
 O(q)\times\prod_{i=1}^r\mathrm{SL}_2,\qquad
 G^\circ=\mathrm{SO}(q)\times\prod_{i=1}^r\mathrm{SL}_2 .
\]
固定同一物理质量、场和源矩阵元后，改变允许的整体局部多项式不改变(DC)。
任意齐次有理消去算子阶数至少为$6+2r$，但完整最小阶数不由(DC)单独决定。
\end{theorem}
\begin{proof}
命题\ref{prop:dispersionpair}将完整物理函数置于实际$M_s$中；
定理\ref{thm:dispersiongysin}给每个所列因子重数至多一，且没有其它无限单值简单因子。
另一方面，命题\ref{prop:dispersiontube}把$R_f$的全部续接轨道置入
完整函数的轨道所生成的系统中。定理\ref{thm:exchangeclassification}已证明
$R_f$恰激发一份$T_X$及所有$r$个椭圆核，故各重数至少一。
有限长度性给(DC)。

所得半单表示与原谱的非平凡半单表示同构，
故群恰为推论\ref{cor:exchangegroup}已由私有阈值和K3输入证明的乘积。
这里保留完整底$S$上的行列式角色；若取有限底覆盖，恒等分支不变，
而$O(q)$的分支数可能改变，不作同一断言。

添加$s$多项式至多添加常值局部系统的子商；
定理\ref{thm:addition}在当前函数类别的证明同样适用，
所以背景外重数保持。完整阶数下界由这些线性独立简单因子维数相加得到。
有限单值成分、常值周期及它们之间的不分裂扩张仍可提高完整阶数。
\end{proof}

\subsection{还没有被分类掉的量：固定周期与非齐次源}
\begin{proposition}[留数湮灭后只剩固定环境周期]\label{prop:dispersionconstants}
令$L\in k(s)[\partial_s]$湮灭实际相对类$[\eta_s^f]$，
而不只是其一个数值积分。则存在固定几何周期$P_1,\ldots,P_b$
及$r_j(s)\in k(s)$，使标准零Taylor规范下
\[
 L\bigl({\cal G}_4/C_4\bigr)=\sum_{j=1}^b r_j(s)P_j .
 \tag{DS}\label{eq:dispersionsource}
\]
所有$P_j$来自同一固定$H^3(U,B_U)$。
这一等式可在相应形式周期中成立。
不加新的算术证明，不能把这些$P_j$全部换成Tate或MZV常数。
\end{proposition}
\begin{proof}
(DR)和留数与Gauss--Manin的相容性给
$\mathrm{res}(L\Omega_s)=L[\eta_s^f]=0$。
由(DG)，$L\Omega_s$提升为固定$H^3(U,B_U)_{\dR}\otimes_k k(s)$中的向量。
取其$k$基$e_j$，写提升为$\sum r_j(s)e_j$。
$\Xi_s$在固定对$(U,B_U)$中的像$\tau$不随$s$变化；
配对得到$P_j=\langle\tau,e_j\rangle$和(DS)。
不同提升的差与$\tau$配对为零。每一步为实际态射、连接和线性关系，
故与定理\ref{thm:ambient-source}一样给形式周期等式。

固定参考纤维$x=s_0$和紧化边界已进入环境；
其常数可以含固定K3或椭圆周期，几何单值为恒等并不意味着算术简单。
因此只知道(DC)不足以恢复完整coaction或扩张类。
以$\mathbb C$为常数域时右侧是一个有理函数，
非零时可由一阶算子$\partial_s-R'/R$湮灭，其中$R=\sum r_jP_j$。
若要求系数仍在$k(s)$，不得偷偷把未知$P_j$添入常数域；
可用全部$r_j$的有限Wronskian消去，阶数不超过它们的常数线性秩。
这些是源的上界，不把$L$的阶数冒称$6+2r$。
\end{proof}

\begin{proposition}[几何精确输出与固定纤维动机输出的边界]\label{prop:dispersionarithmeticboundary}
本节构造了完整减除读出的实际相对周期、留数/管态射和精确的无限单值简单组成。
在每个代数纤维，它还给
\[
 H^3(U,B_U)\longrightarrow M_s\longrightarrow
 H^2(D_s,B_U\cap D_s)(-1)\longrightarrow H^4(U,B_U)
\]
这一实际对象图表，但不由此自动识别$M_s$或其物理最小子商的全部算术组成。
若仅保留最低阶$f$而不匹配在壳数据，上一节允许的sunset背景
仍保证相同受保护下界；(DC)的精确重数不无条件沿用。
\end{proposition}
\begin{proof}
固定纤维中的两个环境项是实际Hodge/Nori对象，
不能按其在族中单值平凡而在这些范畴中置零。
并且单个特殊纤维的子对象不必来自族的平坦子对象；
族的不可约性与特殊点物理最小性之间缺少的正是这一比较。
所以图表、形式管关系和函数分类都不蕴含所述更强算术结论。

未匹配的背景包含$S_{c_i}$及$\partial_{c_i}S_{c_i}$，
本身可以带同类椭圆组成，不能套用“背景外重数”而把这些因子忽略。
正常阈值$225$的变分仍杀掉所有该背景，故上一节的下界不变；
要断言扰动后重数仍恰一，另需共同实际载体的上界或相容匹配。
完整函数的形式最小化和几何函数系统的最小化应明确分开；
这一原则也与经典族周期框架\cite[\S2.4、\S7--8]{brown}一致。
\end{proof}

本节因而把一条完整接受规则接通：
合并树振幅$\longrightarrow$真实全谱$\longrightarrow$实际色散相对对
$\longrightarrow$留数上界与管变分下界$\longrightarrow$完整函数的精确运动学组成。
它没有将全部常数或混合粘合丢弃，也没有将这一个源模型重新定义为整个总目标。
下一主接口仍是更一般物理投影/森林与最小标架之间的实际比较，
而不是给同一模型增加极点数或小数位。



\section{一般多项式角分子的精确投影与K3充要判据}\label{sec:angularpolynomial}
前面的两对不变质量规则没有覆盖交叉标量积，例如$p_1\cdot p_3$。
本节补上这一接口：完整振幅可以含任意次数的标量积多项式，
精确角积分由有限张量缩并完成，并不需要先求其相空间周期。
正性用于识别物理读出；纯核的结论仍须同时使用实际载体和秩一变分。

固定四维、四个有标号的正质量粒子，先写一般$m_i>0$；
涉及六维K3的精确结论时再取$m=(1,2,4,8)$。
记$\mu_i=m_i^2$、$z_{ij}=p_i\cdot p_j$，
\[
 s=\sum_i\mu_i+2\sum_{i<j}z_{ij},\qquad
 a=\mu_1+\mu_2+2z_{12},\quad b=\mu_3+\mu_4+2z_{34}.
\]
系数域为带指定实嵌入的数域$k$，包含质量和耦合。
输入是先合并全部树图后的实标量振幅
\[
 F=\frac{N(s,(z_{ij})_{i<j})}{q(a)r(b)},\qquad
 N\in k[s,(z_{ij})_{i<j}],
 \qquad R_N(s)=\int d\Phi_4^{(4)}\,F^2.
 \tag{AP}\label{eq:angularinput}
\]
允许$q=r=1$；若有极点，采用第\ref{sec:weightedk3}节的阈值下接受条件。
本节的“角分子”限于这些标量积多项式，不包括角依赖传播子、外部自旋投影
或Levi--Civita张量。非零振幅指在质量壳及总动量关系上非零。

\subsection{协变球面矩给有限而精确的角投影}
令$K=p_1+p_2$、$L=p_3+p_4$，$w=K\cdot L=(s-a-b)/2$。
分解为
\[
 p_1=\alpha K+\xi,\quad p_2=(1-\alpha)K-\xi,\qquad
 p_3=\beta L+\zeta,\quad p_4=(1-\beta)L-\zeta,
\]
其中
\[
 \alpha=\frac{a+\mu_1-\mu_2}{2a},\quad
 \beta=\frac{b+\mu_3-\mu_4}{2b},\quad
 K\cdot\xi=L\cdot\zeta=0,\quad
 \xi^2=-\kappa_a^2,\quad \zeta^2=-\kappa_b^2,
\]
\[
 \kappa_a^2=\frac{A(a)}{4a},\quad
 \kappa_b^2=\frac{B(b)}{4b},\qquad
 A(a)=\lambda(a,\mu_1,\mu_2),\quad B(b)=\lambda(b,\mu_3,\mu_4).
\]
以$g=\operatorname{diag}(1,-1,-1,-1)$，
记$h_K^{\mu\nu}=-g^{\mu\nu}+K^\mu K^\nu/a$，$h_L$同理。

\begin{proposition}[角平均的有理下降，不引入新的根式]\label{prop:angularmoments}
在固定$s,a,b$的物理内点，令$\mathcal A$为两个内层二体角的独立归一化平均。
奇矩为零，偶矩严格为
\[
 \mathcal A_K(\xi^{\mu_1}\cdots\xi^{\mu_{2j}})
 =\frac{\kappa_a^{2j}}{(2j+1)!!}
   \sum_{\text{成对配对}}\prod_{\{\ell,t\}}h_K^{\mu_\ell\mu_t},
 \tag{AM}\label{eq:angularmoments}
\]
$\zeta$具有相同公式。于是
\[
 W_N(s,a,b)=\frac{\mathcal A(N^2)}{q(a)^2r(b)^2}
 \in k[s,a,b,a^{-1},b^{-1},q(a)^{-1},r(b)^{-1}],
\]
且实际四粒子谱恰为
\[
 R_N(s)=\frac1{2048\pi^5}\int_{\mathcal R_s}
       W_N(s,a,b)\frac{\sqrt{A(a)B(b)\lambda(s,a,b)}}{sab}\,da\,db.
 \tag{AR}\label{eq:angularreduced}
\]
角投影不新添有限极点；$W_N$在阈值点$(M^2,a_+,b_+)$正则。
\end{proposition}
\begin{proof}
在$K$的静止系，$\xi=(0,\kappa_a n)$，$n$在单位$S^2$上均匀分布。
旋转不变性使偶矩为Kronecker配对和的常数倍。
收缩最后两个指标并用$n\cdot n=1$，其系数满足
$c_j=c_{j-1}/(2j+1)$、$c_0=1$，故$c_j=1/(2j+1)!!$。
反演$n\mapsto-n$使奇矩为零。
协变化得到(AM)；两个衰变角在实际两对相空间分解中独立。

将全部$p_i$作上述代入，先展开完整$N^2$，再逐次用(AM)。
所有缩并最终只含$K^2=a,L^2=b,K\cdot L=w$及质量。
偶矩中的半径仅以$\kappa_a^2,\kappa_b^2$出现，故没有残留平方根；
投影张量只引入$a,b$的逆幂，没有Gram行列式或$\lambda$的逆幂。
这证明所写有理环及极点结论。

剩下的整体空间方向因$N$为Lorentz标量而不影响积分。
每个内层角平均均除以$4\pi$，对应角体积已经包含在
$\Phi_2^{(4)}=\sqrt{\lambda}/(8\pi\,\text{母质量平方})$中。
对三个二体相空间及$da\,db/(2\pi)^2$应用命题\ref{prop:evenphase}，恰得(AR)。
该递归与标准二体角测度见\cite[式(49.9)、(49.12)--(49.18)]{PDGphase}；
本节保持$d\Phi$中含$(2\pi)^4\delta^4$的约定，不重复乘该常数。
\end{proof}

\subsection{物理零方向可以在积分前精确判定}
定义壳上约化多项式
\[
 \overline N(z)=N\left(\sum_i\mu_i+2\sum_{i<j}z_{ij},z\right).
 \tag{AN}\label{eq:angularnormalform}
\]
这里$s$随谱参数变化，不先固定在某一个特殊值。

\begin{theorem}[多项式、角平方和与完整正谱的共同零核]\label{thm:angularzero}
以下三者等价：
\[
 \overline N=0,\qquad W_N\equiv0,\qquad R_N\equiv0.
\]
若$\overline N\ne0$，则任意足够小的正常阈值右邻域内都有正谱点，
并且它不可能在整个邻域恒零。对给定有限多项式输入，这个判定只需
一次线性变量替换和精确系数比较；不存在四个四维动量的额外Gram多项式恒等式。
\end{theorem}
\begin{proof}
$\overline N=0$使物理矩阵元逐点为零，其余两项随之为零。
反过来，固定一个物理内点的$s,a,b$，$N^2$连续非负，
角测度处处为正；故其平均为零恰当$N$在该角纤维处处为零。
再对实际正$a,b$测度应用同一事实，若$R_N$作为函数恒零，
则$N$在全部正能量四粒子物理构型上消失。

证明这些构型的六个$z_{ij}$包含实开集。
四个线性无关的未来类时四向量给Gram矩阵
$G=(p_i\cdot p_j)$，对角固定为$\mu_i$、惯性指数$(1,3)$。
惯性指数在非退化矩阵的小扰动下保持。
因原四向量可取线性无关且均未来，邻近Gram矩阵可由邻近四向量实现，
并保持同一未来分支。对角质量壳约束因此不再限制六个非对角元的小扰动。
其和$P$仍未来类时，可作Lorentz变换到总静止系。
非零实多项式不可能在这一六维开集恒零，故$\overline N=0$。

还可把四个空间动量取任意小而一般、并令其和为零。
相应四向量仍可线性无关；它们的Gram矩阵虽接近阈值退化矩阵，
在每个严格正的小能量处仍有上述开邻域。
所以同一论证适用于任意阈值右邻域。
这里没有对五个独立四向量作错误的Gram计数：
总动量$P$已是四个$p_i$的和，唯一引入的$s$关系正是(AN)。
\end{proof}

\subsection{有限阈值矩，而不是完整求积}
将有理$W_N$通分为$H/D$，其中$H,D\in k[s,a,b]$且
$D(M^2,a_+,b_+)>0$；可取$a,b$的幂及$q^2r^2$为分母。
设$H\ne0$，其阈值平移的最低非零齐次项为$H_\ell$。
令$M_A=m_1+m_2$、$M_B=m_3+m_4$、$M=M_A+M_B$，并写
\[
 h_\ell(x,y)=
 \frac{H_\ell(1,M_Ax/M,M_By/M)}{D(M^2,a_+,b_+)}.
\]
记$\mathfrak D$为三角形$x,y>0,x+y<1$上参数
$(3/2,3/2,3/2)$的归一化Dirichlet平均；对单项式
\[
 \mathfrak D(x^u y^v)=
 \frac{(3/2)_u(3/2)_v}{(9/2)_{u+v}}.
 \tag{AD}\label{eq:angulardirichlet}
\]

\begin{proposition}[角投影后的有限首项严格为正]\label{prop:angularjet}
若$\overline N\ne0$，则$\ell\le\deg H$，
$h_\ell$在开三角形上非负且不恒零，并有
\[
 R_N(M^2+\varepsilon)=
 \frac{\sqrt{\prod_i m_i}}{13440\pi^4M^5}
 \mathfrak D(h_\ell)\,
 \varepsilon^{7/2+\ell}\bigl(1+O(\varepsilon)\bigr),
 \qquad \mathfrak D(h_\ell)>0.
 \tag{AT}\label{eq:angularthreshold}
\]
因此首项由有限多项式平移和(AD)给出，不要求相空间数值积分。
\end{proposition}
\begin{proof}
定理\ref{thm:angularzero}保证$H\ne0$。
平移后的最低项在$H$的有限次数内出现；
齐次非零$H_\ell(t,u,v)$在$t=1$的限制仍是非零多项式，
故$h_\ell$不恒零。
对任意固定开三角形内点，取
$a=a_++(M_A/M)\varepsilon x$、$b=b_++(M_B/M)\varepsilon y$。
足够小$\varepsilon>0$时这是物理内点，而$W_N\ge0$，
所以其最低系数$h_\ell(x,y)$非负。
一个非零非负多项式在正测度下的平均严格为正。

使用内部两对动能及两对之间相对动能作固定单纯形坐标。
三个二体动量因子各给一个平方根，归一化后的极限密度
正是$x^{1/2}y^{1/2}(1-x-y)^{1/2}$。
原$a,b$与这些坐标相差$O(\varepsilon^2)$的解析项；
$D$在阈值非零，使$W_N$的展开在闭单纯形上一致。
密度余项在同一个可积Dirichlet权下为$O(\varepsilon)$，
所以可逐项取得首项和相对误差。
常数由单位权的命题\ref{prop:evenphase}固定，得(AT)。
(AD)由三个Beta积分或逐次Beta递推得到。
这使用正的完整平方平均，不能推广为任意有符号权。
\end{proof}

\subsection{从精确角投影到最小K3因子}
\begin{theorem}[一般标量积多项式的物理K3充要规则]\label{thm:angulark3}
取$m=(1,2,4,8)$以及上述接受域。
则$\overline N\ne0$当且仅当$R_N$的最小函数系统恰含一份$T_X$。
若$q=r=1$，该系统在有限单值简单背景外的半单组成恰为$T_X$；
若有接受的两对极点，则额外无限单值简单因子的上界仍仅为
命题\ref{prop:weightedupper}列出的椭圆曲线，不断言它们全部可见。
任意非零这类读出的齐次有理微分阶数至少为六。
\end{theorem}
\begin{proof}
(AR)是原实际相空间的精确积分下降，所得二形式为$W_N\eta_{2,s}$。
命题\ref{prop:angularmoments}证明它只在原零/无穷远边界和已接受的
两对极点上有极点。正常交叉化的曲面因而与
命题\ref{prop:evenboundary}或命题\ref{prop:weightedupper}相同，
给至多一份$T_X$及所列椭圆候选。
增高极阶不增添除子分量，角次数本身也不复制纯因子。

在正常阈值的局部$D_4$消失区，$W_N$正则，
所以$W_N\eta_2$在三个例外有理曲线上仍积分为零。
实际物理双三角链的类为$c\delta$加这些例外类，$c\ne0$；
这里$c$由链的几何决定，与读出多项式无关。
命题\ref{prop:angularjet}给该周期非零，
故消失方向$\delta$的读出也非零。
新增两对极点仍远离消失区，在阈值邻域同型；
因此命题\ref{prop:evenlocal}的秩一反射及引理\ref{lem:evenvisibility}适用。
已知指定质量射线的$T_X$简单性给至少一份，下界与上界相等。
若$\overline N=0$，函数本身为零，反向蕴含由此成立。

$q=r=1$时角消去新增的$a,b$逆幂只落在有理边界，
其它简单组成为有限单值Artin--Tate背景，故所述半单分类精确。
这里分类的是随$s$变化的函数，不能推出每个特殊$s$的数值非MZV，
也不声称完整微分阶数恰等于六。
\end{proof}

\subsection{一个不能先平均振幅再平方的精确实例}
\begin{proposition}[纯角分子均值为零而物理K3仍保留]\label{prop:angularcounterexample}
令$\nu_a=(a+\mu_1-\mu_2)/2$、$\nu_b=(b+\mu_3-\mu_4)/2$，取多项式
\[
 N_\angle=ab\,p_1\cdot p_3-b\nu_a K\cdot p_3
                 -a\nu_b L\cdot p_1+\nu_a\nu_b K\cdot L
           =ab\,\xi\cdot\zeta.
\]
它的精确角平均为零，但
\[
 \mathcal A(N_\angle^2)
   =\frac{A(a)B(b)}{144}\left(2ab+w^2\right)\ne0.
 \tag{AE}\label{eq:angularexample}
\]
在指定质量上，其谱阈值阶为$\varepsilon^{11/2}$，首项系数为
单位谱常数乘$110592/275$，且最小函数系统恰含一份K3核。
\end{proposition}
\begin{proof}
独立奇矩为零给$\mathcal A(\xi\cdot\zeta)=0$。
二阶矩的收缩则为
\[
 \mathcal A((\xi\cdot\zeta)^2)
 =\frac{\kappa_a^2\kappa_b^2}{9}\,
                   h_K^{\mu\nu}h_{L,\mu\nu}
 =\frac{A(a)B(b)}{144ab}
                   \left(2+\frac{w^2}{ab}\right).
\]
乘$(ab)^2$得到(AE)。
在$(s,a,b)=(225,9,144)$，$w=36$，
$A'(9)=8$、$B'(144)=128$，
故最低项为$27648(a-9)(b-144)$。
代入$a-9=\varepsilon x/5$、$b-144=4\varepsilon y/5$，
得到$h_2=(110592/25)xy$。
由$\mathfrak D(xy)=1/11$和(AT)给所写首项。
定理\ref{thm:angulark3}给K3结论。
因此$\mathcal A(N)^2$不能代替$\mathcal A(N^2)$：
前者把一个非零可观测量及其几何核全部错误消掉。
\end{proof}

\subsection{局部源、完整三圈函数与有限投影核}
\begin{theorem}[任意有限角导数源的完整物理出口]\label{thm:angularsource}
对$q=r=1$，每个实标量积多项式$N$都可由四个不同自由实标量场的
局部正规序复合算符$\mathcal O_N$实现为真空到四粒子的矩阵元。
其完整二点非局部谱为$R_N$，对应有分子三圈banana图。
若$D$为$N$在$s,z_{ij}$中的总次数，取$J=2D+3$，
则减除函数
\[
 G_N(s)-P(s)=(s-s_0)^J\int_{225}^{\infty}
       \frac{R_N(x)\,dx}{(x-s)(x-s_0)^J}
 \tag{AF}\label{eq:angularfull}
\]
绝对收敛，并在$\overline N\ne0$时恰含一份无限单值简单因子$T_X$。
其对应半单几何群为$O(q)$，恒等分支为$\mathrm{SO}(q)$。
固定算符后改变允许的局部二点接触项不改变该组成。

进一步，给定有限算符列表$N_1,\ldots,N_t$，把$\overline N_i$的全部
单项式系数排成矩阵$C$。对于实常数向量$u$，完整对角谱的不可见核恰为
\[
 u^t R(s)u\equiv0
 \quad\Longleftrightarrow\quad Cu=0,\qquad
 R_{ij}(s)=\int d\Phi_4^{(4)}\,N_iN_j.
 \tag{AK}\label{eq:angularkernel}
\]
每个核外组合的对角谱及其完整减除二点函数均保留一份K3。
\end{theorem}
\begin{proof}
单项式$\prod_{i<j}z_{ij}^{n_{ij}}$通过在第$i,j$种场上配对放置
$n_{ij}$个导数并缩并指标实现；每个动量因子来自相应场的导数。
选取实符号系数补偿傅里叶变换的$i$因子。
$s$的幂由整个局部插入上的$-\Box$实现。
四种场不同，正规序四场算符作用于真空只产生每种各一个的四粒子态。
因此任意多项式及全部干涉都可由有限线性组合实现；
不能用作用量中的分部积分把局部源的整体导数直接删掉。

自由Wick定理使二点函数在两个不同点之间恰有四条连接线和两个插入，
所以圈数为$4-2+1=3$。没有其它粒子道。
Wightman谱由完整矩阵元平方给$R_N$。
时间序分布在非重合点已固定，延拓歧义只支撑于两点重合，
故其Fourier变换是局部多项式。
沿用第\ref{sec:fullsourcebridge}节的$G=2\pi\Pi$规范，
其跳跃为$2\pi iR_N$，减除的谱构造给(AF)。

物理区的每个$p_i\cdot p_j$及$s$均为$O(s)$，
故$N=O(s^D)$、$R_N=O(s^{2+2D})$。
$J=2D+3$给无穷远$x^{-2}$尾部，阈值由(AT)可积。
将(AR)和色散变量合并，得到第\ref{sec:dispersioncarrier}节的固定三维模型，
仅把$f(a)^2$换成$W_N(x,a,b)$。
它仍只有原有理水平边界和固定竖直极点；无穷远以同一单位因子拆开。
相对Gysin给完整函数至多一份无限单值简单因子$T_X$，
实际管变分和定理\ref{thm:angulark3}给至少一份。
这还识别所述半单表示及其群；其它常值周期和混合扩张不被丢弃。
局部接触多项式属于有限单值简单背景，不改变这一结论。

最后$u^tR(s)u$就是合并振幅$\sum_i u_iN_i$的完整正谱。
定理\ref{thm:angularzero}把其零性等价于约化多项式为零，
而后者恰为有限线性方程$Cu=0$，证明(AK)。
这里的核指作为$s$函数恒零的常数系数组合；
不声称每个特殊$s$处的谱矩阵有相同秩，也不把各交叉项单独当正谱。
\end{proof}

这一规则将任意有限导数阶的多项式角分子接到真实物理投影和完整三圈函数，
且不需求完整周期的求值。其实现只需多项式代入、有限偶矩缩并、
有限阈值喷射及系数矩阵的线性代数；朴素配对展开可为阶乘规模，不宣称多项式复杂度。
角依赖传播子的积分会改变上述有理环，可能产生新的对数或几何，须另作实际分析。
尤其原$\phi^4$十树图含这类分母，本节没有关闭它的完整四体几何及三圈光学余项。
同样，函数的实际角下降和新色散提升不自动证明原动量/森林载体之间的
形式周期相等；固定纤维算术、混合最小对象及一般图规则仍按原目标继续。



\section{含角传播子的精确旋转推前与四体产生}\label{sec:rotationdenominators}
上一节的多项式角矩不能直接处理含角传播子。
本节在实际四维$\phi^4$四体产生中处理这个缺口：
先保持末态内部构型不变，精确积分它相对入射束的整体方向。
所得规则适用于任意有限个无极点、单交换型线性角分母的完整干涉。
剩余内部形状的积分仍须独立分类；不能由这一步的对数表达给整个截面贴上MPL标签。

\subsection{共同旋转变量与一个统一的角积分核}
令$n\in S^2$，以$d\Omega_n/(4\pi)$为归一化测度。
在带实嵌入的系数域上考虑
\[
 f(n)=c_0+\sum_{i=1}^r\frac{c_i}{D_i(n)},\qquad
 D_i(n)=a_i-b_i\cdot n,\quad a_i>|b_i|,\quad c_i\in\mathbb R.
 \tag{RA}\label{eq:rotationinput}
\]
这里$b_i$是三维Euclid向量；$a_i,b_i,c_i$可代数地依赖保持固定的内部运动学。
条件$a_i>|b_i|$保证正分母；若原传播子为负，将符号吸入$c_i$。
不得给每个传播子选择彼此独立的角变量：所有项依赖同一个$n$。

记$v_i=(a_i,b_i)$为辅助Minkowski向量，其配对为
\[
 A_i=v_i^2=a_i^2-|b_i|^2,\quad
 C_{ij}=v_i\cdot v_j=a_i a_j-b_i\cdot b_j,\quad
 \Delta_{ij}=C_{ij}^2-A_iA_j.
\]
这些辅助向量均为未来类时，故$C_{ij}\ge\sqrt{A_iA_j}>0$。
再引入$v_0=(1,0)$，使$D_0=1$。

\begin{theorem}[单交换角分母的完整平方核]\label{thm:rotationkernel}
对所有上述输入，归一化角平均恰为
\[
 \mathcal A(f^2)=\sum_{i,j=0}^r c_i c_j\,\Gamma_{ij},
 \qquad \Gamma_{ij}=\mathcal K(v_i,v_j),
\]
其中对$v=(a,b),w=(d,e)$置$A=v^2,B=w^2,C=v\cdot w$、
$\Delta=C^2-AB$，有
\[
 \mathcal K(v,w)=
 \begin{cases}
 \displaystyle\frac1{2\sqrt\Delta}
  \log\frac{C+\sqrt\Delta}{C-\sqrt\Delta},&\Delta>0,\\[6pt]
 \displaystyle\frac1C,&\Delta=0 .
 \end{cases}
 \tag{RK}\label{eq:rotationkernel}
\]
实支由$\sqrt\Delta\ge0$及正对数比固定。特别地，
\[
 \Gamma_{ii}=\frac1{A_i},\qquad
 \Gamma_{0i}=
 \frac1{2|b_i|}\log\frac{a_i+|b_i|}{a_i-|b_i|}
\]
而$b_i=0$时$\Gamma_{0i}=1/a_i$。
这是完整干涉矩阵，不能改成各图独立平方之和。
\end{theorem}
\begin{proof}
对正数$D,E$有有限参数恒等式
\[
 \frac1{DE}=\int_0^1\frac{dt}{(tD+(1-t)E)^2}.
\]
它可由一个有理原函数直接验证，$D=E$时取连续极限。
另一方面，沿$b$的方向作极轴，
\[
 \mathcal A\bigl((a-b\cdot n)^{-2}\bigr)
 =\frac12\int_{-1}^1\frac{dz}{(a-|b|z)^2}
 =\frac1{a^2-|b|^2}.
\]
未来类时锥的凸性使中间分母始终为正；连续有界性允许交换积分。
所以
\[
 \mathcal K(v,w)=
 \int_0^1\frac{dt}{A t^2+2Ct(1-t)+B(1-t)^2}
 =\int_0^\infty\frac{du}{A u^2+2Cu+B},
 \tag{RP}\label{eq:rotationparameter}
\]
最后一步用$u=t/(1-t)$。二次式的两个根是
$(-C\pm\sqrt\Delta)/A$；部分分式并取$0,\infty$端点即得(RK)。
当$\Delta=0$，直接积分$(\sqrt A u+\sqrt B)^{-2}$得
$1/\sqrt{AB}=1/C$。
$v=w$及$w=v_0$给所列特殊值。
最后对完整$f^2$作双线性展开，所有交叉项的系数自动正确。
\end{proof}

\begin{proposition}[合并极部后的精确零核]\label{prop:rotationzerokernel}
置$u_i=b_i/a_i$，将所有非零且相同的$u_i$合并，记其不同值为集合$\mathcal U$。
定义
\[
 d_0=c_0+\sum_{u_i=0}\frac{c_i}{a_i},\qquad
 d_u=\sum_{u_i=u}\frac{c_i}{a_i}\quad(u\in\mathcal U).
\]
则
\[
 \mathcal A(f^2)=0
 \quad\Longleftrightarrow\quad
 d_0=0\ \hbox{且所有 }d_u=0.
 \tag{RZ}\label{eq:rotationzero}
\]
包含常数函数的Gram矩阵$\Gamma$之秩恰为$1+|\mathcal U|$。
因此在给定代数实构型上，判零只需归一化、合并相同向量和精确线性求和；
不需求对数值或相空间积分。
\end{proposition}
\begin{proof}
正测度和连续性使$\mathcal A(f^2)=0$等价于$f$在整个实球面上为零。
合并后
\[
 f(n)=d_0+\sum_{u\in\mathcal U}\frac{d_u}{1-u\cdot n},
 \qquad 0<|u|<1.
\]
实球面在复仿射二次曲面$Q:n_1^2+n_2^2+n_3^2=1$中Zariski稠密；
可由球面的有理参数化在一个实开片上直接看出。
故实球面上的恒等式也是$Q$的有理函数恒等式。

每个$1-u\cdot n=0$在$Q$上截出一条光滑不可约仿射圆锥曲线。
实旋转将$u$置于第三轴后，它的方程为
$n_1^2+n_2^2=1-|u|^{-2}\ne0$，故确实光滑不可约。
不同$u$给不同的曲线：一条非退化圆锥曲线的仿射线性包唯一确定其所在平面，
而所取平面方程的常数项均为1。
在一条曲线上避开其它曲线的交点，只有相应一项有简单极点。
有理恒等式迫使该$d_u=0$；逐条处理后$d_0=0$。
反向显然。所得$1+|\mathcal U|$个函数线性无关，
其正$L^2(S^2)$内积的Gram秩即为此数。
\end{proof}

\subsection{在十张真实树图中应用同一个核}
现在回到极点质量$m>0$、相同实标量的四维$\phi^4$理论。
考虑$s>16m^2$的$2\to4$过程，末态有标号动量为$k_1,\ldots,k_4$，
$P=\sum k_i=(\sqrt s,0)$。
总静止系的入射动量写作
\[
 p_\pm=\left(\frac{\sqrt s}{2},\ \pm p\,n\right),\qquad
 p=\frac{\sqrt{s-4m^2}}2.
\]
固定四个$k_i$的内部形状，$n$就是入射方向相对该形状的变量。
十个无序$3|3$树分割及各自权重1沿用原基准；下面给其直接几何重组。
四个类时项之和为
\[
 T=\sum_{i=1}^4\frac1{(P-k_i)^2-m^2}
   =\sum_{i=1}^4\frac1{s-2\sqrt s\,E_i}.
\]
对六个二元子集$J\subset\{1,2,3,4\}$，定义
\[
 K_J=\sum_{j\in J}k_j=(E_J,\mathbf K_J),\quad
 a_J=K_J^2,\quad b_J=(P-K_J)^2,
\]
\[
 \alpha_J=\frac{s-a_J-b_J}{2},\qquad
 B_J=2p\,\mathbf K_J,\qquad
 v_J=(\alpha_J,B_J).
\]
这里大写向量$B_J$与标量$b_J$不同。六个混合树传播子满足
\[
 (p_+-K_J)^2-m^2=-\alpha_J+B_J\cdot n.
\]
于是无耦合的完整分母和是
\[
 S(n)=T-\sum_{|J|=2}\frac1{\alpha_J-B_J\cdot n},
 \qquad \mathcal M_{24}^{(0)}=-g^2S(n).
 \tag{RT}\label{eq:rotationphi4}
\]
$J$和$J^c$都必须保留：它们是两张不同树图，
且$B_{J^c}=-B_J$、$\alpha_{J^c}=\alpha_J$；
这不是把同一个无序$3|3$分割重复计算。

\begin{proposition}[全物理域的接受条件与严格正性]\label{prop:rotationphi4positive}
上述六个$v_J$在每个物理构型上都为未来类时，
并且$\alpha_J-|B_J|>m^2$。
四个类时分母至少为$8m^2$。
对所有有限$s>16m^2$及其每个物理末态形状，都有
\[
 W_{\phi^4}:=\mathcal A(S^2)>0.
 \tag{R+}\label{eq:rotationpositive}
\]
因此首个四体产生截面的系数在整个阈值以上严格为正。
这不要求$S(n)$对每个入射方向都有相同符号。
\end{proposition}
\begin{proof}
$a_J,b_J\ge4m^2$且$\sqrt{a_J}+\sqrt{b_J}\le\sqrt s$，
从而$\alpha_J\ge\sqrt{a_Jb_J}\ge4m^2$。
记$\lambda_J=\lambda(s,a_J,b_J)$，则
\[
 |B_J|^2=\frac{s-4m^2}{4s}\lambda_J,\qquad
 A_J:=v_J^2=a_Jb_J+\frac{m^2}{s}\lambda_J>0.
\]
更强的恒等式
\[
 (\alpha_J-m^2)^2-|B_J|^2
 =(a_J-m^2)(b_J-m^2)+\frac{m^2}{s}(a_J-b_J)^2>0
\]
结合$\alpha_J-m^2>0$，证明严格间隔。
类时项的另一侧是三个未来向质量$m$的粒子，
其总质量至少$3m$，故所述$8m^2$界成立。
这些也独立重证旧基准的无极点条件。

应用命题\ref{prop:rotationzerokernel}。
在任意非零$u=B_J/\alpha_J$组内，每个系数均为$-1/\alpha_J<0$，
所以相同分母的极部相加后绝不会相消。
若所有$B_J=0$，则由
$\mathbf k_1+\mathbf k_2=\mathbf k_1+\mathbf k_3=
\mathbf k_1+\mathbf k_4=0$及总空间动量为零，
可得全部$\mathbf k_i=0$，这只可能在$s=16m^2$。
阈值以上至少一组非恒定分母存在，其非零极部使角平均平方严格为正。
相空间内域非空、正测度且有限；乘以正通量倒数和$1/4!$，
得到完整产生系数的严格正性。
\end{proof}

\begin{theorem}[保留全部干涉的五变量物理推前]\label{thm:rotationfullrate}
记$J_J=\mathcal K(v_0,v_J)$、$K_{IJ}=\mathcal K(v_I,v_J)$。
固定内部形状后的精确读出是
\[
 W_{\phi^4}
 =T^2-2T\sum_J J_J+\sum_J\frac1{A_J}
       +2\sum_{I<J}K_{IJ}.
 \tag{RW}\label{eq:rotationweight}
\]
它包含原十图全部十个平方项和四十五个干涉对。
以$a=(k_1+k_2)^2,b=(k_3+k_4)^2$及两个对内极角余弦$x,y$、
相对方位角$\varphi$为形状坐标，定义
\[
 \rho(s,a,b)=\frac{\sqrt{\lambda(s,a,b)}}{2048\pi^5s}
              \sqrt{1-\frac{4m^2}{a}}\sqrt{1-\frac{4m^2}{b}},
 \quad
 \mathcal R_s=\{a,b\ge4m^2,\ \sqrt a+\sqrt b\le\sqrt s\}.
\]
在原基准的$d\Phi$含$(2\pi)^4\delta^4$约定下，严格有
\[
 \frac{\sigma_{24}^{[g^4]}(s)}{g^4}
 =\frac1{2s\beta\,4!}\int_{\mathcal R_s}\rho(s,a,b)\,da\,db
  \int_{-1}^1\frac{dx}{2}\int_{-1}^1\frac{dy}{2}
  \int_0^{2\pi}\frac{d\varphi}{2\pi}\,W_{\phi^4},
 \quad \beta=\sqrt{1-\frac{4m^2}s}.
 \tag{RF}\label{eq:rotationrate}
\]
式中$\sigma_{24}^{[g^4]}$指截面中阶$g^4$这一整项。
$(T,A_J,C_{IJ})$只需末态不变量：
\[
 C_{IJ}=\alpha_I\alpha_J
 -(s-4m^2)(E_I E_J-K_I\cdot K_J),\qquad
 E_J=\frac{s+a_J-b_J}{2\sqrt s}.
\]
因此公式没有剩余的入射方向变量。
\end{theorem}
\begin{proof}
(RW)是定理\ref{thm:rotationkernel}对(RT)的直接应用。
$T^2$中保留四个平方和六个干涉对，
$T$与六个混合项保留二十四对，
六个混合项内部保留六个平方和十五对；
故恰为原十平方及四十五干涉对，没有使用非物理的逐图正性。

两对相空间递推给$\rho(s,a,b)\,da\,db$乘三个归一化二体角测度，
其$2\pi$约定和二体公式见\cite[式(49.12)--(49.18)]{PDGphase}。
在一般内部形状上，这三个角测度可分解为
\[
 dR_{\mathrm{Haar}}\,
 \frac{dx}{2}\frac{dy}{2}\frac{d\varphi}{2\pi},
 \qquad \int_{\mathrm{SO}(3)}dR_{\mathrm{Haar}}=1.
\]
具体地，外层方向的$d\Omega$与两个内层方位角中的共同方位
组成体积$8\pi^2$的整体旋转；
余下相对方位体积为$2\pi$，两个余弦各长2，
恢复原$(4\pi)^3$。
退化形状是零测集，积分恒等式通过连续性成立。

把整体旋转转移到束方向，
$R^{-1}n_{\mathrm{beam}}$在$S^2$上均匀分布，
故整体旋转积分正是$\mathcal A$，没有额外$4\pi$或$8\pi^2$。
所有分母有前述一致间隔，固定$s$的相空间有限，
所以该平均与形状积分可以交换。
最后保留四个相同末态的$1/4!$及入射通量$2s\beta$。
这与上一节四种不同粒子的源模型不同，不能照搬其不含$1/4!$的公式。
所列$C_{IJ}$只是$\mathbf K_I\cdot\mathbf K_J=E_I E_J-K_I\cdot K_J$的代入。
\end{proof}

\subsection{对数载体、可去退化和完整分类的边界}
\begin{proposition}[旋转纤维的Kummer载体及高幂闭包]\label{prop:rotationkummer}
在$\Delta\ne0$的代数开集上，(RP)是一个明确的相对曲线周期：
\[
 H^1\!\left(\mathbb P^1\setminus\{r_-,r_+\},\,\{0,\infty\}\right),
 \qquad r_\pm=\frac{-C\pm\sqrt\Delta}{A}.
 \tag{RC}\label{eq:rotationcarrier}
\]
在添加$\sqrt\Delta$的有限代数覆盖后，其权分级为
$\mathbb Q(0)$及$\mathbb Q(-1)$；
系数及Tate归一化须保留。旋转积分本身不需要椭圆曲线。

在物理域$\Delta=0$是(RK)的可去退化，而非真实角夹挤：
\[
 \mathcal K=\frac1C\sum_{n\ge0}
       \frac{\Delta^n}{(2n+1)C^{2n}},\qquad |\Delta|<C^2 .
 \tag{RL}\label{eq:rotationlimit}
\]
若同一个单交换角分母出现正整数高幂，则
\[
 \mathcal A(D_i^{-p}D_j^{-q})=
 \frac{(-1)^{p+q-2}}{(p-1)!(q-1)!}
 \partial_{a_i}^{p-1}\partial_{a_j}^{q-1}\mathcal K(v_i,v_j),
 \tag{RD}\label{eq:rotationderivatives}
\]
偏导时固定$b_i,b_j$，先将两份参数视为独立，最后才允许它们相合。
因而这些高幂没有增加对数长度或新的判别式根。
该命题不给含多个不同交换分母乘积的一般振幅同样的闭包。
\end{proposition}
\begin{proof}
二次多项式的根在物理域为负，故路径$u:0\to\infty$不穿过极点。
在$\mathbb P^1$中，$du/(Au^2+2Cu+B)$于无穷远正则；
0及$\infty$是相对端点。
开曲线与端点的上同调正合列给一个单位分级和一个Tate分级，
而部分分式得到的对数比正是其Kummer周期。
这是为已约化核显式给出的载体；若要求它与原完整动量图的形式周期逐一相等，
仍须提交相应比较态射，不能仅凭函数恒等式代替。

在未来类时域$AB>0,C>0$，展开
$\log((1+z)/(1-z))=2\sum_{n\ge0}z^{2n+1}/(2n+1)$
即得(RL)，其极限为$1/C$。
也可直接从(RP)的正分母得到连续解析延拓；
$A=0$或$B=0$等真正边界不在该物理域。
最后利用正间隔在参数小邻域的一致性，对角积分逐次求导，
每次$\partial_a D^{-p}=-pD^{-p-1}$给(RD)。
对数公式的有限代数微分仍在同一代数函数及一次对数生成的域中。
\end{proof}

旋转推前是结构预测的一个明确接受规则：
从图分割提取共同方向上的线性分母，核对正锥条件，
合并重复极部，再用二次式(RP)输出根式与Kummer因子。
对$r$个简单项，至多需要$r(r+1)/2$个双项核及$r$个常数干涉核；
这是代数表达式层面的二次数量界，不是数值精度或比特复杂度结论。
$\phi^4$例有六个混合项；不作互补配对简化时至多二十一个非对角对数核，
其余项有理。

这个结论关闭了完整十图的\emph{整体旋转}接口，也给任意有限$s>16m^2$的严格正性。
它没有关闭(RF)的五个内部形状变量的最小对象。
例如可令$z=\cos\varphi$，
把相对方位测度写为$dz/(\pi\sqrt{1-z^2})$；
再加入boost及内层极角的代数根式和(RC)的区间变量，
得到明确的代数覆盖与Kummer系数推前载体。
但进一步推前可能产生非Tate组成；
“旋转纤维是Kummer”不推出“完整截面是混合Tate”。
原十图还处在等质量射线上，不能把上一节$(1,2,4,8)$的秩六结论直接搬来。

最后，$s\downarrow16m^2$时$T\to1/(2m^2)$、
六个$D_J\to4m^2$，所以$W_{\phi^4}\to1/m^4$；
这恢复旧阈值证明的实际振幅，而非产生一套新数值证书。
本节没有重算四体样本、旧严格粗界或小数，
没有把右端的精确角投影称为三圈前向振幅的独立全切割验证。



\section{二维banana的实际色散比较与最小对象的转移}\label{sec:actualbananadispersion}
本节回到核心二维四传播子图，补齐一个比新增数值更直接的接口：
原正参数周期与正谱色散周期之间存在实际代数映射及有向链关系。
因此这里得到的是形式周期关系，而非由两个解析函数相等猜测动机相等。
经典banana割的递推及色散表示已有长期研究，参见\cite{KreimerBanana}第1--2节；
本节的作用是固定本项目的几何模型、整体符号、归一化及最小标架的传递。
最后把同一构造推广到任意有限条边的二维banana和任意二叉分组树。
这给一整个接受图类的实际比较规则，不宣称已经分类所有高维纯核。

\subsection{正环面与总空间动量为零的质量壳}
固定带实嵌入的数域$k$，$m_i\in k_{>0}$、$\mu_i=m_i^2$，先取$s\in k_{<0}$。
令$\alpha_4=1$，并写
\begin{equation}\label{eq:actualtorusperiod}
 f=\Bigl(\sum_{i=1}^4\mu_i\alpha_i\Bigr)
       \Bigl(\sum_{i=1}^4\alpha_i^{-1}\Bigr),\qquad
 \eta=d\log\alpha_1\wedge d\log\alpha_2\wedge d\log\alpha_3,
 \qquad I(s)=\int_{\alpha_i>0}\frac{\eta}{f-s}.
\end{equation}
这是旧相对比较采用的正物理周期，未另乘圈因子或图对称因子。
令$\tau=(m_1+m_2+m_3+m_4)^2$。
下面的非紧链均取正实取向，在射影紧化及所列代数边界的公共消解后理解为相对链。

\begin{proposition}[代数共同boost及无额外倍数的正链比较]\label{prop:actualboost}
积分\eqref{eq:actualtorusperiod}绝对收敛。
考虑三维代数质量壳
\[
 Y:\quad E_i^2-k_i^2=\mu_i,\qquad \sum_{i=1}^4 k_i=0,
 \qquad z_i=E_i+k_i,
\]
及其正分支$Y_+$：$E_i=\sqrt{k_i^2+\mu_i}>0$。
在相应稠密开集上，映射
\begin{equation}\label{eq:actualboostmap}
 \phi:Y\longrightarrow (\mathbb G_m)^3,\qquad
 \alpha_i=\frac{\mu_4z_i}{\mu_i z_4}\quad(1\leq i\leq3)
\end{equation}
满足$f\circ\phi=S^2$，其中$S=\sum_i E_i$。
其正实限制是有向微分同胚，且
\begin{equation}\label{eq:actualboostform}
 \phi^*\eta=\frac{S}{E_1E_2E_3E_4}\,dk_1\wedge dk_2\wedge dk_3.
\end{equation}
因而正质量壳上的右侧积分与$I^{\mathfrak m}(s)$由实际拉回关系表示同一形式周期。
代数映射一般次数为二，但正链只选择一次，不能在这里乘二。
\end{proposition}
\begin{proof}
由Cauchy不等式$f\geq\tau$。
写$t_i=\log\alpha_i,t_4=0$，则
\[
 f\geq\mu_{\min}\exp(\max_i t_i-\min_i t_i)
       \geq\mu_{\min}\exp(\|t\|_\infty).
\]
因此$s<0$时被积密度由$\mu_{\min}^{-1}e^{-\|t\|_\infty}d^3t$控制。
这直接证明全部无穷远方向的绝对收敛。

在$Y$上有$\mu_i/z_i=E_i-k_i$，所以
$\sum_i z_i=\sum_i\mu_i/z_i=S$。代入\eqref{eq:actualboostmap}即得$f=S^2$。
反之，给定正$\alpha$，置$u_i=\mu_i\alpha_i$、
$A=\sum_i u_i$、$B=\sum_i\mu_i/u_i$，并取
\[
 c=\sqrt{B/A}>0,\qquad z_i=cu_i,\qquad
 E_i=\tfrac12(z_i+\mu_i/z_i),\quad k_i=\tfrac12(z_i-\mu_i/z_i).
\]
则$\sum_i k_i=0$、$E_i>0$，并恢复原$\alpha$。
正平方根的选择唯一；另一根给同时翻转的负能量分支。
这说明该正链比较本身定义于代数二重覆盖，未使用超越boost坐标。

又因$d\log z_i=dk_i/E_i$及$dk_4=-\sum_{j=1}^3dk_j$，
三个$d\log\alpha_i$的系数矩阵为
\[
 \operatorname{diag}(E_1^{-1},E_2^{-1},E_3^{-1})
                  +E_4^{-1}{\bf1}{\bf1}^{\mathsf T}.
\]
其行列式为$S/(E_1E_2E_3E_4)>0$，证明形式及取向公式。
映射图的紧化消解给实际代数态射；推前正链恰为原正链。
加入无穷远及链边界后，形式周期的函子关系即给所述相等。
收敛保证这些模型采用普通收敛相对周期，不需要有限部分正规化。
\end{proof}

\subsection{欧氏动量积分的能量留数比较}
为避免把解析归一化隐藏在符号中，令$Q=\sqrt{-s}>0$，
并在$k'=k(i,Q)$上考虑有向实六维积分
\begin{equation}\label{eq:actualmomentum}
 \mathcal B(-Q^2)=\pi^{-3}\int_{\mathbb R^6}
 \frac{d\omega_1d\omega_2d\omega_3\,dk_1dk_2dk_3}
 {\prod_{j=1}^3(\omega_j^2+E_j^2)
       [(Q-\omega_1-\omega_2-\omega_3)^2+E_4^2]}.
\end{equation}
这里$k_4=-k_1-k_2-k_3$。形式与链同时按所显示的能量优先次序定向。
这是三个圈各除以$\pi$的二维欧氏标量banana。

\begin{proposition}[实际能量留数与Lefschetz归一化]\label{prop:actualenergyresidue}
式\eqref{eq:actualmomentum}绝对收敛，且$\mathcal B(s)=I(s)$。
令$\mathcal M_{\rm raw}^{\mathfrak m}$为去掉$\pi^{-3}$的代数形式周期，
$\mathbb L^{\mathfrak m}$为周期$2\pi i$的Lefschetz元素。
在$k'$上的局部化形式周期代数中，实际留数比较给
\begin{equation}\label{eq:actualmomentumformal}
 (2i)^3(\mathbb L^{\mathfrak m})^{-3}
             \mathcal M_{\rm raw}^{\mathfrak m}=I^{\mathfrak m}(s).
\end{equation}
这条动量比较显式保留基域$k'$；后面的参数--色散比较直接在$k$上完成。
\end{proposition}
\begin{proof}
对$A,B>0$和$q\in\mathbb R$，在上半平面取两个极点$iA,q+iB$，得到
\begin{align}
 \int_{\mathbb R}\frac{d\omega}
 { (\omega^2+A^2)((q-\omega)^2+B^2)}
 &=\pi\left\{\frac1{A((q-iA)^2+B^2)}
             +\frac1{B((q+iB)^2+A^2)}\right\}\notag\\
 &=\frac{\pi(A+B)}{AB[q^2+(A+B)^2]}.
 \label{eq:actualconvolution}
\end{align}
中间两项在重合极点处取合并值；右端在全部正实域正则。
依次卷积三次，能量积分恰为
\[
 \frac{\pi^3S}{E_1E_2E_3E_4(S^2+Q^2)}.
\]
原实密度非负，Tonelli允许先积分能量而不预设有限性；
命题\ref{prop:actualboost}随后证明其有限且等于$I(s)$。

还须说明为何这里有形式周期关系。
先拉到$E_j^2=k_j^2+\mu_j$的代数覆盖，Betti链逐个选择正根；
该覆盖在所选链上次数一。每个能量纤维紧化为$\mathbb P^1$，
实直线连同无穷远点是一条闭合实圆；其有向类是两个上半平面极点小圆的和。
有理一形式在无穷远正则。
对穿孔$\mathbb P^1$的实际留数/Gysin态射，
这个Betti圆类配对为$\mathbb L^{\mathfrak m}$乘相应代数留数和。
式\eqref{eq:actualconvolution}中大括号的一半$i^{-1}$正是该留数和，
故每次卷积的形式周期因子是$\mathbb L^{\mathfrak m}/(2i)$。

这些是带剩余动量参数的相对态射。
正实基上极点不穿过实圆；两上极点重合时使用它们的总留数，
而不分别延拓有极点的两项。其有理合并式的唯一实分母
$AB[q^2+(A+B)^2]$不为零，故总圆类及其配对跨过该重合层。
在极点判别式及无穷远的公共消解上完成该相对构造，
只需细分基链；人工细分面成对抵消，闭纤维不产生新的链端点。
绝对收敛处理剩余无穷远边界。
于是三次实际留数后所得正质量壳标架，
由\eqref{eq:actualboostmap}送回$I^{\mathfrak m}$，证明\eqref{eq:actualmomentumformal}。
整个论证使用代数拉回、相对留数与Betti链关系，没有调用周期映射单射性。
\end{proof}

\subsection{同一有理商的全三维扩展与八张正分支}
令
\[
 A(a)=\lambda(a,\mu_1,\mu_2),\quad B(b)=\lambda(b,\mu_3,\mu_4),\quad
 \lambda(v,w,z)=v^2+w^2+z^2-2vw-2vz-2wz.
\]
以$x$表示谱变量，不与参数$\alpha_i$混用。定义固定三维载体
\begin{equation}\label{eq:actualspectralcarrier}
 V:\ y^2=A(a)B(b)\lambda(x,a,b),\qquad
 \zeta_s=\frac{dx\wedge da\wedge db}{(x-s)y}.
\end{equation}
正谱链$\Xi_+$位于$y>0$，投影区域为
\[
 a\geq(m_1+m_2)^2,\quad b\geq(m_3+m_4)^2,\quad
 x\geq(\sqrt a+\sqrt b)^2.
\]
取$dx\wedge da\wedge db$的正向，记$J(s)=\int_{\Xi_+}\zeta_s$。
不等式所用平方根只是指定半代数实分支；链的代数边界包含于
$A B\lambda=0$和无穷远，因此$\Xi_+$定义$k$上的相对Betti标架。

\begin{theorem}[原参数与色散的实际形式周期等式]\label{thm:actualdispersionformal}
上述正谱积分绝对收敛。在同一$k$上，有
\begin{equation}\label{eq:actualeight}
                   I^{\mathfrak m}(s)=8J^{\mathfrak m}(s).
\end{equation}
该等式由原环面到\eqref{eq:actualspectralcarrier}的次数四有理映射、
正链的八块分割及覆盖对合给出。
次数四、八张正分支和前一命题的正能量二重覆盖是三个不同的计数。
\end{theorem}
\begin{proof}
在原环面上令$u_i=\mu_i\alpha_i$并模去共同尺度，取$u_4=1$。
采用定理\ref{thm:evenquotient}的同一有理变量
\[
 r=u_1/u_2,\quad t=u_3/u_4,\quad z=(u_1+u_2)/(u_3+u_4),
\]
及
\begin{align*}
 a&=(r+1)(\mu_1/r+\mu_2),&
 b&=(t+1)(\mu_3/t+\mu_4),\\
 x&=(z+1)(a/z+b),&
 y&=(\mu_2r-\mu_1/r)(\mu_4t-\mu_3/t)(bz-a/z).
\end{align*}
区别在于现在让$x=f$自由变化，故这是全三维映射$\Pi$。
它定义于$k$，并满足
\begin{equation}\label{eq:actualfullform}
             \frac{\eta}{f-s}=-\Pi^*\zeta_s.
\end{equation}
确实，$da=(\mu_2r-\mu_1/r)d\log r$，$db$同理，
$\partial_{\log z}x=bz-a/z$，而
$\eta=-d\log r\wedge d\log t\wedge d\log z$。
$dx\wedge da\wedge db$中的循环置换为正，因而负号不能删除。

对正谱域的任一严格内点，反解三个二次式
\[
 \mu_2r^2+(\mu_1+\mu_2-a)r+\mu_1=0,\quad
 \mu_4t^2+(\mu_3+\mu_4-b)t+\mu_3=0,\quad
 bz^2+(a+b-x)z+a=0.
\]
它们均有两个不同正根。依次选择三个根，给八个不同的原正环面点。
每一个选择在严格谱域上是一张解析代数分支，
分别由三个括号的符号$\varepsilon_1,\varepsilon_2,\varepsilon_3$标记。
其中四支映到$y>0$，四支映到$y<0$。
二重覆盖$V\to\mathbb A^3$不可约：例如把$x$作为变量，
$\lambda(x,a,b)$有两个不同根，不能由与$x$无关的$AB$变成平方。
故原环面到$V$的一般次数为$8/2=4$，与旧纤维比较一致。

记$\Xi_-$为$y<0$上同样按$dx\wedge da\wedge db$正向的链。
在$y>0$的四块上，\eqref{eq:actualfullform}说明$\Pi$反转正体积取向；
在$y<0$的四块上则保持取向。因此在包含分支及无穷远边界的相对链群中，
\begin{equation}\label{eq:actualchainpush}
                 \Pi_*[\sigma_+]=-4[\Xi_+]+4[\Xi_-].
\end{equation}
覆盖对合$\iota(x,a,b,y)=(x,a,b,-y)$满足
$\iota_*[\Xi_+]=[\Xi_-]$、$\iota^*\zeta_s=-\zeta_s$。
由形式周期的态射关系，$[\Xi_-,\zeta_s]^{\mathfrak m}=-J^{\mathfrak m}$；
再用\eqref{eq:actualfullform}和\eqref{eq:actualchainpush}恰得到$8J^{\mathfrak m}$。

这里链分割的内部面为三个根差等于零的代数面。
它们是正环面内的光滑有穷位置，原形式无极点，
相邻块的有向边界抵消。对映射图、这些面、$x=s$及无穷远作公共消解，
就将上述等式变成实际相对上同调的拉回与链推前关系。
所选实域上$x\geq\tau>0>s$，故不存在色散极点夹挤。
原积分绝对收敛，且每块的拉回密度为相同的正谱密度，
所以$J$也绝对收敛。该论证未通过极限数值等式认定形式周期相等。
\end{proof}

\subsection{物理割的系数及正支的全局解析函数}
用有标号末态的二维相空间规范
\[
 d\Phi_4^{(2)}=(2\pi)^2\delta^{(2)}(P-\textstyle\sum p_i)
                    \prod_{i=1}^4\frac{dk_i}{(2\pi)2E_i},
 \qquad P=(\sqrt x,0).
\]
记$\Phi_4^{(2)}(x)=\int d\Phi_4^{(2)}$；不在此处除以$4!$。

\begin{proposition}[无减除色散与四粒子跳跃的准确常数]\label{prop:actualcutnormalization}
在割平面$\mathbb C\setminus[\tau,\infty)$上，原正支满足
\begin{align}
 I(s)&=8\int_{\tau}^{\infty}\frac{dx}{x-s}
       \int_{\mathcal R_x}\frac{da\,db}{\sqrt{A(a)B(b)\lambda(x,a,b)}}
       \notag\\
     &=32\pi^2\int_{\tau}^{\infty}\frac{\Phi_4^{(2)}(x)}{x-s}\,dx,
       \label{eq:actualphysicaldispersion}\\
 \operatorname{Disc}I(s)&=64\pi^3i\,\Phi_4^{(2)}(s)\qquad(s>\tau),
       \label{eq:actualdisc}
\end{align}
其中$\mathcal R_x=\{a\geq a_+,b\geq b_+,\sqrt a+\sqrt b\leq\sqrt x\}$，
$\operatorname{Disc}I(s)=I(s+i0)-I(s-i0)$。
这识别该二维banana正物理支的全四线割，不是任意图Cutkosky定理。
\end{proposition}
\begin{proof}
在二体质心系直接积分能量$\delta$函数，两根$k=\pm|k|$给
\[
 \Phi_2^{(2)}(x;a,b)=\frac1{\sqrt{\lambda(x,a,b)}}.
\]
每一次插入对子不变量附带$da/(2\pi)$。
两次实际相空间分解因此得到
\[
 \Phi_4^{(2)}(x)=\frac1{4\pi^2}
       \int_{\mathcal R_x}\frac{da\,db}{\sqrt{A B\lambda}}.
\]
两个二体内层的左右根及外层的左右根正是八个离散质量壳分支。
与定理\ref{thm:actualdispersionformal}的八张正参数分支对应，
不再补乘一个方向体积。

收敛已在负实$s$证明。对割平面的任意紧集$K$，存在$C_K$，使
$|x-s|^{-1}\leq C_K(x+1)^{-1}$对$x\geq\tau,s\in K$成立。
由同一个正测度的可积控制，色散式定义该域上的全纯函数，
并与原正参数积分一致。$s>\tau$时谱密度在内点光滑，
局部用$1/(x-s\mp i0)$的差$2\pi i\delta(x-s)$，
远端用上述可积界，即得\eqref{eq:actualdisc}。
所有$2\pi$因子均来自已声明的相空间规范。
这不把上半平面边界值的符号与旧$P/-P$约定互换。
\end{proof}

\subsection{指定固定纤维的完整最小对象与coaction}
这里使用既有核心证明作为输入：在$k=\Q$、$m=(1,2,4,8)$、$s=-64$，
正标架$I^{\mathfrak m}$的最小Nori对象为已经构造的秩八$M$，
并已给出其完整扩张、固定群和比较规范。
这不是从一般参数的纯单值群重新推定固定点群。

\begin{corollary}[固定色散周期继承完整秩八对象]\label{cor:actualdispersionminimal}
在上述指定数据下，$J^{\mathfrak m}$的最小对象与$M$同构。
因此其权重、Hodge层、扩张是否分裂及固定Tannaka群，
均继承既有$I^{\mathfrak m}$的结论，而不仅继承一个纯K3子商。
沿用原规范$ I^{\mathfrak m}=2(C_M^{\mathfrak m})_{8,6}$时，有
\begin{equation}\label{eq:actualdispersioncoaction}
 J^{\mathfrak m}=\tfrac14(C_M^{\mathfrak m})_{8,6},\qquad
 \Delta J^{\mathfrak m}
   =\tfrac14\sum_{j=1}^{8}(C_M^{\mathfrak m})_{8,j}
                         \otimes(g_M^{\dR})_{j,6}.
\end{equation}
右端采用经$I^{\mathfrak m}=8J^{\mathfrak m}$识别后的原对象基；
没有声称它是某组任意选择的谱曲面整周期基。
\end{corollary}
\begin{proof}
有限长度范畴及精确实现函子中，一个形式周期的最小对象由
生成de Rham标架的最小子对象，再除去Betti标架湮灭的最大子对象得到。
这一对象只依赖形式周期关系，见\cite{brown}第2.4节。
乘以非零有理数既不改变生成子对象，也不改变湮灭子对象。
定理\ref{thm:actualdispersionformal}是在$\Q$上的实际态射关系，
所以可直接应用，得到最小对象的同构；不需要周期猜想。
相同对象生成同一Tannaka子范畴，因此群、权重及扩张结构一并转移。
最后把既有矩阵系数及其规范余作用除以八，便得所列式子。

这里没有断言整个$H^3$色散载体与秩八$M$同构；
只断言由当前形式和当前链选择出的最小子商同构。
动量式\eqref{eq:actualmomentumformal}使用$k'$，
不能单凭该扩域后的比较证明$\Q$上的算术下降；
本推论依赖的恰是直接定义于$\Q$的环面--色散映射。
\end{proof}

\subsection{任意有限条边：二叉分组给形式周期保持的有限规则}
上述八分支结构来自二叉分组，因而可以在不积分的条件下推广。
以下仍只处理正质量、$D=2$、单位标量传播子、$s<0$的banana，
不把一般图、任意分子或维数正则化放入接受域。

\begin{theorem}[二维banana的任意二叉树色散规则]\label{thm:actualbinarytree}
设$n\geq2$，取以$n$条带标号边为叶子的任意满二叉树$\mathcal T$。
每个叶子赋值$h_i=\mu_i$，每个内部顶点$v$赋变量$h_v$；根变量记为$x$。
在正域中要求$h_v>(\sqrt{h_L}+\sqrt{h_R})^2$，其中$L,R$是$v$的两个孩子。
令
\[
 Q_{\mathcal T}=\prod_{v\ {
m internal}}\lambda(h_v,h_L,h_R),
 \qquad V_{\mathcal T}:y^2=Q_{\mathcal T},
\]
并以$y>0$、每个变量正体积取向定义
\[
 J_{\mathcal T}(s)=\int
    \frac{dx\prod_{v\ne{\rm root}}dh_v}{(x-s)\sqrt{Q_{\mathcal T}}}.
\]
则正参数banana周期
\[
 I_n(s)=\int_{(\mathbb R_{>0})^{n-1}}
  \frac{\bigwedge_{i=1}^{n-1}d\log\alpha_i}
  {(\sum_{i=1}^n\mu_i\alpha_i)(\sum_{i=1}^n\alpha_i^{-1})-s},
 \qquad \alpha_n=1,
\]
满足实际形式周期关系
\begin{equation}\label{eq:actualbinaryperiod}
                I_n^{\mathfrak m}(s)=2^{n-1}J_{\mathcal T}^{\mathfrak m}(s).
\end{equation}
其有理载体映射的一般次数为$2^{n-2}$；正链在不变量底空间上有$2^{n-1}$张分支。
不同二叉树产生的$J_{\mathcal T}^{\mathfrak m}$相等，
因而具有同一个最小标架对象。
\end{theorem}
\begin{proof}
写$u_i=\mu_i\alpha_i$，对树中每个顶点的叶子集定义
\[
 U_v=\sum_{i\in v}u_i,\qquad V_v=\sum_{i\in v}\mu_i/u_i,
 \qquad h_v=U_vV_v.
\]
在每个内部顶点置$z_v=U_L/U_R$，便有
\begin{equation}\label{eq:actualbinaryquadratic}
 h_v=(z_v+1)(h_L/z_v+h_R),\qquad
 h_Rz_v^2+(h_L+h_R-h_v)z_v+h_L=0.
\end{equation}
给定全部$z_v$和根的共同尺度，依次用
$U_L=z_vU_v/(1+z_v)$、$U_R=U_v/(1+z_v)$恢复所有$u_i$。
所以$z_v$是原射影环面的有理坐标，正域与$z_v>0$一一对应。

按先孩子后父亲的顺序排列内部顶点。
在每次合并中，将两个块内的对数比例坐标与$\log(U_L/U_R)$拼接，
对数体积的Jacobian为正负一：块内列的共同尺度项相消，
剩余尺度比列系数为一。由二叶子的情形归纳，
$\eta_n=\varepsilon\bigwedge_v d\log z_v$，其中$\varepsilon=\pm1$只依赖所选次序。
又因$h_v$只依赖其子树的$z$，$z\mapsto h$的Jacobian为三角形，
对角元为
\[
 \delta_v=h_Rz_v-h_L/z_v,\qquad
 \delta_v^2=\lambda(h_v,h_L,h_R).
\]
故取$y=\prod_v\delta_v$，便得到实际有理映射$\Pi_{\mathcal T}$及
$\Pi_{\mathcal T}^*(\bigwedge_v dh_v/y)=\varepsilon\eta_n$。
把根变量移到最前面只改变同一个可追踪的取向符号。

给定严格正不变量，式\eqref{eq:actualbinaryquadratic}在每个内部顶点
恰有两个不同正根，独立给$2^{n-1}$张分支。
它们各自对正密度的贡献都为$J_{\mathcal T}$。
代数映射到全部$h_v$的一般次数也是$2^{n-1}$：二次式给上界，
这个实开集上的不同正原像达到上界。
$Q_{\mathcal T}$对根变量的两个根不同，故双覆盖不可约，
映到$V_{\mathcal T}$的一般次数为$2^{n-2}$。

对分支按$\delta_v$的符号分割，带取向的链推前与$y\mapsto-y$的作用，
完全如\eqref{eq:actualchainpush}：每一块的实际形式周期都等于正$J_{\mathcal T}$。
内部代数面成对抵消。无穷远收敛仍由
$f\geq\mu_{\min}e^{\|t\|_\infty}$证明，只需将$t$改为$n-1$维。
在公共消解上应用态射关系即得\eqref{eq:actualbinaryperiod}。
所有二叉树都比较到同一个$I_n^{\mathfrak m}$，
所以树间相等也是实际形式周期关系，而非相同解析表达的猜测。
\end{proof}

该规则还直接固定物理相空间的常数：逐个插入非根内部不变量，给
\[
 \Phi_n^{(2)}(x)=(2\pi)^{-(n-2)}
     \int\frac{\prod_{v\ne{\rm root}}dh_v}{\sqrt{Q_{\mathcal T}}},\qquad
 I_n(s)=2^{n-1}(2\pi)^{n-2}
     \int_{(\sum_i m_i)^2}^{\infty}\frac{\Phi_n^{(2)}(x)}{x-s}\,dx.
\]
几何构造只需为$n-1$个内部顶点写出二次式及边界；
若显式列出全部实分支则有$2^{n-1}$项。
这两个有限规模不是识别高维动机全部简单因子的复杂度上界。
特别地，$n$增大后载体维数及分支几何增大，
二次根式表示不能推出mixed Tate、MPL或特定数值属于MZV。

本节实际关闭的是正质量二维banana的原参数--正谱色散标架比较，
以及指定固定纤维的最小对象与coaction转移。
原四维含角传播子的五变量推前、全部三圈前向光学切割、
一般图的固定纤维算术及更广物理链比较仍分别保留。
新增理论没有改变任何旧整格、绝对比较或数值证书。


\section{偶数维banana：能量减除、球面推前与实际谱周期}\label{sec:evenenergymom}

上一节将二维原正参数周期与谱载体作了实际比较。
现在处理升维后的两个障碍：空间动量积分发散，以及必须积分掉的二体角纤维。
本节选定一个明确的能量减除代表，证明其与谱双覆盖之间的形式周期关系。
它给任意有限边数、任意正偶数维的有限规则；不要求先求完整积分。
与另一个独立构造的森林代表之间，解析函数相等和形式周期相等仍分开陈述。
经典banana减除色散的背景见\cite{KreimerBanana}，相空间约定沿用\cite{PDGphase}；
下面的常数、实链和代数推前由本节自行固定。

\subsection{明确的有限代表与最少减除次数}

令$n\geq2$、$D=2k$、$d=2k-1$，质量$m_i>0$，$\mu_i=m_i^2$。
全部粒子有标号，不含$1/n!$。置
\[
 \boldsymbol p_n=-\sum_{i<n}\boldsymbol p_i,\qquad
 E_i=\sqrt{|\boldsymbol p_i|^2+\mu_i},\qquad S=\sum_i E_i,
 \qquad a=(k-1)(n-1),\quad \tau=(\sum_i m_i)^2.
\]
这里$a$只在本节的一般$n,D$公式中表示减除阶数，
最后的四线公式将改用$u,v$作为对子不变量以避免混淆。
取$s,s_0<0$，定义
\begin{equation}\label{eq:energymomkernel}
 K_J(x;s,s_0)=\frac{(s-s_0)^J}{(x-s)(x-s_0)^J},\qquad J\geq0,
 \quad K_0(x;s,s_0)=\frac1{x-s}.
\end{equation}
本节的\emph{能量MOM代表}是下述合并后的单个积分：
\begin{equation}\label{eq:energymomdefinition}
 \mathcal E_{2k,n}^{[J]}(s,s_0)=\pi^{-a}
 \int_{\mathbb R^{d(n-1)}}
       \frac{S}{\prod_i E_i}K_J(S^2;s,s_0)
       \prod_{i<n}d^d\boldsymbol p_i .
\end{equation}
减除前的发散积分不被用来定义右端。若$s=s_0$且$J>0$，约定右端为零。
幂$\pi^{-a}$来自每圈$\pi^{-k}$的欧氏约定：
在固定空间动量时，$n-1$次能量卷积各给一个$\pi$，
留下$\pi^{(1-k)(n-1)}$。这个固定空间动量的计算收敛，
但它本身不允许把Taylor操作穿过一个发散的全动量积分。

\begin{proposition}[精确的收敛阶数]\label{prop:energymomconvergence}
对$k\geq2$及$s\ne s_0$，式\eqref{eq:energymomdefinition}绝对收敛
当且仅当$J\geq a$。$k=1$时$J=0$已经收敛。
这些积分在$s\in\mathbb C\setminus[\tau,\infty)$上给全纯函数；
在$s_0$处，$J>0$时前$J$个Taylor系数为零。
\end{proposition}
\begin{proof}
采用吸收全部动量守恒$(2\pi)^D\delta^{(D)}$的相空间约定，
在总空间动量为零的参考系有
\[
 \Phi_n^{(D)}(x)=
 \frac{1}{2^n(2\pi)^{n(D-1)-D}}
 \int\frac{\prod_{i<n}d^d\boldsymbol p_i}{\prod_i E_i}
                \delta(\sqrt x-S).
\]
因$\int dx\,\delta(\sqrt x-S)K(x)=2SK(S^2)$，得到正测度恒等式
\begin{equation}\label{eq:energyphasespectral}
 \mathcal E_{2k,n}^{[J]}
 =\pi^{-a}2^{n-1}(2\pi)^{n(D-1)-D}
       \int_\tau^\infty\Phi_n^{(D)}(x)K_J(x;s,s_0)\,dx.
\end{equation}
暂可将等式两边同时截断于$S^2\leq R$，所以没有使用尚未证明的收敛。

二体积分直接给
\[
 \Phi_2^{(2k)}(h;u,v)
   =c_{2k}\frac{\lambda(h,u,v)^{k-3/2}}{h^{k-1}},\qquad
 c_{2k}=\frac{\operatorname{vol}(S^{2k-2})}
                   {2^{2k-1}(2\pi)^{2k-2}}.
\]
在物理域$0\leq\lambda(h,u,v)\leq h^2$。
任意二叉树分解含$n-1$个二体因子及$n-2$个内部不变量积分；
每个内部不变量至多为$x$，故$k\geq2$时
\[
                 \Phi_n^{(2k)}(x)\leq Cx^{(k-2)(n-1)+n-2}
                                      =Cx^{a-1}.
\]
下界也有同一幂次。令非根内部变量为$h_v=xr_v$，
在满足全部严格二体不等式的一个固定紧子域上积分。
这样的子域可从根向叶递归选取足够小的正$r_v$；
叶子质量比$\mu_i/x$趋于零，所有内部Källén式在该子域上统一正。
缩放全部二体因子及$n-2$个微分，便得
$\Phi_n^{(2k)}(x)\geq c x^{a-1}$，其中$c>0$，$x$充分大。
$n=2$时不变量积分为空，同一结论由二体公式直接成立。

因此无穷远的绝对被积函数与$x^{a-J-2}$双向可比，
正好在$J\geq a$时可积。有限能量范围的质量壳紧致，
正质量保证原密度无极点，阈值端点不产生发散。
$k=1$时，上一节的质量壳坐标同样适用于任意$n$：
取$z_i=E_i+p_i$及$\alpha_i=\mu_n z_i/(\mu_i z_n)$，
则$f=S^2$，而对数Jacobian矩阵为
$\operatorname{diag}(E_1^{-1},\ldots,E_{n-1}^{-1})+E_n^{-1}\mathbf1\mathbf1^t$，
其行列式为$S/\prod_iE_i$。
故无减除质量壳积分等于上一节的正环面周期，由对数坐标估计收敛；
附加减除核相对于该核有界。
最后，在割平面的任一紧子集上，$|x-s|^{-1}\leq C/(1+x)$，
同一可积控制允许全纯参数积分。
核中显式的$(s-s_0)^J$给所声称的零喷射。
\end{proof}

\subsection{角体积的Lefschetz因子不是数值替换}

设$K\subset\mathbb R$是包含质量、参考点及所取代数运动学点的数域，
$K'=K(i)$，$\mathbb L^{\mathfrak m}$的周期为$2\pi i$。
形式周期陈述先在$K'$上作出。变量形式的等式由同一代数对应在光滑开集上给出。
记$Q_{2r}:\sum_{j=1}^{2r+1}z_j^2=1$的实球面为$\Sigma_{2r}$，
取向使标准代数体积形式$\vartheta_{2r}$在实球面上为正。
$r=0$时取$\vartheta_0=z_1$、$\Sigma_0=[1]-[-1]$，
所以两个定向点各给正贡献，合为二。

\begin{lemma}[实际球面周期]\label{lem:actualsphereperiod}
该带链形式周期满足
\begin{equation}\label{eq:actualsphereperiod}
 [Q_{2r},\vartheta_{2r},\Sigma_{2r}]^{\mathfrak m}
 =\frac{2^{r+1}r!}{(2r)!}\,i^{-r}(\mathbb L^{\mathfrak m})^r.
\end{equation}
等式由代数变量替换、闭圆留数及相对端点关系给出。
\end{lemma}
\begin{proof}
把前两个坐标写为$\sqrt t(\cos\theta,\sin\theta)$，
其余坐标写为$\sqrt{1-t}\,n$，$n\in\Sigma_{2r-2}$。
在相应二次覆盖上这是真实代数映射；$0<t<1$上实链一一参数化，
两个退化端面在原球面上按相对边界处理。
体积形式的拉回为
\[
 \frac12(1-t)^{r-3/2}dt\wedge d\theta\wedge\vartheta_{2r-2}.
\]
置$v=\sqrt{1-t}>0$，区间因子的形式积分变成
$\int_0^1v^{2r-2}dv=1/(2r-1)$，由代数原函数与端点关系成立。
单位圆上令$z=x+iy$，则$d\theta=dz/(iz)$，
其逆时针链给$\mathbb L^{\mathfrak m}/i$。
由实际乘积与态射关系，球面周期$V_r$满足
$V_r=\mathbb L^{\mathfrak m}V_{r-1}/(i(2r-1))$。
从$V_0=2$归纳即得公式。端面可在映射图的共同消解上处理；
证明没有把一个已知的球面体积数值代入形式周期代数。
\end{proof}

\subsection{同一谱双覆盖上的实际推前}

取有标号叶子的任意满二叉树$\mathcal T$，内部顶点集合记为$\mathcal I$。
置$h_i=\mu_i$、$h_{\rm root}=x$，并定义
\[
 Q_{\mathcal T}(h)=\prod_{v\in\mathcal I}\lambda(h_v,h_L,h_R),\qquad
 V_{\mathcal T}:\ y^2=Q_{\mathcal T}(h).
\]
正谱链$\Xi_{\mathcal T,+}$由
$\sqrt{h_v}>\sqrt{h_L}+\sqrt{h_R}$及$y>0$确定，
取$dx\prod_{v\ne{\rm root}}dh_v$的正体积取向。
质量壳积分去掉$\pi^{-a}$的形式周期记为$H_{2k,n}^{[J],\mathfrak m}$，
并定义其Tate归一化
\[
 \mathcal E_{2k,n}^{[J],\mathfrak m}
    =(2i)^a(\mathbb L^{\mathfrak m})^{-a}H_{2k,n}^{[J],\mathfrak m}.
\]
这允许Tate逆扭转；不声称规范化对象仍属于未局部化的有效范畴。

\begin{theorem}[偶数维的实际能量--谱比较]\label{thm:actualevenenergy}
在命题\ref{prop:energymomconvergence}的收敛范围内，
上述两个实际带标架周期在$K'$上满足
\begin{align}\label{eq:actualevenenergy}
 \mathcal E_{2k,n}^{[J],\mathfrak m}
 &=\left(\frac{2(k-1)!}{(2k-2)!}\right)^{n-1}
       [V_{\mathcal T},\omega_{k,J},\Xi_{\mathcal T,+}]^{\mathfrak m},\\
 \omega_{k,J}
 &=K_J(x;s,s_0)
    \frac{Q_{\mathcal T}^{k-1}}{\prod_{v\in\mathcal I}h_v^{k-1}}
    \frac{dx\prod_{v\ne{\rm root}}dh_v}{y}.\nonumber
\end{align}
右端的相对对须包含所有端面、极点及无穷远的适当变换。
不同二叉树给相同形式周期，因而其最小带标架对象在$K'$上同构。
本命题只识别该周期的最小对象，不识别整个环境上同调。
\end{theorem}
\begin{proof}
先构造实际映射。在每个内部顶点的静止系中写
\[
 E_L^*=\frac{h_v+h_L-h_R}{2\sqrt{h_v}},\qquad
 p_v^*=\frac{\sqrt{\lambda(h_v,h_L,h_R)}}{2\sqrt{h_v}},
 \qquad \boldsymbol p_L^*=p_v^* n_v,\quad n_v\in\Sigma_{2k-2}.
\]
右孩子空间动量取相反号。把所有$\sqrt{h_v}$与Källén平方根加入一个
有限代数覆盖$W$，并在实链上各取正根。
从根的静止系递归使用无旋转Lorentz提升，恢复每个叶子的时空动量。
例如对未来总动量$P=(P^0,\boldsymbol P)$、$h=P^2$，提升公式为
\[
 p^0=\frac{P^0 E^*+\boldsymbol P\cdot\boldsymbol p^*}{\sqrt h},\qquad
 \boldsymbol p=\boldsymbol p^*+
 \boldsymbol P\left(\frac{E^*}{\sqrt h}
       +\frac{\boldsymbol P\cdot\boldsymbol p^*}{\sqrt h(P^0+\sqrt h)}\right).
\]
分母在正链内部非零。这给从
$W\times\prod_{v\in\mathcal I}Q_{2k-2}$到原质量壳的代数映射。
反向先取各子树的总动量、平方及其正根，再用逆提升恢复$n_v$；
所以正链内部一一对应。这里没有除以共同旋转群，
所有$n-1$个球面方向都被积分一次。

计算微分的Jacobian可在每个顶点的静止系完成。
若$h=(E_L+E_R)^2$、$p=|\boldsymbol p_L|$，则
\[
 \frac{dh}{dp}=\frac{2hp}{E_LE_R},\qquad
 \frac{E_L+E_R}{E_LE_R}p^{d-1}dp
 =\frac{\lambda(h,h_L,h_R)^{k-3/2}}
             {2^{2k-2}h^{k-1}}\,dh.
\]
递归合并两个孩子时，还须以同一个总质量壳密度匹配中间能量。
等价地，插入总动量恒等式并对其正能量留数，再插入$dh$；
这给每个非根内部顶点的$dh/(2\pi)$以及标准二体密度。
故完整未规范化质量壳形式的推前为
\begin{equation}\label{eq:angularrawconstant}
 \left(\frac1{2^{2k-2}}\right)^{n-1}
 \left(\prod_{v\in\mathcal I}\vartheta_{2k-2}(n_v)\right)
 K_J(x;s,s_0)
 \prod_{v\in\mathcal I}
 \frac{\lambda(h_v,h_L,h_R)^{k-3/2}}{h_v^{k-1}}
 dx\prod_{v\ne{\rm root}}dh_v .
\end{equation}
这里的$2\pi$是说明相空间约定的中间记号，最终Jacobian没有引入超越常数。
也可由\eqref{eq:energyphasespectral}及
$\Phi_n=(2\pi)^{-(n-2)}\int\prod_v\Phi_2\prod dh_v$
逐个核对：所有非角常数恰为$2^{-(2k-2)(n-1)}$。
这些插入是质量壳的代数消元及正分支Jacobian；
$\delta$记法不被作为额外的动机运算。

引理\ref{lem:actualsphereperiod}将每个球面链的实际形式周期换成
$2^{k}(k-1)!i^{-(k-1)}(\mathbb L^{\mathfrak m})^{k-1}/(2k-2)!$。
乘以定义中的$(2i)^a(\mathbb L^{\mathfrak m})^{-a}$后，
全部角纤维Tate因子消去，留下公式中的有理常数。

最后，在$W$上令$y=\prod_v\sqrt{\lambda(h_v,h_L,h_R)}$，
得到实际有限映射$W\to V_{\mathcal T}$。
剩余微分正是$\omega_{k,J}$的拉回；实链选定的正根分支
向$\Xi_{\mathcal T,+}$一一推前。因此有限覆盖的总代数次数不产生额外倍数。
辅助$\sqrt{h_v}$在最终微分中已消失，各个Källén根只剩其乘积$y$。
这是显式映射的链与形式关系，不是凭公式中根式的消失猜测下降。

所有映射都可取图闭包并同时消解。
球面纤维闭合，阈值处塌缩，内部分割面以相反取向抵消；
收敛命题保证无穷远给合法相对链。
因此上述变量替换、球面乘积及推前是实际形式周期关系。
不同树均比较到原来的同一个$H_{2k,n}^{[J],\mathfrak m}$，
最小对象结论由带标架态射关系的函子性得到\cite{brown}。
我们只作$K'$上的结论，不由此补造任意$K$上的算术下降。
\end{proof}

\subsection{接到物理重整化，并保留形式比较的边界}

\begin{proposition}[自由复合源的物理MOM实现]\label{prop:energyphysicalmom}
令$k\geq2$，$\phi_i$是$n$种互异的自由实标量场，
$\mathcal O=:\!\prod_i\phi_i\!:$。
取使欧氏banana每圈度量为$d^{2k}\ell/\pi^k$的整体归一化。
其两点函数在$s_0<0$处实行$a$次Lorentz不变MOM减除后，
解析函数等于$\mathcal E_{2k,n}^{[a]}$。
在相同短距离次数约束下，任意两个局部延拓之间的差为
次数至多$a-1$的$s$多项式；相同MOM喷射唯一固定该差。
\end{proposition}
\begin{proof}
分离点的Wick二点函数恰为$n$条自由传播子的乘积，
互异粒子使其权重为一，真空谱只含各一种粒子的$n$体通道。
在总静止系逐次积分欧氏能量得到
$\pi^{n-1}S/(\prod E_i(S^2-s))$。
在谱变量上施加$1-T_{s_0}^{a-1}$，恰得到$K_a$；
收敛命题保证这是一个无需额外正则化的全纯读出。
固定空间动量的能量卷积和正谱定义同时确定其归一化。

位置空间乘积的短距离缩放次数为$n(2k-2)$。
局部延拓的奇异阶为
\[
             n(2k-2)-2k=2(a-1).
\]
两个满足此次数界的延拓之差支撑于原点，
故为总导数次数至多$2(a-1)$的有限个$\delta$导数的线性组合。
Fourier变换后是动量多项式；Lorentz标量性将它限制为$s$的
次数至多$a-1$的多项式。分散表示给同一个正谱的延拓，
在$s_0$处的$a$个零喷射使这一个局部多项式唯一。
由此得到所声称的解析函数比较。
$J>a$也可使用，但相应增加有限局部喷射，属于另行声明的过减除方案。
\end{proof}

此命题并未证明任意两个预先指定的Nori森林代表形式相等。
形式周期的数值评价是否单射没有被假定。
若要把另一套BPHZ森林对象识别到本节对象，仍须构造实际比较对应，
或者给每一项有限喷射的形式周期恒等式。
经典解析重整化唯一性只解决函数比较；
定理\ref{thm:actualevenenergy}实际解决的是已经明确选定的
能量MOM质量壳代表与谱双覆盖的标架比较。

\subsection{四维四线结果及对统一结构表的含义}

\begin{corollary}[四维的单位比较系数与K3可见性]\label{cor:energyfourline}
对$n=4,k=2,J=3$，取两个对子不变量$u,v$，置
\[
 A(u)=\lambda(u,\mu_1,\mu_2),\quad
 B(v)=\lambda(v,\mu_3,\mu_4),\quad
 y^2=A(u)B(v)\lambda(x,u,v).
\]
在上述$K'$范畴中有实际形式周期等式
\begin{equation}\label{eq:energyfourline}
 \mathcal E_{4,4}^{[3],\mathfrak m}
 =\left[\int_{\Xi_+}
   \frac{(s-s_0)^3 y\,dx\,du\,dv}
        {(x-s)(x-s_0)^3xuv}\right]^{\mathfrak m}.
\end{equation}
因此既有谱载体的上界、相对留数及实际管变分，可转移到这个完整有限能量代表。
在原指定质量$(1,2,4,8)$的光滑运动学族上，
其无限单值简单组成恰保留一份既有秩六K3核。
这个结论不声称其固定纤维最小Nori对象为二维的秩八对象。
\end{corollary}
\begin{proof}
定理\ref{thm:actualevenenergy}的常数在$k=2$时为
$(2\cdot1!/2!)^3=1$，而$Q/y=y$，遂得\eqref{eq:energyfourline}。
这是第\ref{sec:dispersioncarrier}节无交换极点的三维色散对，
区别仅是已明确去除的非零整体规范因子。
其Gysin上界只给一份运动学K3；
第\ref{sec:evendimk3}节正常阈值的非零管变分给匹配下界。
实际标架比较把上下界一起转移到本节代表。
减除只改变有限局部项，不消去正常阈值变分。
纯核的这个函数结论没有识别完整边界扩张，
也没有给出与二维减除阶数、微分和权重都不同的固定对象同构。
\end{proof}

一般$n,k$时，谱载体维数恒为$n-1$，升维只将微分乘以
$Q_{\mathcal T}^{k-1}/\prod h_v^{k-1}$并要求至少$a$次减除。
故由同一边表和一棵二叉树即可写出实际比较、必要减除数和候选载体，
无需完成原$d(n-1)$维积分。
然而$h_v=0$等极点会改变相对边界，微分也会改变可见子商；
相同双覆盖不保证相同混合对象，更不等于所有高维纯核已经分类。
本节没有解决一般含角传播子的五变量物理推前、任意图最小算术分类，
或独立森林代表的全部形式比较；这些仍属于完整总目标。



\section{规则组合的正确性与总目标的关键缺口}
\subsection{有限推导可以保证什么}
每条规则需提交适用谓词、实际模型、形式/链比较、纯核与边界、
以及输出强度。一个节点只能给出：
精确最小对象、几何上界、受保护下界、条件结论或当前规则不适用。
拒绝图不能输出“非MPL”；高阶方程不能直接输出“K3”。

\begin{theorem}[有限组合的可靠性]
若有限推导树的叶节点均有实际输入，每个内部步骤满足其证明前提，
则最终上界、下界与精确输出成立。
简单重数可由相同的上下界确定；完整最小混合对象还需粘合及标架识别。
\end{theorem}
\begin{proof}
逐节点归纳：对象运算用命题\ref{prop:upper}；
割用命题\ref{prop:lower}；
有限方案用定理\ref{thm:ren}；
标架最小化用实际子商定义。
始终保留基域、参数层、链条件，便无未经比较跨到物理对象的步骤。
精确性最后由重数上下界及已证粘合给出。
\end{proof}
这是可靠性而非任意图可计算性的证明。
圈星识别/分支模块及组合模块有有限终止程序；
一般实际子商、扩张非零和物理链仍需有效输入。
有限碰撞分层也不代替可能无限的算术特殊点分类。

\subsection{后续优先闭合的三处主接口}
\begin{enumerate}
\item \textbf{结构预测入口。}
把原覆盖/消元规则接到实际几何，形成接受图类的纯核、边界、分子选择表。
保留一般图方向，不能以banana一例代替。
\item \textbf{割到完整物理标架。}
优先证明更多接受类的链生成、下sector粘合和实际收缩/森林比较。
本篇已给有理树权的实际相空间载体、二体有理角规则与三体椭圆下界；
指定两圈弹性截面另有真实边值、交叉角卷积及唯一椭圆简单因子；
等变Kummer规则将其提升为有限单变量投影类的角色上界和变分精确化；
第\ref{sec:massiveforestperiods}节已给四维正质量欧氏森林的实际周期入口、有限喷射方案比较，
以及消去整个局部多项式空间后的低圈背景；
第\ref{sec:forestparametric}--\ref{sec:sectorboundary}节进一步给合并密度的参数表达，
以扇区共同边界证明该正质量欧氏接受域的任意有限圈完整标架比较；
无质量/IR、一般物理圈边值及推前后最小混合扩张仍须实际输入。
这些决定几何下界是否真正进入可观测量；
一旦给出，本篇方案保持定理立即适用。
\item \textbf{连贯全文与验收。}
按结构表、三个主定理和核心实例组织全文。
本篇另给完整森林方向到临界谱的有限检验，
把固定点漂移、算符混合和非线性共振分开；
实际输入、单谱及正则方案条件不能省略。
局部RG、证书及物理基准各承担对应证明责任，不无限扩展支线。
\end{enumerate}
热质量比较、基准窄误差及其它原任务仍在总目标内；
除非阻断这些主接口，不作为下一批默认工作。
已有Q6、绝对比较、固定群和整格证据保留，不重做。
本文加强了总理论连接，没有宣称任意图分类已经完成。

\section{无质量图的MZV肯定规则：串并联、减除与同拓扑退化}
\label{sec:masslessmzv}

本节补分类表的肯定方向：在完整积分以前，从图的分解证书预报所需常数。
接受域是无质量、非零欧氏外动量的标量两点图；物理对象是维数正规化的
Laurent系数、在明确条件下的局部反项，以及无质量banana的完整重整化函数。
这里不假定每个图都属于同一个可重整化拉格朗日量；图权、耦合和外腿归一化另行乘回。
一圈卷积和单尺度子图替换是经典规则\cite[式(1.2.1)、\S5.2.1]{PanzerMZV}。
以下自给推导它们对分类、物理减除及特殊质量层的具体结论，不主张经典公式的新颖性。

\subsection{物理输入与一圈卷积}
令$D=4-2\varepsilon$，$q^2=x>0$，每圈测度为$d^Dk/\pi^{D/2}$。
边的传播子为$(k_e^2)^{-a_e}$，初始$a_e\in\mathbb Z_{>0}$。
所有外动量只在两个指定端点注入；其余顶点总外动量为零。
若图有$h$圈，使用明确的Euler归一化
\[
 \widehat F_G(x,\varepsilon)
 =e^{h\gamma_E\varepsilon}\mu^{2h\varepsilon}F_G(x,\varepsilon),
 \qquad L=\log(x/\mu^2),\qquad
 \mathcal Z_1=\mathbb Q[\zeta(2),\zeta(3),\ldots]\subseteq\mathcal Z.
\]
此处$\mu>0$是同一减除尺度。更换测度会改变显式$\pi$与$4\pi$因子；
没有先去掉这些规范就不能声称系数属于普通MZV代数。
维数不为整数的整个积分不是一个已构造的有限维Nori对象；本节的算术结论逐个展开系数陈述。

\begin{lemma}[实际卷积核]\label{lem:masslesskernel}
在$0<\operatorname{Re} A,\operatorname{Re} B<D/2$及
$\operatorname{Re}(A+B)>D/2$的实正整数维数收敛域，
\begin{align}
 \int\frac{d^Dk}{\pi^{D/2}(k^2)^A((q-k)^2)^B}
 &=x^{D/2-A-B}\,g_D(A,B),\\
 g_D(A,B)&=
 \frac{\Gamma(D/2-A)\Gamma(D/2-B)\Gamma(A+B-D/2)}
 {\Gamma(A)\Gamma(B)\Gamma(D-A-B)}.\label{eq:masslesskernel}
\end{align}
两边以Schwinger参数形式亚纯延拓到维数和指数的一般复值。
\end{lemma}
\begin{proof}
对两条边引入正参数$t,u$，完成平方后动量高斯积分给
$(t+u)^{-D/2}\exp[-tu x/(t+u)]$。
置$t=\lambda v,u=\lambda(1-v)$，$0<v<1$。
先积$\lambda$得到$\Gamma(A+B-D/2)x^{D/2-A-B}$，
再积$v$得到$B(D/2-B,D/2-A)$，连同原来的
$1/(\Gamma(A)\Gamma(B))$即为所示公式。
收敛域中Tonelli或绝对Fubini合法；复维数版本由该参数公式定义并唯一亚纯延拓，
不是在非整数维度上假想一个通常的Lebesgue空间。
\end{proof}

\subsection{边表证书与全阶常数类}
双端串并联证书是一棵二叉树。叶是有标号边；每个结点的两个子图除所声明端点外不相交。
串联$S$仅把前者末端与后者首端粘合；并联$P$把两对端点分别粘合。
最终叶集须恰为输入边集，端点粘合须恢复原边表，并逐顶点核对无内部动量注入。
这里的并联包含任意边数banana，串联中允许二价顶点。
这是一项充分的接受条件，不是所有MZV图的必要条件。

\begin{theorem}[串并联证书预报全阶单$\zeta$代数]\label{thm:masslesssp}
给定上述证书，设$E$为边数，$h$为圈数，$A_0=\sum_e a_e-2h$。
不需要展开任一多圈积分即可断定
\[
 \widehat F_G(x,\varepsilon)
 =x^{-A_0}e^{-h\varepsilon L}\widehat C_G(\varepsilon),
 \qquad \widehat C_G\in\mathcal Z_1((\varepsilon)).
 \label{eq:masslessspclass}
\]
更精确地，证书给有限有理函数$R_G\in\mathbb Q(\varepsilon)$和有理数$c_k$使
\[
 \widehat C_G(\varepsilon)=R_G(\varepsilon)
 \exp\left(\sum_{k\ge2}\frac{\zeta(k)}k c_k\varepsilon^k\right).
 \label{eq:masslessgammaexp}
\]
若$R_G$在零点的极点阶为$p$，则第$\varepsilon^j$系数可表达为
权重至多$j+p$的单$\zeta$乘积；计$L$的权重为1，同一界适用于去掉$x^{-A_0}$后的系数。
这里是表示的权重上界，不假定$\zeta$值间的代数独立性。
\end{theorem}
\begin{proof}
每个子图记$F=Cx^{-A}$，从叶$(h,A,C)=(0,a_e,1)$开始。
串联时动量相同、环变量互相独立，因此
\[
 h=h_1+h_2,\quad A=A_1+A_2,\quad C=C_1C_2.
\]
并联时分流动量为$k,q-k$，额外增加一次真实卷积，因此
\[
 h=h_1+h_2+1,\quad A=A_1+A_2-D/2,\quad
 C=C_1C_2g_D(A_1,A_2).
\]
从而每个子图均有$A=a+h\varepsilon$、$a\in\mathbb Z$。
恒等式先在各叶指数独立的收敛域证明：指定任一子树的有效指数
$0<A<D/2$，串联可把它分成两个正数；并联可选择两个小于$D/2$的正数，
其和为$A+D/2$。反向归纳给所有叶的非空同时收敛域。
再亚纯延拓至原正整数指数和$D=4-2\varepsilon$，无需非法交换发散积分。

用$\Gamma(z+1)=z\Gamma(z)$把所有Gamma因子移至$\Gamma(1+c\varepsilon)$。
整数位移只产生有理函数。非正整数的分子Gamma具有非零斜率：
并联结点的三条分子斜率为$-(h_1+1),-(h_2+1),h_1+h_2+1$；
所以不会遇到沿整条$\varepsilon$直线恒为极点的分子。
斜率为零的分母Gamma仅来自无圈正指数子树，也不形成此问题。
各并联结点的分子斜率和减去分母斜率和恰为1。
应用经典展开\cite[5.7.3]{DLMFGamma}
\[
 \log\Gamma(1+z)=-\gamma_E z+
 \sum_{k\ge2}\frac{(-1)^k\zeta(k)}k z^k,
\]
恰有$h$个并联结点，故Euler常数项为$-h\gamma_E\varepsilon$，被规定的规范消去。
剩下$c_k=\sum_\nu\sigma_\nu(-1)^kc_\nu^k$，其中
$\sigma_\nu=\pm1$记录分子或分母，$c_\nu\in\mathbb Q$为斜率。
指数展开的$\varepsilon^r$项由总权重$r$的$\zeta$乘积组成；
乘$R_G$及$e^{-h\varepsilon L}$给所述界。
\end{proof}

这个证明也解释为什么需要规范：若保留$e^{-h\gamma_E\varepsilon}$，
展开就含Euler常数；这不能被省略后写成MZV定理。
去掉$x^{-A_0}$后，当$x=\mu^2$时系数属于$\mathcal Z_1$；一般$x$时是
$\mathcal Z_1[L]$，不能把任意质量单位或运动学对数自动叫作MZV常数。
有质量、非整数中心指数或内部动量注入都可能破坏这个接受域。

\begin{proposition}[有限证书的核查成本]\label{prop:masslesscertificate}
给定边表与上述二叉证书，可用不超过$O(E^2)$次集合/端点核对以及
$O(E)$个指数加减操作核实接受条件并生成至多$6h$个Gamma因子的表达树。
无需生成生成树多项式、求积或数值整数关系。
\end{proposition}
\begin{proof}
满二叉树有$2E-1$个结点。逐结点比较子图的带标号边集与粘合端点，
用朴素至多$E$长度列表即可在$O(E^2)$内验证；同时检查外动量标记。
三条递归量$(h,A,C)$每个结点只需常数次操作，每个$P$结点追加六个Gamma记录。
这些是组合操作及符号表达树的界；若把任意大整数位移完全展开成线性因子，
或列举所有$\varepsilon$阶系数，其位复杂度另计，不被这个界隐瞒。
\end{proof}

\subsection{重整化与局部RG的接受条件}
\begin{proposition}[森林数据闭合时的MZV保持]\label{prop:masslessforestclosure}
固定一个有限阶的实际MS森林计算。假设无未分离的IR极点，局部反项规则已成立，
且每个需要的子图和带局部插入的余图在去量纲后，都有上一定理的证书，
或是这类表达的有限有理组合和明确为零的无尺度积分。
则该计算的所有MS反项、有限部分、以及由它们作有限有理代数操作与尺度导数所得的
RG图贡献，系数属于$\mathcal Z_1$；保留尺度时属于$\mathcal Z_1[L]$。
这不声称任一串并联图自动满足所有森林接受条件。
\end{proposition}
\begin{proof}
有限圈阶按子图数递归，用
\[
 C_-(G)=-K\!\left(F_G+
 \sum_{\emptyset\ne\gamma\subsetneq G}C_-(\gamma)F_{G/\gamma}\right),
 \qquad R(G)=(1-K)\overline R(G),
\]
其中求和和局部插入按给定实际森林规则，$K$取负次幂。
归纳假设与接受证书使括号中的各项都属于相应Laurent环。
有限加法、乘法和$K$保持该环。若需要除以常数项为非零有理数的形式级数，
其逆也逐阶留在同一环。$\partial_L$同样保持此环；
局部性另外保证反项的局部系数不含$L$。
没有将包含IR污染的裸极点解释为UV反项，也没有由单图贡献推出完整理论的$\beta$函数。
若有限方案另添任意常数，该结论只在所添常数仍属所声明代数时保持。
\end{proof}

\subsection{实际无质量banana：全阶规则与完整物理读出}
令$B_n$为两点之间的$n\ge2$条单位无质量边，$h=n-1$。
它可以来自$n$种互异自由标量的正规序乘积源两点函数；
剥去Wick归一化和测度常数后，正是这里的有标号标量图。
不加入相同粒子的$n!$，也不把任意$n$的复合源说成同一个可重整化基本场相互作用。

\begin{theorem}[无质量banana的闭式分类与MS有限部分]\label{thm:masslessbanana}
在前述规范下，
\begin{align}
 F_{B_n}(x,\varepsilon)
 &=x^{n-2-(n-1)\varepsilon}
 \frac{\Gamma(1-\varepsilon)^n
       \Gamma(2-n+(n-1)\varepsilon)}{\Gamma(n-n\varepsilon)},
 \label{eq:masslessbananaexact}\\
 \widehat C_{B_n}(\varepsilon)
 &=R_n(\varepsilon)\exp\left[
 \sum_{k\ge2}\frac{\zeta(k)}k
 \bigl(n+(-1)^k(n-1)^k-n^k\bigr)\varepsilon^k\right],\\
 R_n(\varepsilon)&=
 \left[\prod_{j=0}^{n-2}((n-1)\varepsilon-j)
       \prod_{j=1}^{n-1}(j-n\varepsilon)\right]^{-1}.
\end{align}
该实际图在$x>0$无IR发散，所有非空真UV森林的反项插入积分在无质量DR中为零。
因此其MS反项和有限部分分别为
\begin{align}
 C_-(B_n)&=-\frac{r_n}{\varepsilon}x^{n-2},
 &r_n&=\frac{(-1)^{n-2}}{((n-1)!)^2},\\
 F^{\mathrm{MS}}_{B_n}(x)&=r_nx^{n-2}\bigl(b_n-(n-1)L\bigr),
 &b_n&=(n-1)H_{n-2}+nH_{n-1},\label{eq:masslessbananaMS}
\end{align}
其中$H_0=0$、$H_m=\sum_{j=1}^m1/j$。
特别地，去量纲后的减除点值是有理数；所有更高$\varepsilon$系数均在单$\zeta$代数中。
\end{theorem}
\begin{proof}
更一般维数的$n$线公式为
\[
 x^{(n-1)D/2-n}
 \frac{\Gamma(D/2-1)^n\Gamma(n-(n-1)D/2)}
 {\Gamma(n(D/2-1))}.
\]
$n=2$是一圈引理。把$n-1$线子图和一条边并联，
将有效指数$n-1-(n-2)D/2$与1代入卷积核，Gamma因子望远镜消去，得到归纳步。
整数维收敛域的解析指数延拓同前，故所得是实际图的亚纯表达。
把$D=4-2\varepsilon$代入，再把Gamma平移至1，给$R_n$及指数中的系数。

IR方面，非零$q$使全部$n$条线不能同时软化。
任一真子集的$r<n$个软线动量可作为独立的$r$个四维局部坐标；
其奇性至多为$\prod_{j=1}^r(k_j^2)^{-1}$，局部可积。
这同时覆盖移位后某条线变软的层；四维附近仍成立。
因此这里的零点极点是UV问题，不能套用有IR的无尺度原图。

每个真圈子图包含两端顶点和$r<n$条平行边。
把其局部反项缩成一点后，至少一条余边成为无质量蝌蚪；
局部反项只提供有限阶动量多项式。
展开这些多项式后，余图是独立无尺度张量蝌蚪积分的有限和，DR中均为零。
嵌套森林按最外层真子图同样判断；互相重叠的子图不能当作可同时收缩的森林。
所以$\overline R(B_n)=\widehat F_{B_n}$，无需把裸图一般地等同于已减除图。

$R_n$的唯一零点极点来自$j=0$，其留数为$r_n$。
Gamma指数从$\varepsilon^2$开始，而
\[
 R_n(\varepsilon)=\frac{r_n}{\varepsilon}
                 (1+b_n\varepsilon+O(\varepsilon^2)),\qquad
 e^{-h\varepsilon L}=1-h\varepsilon L+O(\varepsilon^2).
\]
二者相乘即给反项及有限部分，$b_n$由有限线性因子的对数导数求出。
这是确切全$n$公式，不依赖数值拟合。
\end{proof}

\begin{corollary}[MOM、正谱与四线显式物理函数]\label{cor:masslessphysical}
令$T_{\mu^2}^{n-2}$取$x=\mu^2$处至$n-2$阶Taylor多项式。
同一源函数满足这些局部MOM条件的代表为
\[
 F^{\mathrm{MOM}}_{B_n}(x)
 =-(n-1)r_n(1-T_{\mu^2}^{n-2})
       \bigl[x^{n-2}\log(x/\mu^2)\bigr].
\]
取$\mathcal A(s\pm i0)=F_{B_n}^{\mathrm{MS}}(-s\mp i0)$，则$s>0$时
\[
 \frac{\operatorname{Disc}\mathcal A(s)}{2\pi i}
 =\frac{s^{n-2}}{(n-1)!(n-2)!}>0.
 \label{eq:masslessspectral}
\]
同一不连续性也适用于上述MOM代表。具体四线有
\begin{align}
 F^{\mathrm{MS}}_{B_4}(x)&=\frac{x^2}{36}
                   \left(\frac{71}{6}-3\log(x/\mu^2)\right),\\
 F^{\mathrm{MOM}}_{B_4}(x)&=-\frac1{12}\left[
 x^2\log(x/\mu^2)-\mu^2(x-\mu^2)-\frac32(x-\mu^2)^2\right],\\
 \frac{\operatorname{Disc}\mathcal A_4(s)}{2\pi i}&=\frac{s^2}{12}.
\end{align}
这些是前述剥去常数的源两点函数，不是未经LSZ和通量处理的散射截面。
\end{corollary}
\begin{proof}
两种局部延拓之差是次数至多$n-2$的多项式；指定Taylor喷射唯一确定它。
MS有限部分中的$b_nx^{n-2}$全部被Taylor减除，得到首式。
对于$s>0$，$\log(-s\mp i0)=\log s\mp i\pi$，
而$(-s\mp i0)^{n-2}=(-1)^{n-2}s^{n-2}$。
代入MS式与$r_n$，符号相消，得到正谱公式。
局部多项式没有割，故MOM同样成立。
最后用$r_4=1/36$、$b_4=71/6$及
$T_{\mu^2}^{2}[x^2\log(x/\mu^2)]=\mu^2(x-\mu^2)+3(x-\mu^2)^2/2$。
\end{proof}

\subsection{同一拓扑为何从K3退化到对数}
\begin{proposition}[四线全零质量层的显式有理化]\label{prop:masslessdegeneration}
固定$s\ne0$，四线谱载体在$m_1=m_2=m_3=m_4=0$时为
\[
 y^2=a^2b^2\lambda(s,a,b),\qquad
 \lambda(s,a,b)=(s-a-b)^2-4ab.
\]
约去重复因子后，在$ab\ne0$的稠密开集上有理坐标
\[
 z=y/(ab),\quad t=\frac{s-a-b-z}{2a},\qquad
 b=\frac{st}{1+t}-at,\quad z=\frac{s}{1+t}-a(1+t).
\]
因此一般有质量的K3载体在此奇异质量层的正规化是有理的；
实际无质量四维函数的对数分类由上面的完整物理公式确定。
\end{proposition}
\begin{proof}
把$t$的定义写成$z=s-a-b-2at$，代入$z^2=(s-a-b)^2-4ab$，
约去非零$4a$即得$(1+t)b=st-at(1+t)$。
这给互逆有理参数化，$t=-1$及$a b=0$只需作为边界保留。
有理化本身不分类全部相对扩张；上一定理和推论另行给完整读出的对数形式。
所以这里没有从环境有理性直接跳到完整周期结论。
\end{proof}

这一结果与分类表K3行保持同一$D=4$、同一四线拓扑，改变的是质量层和已声明的减除输入。
全零质量必须作为M4独立一行，不能从光滑正质量族把K3标签机械外推。
反过来，也不能在正质量积分中直接丢掉质量后保留原答案。
本节独立定义并解析延拓无质量图，不假定质量极限、维数展开及减除三者无条件可交换。
在$D=2$直接全零质量会遇到新的IR问题，因而不在此物理MOM结论内。

表的M3行由\ref{thm:masslesssp}和\ref{prop:masslesscertificate}给全阶肯定证据；
森林物理出口须提交\ref{prop:masslessforestclosure}的附加输入。
M4行由\ref{thm:masslessbanana}--\ref{prop:masslessdegeneration}实际闭合MS、MOM和正谱。
这回答了为什么某些周期能用MZV，以及怎样在积分前给充分判断；
它不提供所有图的必要判据，也不识别每个正规化系数原生的完整最小动机。

\section{原物理标架与实际余作用：旧假设11.2的分项替换}
\label{sec:actualframebridge}
分类表K1必须指向原正参数积分，而不只是一个具有相同微分系统的混合对象。
本节将后续几何构造中的实际箭头集中写出。
旧整合稿的假设11.2确为输入；不能只删除“假设”二字。
替换它需要下述几何来源、扩张比较和求值前的标架等式。
有限Picard、整格及下层非分裂证书沿用各自原范围，不在这里重复计算。

\subsection{已构造的几何输入及其依赖}
取数域$k\subset\mathbb R$、正平方质量$\mu_i$及$s<0$。
记$T=\mathbb G_m^3$，$U_s=\{M\sum_i\alpha_i^{-1}=s\}\subset T$。
原物理对解消为$(P\setminus X,Y\setminus X)$，其有效同调动机为
\[
 E_{\rm env}=H_3^{\rm N}(P\setminus X,Y\setminus X)
 \simeq H^3_{\rm N}(T,U_s)(3).
\]
正链合法性的直接检验是每个Hepp扇区中
$\Phi_s=dt_1dt_2dt_3/(M_pV-sq)$，分母至少为$\mu_{p_0}>0$。
实际相对对偶由正常交叉局部乘积上的开延拓交换态射和Verdier对偶给出。
所用Nori六运算与原普适上同调的识别分别见
\cite[Proposition 2.11]{BridgeIM}及
\cite[Theorem 4.5, Corollaries 4.8,4.20]{BridgeTubach}；
不使用后者猜想性的定理4.22。
Betti保守性只检验已经构造的箭头，不把任意Hodge同构提升为动机箭头。

二十条边界曲线张成的十六维空间有非退化交叉配对。
其逆Gram矩阵给实际Chow投影子；剔除它得到秩六$H_{\rm res}$，
再沿开限制、相对连接和上述对偶嵌入$E_{\rm env}$。
两个核分别由边界类和纯权重四的$H^2(T)$控制，故这一嵌入单射。
正实环面与三个单位圆的乘积恰有一个横截交点；相容定向校准后交数为一，
故原物理链在顶层$\mathbb Q(0)$的坐标为$1$。Tate共轭符号已包含在相对对偶中。
Ayoub--Barbieri-Viale的有效一动机商定理
\cite[Theorem 8.12]{BridgeABV}于是给实际正合列
\[
 0\longrightarrow A=H_{\rm res}(3)\longrightarrow E_{\rm env}
 \longrightarrow K\longrightarrow0,\qquad \operatorname{rk}(A,K)=(6,9).
\]
具体模型、边界及符号表在\cite{BridgeAudit}中定位到原始几何证明。
即使Picard数跳跃，$H_{\rm res}$仍是边界正交因子；
把它称作完整超越核还需另证$\rho=16$。

\begin{lemma}[符号扩张到相对扩张的实际推出]\label{lem:bridgefiber}
设导出态射$\mathcal V\to\mathcal W$满足
$H^0\mathcal V=\mathbb Q(0)$、$H^0\mathcal W=0$，
且$h:\mathbb Q(0)\to\mathcal V$在$H^0$上为恒等。
则从$\operatorname{Fib}(\mathbb Q(0)\to\mathcal W)$到
$\operatorname{Fib}(\mathcal V\to\mathcal W)$的自然态射，
在$H^0$上把核为$H^{-1}\mathcal W$的扩张，
沿
\[
 H^{-1}\mathcal W\twoheadrightarrow
 \operatorname{coker}(H^{-1}\mathcal V\to H^{-1}\mathcal W)
\]
推出成右边的相对扩张。
\end{lemma}
\begin{proof}
两个纤维三角的长正合列都以$\mathbb Q(0)$为商。
自然态射在商上为恒等、在核上为上述满射。
沿此满射推出左边后，两端态射均为恒等，短正合列给中间同构。
纤维的自然态射取在六运算的增强范畴中，避免任意选择三角形之间的箭头。
\end{proof}

将引理应用于$\mathcal V=R\Gamma(T)(3)[3]$、
$\mathcal W=R\Gamma(U_s)(3)[3]$，
$h$为有序坐标单位的杯积。单位圆的乘积校准顶层生成元；
光滑仿射曲面的$H^3$为零。
两次局部化、tame单位及Steinberg关系遂把实际相对扩张送到质量单位。
商去三个环境杯积方向并固定同调余字符为负留数基后，得到
\[
 K=\left[\mathbb Z\longrightarrow\mathbb G_m^8,\quad
 1\longmapsto
 (\mu_1/\mu_4,\mu_2/\mu_4,\mu_3/\mu_4,1,1,1,1,1)\right]_{\mathbb Q}.
\]
这一识别使用完整旗标比较\cite{BridgeAudit}，不仅使用质量对数的数值。
它确定一动机商，不同时证明K3方向下层扩张非零。

\subsection{链提升、形式下降及形式周期恒等}
令$\kappa$为三个质量比在$k^\times\otimes\mathbb Q$中的秩，
$K_{\rm act}$为质量点生成的最小连通子环面的一动机。
取实际拉回$E_{\rm act}=E_{\rm env}\times_KK_{\rm act}$，有
\[
 i:E_{\rm act}\hookrightarrow E_{\rm env},\qquad
 E_{\rm env}/E_{\rm act}\simeq\mathbb Q(1)^{8-\kappa},\qquad
 \operatorname{rk}E_{\rm act}=7+\kappa.
\]
正质量排除了非平凡符号扭转；以上商在实基域上为分裂Tate对象。
定义$M_{\rm phys}=E_{\rm env}^{\vee}$、
$M_{\rm act}=E_{\rm act}^{\vee}$及$r=i^\vee$。

\begin{theorem}[原积分的实际余作用比较]\label{thm:bridgephysical}
正物理链$\sigma$唯一提升为$\sigma_{\rm act}\in(E_{\rm act})_{\rm B}$。
置$\Phi_{\rm act}=r_{\rm dR}[\Phi_s]$，则在同一个Nori形式周期环中
\begin{equation}
 [M_{\rm phys},\sigma,[\Phi_s]]^{\mathfrak m}
 =
 [M_{\rm act},\sigma_{\rm act},\Phi_{\rm act}]^{\mathfrak m}.
 \label{eq:bridgephysical}
\end{equation}
因此两者的规范motivic coaction相同，且
\[
 \Delta I^{\mathfrak m}
 =\sum_{j=1}^{7+\kappa}
 [M_{\rm act},\sigma_{\rm act},e_j]^{\mathfrak m}
 \otimes[M_{\rm act},e_j^\vee,\Phi_{\rm act}]^{\rm dR}
\]
对任意de Rham基$e_j$成立，采用右余作用约定。
\end{theorem}
\begin{proof}
几何复共轭固定正链，而在商$\mathbb Q(1)^{8-\kappa}$上为负恒等。
链的商像同时等于自身与负值，故为零。
Betti实现正合，且$i$单射，给唯一提升$i_{\rm B}\sigma_{\rm act}=\sigma$。
形式沿对偶满射$r$下降。形式周期的态射关系
\[
 [M_{\rm phys},r_{\rm B}^{\vee}\sigma_{\rm act},[\Phi_s]]^{\mathfrak m}
 =
 [M_{\rm act},\sigma_{\rm act},r_{\rm dR}[\Phi_s]]^{\mathfrak m}
\]
就是\eqref{eq:bridgephysical}，参见\cite[\S2.3--2.4]{brown}。
这是求值前等式，未用周期映射单射性；规范余作用保持此关系。
也未选择一个未构造的投影$E_{\rm env}\to E_{\rm act}$。
\end{proof}

\begin{corollary}[指定纤维的最小性及其独立输入]\label{cor:bridgeminimal}
在$m=(1,2,4,8),s=-64,k=\mathbb Q$，
继承指定Picard与下层非分裂证书后，\eqref{eq:bridgephysical}的最小对象秩为八。
\end{corollary}
\begin{proof}
此时$\kappa=1$；$\rho=16$使$A=T(X)(3)$的Hodge实现简单。
质量Kummer扩张非零，下层$0\to A\to W_{-2}E_{\rm act}\to\mathbb Q(1)\to0$
的实Hodge非分裂由既有行列式区间$(877/100,439/50)$给出。
其标记和严格有限输入见\cite{BridgeAudit}，本节不重算该区间。

包含物理链的子对象$R$在$K_{\rm act}$上满射：
顶层非零，而一维Kummer支撑的非分裂性排除真子对象满射顶层。
$R\cap A$只能是$0$或$A$。后一种给$R=E_{\rm act}$；
前一种给$R\simeq K_{\rm act}$，其对偶$F^3=0$，
使物理$F^3$形式与链的配对为零，违反欧氏积分严格正性。
所以链生成全部对象。

若非零子对象$J$被物理形式杀死，简单性及非零全纯留数给$J\cap A=0$。
于是$J$注入$K_{\rm act}$。其非零像或者是$\mathbb Q(1)$，
或者满射顶层而等于$K_{\rm act}$；
两种情况都给下层扩张的截面，与既有非分裂输入矛盾。
故既不能缩小链生成对象，也不能商去形式看不见的子对象。
\end{proof}

规范余作用、最小对象、比较矩阵条目和完整固定群须分别验收。
固定读出$I^{\mathfrak m}=2(C_M^{\mathfrak m})_{8,6}$
仍须使用P8的实际膜、交数及杯积定标；
完整纯群仍须使用端同态域、Frobenius及行列式对象的独立算术证据。
这些都不从定理\ref{thm:bridgephysical}自动推出。
本节闭合的是上述接受域内的原积分标架比较；一般参数的最小性、
特殊化和全局物理支仍按各自证明处理。
非Tate最小对象不推出单个数值$I(-64)$不属于MZV。


\section{物理基准的严格误差预算与独立光学验收}\label{sec:narrowphysics}
在$s/m^2=20$，现有四体严格界仍为
$8.8\times10^{-13}<C_{24}<6.2\times10^{-10}$，
其中$C_{24}=m^2\sigma_{24}^{[g^4]}/g^4$。
普通两种相空间参数化给$C_{24}\simeq5.75125\times10^{-11}$，
尚未认证其展示精度。本节给解析误差预算，不将它写成已算出的窄区间。
严格相空间体积上界$m^{-4}\Phi_4(20m^2)<7.63\times10^{-8}$
沿用已有正权证书\cite{BridgePhysics}。

\subsection{选定点的解析进展：统一角截断}
\subsubsection{真实十图与无量纲化}
采用原物理基准的测度、通量、相同粒子$1/4!$及树振幅规范。
以下把所有动量除以$m$，因而$m=1$、$s=20$；最后所得系数就是$C_{24}$。
固定末态$k_1,\ldots,k_4$的内部形状，令$P=(\sqrt s,\mathbf0)$。
入射方向$n\in S^2$、入射空间动量长度$p=\sqrt{s-4}/2$。
原十个$3|3$树分割各权重为1，分母和为
\[
 S(n)=T-\sum_{|J|=2}\frac1{\alpha_J-B_J\cdot n},
 \quad
 T=\sum_{i=1}^4\frac1{s-2\sqrt s E_i},
\]
\[
 K_J=\sum_{j\in J}k_j,\quad a_J=K_J^2,\quad
 b_J=(P-K_J)^2,\quad
 \alpha_J=\frac{s-a_J-b_J}{2},\quad B_J=2p\,\mathbf K_J.
\]
无量纲振幅满足$m^2\mathcal M_{24}^{(0)}/g^2=-S$。
六个二元子集均计入；$J$与$J^c$是不同树图，
但$\alpha_{J^c}=\alpha_J$、$B_{J^c}=-B_J$。
本项目此前已证明这些真实分母全物理域无极点。
这里给适用于窄认证的更具体统一估计。

\begin{proposition}[20点的收敛比]\label{audit:prop:ratio}
在闭物理相空间上，$\alpha_J\ge4$、$0<T\le1/2$，且
\[
 q_J^2:=\frac{|B_J|^2}{\alpha_J^2}\le\frac49.
\]
\end{proposition}
\begin{proof}
每个二粒子子系统有$a_J,b_J\ge4$且$\sqrt{a_J}+\sqrt{b_J}\le\sqrt{20}$。
故$\alpha_J\ge\sqrt{a_Jb_J}\ge4$，并有$0<20-a_J-b_J\le12$。
利用$\lambda(s,a,b)=(s-a-b)^2-4ab$，
\[
 q_J^2=\left(1-\frac4{20}\right)
 \left(1-\frac{4a_Jb_J}{(20-a_J-b_J)^2}\right)
 \le\frac45\left(1-\frac{64}{144}\right)=\frac49.
\]
另一方面$P-k_i$是三个未来向质量1的粒子的和，
所以$(P-k_i)^2\ge9$，每个类时分母至少为8。
四项相加给$0<T\le1/2$。
\end{proof}

取互补子集的三个代表
$\mathcal J=\{\{1,2\},\{1,3\},\{1,4\}\}$，记
$z_J=(B_J\cdot n)/\alpha_J$，定义
\[
 A=\sum_{J\in\mathcal J}\frac{2}{\alpha_J(1-z_J^2)},\qquad
 A_N=\sum_{J\in\mathcal J}\frac2{\alpha_J}
             \sum_{k=0}^N z_J^{2k},\qquad S_N=T-A_N.
\]
于是$S=T-A$。这一步先将互补图配对再展开，不是逐图抛弃干涉。

\begin{theorem}[保留干涉的全域误差]\label{audit:thm:tail}
对任意整数$N\ge0$，处处有
\[
 0\le A-A_N\le\delta_N:=\frac{27}{10}\left(\frac49\right)^{N+1},
 \qquad |S|,|S_N|\le\frac{27}{10}.
\]
令$\cA(f)=(4\pi)^{-1}\int_{S^2}f(n)\dd\Omega$，
$W=\cA(S^2)$、$W_N=\cA(S_N^2)$，则
\begin{equation}
 |W-W_N|\le\frac{729}{50}\left(\frac49\right)^{N+1}.       \label{eq:audit:1}\end{equation}
\end{theorem}
\begin{proof}
几何级数余项正，且每个配对项的余项不超过
$(2/\alpha_J)(4/9)^{N+1}/(1-4/9)$。
三项求和及$\alpha_J\ge4$给$\delta_N$。
同理$0\le A_N\le A\le27/10$。
结合$0<T\le1/2$，$S,S_N$均落在$[-27/10,1/2]$内。
于是
\[
 |S^2-S_N^2|\le(|S|+|S_N|)|S-S_N|
 \le\frac{27}{5}\delta_N.
\]
归一化正角测度不增大这个上界，得到\eqref{eq:audit:1}。
\end{proof}

角积分$W_N$可以完全消去。若$i_1,\ldots,i_{2r}\in\{1,2,3\}$，
球面均匀矩为
\begin{equation}
 \cA(n_{i_1}\cdots n_{i_{2r}})
 =\frac1{(2r+1)!!}
   \sum_{\text{所有配对}}\ \prod_{\{a,b\}}\delta_{i_ai_b},
 \qquad \cA(\text{奇数次单项式})=0.                       \label{eq:audit:2}\end{equation}
它来自旋转不变性及逐次收缩$\sum_i n_i^2=1$。
因而$W_N$是$T,\alpha_J^{-1},B_I\cdot B_J$的有限有理系数表达式，
无需角采样，也无需在重合分母处数值相减两个对数。
实际实现宜用矩递推，式\eqref{eq:audit:2}不是要求显式列出全部配对。
式\eqref{eq:audit:1}--\eqref{eq:audit:2}是标准级数与球矩方法在本通道的应用，不宣称方法首创。

\begin{corollary}[20阶角截断的截面误差]\label{audit:cor:eta}
令$C_N$以$W_N$代替真实角平均平方而保留其余相空间积分。
使用旧证书的相空间体积上界，有
\[
 |C_{24}-C_N|
 \le\eta_N:=
 \frac{763}{10^{10}}\frac{729}{50\cdot768}
          \left(\frac49\right)^{N+1}.
\]
特别地
\begin{equation}
 \eta_{20}=
 \frac{6400507904}{109932594211669278076171875}
 <6\times10^{-17}.                                    \label{eq:audit:3}\end{equation}
\end{corollary}
\begin{proof}
首项系数为$(2s\beta\,4!)^{-1}\int \dd\Phi_4\,W$，
$\beta=\sqrt{4/5}>4/5$，故$2s\beta\,4!>768$。
正相空间积分、\eqref{eq:audit:1}及旧$\Phi_4<763/10^{10}$给结论。
\eqref{eq:audit:3}的末端不等式可直接交叉相乘检查，未使用浮点尾数。
\end{proof}

注意$C_N$在尚未严格计算。小的$\eta_{20}$不是窄的$C_{24}$区间；
它只隔离了一个已控制的误差来源，剩余五变量认证仍是实际缺口。

\subsubsection{剩余积分的有限证书与验收}
沿此前五变量推前，用$a,b$表示两对质量平方，$x,y$是两对内极角余弦，
$\varphi$是相对方位角，则
\[
 C_N=\frac1{2s\beta\,4!}
 \int_{\mathcal R_s}\rho(s,a,b)\dd a\dd b
 \int_{-1}^1\frac{\dd x}2\int_{-1}^1\frac{\dd y}2
 \int_0^{2\pi}\frac{\dd\varphi}{2\pi}\,W_N,
\]
\[
 \rho=\frac{\sqrt{\lambda(s,a,b)}}{2048\pi^5s}
          \sqrt{1-\frac4a}\sqrt{1-\frac4b}.
\]
平方在相加全部树项之后才取，故所有45个干涉对均保留。
设$Q=\sqrt{20}-4$，以$u,v\in[0,1]$置
\begin{equation}
 A=2+Qu,\quad B=2+Q(1-u)v,\quad a=A^2,\quad b=B^2,
 \quad \left|\frac{\partial(a,b)}{\partial(u,v)}\right|
       =4Q^2(1-u)AB.                                   \label{eq:audit:4}\end{equation}
此处标量$A,B$仅用于变量替换，不是上文的传播子和或向量$B_J$。
\eqref{eq:audit:4}覆盖$\mathcal R_{20}$，边界退化为零测度。
再令$x=2t_1-1,y=2t_2-1,\varphi=2\pi t_3$，得到单位五立方体。

\begin{proposition}[可终止的严格包围路线]\label{audit:prop:boxes}
对任意$\varepsilon>0$，存在有限立方体分割与有限精度的向外区间运算，
给出$L_N\le C_N\le U_N$且$U_N-L_N<\varepsilon$。
于是
\[
 C_{24}\in[L_N-\eta_N,U_N+\eta_N].
\]
这一存在及终止陈述不提供实用时间上界，尚未执行该分割。
\end{proposition}
\begin{proof}
使用两对静止系及沿固定轴的有限Lorentz推动构造末态，
内部空间动量用非负平方根表达。所有分母有$A,B\ge2$、
$\alpha_J\ge4$及类时分母至少8的统一正界；
推动公式无需除以在阈值退化的空间动量长度。
故\eqref{eq:audit:4}后的被积函数在闭立方体连续，包括阈值边界。
式\eqref{eq:audit:2}是有限运算，不产生新的极点。
平方根、三角函数和倒数在这些域上有收敛的包含区间扩张；
平方根参数因区间依赖得到负下端时可与$[0,\infty)$相交。
逐步缩小网格并增加运算精度，区间振幅一致趋零，
正体积乘各格上下包含值给Darboux夹逼。
紧致性与一致连续性保证有限步达到任意给定$\varepsilon$。
\end{proof}

第一选点维持$s/m^2=20$，首个实用验收要求为
\begin{equation}
 0<L<U,\qquad C_{24}\in[L,U],\qquad \frac{U-L}{L}\le10^{-3}. \label{eq:audit:5}\end{equation}
这是完整区间宽度的千分之一标准，不是目前已达到的精度。
初版不追求认证展示的全部六位有效数字。
证书需保存模型/方案、测度、端点、分割、外舍入方式及误差预算；
独立核验从有限分割及包含不等式重建结论，而非仅重读一个中心值。
实际实现若发现五变量界收敛过慢，应优化解析变量或区间模型，
不以无限增加随机样本代替误差证明。

\subsection{包容系数与光学闭合分别验收}
\subsubsection{弹性误差不能遗漏}
在20点写成
\begin{equation}
 \sigma_{\rm incl}
 =\frac{g^2+c_3g^3+c_4g^4}{640\pi m^2}+O(g^5),\quad
 c_4=c_{4,\rm el}+640\pi C_{24}.                         \label{eq:audit:6}\end{equation}
旧近似$c_3\simeq0.0180991989932455$、
$c_4\simeq0.00014250842$已经包含四体；
旧\texttt{elastic\_total}只含弹性，
合并记录位于\texttt{data/phi4\_four\_body/inclusive.json}。
若$C_{24}\in[L,U]$、$c_{4,\rm el}\in[E_-,E_+]$，
則\eqref{eq:audit:6}使用同一正$\pi$区间作向外运算；
$\pi$视作精确常数时结果就是
$[E_-+640\pi L,E_++640\pi U]$。
只认证四体而继续使用未包围的$E_\pm$，不能认证完整$c_4$。
$c_3$也必须单独包围。保留的小数应使整个区间落在同一个舍入单元。

\subsubsection{独立三圈前向计算的目标}
沿用树级弹性振幅为$-g$的约定，写
\[
 \mathcal M_{22}=-g+g^2\mathcal A_1+g^3\mathcal A_2
                      +g^4\mathcal A_3+\cdots.
\]
阶$g^4$幺正性比较是
\begin{equation}
 2\operatorname{Im}\mathcal A_3(s,0)
 =\frac1{2!}\int\dd\Phi_2
       \bigl(|\mathcal A_1|^2-2\operatorname{Re}\mathcal A_2\bigr)
  +\frac1{4!}\int\dd\Phi_4|S_6|^2.                       \label{eq:audit:7}\end{equation}
其右端可由既有弹性与四体振幅构造，但这不构成左端的独立计算。
必须从同一拉格朗日量、同一方案的三圈前向图和反项出发，覆盖允许的二体、四体切割，
核对组合因子及外腿/LSZ处理，再与右端比较。
以\eqref{eq:audit:7}右端定义$\operatorname{Im}\mathcal A_3$属于循环检验。
已有振幅阶$g^3$的检验不能替代\eqref{eq:audit:7}。

最先应完成的是三圈图及反项的完整分析切割恒等式；
数值积分只有在这条结构链固定后才有独立检验意义。
若两端分别有严格区间，区间交叠只表明在该精度未发现矛盾，
并不单独证明全运动学恒等式。全文应分别报告解析证明范围和点值残差范围。


\section{分类表的先行工作与新颖性边界}\label{sec:literaturematrix}
\subsection{文献比较：依据与分级}
截至本稿的审阅于2026年10月14日用题名、作者、banana/sunrise、motivic coaction、
K3与2026交叉检索，追踪到arXiv原文及作者或出版方页面。
只用一手论文支持技术结论。
检索同时检查2025年预印本在2026年发表或修订的版本，不能只搜编号26开头。
网页索引、摘要或未命中的关键词均不能证明不存在先行工作。

比较等级：
\textbf{K}是已知结果，应引用；
\textbf{A}是已有定理在指定输入的应用或重新推导；
\textbf{D}是值得逐式比较的差异候选，尚无新颖性结论；
\textbf{O}是本项目仍开放的验收。
“本文独立推导”仅说明证明来源，不等于优先权。

\subsubsection{直接影响分类表的对照}
\small
\setlength{\tabcolsep}{5pt}
\begin{longtable}{p{.16\textwidth}p{.30\textwidth}p{.36\textwidth}p{.08\textwidth}}
\toprule
项目位置 & 最近先行结果/精确入口 & 本项目可保留的定位与限制 & 级别\\\midrule
\endfirsthead
\toprule 项目位置 & 最近先行结果/精确入口 & 本项目可保留的定位与限制 & 级别\\\midrule
\endhead
E1--E4两圈几何 &
DHV定理5.1、7.5，命题7.6\cite{DHV} &
超椭圆/Hodge上界及适用域内相对上界已有；具体边界、正链与单标量最小性的增量须比较。 & K/D\\
E1任意质量sunrise &
ABW的$D=2$椭圆二重对数；BMW的七主积分$\varepsilon$形式\cite{LitABW,LitBMW} &
椭圆性及表示不是新结果；图族规则须说明分子、注入、边界与特殊点条件。 & K/D\\
E1与细分规则 &
Tapušković一般质量sunrise余作用及边细分\cite{LitTap} &
不能声称首次从细分图得到共轭；需比较本项目零注入/不同质量/分子规则的接受域。 & K/D\\
A1一圈完整对象 &
Mounoud定理1.1、1.2、7.26、7.27\cite{LitMou} &
动量空间Nori对象、切割/收缩及图式余作用已有。特殊参数与本项目标架的映射仍需逐条对照。 & K/D\\
K1--K2四质量banana &
两组2025预印本及2026刊文\cite{LitDMPS,LitPTWW} &
全质量微分系统与计算方法已有；公开纯矩阵明确归原作者。实际固定纤维比较、整格与最小对象是另行比较项。 & K/A/D\\
K类特殊质量层 &
Duhr--Maggio的两参数K3与三质量相等计算\cite{LitDM,LitDMfull} &
某些K3周期可约为椭圆乘积；“有K3”不能写成“每个质量层都需要不可约新函数”。 & K\\
K1指定Picard与纯群 &
Charles定理1、6及EJ定理7.5、7.7\cite{LitCharles,LitEJ} &
一般工具已有，指定纤维算术应用由自身证据承担；不宣布新Picard算法或通用CM排除。 & A/D\\
完整混合群 &
Eskandari三步混合扩张框架\cite{LitEsk} &
一般幺幂群机制已有；项目的指定对象非分裂、表示与交换子必须各自证明。 & A/D\\
物理切割/混合结构 &
Wang--Wang--Wang的Higgs四体积分族，第3--7节\cite{LitWWW} &
真实物理积分族已有MPL/椭圆/K3混合；本$\phi^4$基准不是同模型，不能借其替代\eqref{eq:audit:7}。 & K/D/O\\
O2窄误差 &
第\ref{sec:narrowphysics}节及旧严格相空间界 &
仅角截断新增解析控制；五变量与弹性系数尚无窄证书。 & A/O\\
\bottomrule
\end{longtable}
\normalsize

\subsubsection{DHV：不仅有绝对曲面结论}
DHV定理5.1给$(a,1,c)$图超曲面的超椭圆型混合Hodge上界；
定理7.5在$3D/2\le a+c$条件下还给相对对象的同类上界。
因此将该文笼统描述为“只研究绝对几何”不准确。
对本项目使用的$(2,1,1),D=2$，该不等式恰取等号。
命题7.6给相应绝对$H^2$一般为纯Tate。
本项目若强调其相对边界椭圆组成，需与上述相对上界一起陈述，
真正待比较的是组成的精确重数、原物理标架与最小性。
截至本稿的审阅重读保存PDF第4、5、45页；几何定义及原定理接受条件仍随原文引用。

\subsubsection{ABW/BMW与sunrise余作用}
ABW\cite{LitABW}已给非零任意质量二维sunrise的椭圆二重对数表示。
BMW\cite{LitBMW}第7节式(84)及第9节给模空间上的$\varepsilon$形式及逐阶解，
并区分固定模参数与等质量层。
因此“任意质量需要椭圆结构”“由微分系统构造解”都是背景。
Tapušković\cite{LitTap}第1.2、3、4节的相关命题与证明已进一步对照：
其对象在实现范畴，细分继承原边质量，等质量层与固定质量系数域须区分。
具体质量恒等式、全参数载体及微分秩的对照见第\ref{sec:subdivisioncomparison}节。
该文第5节和全部附录尚未逐式审阅；完整扩张新主张仍须继续比较。

\subsubsection{Charles与EJ：避免误引成无条件万能算法}
Charles定理1的应用需区分端同态域$E$全实或CM以及$\dim_E T$奇偶。
其定理6原文允许算法不终止，并在$X\times X$的余维2 Hodge猜想条件下保证终止；
备注8另列可无条件终止的情形，包括$E=\Q$。
故该篇摘要的概括不能改写为任意K3的无条件有效终止保证。
本项目指定纤维的有限证书是具体应用，不等于解决一般终止问题。

EJ定理7.5的正规性有$X\times X$的Hodge条件；
定理7.7在CM分支使用已知结果满足此条件，并描述分量群及jump character。
它不是对任意K3自动判非CM的定理。
项目以CM为反设时可以调用这个分支，但必须给自己的算术矛盾。
这些核对澄清工具的适用性，并没有重新认证本项目所有旧数值证书。

\subsection{2026年必须纳入的近邻与版本纪律}
\subsubsection{动量空间一圈：实际范畴与物理操作}
Mounoud于2026年5月提交的论文\cite{LitMou}直接使用Nori动机局部系统。
定理1.1连接欧氏一般运动学积分与对象；定理1.2把权重分级写成二次Artin动机的Tate扭转；
第7.10节给图式coaction。
这是本项目“有效对象+切割/收缩+物理读出”的直接近邻，不能仅在背景末尾泛引。
Artin--Tate也不能直接改成普通MZV。
截至本稿的审阅已阅读这些定理和第1.1节；退化层、维数及传播子指数的逐项同构比较仍待完成。

\subsubsection{banana：区分首次预印与发表年}
Duhr等\cite{LitDMPS}的四不等质量系统是项目既有纯矩阵来源；
Pögel等\cite{LitPTWW}的独立路线发表于JHEP 2026年1月。
两者预印本均始于2025年。后者以Baikov过滤与基变换构造$\varepsilon$因子化系统；
它不要求先识别几何，因而也是“无需先完整积分”的直接方法学对照。
本项目不得把任意质量K3系统或其可计算性称作首创。
差异候选应落实到明确固定点的实际相对对象、物理标架、完整扩张和可复核证书，
而不是把微分系统换用动机语言。

三质量相等层的K3周期分解及完整积分研究\cite{LitDM,LitDMfull}
直接约束分类表的特殊点/特殊子族栏。
K3标签与椭圆乘积表示可能同时成立；
可见纯结构、函数表示及整个扩张必须分别比较。
截至本稿的审阅对这两篇完成摘要与范围确认，尚未逐式审阅全部附录。

\subsubsection{实际物理混合与余作用}
2026年7月的Higgs四个有质量bottom夸克积分族\cite{LitWWW}
在四条有质量切割下同时包含MPL、椭圆和K3扇区，并处理扇区混合，
给主积分前两阶$\varepsilon$结果。
其引言明确将积分族解析求解与完整包容衰变率区分。
这要求本项目把“混合结构进入真实物理”表述成具体模型与投影的结论，
而不是泛称此前没有物理实现；也不能把积分族的结构自动等同总可观测量的最小对象。

Kleinschmidt--Porkert--Schlotterer的2026年刊文\cite{LitKPS}
提出genus-one Eisenstein迭代积分余作用并检查其结构性质。
应保留原文“proposal”的证据层次，不能据此撤销或自动证明本项目原积分的实际比较。
Eskandari的2026年8月修订\cite{LitEsk}题名为
\emph{Tannakian fundamental groups of blended extensions}；
旧记录若沿用较早题名，引用须指明版本。

\subsubsection{待扩展范围}
截至本稿的审阅不是截止2026年10月14日所有预印本的穷尽覆盖。
还需系统向前/向后追踪上述论文引文，并补充图多项式约化与线性可约性、
$M_{0,n}$字母/端点、非Tate物理标架、重整化周期与临界响应的直接近邻。
2603.03649的算法会议报告已确认题名、作者组和日期，
截至本稿的审阅不把其摘要宣传当作任意图分类定理。
新近2610.02286工具论文只命中摘要，尚未技术审阅，不进入已证明比较栏。

新颖性最终验收的单位是命题：
同一输入、同一基域、同一对象层次、同一物理标架与同一输出强度逐项比较。
只要其中有未比对项，就标“增量候选/待查新”，不写“首次”“统一解决”。
新证书、应用与综合组织也可有价值，但应按实际贡献命名。


\section{带质量的细分、合流喷射与文献秩数的可比性}
\label{sec:subdivisioncomparison}

文献对照应比较同一物理输入与同一对象，不能只比较“椭圆”“K3”或矩阵阶数。
本节先定位三个直接近邻，再给出质量合流的有限规则。
规则的目标是确定哪些椭圆曲线足以承载给定函数，而不是从环境标签断言每个因子都可见。

\subsection{三个直接近邻及其精确接口}
\label{subsec:directneighbours}
Tapušković\cite{LitTap}在实现范畴$\mathcal H(S)$中工作（第1.2节），
仅处理收敛积分。同质量、零新顶点注入的细分由命题3.3及推论3.4处理。
命题4.1给sunrise载体
\[
 \operatorname{gr}^{W}M
 =H^1(E)(-1)\oplus\Q(-1)^3\oplus\Q(0).
\]
命题4.3及定理4.4给显式余作用；式(51)显示质量比对数，
第4.3.3节式(53)在等质量层给秩三对象。
这些是本项目E1--E2的先行结果，而非仅有积分表示的背景。
实现范畴的结论不自行识别本项目指定Nori固定纤维。

Duhr等\cite{LitDMPS}第2节式(4)、(7)--(8)给$15$个主积分和全四变量系统；
第3节式(24)、(30)、(32)--(33)给六个K3周期与配对。
第4--6节及附录C处理规范基与初值。本项目复用其纯系统，
没有另行首创全质量banana微分方程。

Pögel等\cite{LitPTWW}第2节式(9)、(11)--(14)同样从$15=4+11$主积分出发；
引言和第3节的Baikov过滤把顶扇区分成六个K3方向与五个附加方向。
式(56)和附录B式(95)的基变换含$\varepsilon^{-1}$、$\varepsilon^{-2}$。
这提示必须审阅特殊化过程，不能把规范系统直接置$\varepsilon=0$作为实际几何比较。

\subsection{为什么不同质量的细分不能直接套同质量定理}
\begin{proposition}[细分后的质量项]\label{prop:splitmassidentity}
设连通图$G$的一条边$e$被分成两条串联边，新顶点没有外动量注入。
记其平方质量为$\mu_1,\mu_2$。在参数变换
\[
 \alpha_{e_1}=y\alpha_e,\qquad
 \alpha_{e_2}=(1-y)\alpha_e
\]
下，第一图多项式$U$与无质量运动学部分$F^0$分别拉回$U_G,F_G^0$，而
\begin{equation}\label{eq:splitmassidentity}
 \rho^*F_{G_s}
 =F_G\big|_{\mu_e=\mu_2}
   +(\mu_1-\mu_2)y\alpha_eU_G.
\end{equation}
因而在$\alpha_eU_G\ne0$的稠密开集，右端与$y$无关当且仅当$\mu_1=\mu_2$。
\end{proposition}
\begin{proof}
用$U=\sum_T\prod_{f\notin T}\alpha_f$的约定。
若生成树含原边，细分后树含两条新边；若不含原边，新树恰含其中一条，
两项之和以$\alpha_{e_1}+\alpha_{e_2}$替代$\alpha_e$。
生成二森林同理。唯一额外可能是新顶点单独成为一个分支；其注入动量为零，
所以对$F^0$没有贡献。
最后用$F=F^0+U\sum_f\mu_f\alpha_f$，新边的质量项拉回为
$[\mu_2+(\mu_1-\mu_2)y]\alpha_eU_G$，得到所述公式。
$U_G$非零，故稠密开集上的充要性成立。
\end{proof}

这是图多项式的恒等式；“$y$不再出现”才是此处受质量约束的结论。
它没有证明不同质量永远不存在其它几何对应。
在同一动量流$q$上，不同质量的两传播子改用
\[
 \frac1{(q^2+\mu_1)(q^2+\mu_2)}
 =\frac1{\mu_2-\mu_1}
   \left(\frac1{q^2+\mu_1}-\frac1{q^2+\mu_2}\right).
\]
它产生两个质量特化的读出，而非直接产生原单一质量纤维中的升幂读出。
两类操作都可以研究，不能共享未经核对的“细分保持”结论。

\subsection{质量碰撞后的有限导数规则}
\label{subsec:confluenttheta}
取两端点间三条内部不交的路径，所有内部顶点零注入；仅两个端点有欧氏外动量。
路径$a$上有$r_a$种互异正平方质量$\mu_{aj}$，相应重数为$n_{aj}\ge1$。
令$N_a=\sum_j n_{aj}$，并限定分子为三条路径的动量平方$Z_1,Z_2,Z_3$的多项式$P$，
满足$\deg_{Z_a}P<N_a$。一般夹角/张量分子不在本规则内。
\begin{theorem}[合流质量的有限部分分式证书]\label{thm:confluentcertificate}
定义
\[
 H_{\boldsymbol j}(\boldsymbol Z)=
 \frac{P(\boldsymbol Z)}{
  \prod_{a=1}^3\prod_{\ell\ne j_a}(Z_a+\mu_{a\ell})^{n_{a\ell}}}.
\]
对$1\le k_a\le n_{aj_a}$置
\begin{equation}\label{eq:confluentcoeff}
 C_{\boldsymbol j,\boldsymbol k}=
 \left.
 \left(\prod_{a=1}^3
 \frac{\partial_{Z_a}^{\,n_{aj_a}-k_a}}{(n_{aj_a}-k_a)!}\right)
 H_{\boldsymbol j}(\boldsymbol Z)
 \right|_{Z_a=-\mu_{aj_a}}.
\end{equation}
则有唯一的恒等式
\begin{equation}\label{eq:confluentpartial}
 \frac{P(\boldsymbol Z)}{
       \prod_{a,j}(Z_a+\mu_{aj})^{n_{aj}}}
 =\sum_{\boldsymbol j,\boldsymbol k}
      \frac{C_{\boldsymbol j,\boldsymbol k}}
           {\prod_{a=1}^3(Z_a+\mu_{aj_a})^{k_a}}.
\end{equation}
只需有限个代数域运算和有限阶多项式导数，不需要完整费曼积分。
所有系数为零当且仅当$P=0$。
\end{theorem}
\begin{proof}
单变量真分式在极点$-\mu_j$的主部由
$(Z+\mu_j)^{n_j}P(Z)/\prod_\ell(Z+\mu_\ell)^{n_\ell}$
的前$n_j$项Taylor系数给出；这正是式\eqref{eq:confluentcoeff}。
减去全部主部后，无有限极点且在无穷远趋零，故余项为零。
在其它变量的有理函数域上依次使用这一恒等式。
每一变量的次数条件保证无多项式余项，三个局部系数操作可交换，
从而给出所列混合导数。
局部主部的唯一性逐变量证明展开唯一；有理函数恒零又等价于其非零共同分母的分子恒零。
\end{proof}

当全部$n_{aj}=1$时，上式退化为原网格值公式。
有重复质量时，网格值必须换为有限Taylor喷射。
但$P\ne0$只说明有理被积函数非零；它不是所有积分非零或每个椭圆因子可见的充分条件。
不同导数项之间仍可能发生精确抵消。

\begin{proposition}[积分读出及椭圆载体上界]\label{prop:confluentelliptic}
在$D=2$及上述正质量欧氏条件下，取测度$d^2k\,d^2\ell/\pi^2$。
记$S(s;\mu_1,\mu_2,\mu_3)$为单位边三线sunrise，$s=-p_E^2$。
式\eqref{eq:confluentpartial}的积分等于
\begin{equation}\label{eq:confluentintegral}
 I_P=\sum_{\boldsymbol j,\boldsymbol k}
 C_{\boldsymbol j,\boldsymbol k}
 \frac{(-1)^{|\boldsymbol k|-3}}{\prod_a(k_a-1)!}
 \left.
 \partial_{\mu_1}^{k_1-1}
 \partial_{\mu_2}^{k_2-1}
 \partial_{\mu_3}^{k_3-1}S(s;\boldsymbol\mu)
 \right|_{\mu_a=\mu_{aj_a}}.
\end{equation}
每个导数项有同一个sunrise相对代数对上的周期表示。
若所有对应平面三次曲线$E_{\boldsymbol j}(s)$光滑，便得到至多
$R=r_1r_2r_3$个椭圆载体的表示上界。
在这些相对载体中，非Tate权重组成仅来自各$H^1(E_{\boldsymbol j})(-1)$，
路径重数增大本身不强制引入K3载体。
\end{proposition}
\begin{proof}
每个正质量升幂传播子满足
\[
 (q^2+\mu)^{-k}
 =\frac{(-1)^{k-1}}{(k-1)!}\partial_\mu^{k-1}(q^2+\mu)^{-1}.
\]
在正质量紧邻域内，所有有限阶导数均被一个固定常数乘单位三线sunrise绝对被积函数控制。
后者在二维欧氏区收敛；有限和及受控微分均可通过积分。
次数条件与有限部分分式也给原被积函数的绝对收敛，证明式\eqref{eq:confluentintegral}。

为明确周期表示，令$A=k_1+k_2+k_3$，$U=\alpha_1\alpha_2+\alpha_2\alpha_3+\alpha_3\alpha_1$，
$F=(\sum\mu_a\alpha_a)U-s\alpha_1\alpha_2\alpha_3$。
单位sunrise的参数形式为$\Omega/F$。微分后得到
\begin{equation}\label{eq:confluentform}
 \omega_{\boldsymbol k}=
 \frac{(A-3)!}{\prod_a(k_a-1)!}
 \frac{\prod_a\alpha_a^{k_a-1}U^{A-3}}{F^{A-2}}\,\Omega.
\end{equation}
它也是正整数幂参数公式，Gamma规范与上式一致。
提高极点阶数不增加$F=0$之外的新极点支撑；
在原三角形角点解消上，它是同一族相对上同调的Gauss--Manin导数类。
故每个三元组仍使用其原相对对，而非引入新的曲面。
该对的留数序列把唯一的非Tate权重项放在$H^1(E)(-1)$；
Tate边界和椭圆部分的秩数可由前述命题4.1核对。
有限直和给出载体上界，证明完毕。
\end{proof}

这里得到的是原积分函数的具体有限表示和所用相对载体的上界。
尚未构造整个合流图原生Nori对象到该直和的所有标架态射，
也未求上述和的最小子商。
因此不能写成“原生最小秩等于$6R$”，不能自动套用互异质量行的$2r+1$，
更不能由部分分式项非零认定对应核在最后求和中幸存。
椭圆载体可用的结论也不证明某个特殊数值不属于MZV。

\paragraph{一个不需求积的反例。}
令第一条路径有两条同质量边，$n_{11}=2$、$\mu_{11}=\mu>0$，
另两条路径各有一条边，平方质量为$\nu,\lambda>0$。
取$P=Z_1+\mu$。所有零阶网格值为零，然而
\[
 \frac{Z_1+\mu}{(Z_1+\mu)^2(Z_2+\nu)(Z_3+\lambda)}
 =\frac1{(Z_1+\mu)(Z_2+\nu)(Z_3+\lambda)}.
\]
喷射公式给$C_{(1,1,1),(1,1,1)}=1$，其它系数为零。
物理读出正是非零sunrise；在原E1接受的正常阈值族上，其椭圆核可见。
因此把“零网格值导致消失”外推到质量合流层会直接给错答案。
这同时提供分类表特殊质量栏的可核查测试，而不是一个数值拟合例子。

\subsection{固定质量后的微分秩为什么不是算术秩}
\begin{lemma}[遗忘连接会丢失常数扩张]\label{lem:constantkummerloss}
取$a\in\Q_{>0}$且$a\ne1$。其Kummer一动机$[\mathbb Z\to\mathbb G_m]$，$1\mapsto a$，
拉回任意$s$开集后，单看$s$方向的代数连接是常值的。
其对数读出在$\mathbb C(s)$上的最小齐次微分阶数为一，
但原混合Hodge扩张不分裂。
\end{lemma}
\begin{proof}
所有代数几何数据来自固定纤维，拉回后沿$s$导数为零。
非零常函数$\log a$被$\partial_s$消去，不被非零零阶算子消去。
相应Hodge扩张的类为
\[
 [\log a]\in\operatorname{Ext}^1_{\mathrm{MHS}}(\Q(0),\Q(1))
       =\mathbb C/(2\pi i\Q).
\]
$\log a$为非零实数，而$2\pi i\Q$为纯虚数，故类不为零。
这给出常值连接与非分裂比较数据并存的明确例子；不需假设周期映射单射。
\end{proof}

因此E1的$\mathbb C(s)$秩三与全参数sunrise的实现载体秩六不是相互反驳。
质量对数作为复常数进入系数域，不表示其算术扩张已消失。
等质量使质量比对数实际为零，与把非零对数纳入常数域，是两个不同操作。
E2微分群、K1固定算术群、K2纯几何载体应继续保留各自的对象与基域标签。

\subsection{正则化参数特殊化的必要数据}
\begin{proposition}[奇异基变换不能直接特殊化]\label{prop:epsilonlattice}
令$K=\mathbb C(s)$，$R=K[\varepsilon]_{(\varepsilon)}$。
即使两个自由$R$模上的稳定连接在$K(\varepsilon)$上同构，
它们模$\varepsilon$的连接也可以不同构。
若两个模型由$\mathrm{GL}_r(R)$中的连接相容基变换联系，才可直接诱导同构的特殊纤维。
\end{proposition}
\begin{proof}
在$K(\varepsilon)$上取$e_1,e_2$及
\[
 \nabla e_1=0,\qquad \nabla e_2=e_1\,\frac{ds}{s}.
\]
令$L=Re_1+Re_2$，$L'=Re_1+R(\varepsilon e_2)$。
两者均为自由稳定格，且泛纤维相同。
$L/\varepsilon L$保留非平凡对数扩张；否则$1/s$有$K$中原函数，
与有理函数的导数在$s=0$处留数为零矛盾。
但$L'/\varepsilon L'$中的连接矩阵为零，因
$\nabla(\varepsilon e_2)=\varepsilon e_1 ds/s$。
这两个特殊纤维不同构。若基变换及逆都在$R$上正则，
模$\varepsilon$后仍可逆并保持连接，最后一句立即成立。
\end{proof}

此外，自由秩$r$的$R$格模$\varepsilon$仍为秩$r$。
故把$15$主积分系统与本项目严格二维秩十活动实现联系起来，
需要明确选择正则模型、有限项/极点项、实际子商及物理初值，
不能只做矩阵代入。这不是说文献的$15$维结果错误；两个任务的输入和输出不同。
本节没有构造$15\to10$的实际比较，也没有以维数正规化解代替原始Nori标架比较。

\subsection{本批对分类和查新的实际改变}
同质量细分的几何保持与不同质量的有限和现在有式\eqref{eq:splitmassidentity}作为分界。
合流层由定理\ref{thm:confluentcertificate}和命题\ref{prop:confluentelliptic}给可执行上界规则，
零阶值判据的反例阻止错误外推。
分类表新增E5，只记录这一接受域、椭圆表示上界和未完成的最小性。

文献优先权方面，sunrise显式余作用、同质量细分及全质量banana规范系统均已存在。
本项目的差异候选仍须与同一范畴、同一参数层、同一标架及同一输出强度逐项比较。
本节的部分分式、喷射和格的反例属于标准方法的明确推导与应用，不宣称这些方法本身新颖。
ABW的任意质量椭圆表示、DHV相对上界、Charles/EJ算术条件及2026近邻
继续按前节的来源等级保留；未逐式审阅的文献不因本批读取三篇而升级。
完整窄截面及三圈前向全切割也未由这些代数结论完成。


\begin{thebibliography}{99}\raggedright
\bibitem{brown} F. Brown, \emph{Notes on motivic periods},
\href{https://arxiv.org/abs/1512.06410}{arXiv:1512.06410},
\S2.4, Lemma 2.4.
\bibitem{stacks} The Stacks Project,
\href{https://stacks.math.columbia.edu/tag/02MS}{Tag 02MS},
Lemma 12.10.6. Serre商及精确函子。
\bibitem{BK} F. Brown and D. Kreimer, \emph{Decomposing Feynman rules},
\href{https://arxiv.org/abs/1212.3923}{arXiv:1212.3923}.
实际MOM提升另由内部证明承担。
\bibitem{KR} E. Kani and M. Rosen,
\emph{Idempotent relations and factors of Jacobians},
Math. Ann. 284 (1989), 307--327.
\href{https://eudml.org/doc/164555}{出版记录}。
本篇直接由拉回/迹证明所用分解，未声称本批通读该文。
\bibitem{CarRod} A. Carocca and R. E. Rodr\'iguez,
\emph{Jacobians with group actions and rational idempotents},
\href{https://arxiv.org/abs/math/0305328}{arXiv:math/0305328v1},
\S2式(2.3)、\S5.1 Proposition5.2。
有理群幂等元及商曲线的经典背景；本文带Kummer系数的选择公式另给证明。
\bibitem{HHMZ} J. D. Hauenstein, R. Huang, D. Mehta and Y. Zhang,
\emph{Global Structure of Curves from Generalized Unitarity Cut of Three-loop Diagrams},
\href{https://arxiv.org/abs/1408.3355}{arXiv:1408.3355}.
链式White-house亏格为先行结果，相同亏格不证明星式图与其同构。
\bibitem{muhl} M. M\"uhlbauer,
\emph{Cutkosky's Theorem for Massive One-Loop Feynman Integrals---Part 1},
\href{https://arxiv.org/abs/2206.08402}{arXiv:2206.08402v2},
\S2.3, Theorems17--18及22。局部留数和Picard--Lefschetz背景；
本篇banana正参数阈值与交数直接证明，未调用该文单圈物理定理覆盖多圈。
\bibitem{BKV} S. Bloch, M. Kerr and P. Vanhove,
\emph{A Feynman integral via higher normal functions},
\href{https://arxiv.org/abs/1406.2664}{arXiv:1406.2664},
Theorem2.2.1, \S2.3, \S3.2, \S6.1--6.2。
\bibitem{brownamp} F. Brown,
\emph{Feynman amplitudes, coaction principle, and cosmic Galois group},
\href{https://arxiv.org/abs/1512.06409}{arXiv:1512.06409}。
\bibitem{BKone} S. Bloch and D. Kreimer,
\emph{Feynman Amplitudes and Landau Singularities for 1-loop graphs},
\href{https://arxiv.org/abs/1007.0338}{arXiv:1007.0338v2},
\S7--10, 特别是Proposition10.1。一圈Hodge结构及双对数的经典来源。
\bibitem{SVone} M. Spradlin and A. Volovich,
\emph{Symbols of One-Loop Integrals From Mixed Tate Motives},
\href{https://arxiv.org/abs/1105.2024}{arXiv:1105.2024}。
偶数维多边形的二次型/符号递归背景；本节未宣称实现该算法。
\bibitem{ABWsun} L. Adams, C. Bogner and S. Weinzierl,
\emph{The two-loop sunrise graph with arbitrary masses},
\href{https://arxiv.org/abs/1302.7004}{arXiv:1302.7004},
\S4--5，式(24)--(43)。本篇采用不乘任意质量尺度的归一化；式(37)--(43)的非齐次二阶方程与齐次三阶上界为该文已有结果。
\bibitem{BKVsun} S. Bloch, M. Kerr and P. Vanhove,
\emph{Local mirror symmetry and the sunset Feynman integral},
\href{https://arxiv.org/abs/1601.08181}{arXiv:1601.08181v3},
\S3--4。该文参数$s$是外部动量平方的倒数，本篇不沿用它的变量约定。
\bibitem{DHV} C. F. Doran, A. Harder, P. Vanhove, with appendix by E. Pichon-Pharabod,
\emph{Motivic geometry of two-loop Feynman integrals}.
\href{https://arxiv.org/abs/2302.14840}{arXiv:2302.14840}.
本篇所用保存版本的定理5.1、7.5与命题7.6。
\bibitem{Singer} M. F. Singer,
\emph{Introduction to the Galois Theory of Linear Differential Equations},
\href{https://arxiv.org/abs/0712.4124}{arXiv:0712.4124}；
\href{https://singer.math.ncsu.edu/papers/LMSlectures.pdf}{作者讲义}，
Theorem1.2.8、Proposition1.3.27（讲义第20页）、\S1.5的张量与对称幂。
本节应用正则奇异系统的单值稠密定理，不作该定理的新颖性宣称。
\bibitem{DM} P. Deligne and J. S. Milne,
\emph{Tannakian Categories},
\href{https://www.jmilne.org/math/xnotes/tc.pdf}{作者修订讲义}，
Proposition2.21（第25页）。群商与满Tannaka子范畴的经典对应。
\bibitem{BWpoly} C. Bogner and S. Weinzierl,
\emph{Feynman graph polynomials},
\href{https://arxiv.org/abs/1002.3458}{arXiv:1002.3458}，
\S2式(6)--(9)、\S4式(32)--(35)。图多项式与参数归一化的经典来源。
\bibitem{KreimerBanana} D. Kreimer,
\emph{Bananas: multi-edge graphs and their Feynman integrals},
Lett. Math. Phys. 113, 38 (2023),
\href{https://doi.org/10.1007/s11005-023-01660-4}{doi:10.1007/s11005-023-01660-4}，
\S1.1--1.3、\S2.3及Theorem2.2。相空间递推与色散为已有研究，
本篇二维常数规范及分支模型另行推导。
\bibitem{HuyK3} D. Huybrechts, \emph{Lectures on K3 Surfaces},
\href{https://www.math.uni-bonn.de/people/huybrech/K3Global.pdf}{作者讲义}，
Chapter1的双覆盖与K3定义、Chapter14的有理双点与ADE例外格。
\bibitem{Aleksejevs} A. Aleksejevs,
\emph{Dispersion Approach in Two-Loop Calculations},
Phys. Rev. D 98, 036021 (2018),
\href{https://arxiv.org/abs/1804.08914}{arXiv:1804.08914v6}，
\S II.A，式(15)，原PDF第8页。
色散子圈插入与在壳平方因子消去是已有方法；
本篇标量规范、相对模型、收敛范围与椭圆可见性逐项另证。

\bibitem{PDGphase} Particle Data Group, \emph{Kinematics},
Review of Particle Physics,
\href{https://pdg.lbl.gov/2025/reviews/rpp2025-rev-kinematics.pdf}{2025 review}.
本文的$d\Phi_n$包含$(2\pi)^4$；同种末态阶乘另行声明。
\bibitem{DDphase} A. I. Davydychev and R. Delbourgo,
\emph{Explicitly symmetrical treatment of three-body phase space},
J.\ Phys.\ A \textbf{37} (2004), 4871--4886,
\href{https://arxiv.org/abs/hep-th/0311075}{arXiv:hep-th/0311075}.


\bibitem{WeinEll} S. Weinzierl,
\emph{Iterated integrals related to Feynman integrals associated to elliptic curves},
\href{https://arxiv.org/abs/2012.08429}{arXiv:2012.08429} (2020).
此综述提供属零/属一迭代积分的一般背景，不承担本文特定截面的不相消证明。

\bibitem{GZcrit} L. Guo and B. Zhang,
\emph{Differential Birkhoff decomposition and the renormalization of multiple zeta values},
\href{https://arxiv.org/abs/0710.0432}{arXiv:0710.0432}, Theorem3.2。
本文平方零森林方向式沿用其经典递归，并单独推导。
\bibitem{DScrit} L. Di Pietro and E. Stamou,
\emph{Operator mixing in the $\epsilon$-expansion: Scheme and evanescent-operator independence},
Phys.\ Rev.\ D \textbf{97} (2018), 065007,
\href{https://arxiv.org/abs/1708.03739}{arXiv:1708.03739v1},
\S3式(3.3)--(3.8)。有限方案在固定点的相似变换及算符混合限制；
本文不以该文证明任意有限截块闭合。
\bibitem{CMVcrit} T. Carletti, A. Margheri and M. Villarini,
\emph{Normalization of Poincar\'e Singularities via Variation of Constants},
\href{https://arxiv.org/abs/math/0408021}{arXiv:math/0408021v1},
\S1的共振条件及经典Poincar\'e--Dulac背景。
本文只使用并直接证明形式逐次消元，不调用该文的解析收敛结论。

\bibitem{CPVprod} P. Calabrese, A. Pelissetto and E. Vicari,
\emph{Multicritical phenomena in O($n_1$)$\oplus$O($n_2$)-symmetric theories},
Phys.\ Rev.\ B \textbf{67} (2003), 054505,
\href{https://arxiv.org/abs/cond-mat/0209580}{arXiv:cond-mat/0209580v2}.
原PDF第2页式(1.1)、第5--6页插入空间及第12页式(4.1)；
HTML重排相应尺度公式为式(23)。乘积尺度关系为经典结果。
\bibitem{KPprod} M. V. Kompaniets and E. Panzer,
\emph{Minimally subtracted six loop renormalization of O(n)-symmetric phi4 theory
and critical exponents},
\href{https://arxiv.org/abs/1705.06483}{arXiv:1705.06483v2},
式(16)--(19)。六圈积分及原始系数归于该文。
\bibitem{BrownMZVprod} F. Brown, \emph{On the decomposition of motivic multiple zeta values},
\href{https://arxiv.org/abs/1102.1310}{arXiv:1102.1310v2}.
本项目的$\partial_3,\partial_5$投影约定与前文RG篇保持一致。


\bibitem{ZimMass} W. Zimmermann,
\emph{Convergence of Bogoliubov's method of renormalization in momentum space},
Commun. Math. Phys. \textbf{15} (1969), 208--234,
\href{https://doi.org/10.1007/BF01645676}{doi:10.1007/BF01645676}.
本节调用其经典有质量收敛定理，不重新证明之；
此次原扫描访问受限，操作公式按下一条的作者原页核对。
\bibitem{BGHMass} D. N. Blaschke, F. Gieres, F. Heindl, M. Schweda and M. Wohlgenannt,
\emph{BPHZ renormalization and its application to non-commutative field theory},
Eur. Phys. J. C \textbf{73} (2013), 2566,
\href{https://arxiv.org/abs/1307.4650}{arXiv:1307.4650v2}.
仅用第二节交换理论，原PDF第8页路由、第11--12页式(31)--(41)及第15页局部性；
不将其非交换模型纳入本节周期定理。
\bibitem{HMSMass} A. Huber and S. M\"uller-Stach,
\emph{Periods and Nori Motives}, Springer (2017),
\href{https://home.mathematik.uni-freiburg.de/arithgeom/preprints/buch-errata/buch-errata.pdf}{作者2025年7月1日勘误版}.
Theorem12.2.1、Lemma12.2.4，PDF第295及299--302页；
同时采用修正后的naive periods定义，不忽略实部/虚部及基域区别。

\bibitem{BEKforest} S. Bloch, H. Esnault and D. Kreimer,
\emph{On Motives Associated to Graph Polynomials},
Commun. Math. Phys. \textbf{267} (2006), 181--225,
\href{https://page.mi.fu-berlin.de/esnault/preprints/helene/79-motives.pdf}{作者存档出版版},
\href{https://arxiv.org/abs/math/0510011}{arXiv:math/0510011}.
第6节与第10节，尤其Proposition10.2及Remark10.3：
解析参数化、几何对应和对应是否同构是不同问题。
本文裸欧氏测度常数独立推导，不直接搬用其射影实链规范。

\bibitem{PanzerSector} E. Panzer,
\emph{Hepp's bound for Feynman graphs and matroids},
Ann. Inst. H. Poincar\'e Comb. Phys. Interact. \textbf{10} (2023), 31--119,
\href{https://arxiv.org/abs/1908.09820}{arXiv:1908.09820v3}.
原PDF第11页Lemma2.8、Proposition2.9及第13页配置多项式脚注。
Hepp扇区及贪心基为先行方法，本文同步动量缩放的系数另行保留。
\bibitem{CKNSector} L. de la Cruz, D. A. Kosower and P. P. Novichkov,
\emph{Finite integrals from Feynman polytopes},
Phys. Rev. D \textbf{111} (2025), 105013,
\href{https://doi.org/10.1103/PhysRevD.111.105013}{doi:10.1103/PhysRevD.111.105013},
\href{https://arxiv.org/abs/2410.18014}{arXiv:2410.18014}.
作者预印本第27页Appendix B式(B2)--(B6)；
欧氏正核与Minkowski限制分别保留，不宣称解决其一般分子多面体猜想。

\bibitem{WahlD4} J. M. Wahl, \emph{Simultaneous resolution of rational singularities},
Compositio Math. \textbf{38} (1979), 43--54,
\href{https://www.numdam.org/item/CM_1979__38_1_43_0.pdf}{原文}，
第43--44页的ADE例外配置及Brieskorn同时消解定理。
同时消解与局部根格为经典输入，具体三重分支、三个固定例外类及物理变分另行论证。

\bibitem{DeligneWeight} P. Deligne, \emph{Théorie de Hodge, II},
Publ. Math. IHÉS \textbf{40} (1971), 5--57,
\href{https://www.numdam.org/item/PMIHES_1971__40__5_0.pdf}{原文}，
第3.2节，尤其Corollaires3.2.13与3.2.17（印刷第38--39页）。
正常交叉边界的权谱序列及紧化像为经典输入；
本文另行识别实际极点曲线、物理平方权和阈值变分。

\bibitem{LvovskiEll} S. Lvovski,
\emph{On monodromy in families of elliptic curves over $\mathbb C$},
\href{https://arxiv.org/abs/1705.02129}{arXiv:1705.02129v5} (2018).
Proposition4.2与Corollary4.5，原PDF第11--12页：
非等平凡光滑属一族的单值在$\mathrm{SL}_2(\mathbb Z)$中有有限指数。
本文用此经典输入识别交换曲线的简单性；
具体私有阈值、物理链续接和留数读出的非零性另行证明。

\bibitem{GKPsource} S. Groote, J. G. K\"orner and A. A. Pivovarov,
\emph{A new technique for computing the spectral density of sunset-type diagrams:
integral transformation in configuration space},
\href{https://arxiv.org/abs/hep-ph/9805224}{arXiv:hep-ph/9805224v3} (2000).
原PDF第3--4页的配置空间乘积与谱关系、第6页的局部减除讨论。
本文沿用这一经典联系；有限阶源归一化、全部中间态比较及几何下界另证。

\bibitem{internal} 本项目：
\nolinkurl{FINITE_MOM_MOTIVIC_COMPARISON}、
\nolinkurl{CUT_RESIDUE_VISIBILITY}、
\nolinkurl{OPERATIONAL_MOTIVIC_INTERFACE}、
\nolinkurl{conic_star_classification}、
\nolinkurl{THEORY_MANUSCRIPT}。
本篇加性Serre比较及方案推论完整另证；
新颖性尚未作独立系统检索。

\bibitem{ClassWidth} F. Brown, \emph{On the periods of some Feynman integrals}, arXiv:0910.0114v2. Theorems 1--2,118; \url{https://arxiv.org/abs/0910.0114}.
\bibitem{ClassZigzag} F. Brown and O. Schnetz, \emph{Proof of the zig-zag conjecture}, arXiv:1208.1890. Theorem 1.1; \url{https://www.ihes.fr/~brown/ZigZag.pdf}.
\bibitem{ClassKthree} F. Brown and O. Schnetz, \emph{A K3 in phi4}, arXiv:1006.4064; final author version, \S6.2, equation (30), \S7. \url{https://www.ihes.fr/~brown/K3inphi4.pdf}.
\bibitem{ClassFraming} F. Brown and D. Doryn, \emph{Framings for graph hypersurfaces}, arXiv:1301.3056. Theorem 5 and the following scope discussion; \url{https://www.ihes.fr/~brown/Framings.pdf}.
\bibitem{ClassMzero} F. Brown, \emph{Multiple zeta values and periods of moduli spaces M0,n}, Ann. Sci. Ec. Norm. Super. 42 (2009), 371--489. \url{https://arxiv.org/abs/math/0606419}.

\bibitem{PanzerMZV} E. Panzer, \emph{Feynman integrals and hyperlogarithms},
arXiv:1506.07243; author thesis, equation (1.2.1), \S5.2.1.
\url{https://people.maths.ox.ac.uk/panzer/paper/phd.pdf}.
One-loop convolution and one-scale insertion rules are classical input; the present
classification certificate and the displayed normalization are derived explicitly.
\bibitem{DLMFGamma} NIST Digital Library of Mathematical Functions, \S5.7, equation 5.7.3,
\url{https://dlmf.nist.gov/5.7.E3}.

\bibitem{LitABW} L. Adams, C. Bogner, S. Weinzierl,
\emph{The two-loop sunrise graph in two space-time dimensions with arbitrary masses in terms of elliptic dilogarithms}.
\href{https://arxiv.org/abs/1405.5640v2}{arXiv:1405.5640v2} (2014).

\bibitem{LitBMW} C. Bogner, S. Müller-Stach, S. Weinzierl,
\emph{The unequal mass sunrise integral expressed through iterated integrals on
$\overline{\mathcal M}_{1,3}$}.
\href{https://arxiv.org/abs/1907.01251v2}{arXiv:1907.01251v2} (2020).

\bibitem{LitTap} M. Tapušković,
\emph{The cosmic Galois group, the sunrise Feynman integral, and the relative completion of $\Gamma_1(6)$}.
\href{https://arxiv.org/abs/2303.17534}{arXiv:2303.17534} (2023).
当前核对第1.2、3、4节相关定义、命题及证明；第5节与全部附录未作逐式审阅。

\bibitem{LitMou} U. Mounoud,
\emph{Motivic Galois theory for one-loop Feynman integrals in momentum space}.
\href{https://arxiv.org/abs/2605.20106v1}{arXiv:2605.20106v1} (2026).

\bibitem{LitDMPS} C. Duhr, S. Maggio, F. Porkert, C. Semper, S. F. Stawinski,
\emph{Three-loop banana integrals with four unequal masses}.
\href{https://arxiv.org/abs/2507.23061}{arXiv:2507.23061}, JHEP 12 (2025) 034.

\bibitem{LitPTWW} S. Pögel, T. Teschke, X. Wang, S. Weinzierl,
\emph{The unequal-mass three-loop banana integral}.
\href{https://arxiv.org/abs/2507.23594v2}{arXiv:2507.23594v2},
\href{https://doi.org/10.1007/JHEP01(2026)021}{JHEP 01 (2026) 021}.

\bibitem{LitDM} C. Duhr, S. Maggio,
\emph{Feynman integrals, elliptic integrals and two-parameter K3 surfaces}.
\href{https://arxiv.org/abs/2502.15326}{arXiv:2502.15326}, JHEP 06 (2025) 250.

\bibitem{LitDMfull} C. Duhr, S. Maggio,
\emph{Three-loop banana integrals with three equal masses}.
\href{https://arxiv.org/abs/2511.19245}{arXiv:2511.19245} (2025).

\bibitem{LitCharles} F. Charles,
\emph{On the Picard number of K3 surfaces over number fields}.
\href{https://arxiv.org/abs/1111.4117}{arXiv:1111.4117}.
本篇核对定理1、6和备注8。

\bibitem{LitEJ} A.-S. Elsenhans, J. Jahnel,
\emph{On the component group of the algebraic monodromy group of a K3 surface}.
\href{https://arxiv.org/abs/2304.07585v2}{arXiv:2304.07585v2} (2024).
本篇核对定理7.5、7.7及其证明。

\bibitem{LitEsk} P. Eskandari,
\emph{Tannakian fundamental groups of blended extensions}.
\href{https://arxiv.org/abs/2407.01379v3}{arXiv:2407.01379v3}, 2026年8月5日修订。

\bibitem{LitWWW} J. Wang, X. Wang, Y. Wang,
\emph{Analytic result of a three-loop integral family in the Higgs decay to four massive bottom quarks}.
\href{https://arxiv.org/abs/2607.03398v1}{arXiv:2607.03398v1} (2026).

\bibitem{LitKPS} A. Kleinschmidt, F. Porkert, O. Schlotterer,
\emph{Towards motivic coactions at genus one from zeta generators}.
\href{https://arxiv.org/abs/2508.02800}{arXiv:2508.02800},
\href{https://doi.org/10.1007/JHEP05(2026)105}{JHEP 05 (2026) 105}.

\bibitem{BridgeIM} F. Ivorra and S. Morel,
\emph{The four operations on perverse motives}, Proposition 2.11, v3.
\url{https://arxiv.org/pdf/1901.02096v3}.

\bibitem{BridgeTubach} S. Tubach,
\emph{On the Nori and Hodge realisations of Voevodsky motives},
Theorem 4.5, Corollaries 4.8 and 4.20, v4.
\url{https://arxiv.org/html/2309.11999v4}.

\bibitem{BridgeABV} J. Ayoub and L. Barbieri-Viale,
\emph{Nori 1-motives}, Theorem 8.12.
\url{https://arxiv.org/pdf/1206.5923}.

\bibitem{BridgeAudit} 费曼周期研究项目，\emph{从原始费曼积分到实际motivic coaction}，
\texttt{HYPOTHESIS\_11\_2\_AUDIT.tex}（2026）。
逐项定位\texttt{GEOMETRIC\_COMPARISON}、\texttt{K3\_KERNEL}、
\texttt{KUMMER\_BOUNDARY}、\texttt{MINIMALITY}及\texttt{LOWER\_REGULATOR}。
固定群与P8另依赖\texttt{FIXED\_PURE\_GROUP}和\texttt{RELATIVE\_TOP\_AUDIT}。
内部可审阅证明稿及既有证书，不表示新增独立重放或外部同行认证。
\bibitem{BridgePhysics} 费曼周期研究项目，
\texttt{PHI4\_FOUR\_BODY.md}及\texttt{data/phi4\_four\_body}中的
正权包围、独立重放和包容系数记录；\emph{物理窄误差认证与系统文献比较}，
\texttt{PHYSICS\_CERTIFICATION\_LITERATURE\_AUDIT.tex}（2026）。
前者的粗严格界与普通数值分开使用；后者仅新增角截断理论界，不给窄总率证书。

\end{thebibliography}
\end{document}

